\documentclass[11pt,a4paper]{ltjsarticle}
\usepackage{amsmath,amssymb,amsthm,mathtools,mathrsfs,array,longtable}
\usepackage[hidelinks]{hyperref}
\hypersetup{pdftitle={三つの有理三次元多様体の Poisson 成分と非可換双有理模型},
 pdfsubject={Poisson structures and noncommutative birational geometry}}
\newtheorem{theorem}{定理}[section]
\newtheorem{proposition}[theorem]{命題}
\newtheorem{lemma}[theorem]{補題}
\newtheorem{corollary}[theorem]{系}
\theoremstyle{definition}
\newtheorem{definition}[theorem]{定義}
\newtheorem{remark}[theorem]{注意}
\newcommand{\PP}{\mathbb P}
\newcommand{\CC}{\mathbb C}
\newcommand{\Frac}{\operatorname{Frac}}
\newcommand{\Qgr}{\operatorname{Qgr}}
\newcommand{\qgr}{\operatorname{qgr}}
\newcommand{\AD}{\operatorname{AD}}
\title{有理三次元多様体の Poisson 構造と非可換双有理幾何\\\large 成分分類・量子化・blow-up・半直交分解}
\author{}
\date{2026年11月8日\quad 統合・再検証第12版}
\begin{document}
\maketitle
\begin{abstract}
$\PP^1\times\PP^2$ と $(\PP^1)^3$ の Poisson 概型の既約成分を分類し、
$\PP^3$ を含む三模型の非可換射影量子化と双有理変換を調べる。
前者は九成分を持ち、Riccati 成分の生成点の長さは2、他は1である。
後者は八成分を持ち、全成分は生成点で被約である。
成分の網羅性、次元、生成点の構造、および指定 del Pezzo 因子への延長を本文で証明する。
さらに射影空間の任意の有限積上の Poisson 方向について、
有限型の曲線上の smooth proper dg 圏の代数的量子化を構成する。
$\PP^1\times\PP^2$ の七成分が
log-canonical な有理 Poisson 関数体に帰着し、残る二成分には
楕円型および Riccati 型の模型が得られる。
Riccati 型と非零 log-canonical 型を分離する次数評価と
非可換分数体のスカラー交換子不変量を与える。
二平面と二次曲面の対数型、line--conic 型、および三重積の対数型について、
共通の量子トーラス分数体を持つ射影量子化を明示する。
楕円型には、二点 evaluation による有限代数から二つの収縮を持つ
相対 smooth proper 圏を構成し、quadric の射影環と共通関数斜体を復元する。
この共通圏の反標準切断環は全次数で平坦であり、
Sklyanin 平面の二点 blow-up 環として表示される。
三重積の最大 pencil 成分には、八生成元・三十七二次関係式の
Segre 斉次環の代数的量子化を構成する。
また楕円曲線を中心とする共通圏に二つの半直交分解を与え、
両端点の一次 Poisson 構造と共通関数斜体を同定する。
三模型のトーラス型には同じ双指標で変形した共通 Rees 模型を構成し、
全 blow-up 段階の Orlov 型半直交分解を証明する。
Riccati 型には $S(2,3)$ 内の特殊族を与える通常の射影三空間の量子化を構成する。
さらに例外型の Weyl 関数体と、Sklyanin pencil の
中心部分体の有限拡大後の比較を記述する。
古典的な対応を滑らかな中心の blow-up と blow-down で表示し、
三次元にも適用できる非可換 blow-up と半直交分解の先行結果の仮定を明示する。
楕円 quadric の一点 blow-up では hereditary 部分圏の消滅補題により
圏論的 blow-up と切断環の射影圏を比較する。
中央正則列の Rees 環・随伴次数環を証明し、中央二変数環上の自由性と、
アフィン blow-up の基底変換・Orlov 型半直交分解を構成する。
射影圏への降下には非 perfect な頂点 torsion の問題が残る。
反標準 section functor からの圏復元は Artin--Zhang の ample 条件の下で導き、
条件を外せない可換 weak del Pezzo の反例を与える。
関数斜体の回復を零ランク対象による局所化として記述し、
同じ型の半直交成分を持つ量子平面の関数斜体が異なる例を与える。
さらに、非可換射影三空間への斜体的双有理性を AS-有理性として定式化し、
量子トーラス型と Weyl 型の異なる非可換 $\PP^3$ 関数斜体を証明する。
全ての三次元トーリック双指標量子化の AS-有理性を与える四次元量子多項式環を構成し、
有理 Casimir 体と中央元をもつ曲面模型の安定化を記述する。
\end{abstract}
\setcounter{tocdepth}{2}
\begingroup
\footnotesize
\makeatletter
\renewcommand{\@pnumwidth}{2em}
\makeatother
\tableofcontents
\endgroup
\clearpage

\section{目的、先行研究と結果の範囲}

三つの可換な有理三次元多様体を
$X_3=\PP^3$、$X_{12}=\PP^1\times\PP^2$、$X_{111}=(\PP^1)^3$ と書く。
各 $X$ の Poisson 概型は
\[
 \operatorname{Pois}(X)
 =\{\pi\in H^0(X,\wedge^2T_X):[\pi,\pi]=0\}
\]
である。成分は自己同型群とスカラー倍で割る前のものを用いる。
本書の問題は、成分の一般元を量子化した非可換多様体の間で、
分数体の同型と共通 blowup を構成することである。
古典的なファイバーと一次の Poisson 括弧を固定し、
量子化の高次項は双有理比較と両立するように選ぶ。

$X_3$ の六成分とそれらの一般元の斉次量子化は
Pym \cite[定理1.1、命題2.9、第3節]{Pym2015} による。
$X_{12}$ の九成分は第\ref{p12:classification}節、
$X_{111}$ の八成分は第\ref{cube:sec:classification}節で、
正規形から網羅性と生成点の概型構造まで証明する。
三次曲線補集合の基本群などの一般定理は、その仮定と参照箇所を明記して用いる。
古典的 Poisson 双有理変換の共通 blow-up は第\ref{sec:classical-factorizations}節で構成する。
全 Poisson 方向の代数的な圏量子化の一般存在定理は本書第\ref{sec:tilting-quantization}節に証明を記す。

非可換曲面の既知の双有理定理は
Presotto--Van den Bergh \cite{PVdB2016}、Presotto \cite{Presotto2018}、
Rogalski--Sierra--Stafford \cite{RSS2023} による。
とくに quadratic Sklyanin 平面と cubic $\mathbb Z$-代数で表される quadric の間には、
楕円曲線上の直線束と平行移動を照合した関数体の同型がある。
本書ではこの比較を、指定した古典的な曲面リンクを持つ代数的な族に組み込む。
Poisson 関数体の障害と整数係数不変量には
Goodearl--Launois \cite{GoodearlLaunois}、
非可換分数体のスカラー交換子には
Goodearl--Yakimov \cite{GoodearlYakimov} を参照する。
明示的な Ore 構成とその適用、相対的な構成、および本書の直接計算の
新規性・優先関係は未確定である。

blow-up と SOD に関しては、Van den Bergh の非可換曲面論、
滑らかな中心に沿う order の blow-up、および非可換射影束の定理を用いる。
第\ref{sec:relative-common-blowup}節では、平面の Beilinson 代数に
楕円曲線上の二点の evaluation 加群を加え、階数23の有限代数から共通圏を作る。
exceptional objects の \'etale 局所持上げと四頂点の quadric 再構成により、
古典点を含む二つの SOD と両端点の射影環が得られる。
関数斜体の比較は付値環上の有限複体の消去により証明する。
$\PP^1$ との積は、三次元の相対共通模型と二種類の SOD を与える。
第\ref{sec:anticanonical-common-ring}節では、その曲面共通圏の
反標準 orbit 環を計算する。restriction の Serre 三角と楕円曲線上の消滅により
全次数の平坦性を証明し、mutation の自然射と RSS の Hilbert 関数を用いて
平面側の二点 blow-up 環へ同定する。

第\ref{sec:toric-nc-blowups}節では、
Cirio--Landi--Szabo \cite{ToricCLS} の量子トーリック構成と
Orlov の同変性 \cite{ToricElagin} を用い、三模型の共通 blow-up を構成する。
全重みの Hom と高次積を双指標で変形し、非同変導来圏の生成性を証明する。
共通模型は長さ24の full exceptional collection を持つ。
第\ref{sec:nc-blowups}節ではさらに、
Van den Bergh の切断環比較を quadric に適用する消滅補題を証明し、
固定された一般の Sklyanin quadric について RSS の共通環に二つの SOD を与える。
第\ref{sec:sod-functions}節では、関数斜体を保持する局所化との両立条件を示す。
全ての成分間の対応を中心付きの非可換 blow-up で表示する問題は残る。


第\ref{sec:native-pencil}節では Rains \cite{PencilRainsFiltered} の
filtered deformation の族のうち、次数六の $\mathbb P^1$ ruling を用いる。
同論文の非可換 Fano 三次元多様体の構成方針に対し、本書では
一般の $(1,1,1)$ pencil の標識、相対 ruling、二重中央拡大の全次数 PBW 性、
一次 Poisson 類、および無限位数ファイバーの中心を具体的に同定する。
第\ref{pcat:sec}節では Orlov \cite{PencilOrlovGluing} の dg gluing を用い、
楕円曲線のグラフを中心とする相対共通圏を構成する。
二回の mutation により階数64の三重積の量子化と第二の楕円成分を得る。
これらの具体的な主張の新規性は未確定である。

\subsection{主張の状態と適用範囲}
本書の基本定理は、先行定理を仮定つきで使用した結果と、
ここで直接証明した補題を区別して述べる。
射影空間の任意の有限積の量子化は第\ref{sec:restored-finite-products}節で、
大域 formality と Artin 近似に依存する存在定理として示す。
反標準環からの共通圏の復元は第\ref{sec:restored-anticanonical}節の
ample 性を含む明示的な十分条件のもとで成立し、
本研究の特定の共通圏でその条件はまだすべて確認できていない。
中央 pencil の Rees 環から例外随伴次数環を構成し、
制限関手に対応する dg gluing の SOD は第\ref{sec:restored-rees-sod}節で証明する。
アフィン圏では Rees 模型と gluing が同値であるが、
非可換射影圏への導来 torsion 局所化を伴う降下は未証明である。
2026年の二文献の定理は第\ref{sec:restored-primary}節で
著者の定理として正確な適用範囲を記す。

\subsection{成分ごとの量子化模型}

$\PP^3$ について、既知の六族は次の通りである。
\[
\begin{array}{c|l}
\text{Poisson 成分}&\text{一般元の非可換射影模型}\\\hline
L(1,1,1,1)&\text{四生成元 skew-polynomial 環}\\
L(1,1,2)&\text{二次の反復 Ore 拡大}\\
R(2,2)&\text{四次元 Sklyanin 代数}\\
R(1,3)&\text{三次元 Sklyanin 代数の中心拡大}\\
S(2,3)&\text{可換三変数環の微分 Ore 拡大}\\
E(3)&\text{二次元 affine Lie 代数の CGG 量子化}
\end{array}
\]
$R(1,3)$ の一般元の関数体の比較は、本書の残る課題に含める。

三重積の八成分は次の四種類になる。
\[
\begin{array}{c|c|c|l}
\text{成分}&\text{個数}&\text{次元}&\text{本書での比較}\\\hline
\mathcal T&1&9&\text{三次元量子トーラス}\\
\mathcal S_{ij}&3&9&\text{非可換 quadric と中央 }\PP^1\text{ の積}\\
\mathcal L_i&3&10&\text{二つの交換パラメーターを持つ量子トーラス}\\
\mathcal R&1&13&\text{中央 pencil を持つ Segre 環と楕円中心の共通圏}
\end{array}
\]
$\mathcal R$ の斉次環は第\ref{sec:native-pencil}節、
二つの収縮を持つ相対共通圏は第\ref{pcat:sec}節で構成する。
両構成の導来同値と分数斜体の同定は追加の比較問題である。
$X_{12}$ の九成分の具体的な関数体は第\ref{sec:p1p2-models}節で一覧にする。

\subsection{証明した対応}

次の表では非可換関数体の同型と、中心の拡大を伴う比較を明記する。
\begin{center}
\small
\begin{tabular}{p{.28\textwidth}|p{.28\textwidth}|p{.34\textwidth}}
一方の模型&他方の模型&得られる対応\\\hline
三模型のトーラス型&同じ $q_{ij}$ を持つトーラス型&共通 Rees 模型、各段階の SOD、同じ分数体\\
$X_3$ の一般の $L(1,1,2)$&$X_{12}$ の一般の line--conic 型&明示的な Ore 分数式による同型\\
$X_{111}$ の一般の $\mathcal L_i$&一係数零のトーラス型&generalized Weyl 模型の分数体の同型\\
$X_{12}$ の一般の垂直型&$X_{111}$ の一般の二因子型&相対共通圏、二つの SOD、同じ分数体\\
$X_{111}$ の一般の最大 pencil 型&同じ楕円曲線を持つ二因子型&楕円中心の共通圏、二つの SOD、同じ斜体 $D(u)$\\
$X_3$ の一般の $E(3)$&垂直型の三重直線退化&中央な一変数上の Weyl 分数体\\
$X_{12}$ の一般の Riccati 型&$X_3$ の $S(2,3)$ 内の特殊族&微分 Ore 分数体の同型\\
四次元 Sklyanin pencil&三次元 Sklyanin 曲面&中心部分体の有限拡大後の同型
\end{tabular}
\end{center}

『一般』は、明記した代数的な非退化条件を満たす点を意味する。
平行移動の無限位数や係数の $\mathbb Q$-線形独立性を用いる箇所では、
その条件を追加して述べる。これらの条件を複素閉点の Zariski 開条件とは扱わない。

\part{基礎、変形理論と量子化}

\section{一般の射影空間の積における矩形 tilting 束と消滅}
\label{sec:products-foundation}
この節は以後の量子化の共通する可換幾何学的入力を、任意の有限積について証明する。
この節単独では量子化を主張しないが、次節で代数的存在定理へ進む。
\begin{proposition}\label{prop:products-tilting-rank}
$X=\prod_{i=1}^r\mathbb P^{n_i}$ とし、$n_i\ge1$ とする。
\[
\mathcal E=\bigoplus_{0\le a_i\le n_i}
\mathcal O_X(a_1,\ldots,a_r),\qquad A=\operatorname{End}_X(\mathcal E).
\]
このとき $\mathcal E$ は full strong tilting 束であり、
\[
\dim_{\mathbb C}A=\prod_{i=1}^r\binom{2n_i+2}{n_i}.
\]
また、全ての $p\ge0$ と $q>0$ について
\[
H^q(X,\wedge^pT_X)=0.
\]
従って HKR 同型の下で $HH^m(X)$ は次数 $m$ の大域 polyvector に一致する。
\end{proposition}
\begin{proof}
各 $\mathbb P^{n_i}$ の Beilinson の full strong exceptional collection
$\bigl(\mathcal O,\ldots,\mathcal O(n_i)\bigr)$ の外積は
$\mathcal E$ の直和分解を与え、Kunneth の公式と生成性から full strong tilting となる。
単一因子における自己準同型代数の次元は
\[
\sum_{0\le a\le b\le n}h^0(\mathbb P^n,\mathcal O(b-a))
=\sum_{d=0}^n(n+1-d)\binom{n+d}{n}
=\binom{2n+2}{n}.
\]
最後の等式は binomial coefficient の Pascal 恒等式を二度適用すれば従う。
全積では Hom の Kunneth 公式によって次元が積になる。
一方 $T_X=\bigoplus_i\mathrm{pr}_i^*T_{\mathbb P^{n_i}}$ より
$\wedge^pT_X$ は外冪の外積の直和に分解する。
$\wedge^jT_{\mathbb P^n}\simeq\Omega^{n-j}_{\mathbb P^n}(n+1)$ であり、
Bott の消滅（または捻った Euler 複体）から高次コホモロジーは消える。
再び Kunneth の公式を適用すればよい。HKR 分解の $q>0$ の項はすべて消える。
\end{proof}
\begin{remark}
この命題は $X_{12}$ で $\dim A=60$、$X_{111}$ で $\dim A=64$ を与える。
この入力から第\ref{thm:restored-product-quant}定理で、
各 Poisson 方向に沿う代数的圏量子化を構成する。
\end{remark}
\section{任意の射影空間の有限積に対する代数的圏量子化}
\label{sec:restored-finite-products}

本節は第\ref{sec:products-foundation}節と第\ref{sec:tilting-quantization}節の
議論の適用範囲を拡張する。以下の結論は、\emph{圏の量子化}についてのものであり、
指定した多重偏極の斉次環の量子化を意味しない。
また、異なる高次項をもつ量子化が一意に定まるという主張でもない。

\begin{theorem}[有限積への代数的量子化の拡張]\label{thm:restored-product-quant}
$k=\mathbb C$、$X=\prod_{i=1}^r\mathbb P^{n_i}$、$d=\sum_i n_i$ とおき、
$\pi\in H^0(X,\wedge^2T_X)$ が $[\pi,\pi]=0$ を満たすとする。
$\mathcal E=\bigoplus_{0\le a_i\le n_i}\mathcal O_X(a_1,\ldots,a_r)$、
$A_0=\operatorname{End}_X(\mathcal E)$ とする。
すると、滑らかな点付き有限型代数曲線 $(B,b)$ と
$B$ 上の有限局所自由な結合代数 $\mathcal A$ が存在して、次を満たす。
\begin{enumerate}
\item $\mathcal A_b\simeq A_0$ および
$\operatorname{Perf}(\mathcal A_b)\simeq\operatorname{Perf}(X)$。
\item $\operatorname{rank}_{\mathcal O_B}\mathcal A
=\prod_i\binom{2n_i+2}{n_i}$。
$B$ を縮めれば、すべての幾何学的ファイバーの大域次元は高々 $d$、
かつ $\mathcal A$ は相対双加群として perfect である。
従って対応する perfect dg 圏は相対的に smooth proper である。
\item 接ベクトル $v\in T_bB$ の Kodaira--Spencer 類は
$\operatorname{HH}^2(X)$ の HKR 分解の下で $\pi$ である。
\end{enumerate}
さらに $A_0$ の形式的量子化 $\widehat\mu$ を一つ固定すれば、
任意の整数 $N\ge1$ に対して、この代数曲線と族を選び直し、
形式的な基底・corner 表示で結合積の $N$-jet を
$\widehat\mu$ の $N$-jet に一致させられる。
ここでの最後の主張は、\emph{各有限 $N$ ごと}の代数化であり、
一つの代数的な族が任意の $N$ に対して同じ指定形式解を与えるという意味ではない。
\end{theorem}
\begin{proof}
第\ref{prop:products-tilting-rank}命題により、矩形 tilting 束の高次自己拡大は零で、
$\operatorname{Perf}(X)$ と $\operatorname{Perf}(A_0)$ は Morita 同値である。
$T_X$ の外冪の高次コホモロジーはすべて消えるため、
大域 polyvector の Schouten dg Lie 代数
\[
\mathfrak g=\bigoplus_{p\ge0}H^0(X,\wedge^pT_X)[1-p],\qquad d_{\mathfrak g}=0
\]
は derived global polyvector dg Lie 代数へ準同型となる。
Van den Bergh の大域 formality と tilting の Hochschild 変形理論
\cite{VdBGlobal,BW2024} によって、これは $A_0$ の
normalized Hochschild cochain dg Lie 代数と $L_\infty$ 同値である。
$h\pi$ は Schouten Maurer--Cartan 方程式を満たすので、この同値で
$A_0[[h]]$ 上の形式的な結合積 $\widehat\mu$ に移る。
その一次 Hochschild 類は $\pi$ に対応する。

$A_0$ の原始直交冪等元 $e_a$ は有限個で、$E=\bigoplus_a ke_a$ は分離的である。
形式変形では冪等元が持ち上がり、$E$-relative な normalized cochain 複体を用いると、
各 $e_bA_0e_a$ の階数を固定した corner の積に移せる。
半順序 $a\le b$ を満たさない corner は零のままである。
その積は、有限個の構造定数が満たす有限個の結合律の多項式方程式に等しい。
この方程式の形式解 $\widehat\mu$ に Artin 近似を適用すると、
任意の $N$ について $k[h]_{(h)}$ の Hensel 化上の解 $\mu_N$ が得られ、
$\mu_N\equiv\widehat\mu\pmod {h^{N+1}}$ とできる
\cite{StacksArtin}。
Hensel 化は有限型の étale 近傍のフィルター付き極限であり、
有限個の構造定数はその一つに降下する。
従って $(\mathbb A^1,0)$ の滑らかな étale 近傍
$(B,b)$ 上で、要求した結合代数の族が得られる。
$N\ge1$ としたので一次 Kodaira--Spencer 類が保持される。

残る滑らかさと大域次元を確認する。
矩形の半順序の狭義鎖の長さは高々 $d$ である。
$J=\bigoplus_{a<b}e_b\mathcal Ae_a$ は
$J^{\otimes_E(d+1)}=0$ を満たす。
$E$ は基底環上の分離的代数 $\mathcal O_B^{\oplus |I|}$ であるから、
\[
 0\to \mathcal A\otimes_EJ^{\otimes_Ed}\otimes_E\mathcal A
 \to\cdots\to\mathcal A\otimes_E\mathcal A\to\mathcal A\to0
\]
という相対 normalized bar 分解があり、全項は
$\mathcal A$-双加群として相対射影的で、全長は $d$ である。
同じ構成を左加群に適用して各幾何学的ファイバーの
大域次元が高々 $d$ であることが従う。
$\mathcal A$ は有限局所自由なので properness も従う。
階数は第\ref{prop:products-tilting-rank}命題で計算済みである。
\end{proof}

\begin{remark}[制約と証明の所在]
形式量子化を代数的曲線へ写す上の主張は、
大域 formality と Artin 近似という引用定理に依存する。
$\operatorname{Pois}(X)$ の全ての成分の概型同型や、
一意な代数化、多重偏極の保存、AS-regular な ambient の存在は主張しない。
特に任意の有限積の場合に、ある一つの連結した代数的族で
すべての Poisson 方向を同時にパラメータ化することまではここでは証明していない。
\end{remark}

\section{量子化の定式化と先行研究}

基礎体を $\CC$ とする。点付き代数曲線 $(B,b_0)$ 上の非可換変形が
$\pi\in H^0(X,\wedge^2T_X)$ を量子化するとは、古典的なファイバーを
$X$ と同定したときの一次変形類が HKR 同型を通して $\pi$ になることをいう。
斉次座標環 $R$ を変形する場合は、局所関数について
\[
 \left.\frac{\widetilde f\widetilde g-\widetilde g\widetilde f}{t}
 \right|_{t=0}=\{f,g\}_{\pi}
\]
という通常の半古典極限で同じ条件を指定できる。
ここで $t$ は $b_0$ で一次の零点を持つ局所パラメーターである。

第\ref{sec:tilting-quantization}節の一般定理では、非可換多様体を
有限次元代数の perfect dg 圏によって表す。
古典的な点の近傍で、全ての Poisson 方向を含む有限型代数的な族が存在する。
この族は選んだ tilting bundle と古典的ファイバーの標識を伴う。
第\ref{sec:toric-quantization}節と第\ref{sec:surface-quantization}節では、
次数ごとに平坦な connected graded algebra とその $\operatorname{qgr}$ を用いる。
得られた範囲を次に記す。
\[
\begin{array}{c|c|l}
\text{模型}&\text{Poisson 構造}&\text{代数的な構成}\\\hline
\PP^1\times\PP^2, (\PP^1)^3&\text{全て}&
     \text{tilting 代数の \'etale 局所変形}\\
\PP^1\times\PP^2, (\PP^1)^3&\text{トーラス型の一般元}&
     \text{Cox 環の双指標 twist}\\
\PP^1\times\PP^2&V_{\mathrm{vert}}\text{ の全て}&
     \text{quadratic AS-regular 曲面との積}\\
(\PP^1)^3&\mathcal S_{ij}\text{ の全て}&
     \text{非可換 quadric との積}\\
(\PP^1)^3&\mathcal L_i\text{ の一般元}&
     \text{同変 quadric と Ore 拡大}\\
(\PP^1)^3&\mathcal R\text{ の一般元}&
     \text{次数六楕円代数の二重中央拡大}
\end{array}
\]
成分の分類と記号は本書の成分分類 による。
一般存在定理は、これらの成分分類の網羅性を仮定しない。

大域的量子化には Kontsevich \cite{KontsevichAlgebraic} と
Van den Bergh \cite{VdBGlobal} の formality を用いる。
tilting bundle による変形理論の記述は
Barmeier--Wang \cite[\S3.2]{BW2024} に詳しい。
本稿の一般存在証明は、この形式変形理論を有限次元の結合律の概型に移し、
Artin 近似 \cite{StacksArtin} で有限型の基底に置き換えるものである。
矩形の半順序を用いることで、全ファイバーの相対 smoothness と
大域次元の上界を直接証明できる。

明示的な構成は、量子トーリック多様体 \cite{CLS}、
三次元 quadratic AS-regular 代数 \cite{MoriSmith}、
Van den Bergh の非可換 quadric \cite{VdBQuadrics} に依存する。
二因子型では cubic $\mathbb Z$-代数の二周期的な表示を用い、
対数型ではその同変性を利用する。
先行研究の一般理論と今回の具体化の関係は各証明中に明記した。
この具体化の新規性および優先関係は未確定である。

\section{全ての Poisson 方向の代数的圏量子化}
\label{sec:tilting-quantization}

この節では、滑らかで固有な dg 圏の変形という意味での非可換代数多様体を扱う。
半古典極限は
\[
 HH^2(X)\simeq
 H^0(X,\wedge^2T_X)\oplus H^1(X,T_X)\oplus H^2(X,\mathcal O_X)
\]
による Kodaira--Spencer 類で指定する。
以下の二つの多様体では後の二項は零である。
標識した古典的ファイバーの変形底を用い、自己同値による大域的な商は別に考える。

\subsection{有限次元の tilting 代数}

$X=\mathbb P^1\times\mathbb P^2$ の場合には
\[
 \mathcal E_X=\bigoplus_{\substack{0\leq a\leq1\\0\leq b\leq2}}\mathcal O_X(a,b),
 \qquad
 A_X=\operatorname{End}_X(\mathcal E_X)
\]
とし、$X=(\mathbb P^1)^3$ の場合には
\[
 \mathcal E_X=\bigoplus_{a,b,c\in\{0,1\}}\mathcal O_X(a,b,c),
 \qquad
 A_X=\operatorname{End}_X(\mathcal E_X)
\]
とする。接束を $T_X$ と書く。
射影空間上の標準的な full strong exceptional collection の外積を用いているので、
\[
 \operatorname{Perf}(X)\simeq\operatorname{Perf}(A_X).
\]
$B_n=\operatorname{End}_{\mathbb P^n}
(\bigoplus_{i=0}^{n}\mathcal O(i))$ と置けば、それぞれ
$A_X=B_1\otimes B_2$、$A_X=B_1^{\otimes3}$ である。

直和因子を添字集合 $I$ の元 $\alpha$ で表す。
$I$ は $\{0,1\}\times\{0,1,2\}$ または $\{0,1\}^3$ であり、
座標ごとの半順序を入れる。
\[
 V_{\beta\alpha}
 =\operatorname{Hom}_X(\mathcal O_X(\alpha),\mathcal O_X(\beta))
\]
は $\alpha\leq\beta$ の場合だけ非零である。
$X=\prod_i\mathbb P^{n_i}$ と書くと
\[
 \dim V_{\beta\alpha}
 =\prod_i\binom{n_i+\beta_i-\alpha_i}{n_i}.
\]
対角部分 $E=\bigoplus_{\alpha\in I}\mathbb C e_\alpha$ と、
狭義上三角部分 $J=\bigoplus_{\alpha<\beta}V_{\beta\alpha}$ を固定する。
各 $V_{\beta\alpha}$ に内部次数 $|\beta|-|\alpha|$ を与える。

\begin{lemma}\label{lem:directed-quantization}
$E$ の積、$E$ による左右作用、および上記の corner 空間を保つ任意の結合的積に対し、
$A=E\oplus J$ は有限次元であり、$A$ の大域次元は $3$ 以下である。
相対 normalized bar resolution は長さ $3$ の有限生成射影 $A$-双加群分解となる。
同じ構成は任意の可換基底環上で成り立ち、相対的な smoothness と properness を与える。
\end{lemma}

\begin{proof}
半順序集合 $I$ の狭義鎖の長さは高々 $3$ であるから
$J^{\otimes_E4}=0$ である。
射影 $A\to E$ は代数準同型であり、相対 normalized bar resolution は
\[
 0\longrightarrow A\otimes_EJ^{\otimes_E3}\otimes_EA
 \longrightarrow A\otimes_EJ^{\otimes_E2}\otimes_EA
 \longrightarrow A\otimes_EJ\otimes_EA
 \longrightarrow A\otimes_EA
 \longrightarrow A\longrightarrow0
\]
となる。この複体の完全性は通常の bar contraction による。
各項は $Ae_\beta\otimes_{\mathbb C}e_\alpha A$ の有限直和であるから
$A$-双加群として射影的である。
同様の左加群 bar resolution を任意の $A$-加群に適用すれば大域次元の上界を得る。
可換基底環上でも corner 空間を有限自由加群に置き換えることで同じ証明が使える。
\end{proof}

古典的な積は Koszul である。Koszul 条件は、この有限 bar 複体について
$i\ne j$ に対する $\operatorname{Tor}^{A}_i(E,E)_j=0$ と書ける。
内部次数は高々 $3$ なので有限個の消失条件であり、古典的な積の Zariski 開近傍で保たれる。
Koszul 性のもとでは、corner Hilbert 行列の逆行列から
生成元、二次関係式、三次シジジーの数が決まる。
\[
\begin{array}{c|r|r|r|r|r}
 X&\dim A_X&\text{頂点}&\text{矢}&\text{二次関係式}
      &\text{三次シジジー}\\ \hline
 \mathbb P^1\times\mathbb P^2&60&6&18&18&6\\
 (\mathbb P^1)^3&64&8&24&24&8
\end{array}
\]
最後の列は極小頂点と極大頂点の間の $\operatorname{Tor}_3$ の次元である。
従ってこの開近傍では大域次元はちょうど $3$ となる。
数値は \texttt{check\_tilting.py} で整数演算により再現できる。

\subsection{Jacobi 概型と半普遍変形底}

\[
 C_X=\operatorname{Pois}(X)
 =\{\pi\in H^0(X,\wedge^2T_X)\mid[\pi,\pi]=0\}
\]
を Jacobi 方程式が定義するアフィン概型とする。
この記号にはその非被約構造も含める。

\begin{proposition}\label{prop:poisson-hull}
$A_X$ の結合代数変形の代数的な半普遍底 $(S_X,s_0)$ を取り、
古典的な点での標識を固定できる。
このとき
\[
 \widehat{\mathcal O}_{S_X,s_0}
 \simeq
 \widehat{\mathcal O}_{C_X,0}
\]
という完備局所環の同型があり、接空間上の同型は HKR 同型
$HH^2(A_X)\simeq H^0(X,\wedge^2T_X)$ と整合する。
\end{proposition}

\begin{proof}
$\mathbb P^1$ 上の外積接束は $\mathcal O,\mathcal O(2)$ であり、
$\mathbb P^2$ 上では $\mathcal O,T_{\mathbb P^2},\mathcal O(3)$ である。
Euler 完全系列と K\"unneth の公式から
\[
 H^q(X,\wedge^pT_X)=0\qquad(q>0)
\]
を得る。従って global polyvector field の dg Lie 代数
\[
 \mathfrak g_X^{p-1}=H^0(X,\wedge^pT_X),\qquad d=0
\]
は derived global polyvector field を計算する。
Van den Bergh の大域的 formality 定理
\cite[定理1.1および式(1.3)]{VdBGlobal} と tilting による変形理論の同値
\cite[定理3.18、命題3.21]{BW2024} により、
$\mathfrak g_X$ は $A_X$ の Hochschild dg Lie 代数と $L_\infty$ 準同型の意味で同値である。
$E$ は分離的なので、$E$-relative normalized Hochschild 複体を用いてよい。
実際、形式変形の中で $E$ の原始冪等元は持ち上がり、その持ち上げは
古典的なファイバー上で恒等な共役を除いて一意である。
平坦性によって各 corner の階数も保たれるので、全ての結合代数の形式変形を
この固定した $E$ と corner 空間を持つ表示に移せる。

この同値を、有限型の変形底に結び付ける。
$C_E^n=\operatorname{Hom}_{E\text{-}E}(J^{\otimes_E n},A_X)$ とし、
$\mu_0$ を古典的な積、$d=[\mu_0,-]$ と書く。
$\mu_0+C_E^2$ の中で結合律を満たす積の概型を $M_X$ とする。
これは有限個の二次方程式が定義するアフィン概型である。
$B_E^2=dC_E^1$ の線形補空間 $L\subset C_E^2$ を選び、
\[
 S_X=M_X\cap(\mu_0+L)
\]
と置くと、$T_{s_0}S_X\simeq HH^2(A_X)$ となる。

この切片が hull を与えることを確認する。
corner ごとの基底変換群が $M_X$ に作用し、その無限小作用の像は $B_E^2$ である。
$C_E^1$ 内で $\ker d$ の補空間を選ぶと、この補空間から $B_E^2$ への微分は同型である。
形式的陰関数定理によって、古典的な積の形式近傍で基底変換方向を取り除ける。
小拡大に沿って基底変換を持ち上げ、$B_E^2$ 成分を逐次消せるので、
$\widehat{S_X}_{s_0}$ から結合代数の変形関手への射は形式的に滑らかであり、
接空間上で同型を与える。

一方、$\mathfrak g_X$ の微分は零であるため、その Maurer--Cartan 方程式は
まさに $[\pi,\pi]=0$ である。
gauge 群は小拡大に沿って持ち上がり、原点の一次接空間には自明に作用するので、
$\widehat{C_X}_0$ も対応する変形関手の hull である。
ここで用いる gauge 群は、基底の極大イデアルを $\mathfrak m$ として
$\exp(\mathfrak m\otimes\mathfrak g_X^0)$ であり、標識した古典的ファイバー上で恒等となる。
上記の $L_\infty$ 同値と hull の一意性
\cite[命題2.9]{Schlessinger} によって所要の完備局所環の同型を得る。
\end{proof}

\begin{theorem}\label{thm:algebraic-quantization-all}
$X=\mathbb P^1\times\mathbb P^2$ または $(\mathbb P^1)^3$ とする。
点付き有限型アフィン概型 $(U,u_0)$ と、両方とも \'etale な射
\[
 (C_X,0)\xleftarrow{\ p\ }(U,u_0)
 \xrightarrow{\ q\ }(S_X,s_0)
\]
および $U$ 上の有限自由な結合代数 $\mathcal A$ が存在し、次を満たす。
\begin{enumerate}
 \item $\mathcal A_{u_0}\simeq A_X$ であり、
 $\operatorname{Perf}(\mathcal A_{u_0})\simeq\operatorname{Perf}(X)$ である。
 \item $\mathcal A$ の階数はそれぞれ $60,64$ である。
 $U$ を縮小すると全ての幾何学的ファイバーは Koszul、大域次元 $3$ となり、
 $\operatorname{Perf}(\mathcal A)$ は相対的に smooth かつ proper である。
 \item 古典的な点の Kodaira--Spencer 写像は同型
 \[
 T_{u_0}U\simeq H^0(X,\wedge^2T_X)
       \xrightarrow{\ \sim\ }HH^2(A_X)
 \]
 である。その次元はそれぞれ $34,27$ である。
 \item 任意の $\pi\in C_X(\mathbb C)$ に対し、滑らかな点付き代数曲線
 $(B,b_0)$ 上の $\mathcal A$ の引戻しで、
 局所座標 $t$ に関する一次変形類が $\pi$ となるものが存在する。
\end{enumerate}
\end{theorem}

\begin{proof}
命題 \ref{prop:poisson-hull} の完備局所環の同型に Artin 近似を適用する。
具体的には \cite[補題37.39.2、Tag 0CAT]{StacksArtin} を $N=2$ で用いる。
これにより、与えた同型と一次まで整合する共通の \'etale 近傍 $U$ を得る。
$S_X\subset M_X$ の普遍的な積を $U$ に引き戻して $\mathcal A$ を定める。
その基礎加群は $A_X\otimes_{\mathbb C}\mathcal O_U$ である。
補題 \ref{lem:directed-quantization} とその後の Koszul 開条件が第二項を与える。

HKR 次元は
\[
 h^0(\wedge^2T_{\mathbb P^1\times\mathbb P^2})
 =3\cdot8+10=34,\qquad
 h^0(\wedge^2T_{(\mathbb P^1)^3})=3\cdot3\cdot3=27
\]
である。Artin 近似を二次の極大イデアルまで整合させたので、
Kodaira--Spencer 写像は指定した HKR 同型になる。

最後に直線 $\mathbb A^1\to C_X$, $t\mapsto t\pi$ に沿って $U$ を基底変換する。
点 $(u_0,0)$ の近傍 $B$ は $\mathbb A^1$ 上 \'etale なので滑らかな代数曲線である。
この曲線に沿う一次変形類は $\pi$ である。
\end{proof}

\begin{remark}\label{rem:quantization-scope}
定理は、全ての Poisson 成分を含む、古典的な点の \'etale 局所的な代数的圏量子化を与える。
命題 \ref{prop:poisson-hull} は Jacobi 概型の非被約構造を保持するので、
形式局所的な既約成分と冪零方向も半普遍変形底に引き継がれる。
特定の形式的 star product の全係数との一致を指定する場合には、
その形式族自体の代数化を追加して調べる必要がある。
また、偏極、多重斉次座標環、あるいは指定した connected AS-regular ambient を伴う構成では、
それらの追加構造も変形の対象に含める必要がある。
この節で得た代数の表示は、有限次元の tilting 代数とその有限本の関係式によるものである。
\end{remark}


\section{トーラス型の明示的な代数的量子化}
\label{sec:toric-quantization}

$X=\prod_{a=1}^r\PP^{n_a}$ とし、$(n_1,\ldots,n_r)=(1,2)$ または
$(1,1,1)$ とする。稠密トーラスの座標を $x_1,x_2,x_3$ と書く。
\[
 \pi_\lambda=\sum_{1\leq i<j\leq3}\lambda_{ij}
 (x_i\partial_{x_i})\wedge(x_j\partial_{x_j})
\]
の量子化は有限個の二次関係式で記述できる。

各因子の Cox 変数を $z_{a,0},\ldots,z_{a,n_a}$ とする。
$\deg z_{a,j}=e_a\in\mathbb Z^r$ とし、トーラス重みは
$\operatorname{wt}(z_{a,0})=0$、残りの三変数の重みは
$M=\mathbb Z^3$ の基底に取る。
\[
 B=\CC[q_{12}^{\pm1},q_{13}^{\pm1},q_{23}^{\pm1}],\qquad
 c(m,n)=\prod_{i<j}q_{ij}^{m_i n_j}
\]
と置く。通常の Cox 環の基底変換
$B\otimes_\CC\CC[z_{a,j}]$ に対し、重み斉次な元の積を
\[
 f\star g=c(\operatorname{wt}f,\operatorname{wt}g)fg
 \tag{*}
\]
で定め、この多重斉次環を $A_B$ と書く。

\begin{proposition}\label{prop:toric-algebraic}
$A_B$ は $B$ 上平坦な多重斉次環であり、各幾何学的ファイバーは
Cox 変数の個数を大域次元とする量子多項式環である。
対角部分環
\[
 R_B=\bigoplus_{d\geq0}(A_B)_{(d,\ldots,d)}
\]
は次数 $1$ で生成される noetherian 斉次環であり、各次数は $B$ 上自由である。
その階数は
\[
 \operatorname{rank}_B(R_B)_d=
 \begin{cases}
 (d+1)\binom{d+2}{2}&X=\PP^1\times\PP^2,\\
 (d+1)^3&X=(\PP^1)^3.
 \end{cases}
\]
$q_{12}=q_{13}=q_{23}=1$ のファイバーは $X$ の Segre 座標環である。
各幾何学的ファイバー $R_q$ の tails は完全強例外列をもち、
$D^b(\operatorname{qgr}R_q)$ は滑らかで固有な三次元の非可換射影多様体を与える。
より具体的には、大域次元 $3$ の有限次元代数の導来圏に同値である。
\end{proposition}

\begin{proof}
双指標の恒等式
\[
 c(m,n)c(m+n,p)=c(n,p)c(m,n+p)
\]
から積の結合性が従う。順序を固定した Cox 単項式は $B$ 基底であり、
二変数の交換は可逆なスカラーを掛ける操作である。
従って $A_B$ は反復 skew polynomial extension である。
各体上のファイバーの通常の Koszul 複体は量子 Koszul 複体に変形し、
Cox 変数の個数を長さとする最小分解を与える。

対角次数の単項式を選べば自由性と階数公式が従う。
各ブロックから一変数ずつ選んだ積は $(1,\ldots,1)$ 次であり、
それらが $R_B$ を生成する。それらの間にもスカラー交換関係があるので、
$R_B$ は有限生成量子多項式環の商である。従って noetherian である。

すべての Cox 変数は正規元であり、上記 Segre 単項式も正規元である。
各 Segre 単項式の Ore 局所化の次数零部分は、各ブロックの選んだ変数を
逆にした Cox 環の多重次数零部分と一致する。
これによって $\operatorname{qgr}R_q$ は Cox 環の多重射影化に一致する。
量子トーリック多様体を正規単項式の局所化で貼り合わせる方法は
\cite[\S\S3--5]{CLS} にも記述されている。

この圏の多重シフトを $\mathcal O(\mathbf a)$ と書く。
Cox の \v Cech 複体を各トーラス重みに分けると、可換模型との相違は
非零スカラーのみであり、双指標恒等式による基底変更で可換複体と同型になる。
従って
\[
 \dim\operatorname{Ext}^k\bigl(\mathcal O(\mathbf a),
       \mathcal O(\mathbf b)\bigr)
 =\dim H^k\bigl(X,\mathcal O(\mathbf b-\mathbf a)\bigr).
\]
特に
\[
 T_q=\bigoplus_{0\leq a_i\leq n_i}\mathcal O(\mathbf a)
\]
は強例外列の直和になる。各ブロックの変数の量子 Koszul 複体は、
そのブロックの零集合に台をもつ余核をもち、この余核は多重射影化で零になる。
従って各方向の任意のシフトは上記有限個のシフトから生成される。
$A_q$ の有限大域次元により全ての有限生成多重斉次加群は有限の自由分解をもち、
$T_q$ の生成性が従う。

$\operatorname{End}(T_q)$ は可換模型の矩形 Beilinson 代数の
トーラス重みによる双指標 twist である。
可換模型の各頂点単純加群の最小射影分解は重み斉次に取れ、
同じ twist はこの分解の完全性と最小性を保つ。
従って全ての頂点単純加群の射影次元は保たれ、
$\operatorname{gldim}\operatorname{End}(T_q)=\sum_i n_i=3$ となる。
この有限次元代数の有限大域次元から導来圏の滑らかさと固有性が従う。
\end{proof}

\begin{corollary}\label{cor:toric-semiclassical}
任意の $\lambda_{ij}\in\CC$ に対し
\[
 q_{ij}=1+h\lambda_{ij},\qquad
 B_\lambda=\CC\left[h,
 \left(\prod_{i<j}(1+h\lambda_{ij})\right)^{-1}\right]
\]
と基底変換すると、$\pi_\lambda$ を半古典極限にもつ代数的な族を得る。
$\lambda_{ij}$ 自身も独立変数としてよい。
\end{corollary}

\begin{proof}
重み斉次な $f,g$ の重みを $m,n$ とすると
\[
 \left.\frac{f\star g-g\star f}{h}\right|_{h=0}
 =\left(\sum_{i<j}\lambda_{ij}(m_i n_j-m_j n_i)\right)fg.
\]
これは $\pi_\lambda$ の括弧である。
Cox の次数零局所化でも同じ式が成り立つ。
上の基底は $h=0$ を含む有限型の滑らかなアフィン曲線であり、
すべての積と関係式はこの曲線上で代数的である。
\end{proof}

この構成に $\operatorname{Aut}(X)$ による座標変換を合わせると、
本書 のトーラス成分の開稠密な族を量子化できる。
三重積の一般のトーラス型については、三つの係数と
$\operatorname{PGL}_2^3$ の座標変換を合わせた $9$ 次元の Poisson 族を覆う。
$\PP^1\times\PP^2$ では同様に $11$ 次元の成分を覆う。
本節の明示式が直接扱うのは最大トーラスに共役な作用であり、
閉包の全ての境界点に対する量子化は一般存在定理による。


\section{曲面の量子化から得られる斉次環}
\label{sec:surface-quantization}
\subsection{平面の三次式からの垂直型の量子化}
\label{subsec:quant-plane}

三次式 $f\in\mathbb C[X_0,X_1,X_2]_3$ に対応する
$\mathbb P^2$ の Poisson 構造を $\sigma_f$ と書く。
アフィン座標 $u=X_1/X_0$、$v=X_2/X_0$ では
$\sigma_f=f(1,u,v)\partial_u\wedge\partial_v$ と規格化する。
二次式の対称化を $g\mapsto\widehat g$ とし、
\[
 A_{f,t}=\mathbb C[t]\langle X_0,X_1,X_2\rangle\Big/
 \left(
 [X_0,X_1]-\frac t3\widehat{f_{X_2}},\quad
 [X_1,X_2]-\frac t3\widehat{f_{X_0}},\quad
 [X_2,X_0]-\frac t3\widehat{f_{X_1}}
 \right)
 \tag{E1}
\]
と置く。Mori--Smith \cite[\S1.7.3、定理1.4、\S1.8.2]{MoriSmith}
の構成と正則性判定により、$0\in U\subset\mathbb A^1_t$ なる
Zariski 開集合上で、これは次数ごとに平坦な三次元 quadratic AS-regular
代数の族である。各次数 $n$ の階数は $\binom{n+2}{2}$ である。
$f=0$ および特異な零三次曲線も許される。

正則性の開性はここでは有限の条件で判定できる。
三つの関係式は、交代三次テンソルに $-(t/3)\widehat f$ を加えた
superpotential の巡回偏微分である。
関係式の独立性を保つ開集合では標準性を満たし、AS 正則性は
対応する $3\times3$ 行列の $2\times2$ 小行列式が
$\mathbb P^2$ に共通零点を持たないことと同値である。
この条件は射影性から開条件であり、$t=0$ の多項式環で成立する。
各次数の商の次元がこの開集合上で一定なので、各斉次部分は局所自由である。

関係式から、アフィン錐上の半古典極限は
\[
 \{X_0,X_1\}=\frac13 f_{X_2},\qquad
 \{X_1,X_2\}=\frac13 f_{X_0},\qquad
 \{X_2,X_0\}=\frac13 f_{X_1}.
\]
Euler の恒等式を使うと
\[
 \{u,v\}
 =\frac{X_0 f_{X_0}+X_1f_{X_1}+X_2f_{X_2}}{3X_0^3}
 =f(1,u,v).
\]
従って $\operatorname{qgr}A_{f,t}$ は指定した $\sigma_f$ を量子化する。
さらに $P=\mathcal O_U[Y_0,Y_1]$ との Segre 積
\[
 C_{f,t}=A_{f,t}\mathbin{\circ}P
   =\bigoplus_{n\geq0}(A_{f,t})_n\otimes_{\mathcal O_U}P_n
 \tag{E2}
\]
を取ると、$t=0$ は $\mathbb P^2\times\mathbb P^1$ の
$\mathcal O(1,1)$ に関する切断環である。
半古典極限の混合括弧は零であり、垂直型成分の任意の元を得る。
次数 $n$ の階数は $(n+1)\binom{n+2}{2}$ である。

この積の圏は、$\mathbb P^1_U$ 上の定数族 $\operatorname{qgr}A_{f,t}$
としても記述できる。アフィン開集合ごとには中央多項式拡大を取り、
その間を通常の $\mathbb P^1$ の座標変換で貼り合わせる。
各因子の ample sequence の外積から、同じ圏が上の Segre 切断環で復元される。
Segre 積と二重次数付き商圏の同値については
\cite[\S3.2、定理3.32]{BallardFarman} を用いる。

\subsection{双二次式からの二因子型の量子化}
\label{subsec:quant-quadric}

二次元ベクトル空間 $A,B$ を固定する。
$\mathbb P(A^*)\times\mathbb P(B^*)$ の Poisson 構造は、
体積形式を選ぶと $D\in\operatorname{Sym}^2 A\otimes\operatorname{Sym}^2 B$
で表される。$A$ の基底を $a_0,a_1$、$B$ の基底を $b_0,b_1$ とし、
$A\otimes B\otimes A\otimes B$ において
\begin{align*}
 w_0={}&a_0b_0a_1b_1-a_1b_0a_0b_1\\
       &-a_0b_1a_1b_0+a_1b_1a_0b_0
\end{align*}
と置く。ここで並置はテンソル積を意味する。
$D$ を第一・第三因子および第二・第四因子に関して対称化した
テンソルを $\widetilde D$ と書き、
\[
                   w(t)=w_0-t\widetilde D
                   \tag{E3}
\]
と定める。

$V_i=A$ または $B$ を $i$ の偶奇で定め、$w_i(t)$ を $w(t)$ の
$i$ 回の巡回移動とする。$V_{i+3}^*$ による $w_i(t)$ の縮約の像を
三次の関係式とし、これらでテンソル $\mathbb Z$-代数を割ったものを
$\mathcal A_{D,t}$ と書く。
Van den Bergh の geometric quintuple の判定と分類定理
\cite[\S5.2、定理5.2.2、定義5.3.1、定理5.5.10]{VdBQuadrics} により、
$0$ を含む Zariski 開集合上で全てのファイバーが
三次元 cubic AS-regular $\mathbb Z$-代数になる。
実際 geometric 条件の失敗は、隣接二因子の非零共変ベクトルによる
縮約が零になるという射影的 incidence の像であり閉条件である。
$w_0$ は linear quadric のテンソルなのでこの閉集合の外にある。
定理によって全ファイバーの各斉次部分の次元は一定であり、
関係式の行列の階数も一定になるため、族は次数ごとに平坦である。

$w(t)$ は二回の巡回移動で不変なので、$\mathcal A_{D,t}$ は
選んだ $V_{i+2}=V_i$ によって二周期的である。
従って偶数 Veronese
\[
 B_{D,t}=\bigoplus_{n\geq0}(\mathcal A_{D,t})_{0,2n}
 \tag{E4}
\]
には通常の $\mathbb N$-次数付き代数構造が入り、
\[
 \dim(B_{D,t})_n=(n+1)^2,\qquad
 B_{D,0}=\bigoplus_{n\geq0}
 H^0(\mathbb P^1\times\mathbb P^1,\mathcal O(n,n)).
\]
この Veronese は同じ非可換 quadric の圏を与える。

この族の Kodaira--Spencer 写像は $D$ に対応する Poisson 構造である。
ここではその規格化も確かめる。
写像は $\operatorname{SL}(A)\times\operatorname{SL}(B)$-同変であり、
始域および
\[
 HH^2(\mathbb P^1\times\mathbb P^1)
 =H^0(\mathcal O(2,2))
 =\operatorname{Sym}^2A\otimes\operatorname{Sym}^2B
\]
は同じ既約表現である。従って最高ウェイトベクトル
$D=a_0^2b_0^2$ で確認すればよい。
全ての位置の $a_0,b_0$ に対応する矢を universal localization で可逆にし、
それらによって各対象を同定する。
各位置に残る矢を $u_k,v_k$ とすると、三次関係式のうち二つは
$u_{k+1}=u_k$ および $v_{k+1}=v_k$ を与える。
残る関係式は、共通の座標 $u,v$ に対し
\[
                         uv-vu=t
\]
となる。この局所表示は Weyl 代数の族であり、PBW 基底 $u^iv^j$ を持つ。
古典的な点では $u=a_1/a_0$、$v=b_1/b_0$ なので、半古典極限は
$D(1,u,1,v)\partial_u\wedge\partial_v$ である。

以上から
\[
 C_{D,t}=B_{D,t}\mathbin{\circ}\mathcal O_U[Y_0,Y_1]
 \tag{E5}
\]
は三重積の任意の二因子型 Poisson 構造を量子化する代数的な族である。
各次数の階数は $(n+1)^3$、$t=0$ の切断環は
$\mathcal O(1,1,1)$ に関する Segre 環である。
$D$ の零曲線の滑らかさは仮定しない。

\subsection{同変性と対数型への適用}
\label{subsec:quant-elliptic-equivariance}

式 (E3) は二因子の線形座標変換と両立する。
特に代数群が $w_0$ と $\widetilde D$ を同じ指標で変換するなら、
全ての $B_{D,t}$ にその群の次数付き作用がある。
三重積の一般の対数型では、二つの $(1,1)$ 因子の交点を
$(0,0)$ と $(\infty,\infty)$ に送って
\[
 F=f x+g z,\qquad G=h x+k z
\]
と書ける。本書 の対数型の表示を用いると
\[
 \pi=\sigma+y\partial_y\wedge V,\qquad
 \sigma=aFG\partial_z\wedge\partial_x,\qquad
 V=\kappa(x\partial_x+z\partial_z),\quad
 \kappa=b(fk-gh)
 \tag{E6}
\]
である。$(u,v)=(x,z)$ に対して $D=-aFG$ と取る。
$\sigma$ は radial $\mathbb G_m$ 作用で不変である。
斉次座標の重みを $(0,1)$ と $(0,1)$ に取ると、
$w_0$ と $\widetilde D$ はともに重み $2$ を持つため、
上の量子化は同じ $\mathbb G_m$ 作用 $\rho$ を保つ。

$q=1+t\kappa$ を可逆にした開集合に制限し、
\[
 T=B_{D,t}[Y_0][Y_1;\rho_q],\qquad
 Y_1c=\rho_q(c)Y_1,\qquad Y_0\text{ は中央}
 \tag{E7}
\]
という二重次数付き Ore 拡大を作る。
次数は $\deg c=(\deg_B c,0)$、$\deg Y_i=(0,1)$ とする。
その対角部分 $\bigoplus_{n\geq0}T_{(n,n)}$ は次数ごとに平坦であり、
各次数の階数は $(n+1)^3$ である。
アフィン座標 $y=Y_1/Y_0$ について
\[
 \{y,c\}=\kappa y(x\partial_x+z\partial_z)(c)
\]
なので、半古典極限は式 (E6) と一致する。
従って各対数型成分の一般元についても、代数的な
Segre 切断環の族が得られる。

\subsection{pencil 成分のパラメーターと非可換 quadric}
\label{subsec:quant-pencil-parameters}

三重積の $(1,1,1)$-pencil の基底楕円曲線 $E$ は、
三つの射影から次数 $2$ の直線束 $L_0,L_1,L_2$ を持つ。
一般の場合、この四つ組は Van den Bergh の elliptic quadruple の条件を満たす。
また pencil の二つの生成元を第四の二次元因子にまとめると、
geometric quintuple の四重テンソルが得られる。
これは、本書 の pencil 成分の射影 GIT 商
$\mathbb P(1,2,2,3)$ と、Okawa--Ueda
\cite[定理1.1]{OkawaUedaQuadrics} の非可換 quadric のコンパクトモジュライ
$\mathbb P(2,4,4,6)$ が同じ重み付き射影空間になることを説明する。

この四重テンソルは非可換 quadric 曲面を与える。
本書 の「最大 pencil 成分の Segre 射影量子化」では、
$L_E=L_0\otimes L_1\otimes L_2$ に対する filtered deformation の
相対 $\mathbb P^1$ ruling から三次元の斉次環を構成する。
二つの次数一の中央元 $s,t$ を持ち、各次数の階数は $(n+1)^3$ である。
さらに同論文の「楕円曲線を中心とする基本変換の相対共通圏」では、
楕円曲線のグラフに沿う gluing と二回の mutation により、
元の三重積の量子化を一方の収縮とする二つの SOD を得る。


\section{反標準偏極による半形式的な射影量子化}
\label{sec:anticanonical-quantization}

全ての Poisson 構造について、偏極を伴う半形式的な射影量子化も存在する。
$\PP^1\times\PP^2$ では $L=\mathcal O(2,3)=-K_X$、
$(\PP^1)^3$ では $L=\mathcal O(1,1,1)$ を用いる。

任意の Poisson 多様体の標準束 $K_X$ は標準的な平坦 Poisson 接続を持つ。
局所体積形式 $\omega$ を選ぶと、Hamiltonian ベクトル場の発散が定める
modular vector field $\operatorname{div}_{\omega}\pi$ を用いて書ける。
局所体積形式の変更に伴う変換則は $K_X$ の遷移関数の対数微分と一致するので、
$K_X^{\times}$ 上の $\mathbb G_m$-不変 Poisson 双ベクトルに貼り合わさる。
双対を取ると $-K_X$ に同じ性質があり、直線束の $m$ 乗根が存在すれば、
接続の局所係数を $m$ で割ることでその根にも同じ構造が入る。
この構成の Poisson 加群としての定式化は
Polishchuk \cite[\S1]{PolishchukQuant} を参照されたい。

上の $L$ は ample であり、$H^1(X,\mathcal O_X)=H^2(X,\mathcal O_X)=0$ だから、
Kontsevich \cite[\S3.1、構成2]{KontsevichAlgebraic} により
十分大きい $n_0$ に対する切断環
\[
 \CC\oplus\bigoplus_{n\geq n_0}H^0(X,L^n)
\]
の次数ごとに自由な noetherian $\CC[[t]]$-代数変形が得られる。
これは全ての $\pi$ に対する偏極付き半形式的量子化を与える。
この節の基底は $\operatorname{Spec}\CC[[t]]$ である。
第\ref{sec:tilting-quantization}節の有限型基底上の圏の族に、
この偏極と全次数の切断環を同時に実現する構成は追加の課題である。

\section{成分の局所幾何と残る課題}

定理\ref{thm:algebraic-quantization-all}では、Jacobi 概型と
非可換圏の半普遍底は古典的な点の共通 \'etale 近傍を持つ。
従って本書 の三重積の分類から、標識付き半普遍底の原点における
既約分枝は八つで、その次元は $13,10,10,10,9,9,9,9$ となる。
元の成分は斉次錐なので、原点での完備化によって既約性が保たれる。
この主張は半普遍変形底についてのものであり、
自己同値をさらに同一視する大域モジュライを定めるにはその作用を指定する。

$\PP^1\times\PP^2$ についても本書 の分類から、
同様に九つの分枝と Riccati 分枝の生成点の長さ2が対応する。
その分類の網羅性と生成点の長さの証明は本書に収めている。
本稿の量子化の存在と完備局所環の同型は、この数値的帰結と独立に証明した。

三重積の八成分の一般元について、元の Segre 偏極を持つ
平坦斉次量子化が得られた。最大 pencil の証明は本書
本書 の「最大 pencil 成分の Segre 射影量子化」にある。
別の節では楕円中心の相対共通圏と二つの SOD も構成した。
この特定の斉次環と圏の導来同値、$\PP^1\times\PP^2$ の残る成分の
具体的な関係式、および全成分での偏極付き双有理比較を引き続き調べる。

\texttt{check\_tilting.py} は有限次元代数の corner の次元、
Hilbert 行列の逆行列、鎖の長さと Hochschild cochain の次元を整数演算で確認する。
\texttt{check\_quantization.py} は明示的な族の半古典極限と
四重テンソルの二周期性を確認する。
formality、Artin 近似、および AS 正則性の分類定理は、上記の文献に依存する。


\part{Poisson 概型の既約成分}
\section{\texorpdfstring{$X_{12}$}{X12} の Poisson 概型の分類}
\label{p12:classification}
\subsection{設定、主結果と先行研究}
$X=\PP^1\times\PP^2$ とする。
Polishchuk~\cite[定理6.1, 式(16)]{Polishchuk1997} は、階数2束の行列式の Poisson 加群構造を固定すると、
その射影束上で基底の双ベクトルと両立する Poisson 構造を、束の Poisson 加群構造で記述できることを示した。
本稿では $Y=\PP^2$、束 $\mathcal O_Y^{\oplus2}$ にこの対応を適用する。
平面の退化三次曲線の補集合で得られる平坦 $\mathrm{SL}_2$ 接続の局所モノドロミーと、
混合項の次数制約を組み合わせることが、各ファイバーの分類の要点である。
Loray--Pereira--Touzet~\cite[定理9.1]{p12-LPT} は Picard 数1の Fano 三次元多様体上の Poisson 構造を分類した。
Belmans--Fatighenti--Tanturri~\cite[\S2.2]{p12-BFT} は高い Picard 数での分類問題と多重ベクトル場の次元を論じている。
本節で扱う $X$ の Picard 数は2である。
Pym~\cite{p12-Pym} は1形式や葉層を用いる手法を概説している。
Okumura~\cite[定理1.2]{p12-Okumura} の積に関する結果は偶数次元で各因子の次元が3以上の場合の
SNC log symplectic 構造を扱う。以下の個々の成分について既存文献との優先関係は未確定である。

接束の分解と $\bigwedge^2T_{\PP^2}\cong\mathcal O_{\PP^2}(3)$ より
\[
 H^0(X,\textstyle\bigwedge^2T_X)
 \cong (\mathfrak{sl}_2\otimes\mathfrak{sl}_3)\oplus H^0(\PP^2,\mathcal O(3)),
 \qquad \dim=24+10=34.
\]
また $-K_X=\mathcal O(2,3)$ より $h^0(X,\bigwedge^3T_X)=30$ である。

\begin{definition}\label{p12:def:scheme}
$V=H^0(X,\bigwedge^2T_X)$ と $W=H^0(X,\bigwedge^3T_X)$ を置く。
本節の Poisson 概型は、二次写像
\[
 V\longrightarrow W,\qquad \pi\longmapsto\tfrac12[\pi,\pi]
\]
の零ファイバー $\operatorname{Pois}(X)$ である。
したがって、その座標環は $\operatorname{Sym}(V^\vee)$ を
$\tfrac12[\pi,\pi]$ の全座標二次式で割った環である。
既約成分 $Z$ の生成点 $\eta$ における重複度は
$\operatorname{length}\mathcal O_{\operatorname{Pois}(X),\eta}$ と定める。
以下の分類では固定した $X$ 上の双ベクトルそのものをパラメータとする。
\end{definition}


本節で分類する成分の一覧を掲げる。次元は
$\operatorname{Aut}(X)$ で割る前のアフィン Poisson 概型の次元を表す。
\[
\begin{array}{c|c|l}
\text{記号}&\text{次元}&\text{構成または特徴}\\\hline
V_{\mathrm{vert}}&10&\PP^2\text{ 上の Poisson 構造の引き戻し}\\
C_{\mathrm{tor}}&11&\text{最大トーラスの双ベクトル}\\
C_{\mathrm{fib}}&11&H\omega,\ H\in H^0(\mathcal O(2,1)),\ \omega\in H^0(\Omega^1_{\PP^2}(2))\\
C_{\mathrm{line\text{-}conic}}&11&aLQ\,du+h(u)(L\,dQ-2Q\,dL)\\
C_{(1,1)^2(0,1)}&11&\text{$(1,1)$ 因子を葉とする pencil}\\
C_{(1,1)(1,0)(0,2)}&12&\text{$(1,1)$ 因子を葉とする logarithmic 型}\\
C_{\mathrm{Ric}}&12&\text{共点三直線と直線束上の Riccati 方程式}\\
C_4&12&\text{四因子の対数形式、二つの自由な留数}\\
C_{(1,2)(1,0)(0,1)}&14&\text{$(1,2)$ 因子を葉とする logarithmic 型}
\end{array}
\]
\begin{theorem}\label{p12:th:complete}
$\mathbb C$ 上で、上表の九つが $\operatorname{Pois}(X)$ の全既約成分である。
$C_{\mathrm{Ric}}$ の生成点での長さは2であり、他の八つの生成点での長さは1である。
\end{theorem}
網羅性は射影三次式による層別化から従う。証明を第\ref{p12:sec:conclusion}節でまとめる。
成分の一般点の重複度と、二つの指定因子への制限ファイバーの概型構造は別の問題であり、
後者は定理\ref{p12:th:extension5}, \ref{p12:th:extensionfibers} と命題\ref{p12:prop:zero-restriction} で扱う。

\subsection{アフィン表示と Jacobi 方程式}
標準アフィン開集合の座標を $u$（$\PP^1$）および $v,w$（$\PP^2$）とし、
\[
 \pi=A\,\partial_u\wedge\partial_v+B\,\partial_u\wedge\partial_w+C\,\partial_v\wedge\partial_w.
\]
ここで $C=C(v,w)$ は次数 $\le3$ の多項式であり、$(A,B)$ は $\mathfrak{sl}_2\otimes\mathfrak{sl}_3$ の元を表す。$\mathfrak{sl}_2$ 側の基底は $1,u,u^2$、$\mathfrak{sl}_3$ 側の8つのベクトル場の基底は
\[
 \partial_v,\quad\partial_w,\quad v\partial_v,\quad w\partial_v,\quad v\partial_w,\quad w\partial_w,
 \quad v^2\partial_v+vw\partial_w,\quad vw\partial_v+w^2\partial_w.
\]
Schouten 条件 $[\pi,\pi]=0$ は、Jacobi 恒等式により次の多項式の恒等的消滅と同値である。
\begin{equation}\label{p12:eq:jacobi}
 J(A,B,C):=A(C_v+B_u)+B(C_w-A_u)-C(B_w+A_v)=0.
\end{equation}
$J$ は $u$ に関して次数 $\le2$、$v,w$ に関して全次数 $\le3$ であり、係数比較によって30個の斉次二次式を与える。

\subsection{垂直成分}
\begin{theorem}\label{p12:th:vertical}
 $V_{\mathrm{vert}}=\{(A,B,C):A=B=0\}\cong\mathbb{A}^{10}$ は $\operatorname{Pois}(X)$ の既約成分であり、滑らかな平面三次曲線を表す一般点で $\operatorname{Pois}(X)$ は滑らかである。
\end{theorem}
\begin{proof}
 $A=B=0$ なら \eqref{p12:eq:jacobi} は恒等的に零である。$C$ が滑らかな平面三次曲線を定める点での一次変形を $\dot\pi=\dot C+\eta$ ($\eta\in\mathfrak{sl}_2\otimes\mathfrak{sl}_3$) とする。一次 Jacobi 条件は
 $[C,\eta]=0$ である。$\eta=\sum_i a_i\wedge b_i$ と表せば、これは
 \[
 \sum_i a_i\otimes\mathcal L_{b_i} C=0
 \]
 となる（符号規約は零条件に影響しない）。滑らかな平面三次曲線の射影自己同型群は有限であるから、$b\mapsto\mathcal L_bC$ は $\mathfrak{sl}_3\to H^0(\PP^2,\mathcal O(3))$ の単射である。したがって $\eta=0$。接空間は $\dot C$ の10次元空間に一致する。$V_{\mathrm{vert}}$ 自身の次元も10なので、この点は $\operatorname{Pois}(X)$ の滑らかな点であり、$V_{\mathrm{vert}}$ は既約成分である。
\end{proof}

\subsection{トーラス由来の成分}
$E_u=u\partial_u$, $E_v=v\partial_v$, $E_w=w\partial_w$ は $\mathrm{Aut}(X)$ の最大トーラスの大域ベクトル場である。
\[
 \pi_{\lambda,\mu,\nu}=\lambda E_u\wedge E_v+\mu E_u\wedge E_w+\nu E_v\wedge E_w
\]
は任意の $\lambda,\mu,\nu\in\mathbb C$ について Poisson 構造をなす。$G=\mathrm{PGL}_2\times\mathrm{PGL}_3$ とし、$C_{\mathrm{tor}}$ を
\[
 G\cdot\{\pi_{\lambda,\mu,\nu}\mid(\lambda,\mu,\nu)\in\mathbb C^3\}
\]
の Zariski 閉包と定める。
\begin{theorem}\label{p12:th:toric}
 $C_{\mathrm{tor}}$ は $\operatorname{Pois}(X)$ の11次元既約成分である。$C_{\mathrm{tor}}\ne V_{\mathrm{vert}}$ である。
\end{theorem}
\begin{proof}
 $G\times\mathbb A^3$ の既約性から $C_{\mathrm{tor}}$ は既約である。点
 \[
 \pi_0=2uv\,\partial_u\wedge\partial_v+3uw\,\partial_u\wedge\partial_w+5vw\,\partial_v\wedge\partial_w
 \]
 を考える。$G$ の Lie 代数のうち最大トーラスに属する3次元部分は $\pi_0$ を固定する。残りの root direction は $\mathfrak{sl}_2$ で2つ、$\mathfrak{sl}_3$ で6つあり、$\mathcal L_V\pi_0$ はそれぞれ互いに異なる非零トーラス重みを持つ（これらが非零であることは上の座標表示への代入で直接確認できる）。したがって軌道の接空間は8次元である。$\lambda,\mu,\nu$ に関する3つの変分は互いに独立な重み0のベクトルなので、$C_{\mathrm{tor}}$ の $\pi_0$ における局所次元は少なくとも11である。

 一方、\eqref{p12:eq:jacobi} を $\pi_0$ で線形化する。$\dot\pi$ の30係数への線形写像の rank が23以上であることを示す。トーラス重みによって各項を分解する。$\dot C=v^i w^j$（$i+j\le3$）とすると、線形化された Jacobi 多項式は
 \[
 \dot J=u\bigl(2(i-1)+3(j-1)\bigr)v^i w^j.
 \]
 唯一 $(i,j)=(1,1)$ のとき係数が消え、残り9つは独立な重みの非零像となる。

 次に $\dot C=0$ とし、$(\dot A,\dot B)=u^k(P,Q)$（$k=0,2$）で $P\partial_v+Q\partial_w$ を $\mathfrak{sl}_3$ の前節の基底に走らせる。$k=0$ では像は順に
 \[
 8w,\ 3v,\ 3vw,\ 8w^2,\ 3v^2,\ -2vw,\ -4v^2w,\ -4vw^2,
 \]
 $k=2$ では順に
 \[
 u^2\bigl(2w,\ 7v,\ -3vw,\ 2w^2,\ 7v^2,\ 2vw,\ -6v^2w,\ -6vw^2\bigr).
 \]
 各 $k$ について7つの異なる単項式が得られ、合計14次元の像を作る。先の9つの像は $u$ 次数が1であるため、これらと独立である。ゆえに rank $\ge9+14=23$、すなわち
 \[
 \dim T_{\pi_0}\operatorname{Pois}(X)\le34-23=11.
 \]
 これと11次元の族の存在より $\dim T_{\pi_0}\operatorname{Pois}(X)=11$、かつ $\pi_0$ は滑らかな点である。よってその点を通る11次元既約閉集合 $C_{\mathrm{tor}}$ は既約成分である。$\pi_0$ は mixed term を持つので $V_{\mathrm{vert}}$ とは異なる。
\end{proof}


\subsection{既約三次曲線上の剛性と零ファイバー}
$Y=\PP^2$ とし、射影に伴う線形写像
$q:\operatorname{Pois}(X)\to H^0(Y,-K_Y)$ で平面成分を取り出す。
$D=(C=0)$ が被約既約な三次曲線である $C$ の全体を $V^{\mathrm{irr}}$ と書く。
これは原点と直線・二次式の積からなる閉集合の補集合である。

\begin{theorem}\label{p12:th:irreducible-cubic}
$q^{-1}(V^{\mathrm{irr}})\to V^{\mathrm{irr}}$ は概型の同型であり、
その逆は混合項を零とする切断である。特に $D$ が滑らかな場合のほか、
既約な結節三次曲線および尖点三次曲線の場合にも非零混合項は存在しない。
\end{theorem}
\begin{proof}
$C$ に対応する平面の双ベクトルを $\sigma$、その非退化開集合を
$U=Y\setminus D$ とする。$\pi_1(U)=\mathbb Z/3$ は三つの型すべてについて成り立つ
\cite[\S2.3.2, (3e)]{p12-Cogolludo}。滑らかな場合は
$\{C=1\}\to U$ が $\mu_3$ 被覆で、斉次孤立特異点の Milnor
ファイバーが単連結であることからも従う~\cite[\S\S5--6]{p12-Milnor}。

$\pi$ の混合項を $\sum_{i=1}^3 e_i\wedge V_i$ と書く。
ここで $e_i$ は $H^0(\PP^1,T_{\PP^1})$ の基底、$V_i$ は $Y$ の射影ベクトル場である。
$\PP^1\times U$ 上の特性分布は $TU$ に同型に射影するので、
自明な $\PP^1$ 束の接続を定める。水平持ち上げの垂直成分は、符号の取り方を除いて
$\sum_i (\iota_{V_i}\sigma^{-1})\,e_i$ である。Jacobi 恒等式により水平分布は可積分である。
射影ベクトル場の Lie 環 $\mathfrak{pgl}_2$ を $\mathfrak{sl}_2$ と同一視すると、
これは自明な階数2束上の平坦な接続 $d+\Omega$ を与える。
すなわち $\Omega$ を水平持ち上げから符号を合わせて取ると、曲率の射影像が零であり、
曲率はトレース零だから曲率自身も零である。
この対応は Poisson 加群の射影化~\cite{Polishchuk1997} の局所表示でもある。

$D$ の滑らかな点の近傍で $D=(x=0)$ となる座標 $(x,y)$ を取る。
$\sigma^{-1}$ は1位の極しか持たないので
$\Omega=(A\,dx+B\,dy)/x$ と書ける。$A,B$ は正則なトレース零行列である。
平坦性の $x^{-2}$ の係数は、$A_0=A(0,y),\,B_0=B(0,y)$ として
\[
 [A_0,B_0]=B_0
\]
となる。ある点で $B_0\ne0$ なら $\operatorname{ad} A_0$ が固有値1を持ち、
$A_0$ の二つの固有値は $1/2,-1/2$ である。この点で $y$ を固定し、
$x=r e^{i\theta}$ に沿った平行移動を考える。その行列の微分方程式の係数は
$-i A(r e^{i\theta},y)$ であり、$r\to0$ で一様に $-i A_0$ に収束する。
円周のホロノミー行列も $\exp(-2\pi i A_0)$ に収束する。
各ホロノミー行列は3乗が単位行列であるが、極限の固有値は $-1$ となり矛盾する。
したがって $B_0=0$ で、$\Omega$ は $D$ の滑らかな点で対数極を持つ。
$\sigma=x\,h\partial_x\wedge\partial_y$ ($h$ は単元) と書けば、
$\iota_{V_i}\sigma^{-1}$ の $dy$ 成分が正則であることは
$V_i(x)\in(x)$ と同値である。よって $V_i$ は $D$ を保存する。

$\mathcal L_V\sigma=0$ を満たす射影ベクトル場は零である
\cite[Prop.~15.1]{Polishchuk1997}。実際、$V$ は $D$ に接し、
$\alpha=\iota_V\sigma^{-1}$ は $U$ 上で閉じており、$D$ の滑らかな点で
対数極を持つ。$\pi_1(U)$ が有限なので任意の閉路に沿う周期は零である。
とくに局所メリディアンの周期が零だから対数留数が消え、$\alpha$ は正則に延びる。
残る有限個の特異点にも Hartogs の定理で延び、
$H^0(Y,\Omega_Y^1)=0$ より $V=0$ となる。
一方、$D$ を保存するベクトル場 $V$ は $\mathcal L_V\sigma=c_V\sigma$
を満たす。実際、左辺は $D$ 上消える $-K_Y$ の切断である。
いま示したことから $V\mapsto c_V\in\mathbb C$ は単射なので、
$D$ を保存する射影ベクトル場の空間は高々1次元である。
従って混合項が非零なら $e\wedge V$ と書けるが、このとき
式~\eqref{p12:eq:jacobi} は $\mathcal L_V\sigma=0$ を与え、矛盾する。

これで閉点の主張を得た。各閉点 $(0,C)$ での一次 Jacobi 条件は
各 $i$ について $\mathcal L_{\dot V_i}\sigma=0$ であり、混合方向は零である。
垂直な10次元部分空間が通り、接空間も10次元なので局所環は正則である。
従ってこの開部分概型は被約で、上の切断による同定は概型として成立する。
\end{proof}

\begin{proposition}\label{p12:prop:zero-cubic}
$q^{-1}(0)_{\mathrm{red}}$ の既約成分は次の二つである。
第一はテンソル階数が高々1の10次元錐であり $C_{\mathrm{tor}}$ に含まれる。
第二は $p\in\PP^2$ を動かして、$p$ を通る直線束に接する三つの射影ベクトル場を
任意に選ぶ11次元族であり、$C_{\mathrm{fib}}$ に一致する。
\end{proposition}
\begin{proof}
$\partial_u\wedge (V_0+uV_1+u^2 V_2)$ の Jacobi 条件は
$(V_0+uV_1+u^2 V_2)\wedge(V_1+2uV_2)=0$ であり、
各 $V_i\wedge V_j=0$ と同値である。階数が2以上なら、
どの非零の $V\in\langle V_i\rangle$ に対しても、これと独立な $W$ が取れ、
$W=fV$ を満たす有理関数 $f$ がある。$V$ に因子的な零点がなければ
$f$ は余次元1の全点で正則、従って定数であり矛盾する。
行列 $M$ が定める射影ベクトル場の零点は固有空間の射影化の和である。
直線上で消える非零射影ベクトル場には、スカラー行列を引いて
階数1の表示 $M=p\otimes\ell$ がある。その一般点での方向は
その点と $[p]$ を結ぶ直線の方向であるから、すべての $V_i$ が同じ中心を持つ。
逆に、中心 $[p]$ に対し
$W_p=\{p\otimes\ell\bmod\mathbb C I:\ell\in(\mathbb C^3)^*\}$
内の任意の三つ組は条件を満たす。
この族は $\PP^2$ 上の階数9ベクトル部分束の全空間からの射影像なので
既約閉であり、階数2以上では中心は一意なので次元は11である。
中心を固定した直線束の1形式に $H\in H^0(\mathcal O(2,1))$ を掛けたものが
ちょうどこの族である。

第一の族は Segre 錐なので次元 $3+8-1=10$ である。その一般点では
両因子のベクトル場が半単純で、最大トーラスに属するから閉包全体が
$C_{\mathrm{tor}}$ に入る。一般の $\mathfrak{sl}_3$ の元は相異なる三つの
固有値を持ち、因子的な零点を持たないので第二の族には入らない。
第二の族には階数3のテンソルが存在する。この二つは互いに包含されず、
零ファイバーの既約成分を尽くす。
\end{proof}


\subsection{二つの del Pezzo 因子の幾何}
$H_1,H_2$ を $X$ の二つの因子から引いた超平面類とする。
$-K_X=2H_1+3H_2$ であり、$(1,a)$ 型の滑らかな因子 $S$ に対し
$-K_S=(H_1+(3-a)H_2)|_S$ が成り立つ。その次数は
\[
((-K_X-S)^2\cdot S)_X=
\begin{cases}8&(a=1),\\5&(a=2)\end{cases}
\]
と直接交点積を取ればよい。
射影 $S\to\PP^2$ は、二つの $a$ 次形式の基点を爆発したものだから、
一般の $a=1$ では1点を爆発して $S\cong\mathbb F_1$（8次）、
$a=2$ では一般位置の4点を爆発して5次 del Pezzo となる。
それぞれ
\[
H^0(S,-K_S)\cong
\begin{cases}
H^0(S,\mathcal O_S(1,2)),\quad h^0=9 &(a=1),\\
H^0(S,\mathcal O_S(1,1)),\quad h^0=6 &(a=2).
\end{cases}
\]
曲面では正則双ベクトルはすべて Poisson である。

\subsubsection{Poisson 因子の延長問題}
滑らかな因子 $S\subset X$ を Poisson 部分多様体とするための条件は、
イデアル層 $I_S$ が Poisson イデアルとなること、すなわち
\[
\{I_S,\mathcal O_X\}\subset I_S
\]
である\cite[\S4.1]{p12-Pym}。これは $\pi|_S\in\bigwedge^2T_S$ と同値である。
この条件を満たす双ベクトルの層を $\bigwedge^2T_X(-\log S)$ と表すと、
\begin{equation}\label{p12:eq:log-exact}
0\longrightarrow \bigwedge^2T_X(-S)
\longrightarrow \bigwedge^2T_X(-\log S)
\longrightarrow \bigwedge^2T_S\longrightarrow0
\end{equation}
が成り立つ（局所座標 $S=\{z=0\}$ で直接検証できる）。

\begin{theorem}\label{p12:th:extension5}
$S$ を滑らかな5次 del Pezzo 型の $(1,2)$ 因子とする。
任意の $\beta\in H^0(S,-K_S)$ は、$S$ に接する大域双ベクトル
$\widetilde\beta\in H^0(X,\bigwedge^2T_X)$ に\emph{一意に}延長する。
ただし $\widetilde\beta$ 自体が Poisson になる必要十分条件は
$\beta$ に対応する $H\in H^0(X,\mathcal O_X(1,1))\cong\mathbb C^2\otimes\mathbb C^3$
のテンソル階数が1以下であることである。
従って、延長可能な $\beta$ の集合は4次元の Segre 錐である。
\end{theorem}
\begin{proof}
接束の分解を用いると、$S\sim H_1+2H_2$ に対し
\[
\bigwedge^2T_X(-S)
=\mathcal O(-1,1)\oplus
\bigl(\mathcal O_{\PP^1}(1)\boxtimes T_{\PP^2}(-2)\bigr).
\]
$H^0(\PP^2,T_{\PP^2}(-2))=H^1(\PP^2,T_{\PP^2}(-2))=0$ は
捻った Euler 列から従う。$\PP^1$ 上の $\mathcal O(-1)$ も $H^0,H^1$ が零なので、
$H^0,H^1$ ともに $\bigwedge^2T_X(-S)$ では零になる。
よって \eqref{p12:eq:log-exact} の大域切断により
\[
H^0(X,\bigwedge^2T_X(-\log S))\ \xrightarrow{\sim}\ H^0(S,-K_S).
\]
一方、完全列 $0\to\mathcal O_X(0,-1)\to\mathcal O_X(1,1)\to\mathcal O_S(1,1)\to0$
から、$H^0(X,\mathcal O_X(1,1))\xrightarrow{\sim}H^0(S,-K_S)$ も従う。

$F\in H^0(X,\mathcal O(1,2))$ を $S$ の定義式、
$H\in H^0(X,\mathcal O(1,1))$ を $\beta$ に対応する切断とする。
積多重次数を持つアフィン錐上で、二つの射影因子の方向の外微分を $d_1,d_2$ と書くと、
\begin{equation}\label{p12:eq:formdp5}
\alpha(F,H)=H(d_1F+d_2F)-F(d_1H+2d_2H)
\end{equation}
は両方の Euler ベクトル場との縮約が零なので、
$\Omega^1_X(2,3)\cong\bigwedge^2T_X$ の大域切断に降下する。
$S$ への制限では $\alpha|_S=H\,dF$ であり、必要な $\widetilde\beta$ に対応する。

$S$ が5次 del Pezzo であるとき、四つの一般位置の基点は $\mathrm{PGL}_3$ の作用で標準形へ移せる。
従って $X$ の自己同型により、アフィン座標 $(u,v,w)$ で
\[
F=uvw+(v-1)(w-1)
\]
としてよい。$H$ を
\[
H=(a_0+a_1v+a_2w)+u(a_3+a_4v+a_5w)
\]
とし、$\alpha=C\,du-B\,dv+A\,dw$ と置く。
\eqref{p12:eq:jacobi} に代入して整理すると、
\begin{align}\label{p12:eq:minoridentity}
J(A,B,C)=F\bigl[&M_{01}((u+1)v-1)
-M_{02}((u+1)w-1)\\
&-M_{12}(v+w-2)\bigr],\nonumber
\end{align}
ただし
\[
M_{01}=a_0a_4-a_1a_3,\quad
M_{02}=a_0a_5-a_2a_3,\quad
M_{12}=a_1a_5-a_2a_4.
\]
従って $J=0$ は三つの $2\times2$ 小行列式 $M_{01}=M_{02}=M_{12}=0$
と同値である。このアフィン代数集合は $2\times3$ 行列の階数1以下の錐、
すなわち $\operatorname{Cone}(\PP^1\times\PP^2)$（次元4）である。
式 \eqref{p12:eq:minoridentity} は付属コードでも記号的に検証している。
\end{proof}

\begin{remark}
上の $H=U\cdot L$ が分解可能な場合、\eqref{p12:eq:formdp5} は
\[
\alpha=UL\,dF-FL\,dU-2FU\,dL
\]
となる。これは $F/(UL^2)$ を有理第一積分とする Poisson 構造であり、
固定した $S$ に限らず、その pencil に含まれる一般の5次 del Pezzo 因子を Poisson 葉の閉包とする。
\end{remark}

\subsubsection{一般の del Pezzo Poisson 構造の延長障害}
次は5次・8次の両方に有効な制約である。零曲線が特異な既約曲線であっても成立する。
\begin{proposition}\label{p12:prop:flatobstruction}
$S\subset X$ を滑らかな $(1,1)$ または $(1,2)$ 型 del Pezzo 因子とする。
$\beta\in H^0(S,-K_S)$ の零因子 $E=\{\beta=0\}$ が既約なら、
$\beta$ を誘導して $S$ を Poisson 部分多様体とする ambient Poisson 構造は存在しない。
\end{proposition}
\begin{proof}
もし延長 $\pi$ が存在すれば、Poisson イデアル $I_S$ の商
$I_S/I_S^2\cong N_{S/X}^{\vee}$ は $S$ 上の Poisson module になる。
$S\setminus E$ は正則シンプレクティックなので、$N_{S/X}^{\vee}|_{S\setminus E}$
は平坦な代数的接続を持ち、従ってその de Rham 第一 Chern 類は零である。
$E$ が既約なら、局在長完全列と Borel--Moore 双対性により
\[
\ker\bigl(H^2_{\mathrm{dR}}(S)\to H^2_{\mathrm{dR}}(S\setminus E)\bigr)
=\mathbb C[E].
\]
従って $c_1(N_{S/X})$ は $[E]=c_1(-K_S)$ に比例しなければならない。
しかし $N_{S/X}=\mathcal O_S(1,a)$ と $-K_S=\mathcal O_S(1,3-a)$ は
$a=1,2$ のいずれについても独立な Picard 類である。これは矛盾する。
Poisson module と平坦性の標準的背景は\cite{p12-GualtieriPym,p12-Pym} を参照。
\end{proof}


\subsubsection{固定した8次 del Pezzo 因子での完全な延長判定}
$S$ を滑らかな $(1,1)$ 因子とする。自己同型で
\[
 F=uv+w=0
\]
とおける。$p=(v,w)=(0,0)$ とすると
$S=\operatorname{Bl}_{p}\PP^2$ であり、$S$ 上の反標準切断は
$p$ を通る平面三次曲線の方程式 $P(v,w)$ に対応する。

\begin{theorem}\label{p12:th:fixed8}
$S\subset X$ を滑らかな $(1,1)$ 因子とし、
$\beta\in H^0(S,-K_S)$ を非零切断とする。
$\beta$ が $S$ を Poisson 部分多様体とする
$X$ 上の正則 Poisson 双ベクトルへ延長するための必要十分条件は、
対応する平面三次式 $P$ が一次式と二次式の積に分解することである（重複因子を許す）。
従って、延長可能な切断の集合は、$H^0(S,-K_S)\simeq\mathbb{A}^9$ 内の
相異なる二つの7次元既約閉錐
\[
 \mathcal R_{\ell_p q}=\{\ell_p q\mid p\in\ell_p,\deg q=2\},
 \qquad
 \mathcal R_{\ell q_p}=\{\ell q_p\mid p\in q_p,\deg \ell=1\}
\]
の和集合である（$0$ も含む）。
\end{theorem}
\begin{proof}
$S\to\PP^2$ が同型となる稠密開集合で見ると、平面への射影 $q(\pi)$ は、
指定した $\beta$ に対応する三次式 $P$ に一致する（標準化での非零定数倍を許す）。
$P$ が既約なら定理\ref{p12:th:irreducible-cubic} より $\pi$ は平面方向のみである。
$P\ne0$ となる点でその像は平面全体の接方向であり、
非定数有理関数 $u=-w/v$ のグラフ $S$ に接することはできない。
これで必要性が従う。

後で制限ファイバー自体を計算するため、接する一般の双ベクトルを記述する。
$\bigwedge^2T_X\simeq\Omega^1_X(2,3)$ の対応で、
$S$ に接する一般の双ベクトルは以下の1形式で表せる。
\begin{equation}\label{p12:eq:normal8}
 \alpha=2H\,dF-F(2d_1H+d_2H)+F\gamma,
\end{equation}
ただし $d_i$ はそれぞれの射影因子の斉次微分であり、
\begin{align*}
H={}&a_0+a_1u+v(b_0+b_1u)+b_2uw
+v^2(c_0+c_1u)+c_2uvw+c_3uw^2,\\
\gamma={}&(x_0+x_1u)\,dv+(y_0+y_1u)\,dw
+(z_0+z_1u)(v\,dw-w\,dv).
\end{align*}
この表示が全切断を尽くすことを示す。接する双ベクトルの完全列
\eqref{p12:eq:log-exact} の核は
\[
H^0(X,\bigwedge^2T_X(-S))
 \cong H^0(\PP^1,\mathcal O(1))\otimes
 H^0(\PP^2,T_{\PP^2}(-1)),\qquad\dim=6
\]
であり、$H^1$ は消える。
従って大域切断空間の次元は $6+9=15$ である。
\eqref{p12:eq:normal8} は15個の独立な係数を含み、核と反標準切断の
両方を尽くす（付属コードで直接の線形代数検証も行う）。

$\alpha=C\,du-B\,dv+A\,dw$ として
\eqref{p12:eq:jacobi} に代入すると、$J(A,B,C)$ は $F$ で割れ、
\[
 J(A,B,C)=F\sum_{i,j,k}p_{ijk}u^iv^jw^k,
\]
ただし指数は $0\le i\le1$, $j,k\ge0$, $j+k\le2$ を満たす12組にわたる。
$\alpha|_S=2H\,dF$ であり、平面三次曲線のアフィン方程式は
\begin{align}\label{p12:eq:cubic8}
P(v,w)={}&a_0v-a_1w+b_0v^2-b_1vw-b_2w^2\nonumber\\
 &+c_0v^3-c_1v^2w-c_2vw^2+c_3w^3.
\end{align}
逆に $C$ が可約なら $C=\ell q$ と因数分解し、$p\in C$ から
$p\in\ell$ または $p\in q$ である。
前者では $\ell$ に対応する $U\in H^0(\mathcal O_X(1,0))$ と
$q$ に対応する $Q\in H^0(\mathcal O_X(0,2))$ を取り、
\[
 \alpha=2UQ\,dF-2FQ\,dU-FU\,dQ
\]
を用いる。後者では $q$ の狭義変換が
$\mathcal O_S(1,1)$ の切断に対応するので
$G\in H^0(\mathcal O_X(1,1))$ と $L\in H^0(\mathcal O_X(0,1))$
を取り、
\[
 \alpha=L(G\,dF-F\,dG)
\]
を用いる。いずれも可積分で $S$ に所望の Poisson 構造を誘導する
（スカラー倍は調整できる）。
最後に、第一の可約三次曲線の射影パラメータ空間は
$\PP^1\times\PP^5$、第二は $\PP^2\times\PP^4$ なので、
両錐はそれぞれ7次元であり、一般点の因数分解の型が違うから相異なる。
\end{proof}

\begin{proposition}\label{p12:prop:zero-restriction}
固定した滑らかな $(1,1)$ 因子 $S$ に Poisson 因子として接し、
$S$ 上では零 Poisson 構造を誘導する ambient Poisson 構造は、
\eqref{p12:eq:normal8} で $H=0$ とした $\alpha=F\gamma$ のうち
\[
\operatorname{rank}\begin{pmatrix}x_0&y_0&z_0\\x_1&y_1&z_1\end{pmatrix}\le1
\]
を満たすものに正確に一致する。これは6次元線形空間中の4次元 Segre 錐である。
\end{proposition}
\begin{proof}
$H=0$ なら $\alpha=F\gamma$ であり、
$\alpha\wedge d\alpha=F^2\gamma\wedge d\gamma$。
$\gamma=\gamma_0+u\gamma_1$、
$\gamma_i\in H^0(\PP^2,\Omega^1(2))$ と書くと、
$\gamma\wedge d\gamma$ が零であることは
$\gamma_0\wedge\gamma_1=0$ に等しい。
自然な写像
$\bigwedge^2H^0(\PP^2,\Omega^1(2))\longrightarrow
H^0(\PP^2,\Omega^2(4))$ は3次元空間の同型であるため、
$\gamma_0,\gamma_1$ は比例する。これは表示した階数条件と同値である。
\end{proof}

\subsubsection{8次 del Pezzo 因子での延長の個数とスキーム構造}
\label{p12:sec:markedfibers}
滑らかな $(1,1)$ 因子 $S=\operatorname{Bl}_p\PP^2$ を固定する。
$\mathcal E_\beta$ は、$S$ に接し、指定した切断
$\beta\in H^0(S,-K_S)$ に制限される大域双ベクトルのアフィン空間の中で、
Jacobi 方程式の全係数を課して得られる閉部分スキームとする。
スカラー倍も自己同型も同一視しない。
以下、$D_2=\operatorname{Spec}\mathbb C[\epsilon]/(\epsilon^2)$、
$N=\operatorname{Spec}\mathbb C[s,t]/(st)$、$\mathrm{pt}=\operatorname{Spec}\mathbb C$ とする。

\begin{theorem}\label{p12:th:extensionfibers}
非零 $\beta$ に対応する三次曲線 $C$ が可約のとき、
$\mathcal E_\beta$ のスキーム同型型は以下の通りである。
$L$ は直線、$Q$ は滑らかな二次曲線、$M$ は $L$ と異なる直線を表す。
\[
\begin{array}{lll}
C & p\text{ の位置} & \mathcal E_\beta\\ \hline
L+Q & \text{一方の成分だけに属する} & \mathrm{pt}\\
L+Q & L\cap Q & \mathrm{pt}\sqcup\mathrm{pt}\\
\text{三角形} & \text{辺の滑らかな点} & \mathbb A^1\\
\text{三角形} & \text{頂点} & N\\
\text{共点の相異なる三直線} & \text{共点以外} & \mathbb A^1\\
\text{共点の相異なる三直線} & \text{共点} & \mathbb A^2\\
2L+M & L\setminus M & \mathbb A^1\times D_2\\
2L+M & M\setminus L & \mathbb A^2\\
2L+M & L\cap M & \mathbb A^1\times N\\
3L & L & Z_3
\end{array}
\]
ここで $Z_3=\operatorname{Spec}\mathbb C[A,B,C',D]/(AD-2BC',B^2,BD)$ である。
最初の二行は $L$ が割線の場合にも接線の場合にも成立する。
特に定理\ref{p12:th:fixed8} の二つの7次元錐のいずれでも一般の切断の
延長は一意で、そのファイバーは被約である。最後の行の台は二つの平面の和であり、
$D=0$ の平面の一般点で重複度2を持つ。
\end{theorem}

\begin{proof}
$S\to\PP^1$ は $p$ を中心とする直線束への射影であるため、
標識点を保つ平面の射影変換はその直線束上の変換と合わせて $(X,S)$ の
自己同型に持ち上がる。さらに $\beta$ を非零定数倍すると、
すべての延長を同じ定数倍するスキーム同型が得られる。
従って曲線と点の対 $(C,p)$ の射影同値類ごとに計算すればよい。

$L+Q$ の場合は $L$ が割線か接線かを区別する。
二次曲線の自己同型群 $\operatorname{PGL}_2$ の中で、
割線の安定化群は交点二点の対を保存し、その補集合である
$Q$ 上の開軌道にも $L$ 上の開軌道にも推移的に作用し、二交点を交換する。
実際、二点を $0,\infty$ とすると対角群が両方の開軌道を与える。
接線の場合も一点の安定化群が $Q$ と接線それぞれの接点の補集合で推移的である
（標準二次 Veronese 表現で平行移動の作用により確認できる）。
三角形では頂点を座標点に取れば、対角群と座標の置換が表の二つの軌道を与える。
共点の三直線では、直線束の射影変換で任意の三方向を互いに移せる。
さらに、共点を $[0:0:1]$ とすると、
$[x:y:z]\mapsto[x:y:z+ax+by]$ が各直線上の共点の補集合で推移的である。
二直線の配置も、交点と二直線を座標位置に取れば、各直線上で交点の補集合が一軌道となる。
三重直線ではその安定化群が直線上で推移的である。これで各場合を尽くす。

以下の有理係数計算で非被約構造も含めた同型型を決定する。
$p=(0,0)$ とし、アフィン式を全次数3に斉次化して $C$ を復元する。
最後の行以外は、消去後の自由変数を必要に応じて $s,t,r$ と付け直した。
「環」欄には $\mathbb Q$ 上で計算した座標環を記す。$k=\mathbb Q$ とおく。
\[
\begin{array}{lll}
\text{型と点} & P(v,w)&\text{環}\\ \hline
\text{割線上のみ} &v(1+v^2+w^2)& k\\
\text{割線型、二次曲線上のみ} &v+w+vw& k\\
\text{割線型、交点} &v(w+v^2+w^2)& k[t]/(t^2-1)\\
\text{接線上のみ} &v(v+(w+1)^2)&k\\
\text{接線型、二次曲線上のみ} &(v-1)(v+w^2+2w)&k\\
\text{接線型、接点} &v(v+w^2)&k[t]/(t^2-1)\\
\text{三角形、辺} &v(1+w)(1+v)&k[t]\\
\text{三角形、頂点} &vw(1+v+w)&k[s,t]/(st)\\
\text{共点型、辺} &v(v+w-1)(w-1)&k[t]\\
\text{共点型、共点} &vw(v+w)&k[s,t]\\
2L+M,\ p\in L\setminus M & v^2(1+w)&k[s,t]/(t^2)\\
2L+M,\ p\in M\setminus L & v(1+w)^2&k[s,t]\\
2L+M,\ p\in L\cap M &v^2w&k[r,s,t]/(st)\\
3L &v^3& k[A,B,C',D]/(AD-2BC',B^2,BD)
\end{array}
\]
具体的には \eqref{p12:eq:cubic8} で $P$ の九係数を読み、
\eqref{p12:eq:normal8} の $H$ に代入する。このアフィン空間では
$J/F$ の十二係数が Jacobi イデアルを生成する。
まず $F$ 倍を除いても、すべての係数が生成するイデアルは変化しない。
実際 $F=w+uv$ は $w$ についてモニックなので、積の係数から商の係数を
最高 $w$ 次数から順に回収できる。六変数
$(x_0,x_1,y_0,y_1,z_0,z_1)$ の辞書式順序を用いる。
例えば割線型の交点の行のイデアルは
\[
 (x_0,\ x_1+z_1,\ 2y_0+z_1-1,\ y_1,\ z_0,\ z_1^2-1),
\]
三重直線の行では
\[
 (x_0z_1-2y_0z_0,\ x_1-2y_0,\ y_0^2,\ y_0z_1,\ y_1)
\]
である。後者では $A=x_0,B=y_0,C'=z_0,D=z_1$ とおけばよい。
全行について、元の十二係数と消去表示を与える生成系のそれぞれから被約
Gröbner 基底を $\mathbb Q$ 上で計算し、一致を照合する完全なコードを
\path{p1p2/verify_extension_fibers.py} に収録した。
$\mathbb C$ への基底変換で主張のスキームを得る。
\end{proof}


\subsection{del Pezzo 因子から得られる三つの既約成分}
以下の $dF$ などの外微分は各因子のアフィン錐上で計算する。
表示した1形式は多重次数が $-K_X=(2,3)$ で、両 Euler 場を消すので
大域切断 $\Omega^1_X(2,3)$ を与える。
また、積を外に出して書く logarithmic 1形式は閉形式なので可積分である\cite[\S3.3]{p12-Pym}。

\begin{theorem}\label{p12:th:threecomponents}
次の三つの族の Zariski 閉包は、$\operatorname{Pois}(X)$ の互いに異なる既約成分を与える。
\begin{enumerate}
\item $F\in H^0(\mathcal O(1,2))$、$U\in H^0(\mathcal O(1,0))$、
$L\in H^0(\mathcal O(0,1))$ に対し
\[
\alpha_{5}=UL\,dF-FL\,dU-2FU\,dL,
\qquad \dim C_{(1,2)(1,0)(0,1)}=14.
\]
一般の $\{F=0\}$ は5次 del Pezzo で、その上の Poisson 構造は $(UL)|_S$ に対応する。
\item $F\in H^0(\mathcal O(1,1))$、$U\in H^0(\mathcal O(1,0))$、
$Q\in H^0(\mathcal O(0,2))$ に対し
\[
\alpha_{8a}=2UQ\,dF-2FQ\,dU-FU\,dQ,
\qquad \dim C_{(1,1)(1,0)(0,2)}=12.
\]
一般の $\{F=0\}$ は8次 del Pezzo で、その上の Poisson 構造は $(2UQ)|_S$ に対応する。
\item $F,G\in H^0(\mathcal O(1,1))$、$L\in H^0(\mathcal O(0,1))$ に対し
\[
\alpha_{8b}=L(G\,dF-F\,dG),
\qquad \dim C_{(1,1)^2(0,1)}=11.
\]
一般の $\{F=0\}$ は8次 del Pezzo で、その上の Poisson 構造は $(LG)|_S$ に対応する。
\end{enumerate}
\end{theorem}
\begin{proof}
各表示は、非零な積切断で割るとそれぞれ
\[
\frac{dF}{F}-\frac{dU}{U}-2\frac{dL}{L},\qquad
2\frac{dF}{F}-2\frac{dU}{U}-\frac{dQ}{Q},\qquad
\frac{dF}{F}-\frac{dG}{G}
\]
となり、閉1形式である。従って $\alpha\wedge d\alpha=0$、すなわち $[\pi,\pi]=0$。
第一の族では $\alpha|_{F=0}=UL\,dF$、第二では $\alpha|_{F=0}=2UQ\,dF$、
第三では $\alpha|_{F=0}=LG\,dF$ であり、指定された曲面は Poisson 因子である。

成分性は次の線形代数の証明書で確立する。$\operatorname{Pois}(X)$ は34変数、30二次式で定義されるので、
ある Poisson 点で Jacobi の微分の階数が $34-d$ 以上、族をパラメーター表示する写像の微分の階数が $d$ 以上ならば、
族の閉包の局所次元は $d$ 以上、Poisson スキームの接空間次元は $d$ 以下となる。
よってその点は滑らかで、閉包は $d$ 次元の既約成分である。

アフィンチャート $([u:1],[v:w:1])$ における次の点を選ぶ。
\[
\begin{array}{c|l}
\mathrm{(5)}&F=uvw+(v-1)(w-1),\quad U=2u+3,\ L=v+2w+1\\
\mathrm{(8a)}&F=uv+w,\quad U=2u+3,\ Q=2v^2+vw+3w^2+5v+7w+11\\
\mathrm{(8b)}&F=(u+1)(v+2w+1)+(2u+3)(2v-w+3),\\
&G=(3u+1)(v+3w+2)+(u+4)(2v+w+5),\quad L=2v+3w+5.
\end{array}
\]
\eqref{p12:eq:jacobi} を用いて $30\times34$ の Jacobian 行列を、
全ての定義因子の係数で微分して、対応する parameter differential 行列を作る。
両行列は整数係数である。$\mathbb F_{101}$ へ還元した後、
それぞれ次の階数以上を持つことが、付属の独立した消去コードで確かめられる。
\[
\begin{array}{c|r|r|r}
\text{点}&\mathrm{rank}(dJ)&\mathrm{rank}(d\Phi)&34-\mathrm{rank}(dJ)\\\hline
\mathrm{(5)}&\ge20&\ge14&\le14\\
\mathrm{(8a)}&\ge22&\ge12&\le12\\
\mathrm{(8b)}&\ge23&\ge11&\le11
\end{array}
\]
全ての不等号は $\mathbb F_{101}$ 上で抽出した非零小行列式を有理数上に持ち上げることで得られる。
具体的には、行・列を辞書順に走査して独立な極大 minor を選ぶと、
Jacobi 行列の minor の行列式は mod $101$ で順に $71,20,54$、
パラメーター微分行列では $76,75,70$ となる。
各行列の作り方、列順、上記数値および必要な rank は同梱の
\path{p1p2/verify_del_pezzo.py} に完全に再現される。
この有限体上の非零 minor は特性0での非零を保証する。
よって $d=14,12,11$ として上記の成分性が得られる。

相異なる次元の成分は当然異なる。また11次元の $C_{(1,1)^2(0,1)}$ は
一般点で $\{L=0\}$ に沿って双ベクトルが零になる一方、
$C_{\mathrm{tor}}$ の一般点では零集合に余次元1の成分がない。
従って、この二つの11次元既約成分も互いに異なる。
\end{proof}

\begin{remark}
固定した $(1,1)$ と $(1,2)$ 因子への延長可能集合は定理\ref{p12:th:fixed8} と\ref{p12:th:extension5} で集合論的に決定した。
\end{remark}


\subsection{直線束と直線・二次曲線からの成分}

\subsubsection{射影平面の直線束からの成分}
$\Omega^1_X(2,3)=
(\Omega^1_{\PP^1}(2)\boxtimes\mathcal O_{\PP^2}(3))
\oplus(\mathcal O_{\PP^1}(2)\boxtimes\Omega^1_{\PP^2}(3))$
と分解する。
$H\in H^0(X,\mathcal O(2,1))$ と
$\omega\in H^0(\PP^2,\Omega^1_{\PP^2}(2))$ に対して
$\alpha=H\,\mathrm{pr}_2^*\omega$ は $\Omega^1_X(2,3)$ の切断である。
$\omega$ は平面から来るので $\omega\wedge d\omega=0$、従って
$\alpha\wedge d\alpha=0$ である。
\begin{theorem}\label{p12:th:fibration}
$\alpha=H\omega$ の全体の Zariski 閉包
$C_{\mathrm{fib}}$ は $\operatorname{Pois}(X)$ の11次元既約成分である。
その一般の切断は $(2,1)$ 型の零因子を持つ。
\end{theorem}
\begin{proof}
$H^0(X,\mathcal O(2,1))$ は9次元、
$H^0(\PP^2,\Omega^1(2))$ は3次元で、
パラメータ空間は既約である。
$\omega=dw$ および
\[
H=(1+u+u^2)+(2+3u+5u^2)v+(7+11u+13u^2)w
\]
に対応する点 $\pi=(H,0,0)$ では、
Jacobi 方程式の微分の階数が少なくとも23、
族のパラメータ微分の階数が少なくとも11である。
実際、34変数を本文の基底順に並べ、30式の微分から
$23\times23$ の非零小行列式を、族の微分から $11\times11$ の
非零小行列式を抽出すると、それぞれ
$\mathbb F_{101}$ 上で行列式 $6,94$ が得られる。
再現スクリプト \path{p1p2/verify_new_components.py} により
有限回の整数演算として検証できる。
従ってこの点は Poisson 概型の滑らかな11次元点であり、
族の閉包は既約成分である。
一般の $H$ は $(2,1)$ 型の既約因子を定め、$\omega$ の零集合は
$\PP^2$ の一点だから、共通の因子として $H$ が現れる。
\end{proof}

\subsubsection{直線と滑らかな二次曲線からの成分}
$L\in H^0(\PP^2,\mathcal O(1))$ と
$Q\in H^0(\PP^2,\mathcal O(2))$ を取り、
$h\in H^0(\PP^1,\mathcal O(2))$ および $a\in\mathbb C$ に対して
\begin{equation}\label{p12:eq:lineconic}
\alpha=aLQ\,du+h(u)(L\,dQ-2Q\,dL)
\end{equation}
を定める。$du$ はアフィン座標で表した
$\Omega^1_{\PP^1}(2)\simeq\mathcal O_{\PP^1}$ の標準生成切断である。
分子 $L\,dQ-2Q\,dL$ は Euler 縮約が零の
$\Omega^1_{\PP^2}(3)$ の切断なので、
\eqref{p12:eq:lineconic} は大域切断を定める。
$LQ\ne0$ では
\[
\frac\alpha{LQ}=a\,du+h(u)\,d\log\!\left(\frac Q{L^2}\right),
\]
であり、これは Frobenius 条件を満たす。
\begin{theorem}\label{p12:th:lineconic}
\eqref{p12:eq:lineconic} の全体の Zariski 閉包
$C_{\mathrm{line\text{-}conic}}$ は $\operatorname{Pois}(X)$ の
11次元既約成分である。一般の切断には余次元1の零成分がなく、
$C(v,w)$ は滑らかな二次曲線と直線の積を表す。
\end{theorem}
\begin{proof}
$L$、$Q$、$h$、$a$ の選択空間は既約である。
アフィン座標において
\begin{align*}
 L&=2v+3w+5,\\
 Q&=2v^2+vw+3w^2+5v+7w+11,\\
 h&=1+2u+3u^2,\qquad a=7
\end{align*}
と取る。対応する $\pi$ の Jacobi 方程式の微分階数は
$\mathbb F_{101}$ 上で23、
13個のパラメータに関する微分階数は11である。
抽出した極大小行列式はそれぞれ $6,22\pmod{101}$ で非零であり、
有理数体上でも同じ階数下界が成り立つ。
定理\ref{p12:th:fibration} と同様に11次元既約成分となる。

$C=\alpha_u=aLQ$ は直線と滑らかな二次曲線の積である。
共通の因子が存在すれば $L$ または $Q$ が
$\alpha_v,\alpha_w$ をも割らねばならないが、
$L=0$ では $-2hQ\,dL$、$Q=0$ では $hL\,dQ$ となり、
一般の $L,Q$ ではどちらも恒等的には零でない。
従って一般の $\alpha$ は余次元1の零成分を持たない。
\end{proof}

\subsection{四因子の対数形式と成分の交わり}
\label{p12:sec:four}
$F,U,L,M$ の双次数を順に $(1,1),(1,0),(0,1),(0,1)$ とする。
$a,b\in\mathbb C$ に対し
\begin{equation}\label{p12:eq:four}
 \alpha=FULM\left(a\,d\log\frac{F}{UM}
                    +b\,d\log\frac L M\right)
\end{equation}
を考える。両商は次数零であり、この式は大域的な捻れ1形式に延びる。
非零の因子積で割ると閉形式になるので可積分である。
\begin{theorem}\label{p12:th:four}
\eqref{p12:eq:four} の像の閉包 $C_4$ は $\operatorname{Pois}(X)$ の12次元既約成分である。
$a\ne0$ のとき $F=0$ は Poisson 因子で、そこでの制限は
$a(ULM)|_{F=0}$ に対応する。$F$ が滑らかならこの因子は8次 del Pezzo 曲面である。
\end{theorem}
\begin{proof}
パラメータ空間は既約である。一般の四因子の配置では
$U=1,\ L=v,\ M=w,\ F=u+v+w$ をアフィン表示とできる。
実際、$U$ を無限遠に置き $F=uD+E$ とする。平面の三つの一次形式 $L,M,D$
を座標に取り、$D=1$ 上で $u$ の平行移動と $v,w$ の定数倍を用いる。
この配置で $a=1,\ b=r$ と取ると、係数は
\[
 A=vw-(1+r)vF,\qquad B=-vw-rwF,\qquad C=vw.
\]
$r=2$ における30行34列の Jacobi 微分の階数は22であり、
族の微分の階数は12である。いずれも有理数上の正確な消去である。
後者の列は、各 $F,U,L,M$ をその切断空間の単項式基底の元で
一つずつ置換したものと $(a,b)$ の二つの基底変分で尽くされる。
この計算と下記の不変因子は \path{p1p2/verify_residues.py} で再現できる。
従って像の閉包の次元と接空間次元はこの点で一致する。
$F=0$ 上では $\alpha=aULM\,dF$ であり、制限の主張も従う。
\end{proof}

12次元の既出成分との相異性は $C$ の因数分解でわかる。
\eqref{p12:eq:four} の $du$ 係数はスカラーを除き
$LM(U\,\partial_u F-F\,\partial_u U)$ で、つねに完全分解する。
定理\ref{p12:th:threecomponents} の第2族では一般に滑らかな二次曲線が因子となる。
したがってこの二つの12次元成分は異なる。

留数比に関する跳躍も決定できる。上記一般配置を固定し、
$r=b/a$ と置くと、$dJ$ の $\mathbb Q[r]$ 上の非零 Smith 不変因子は単元倍を除いて
\[
 \underbrace{1,\ldots,1}_{19},\quad f,f,\quad
 f\,r(r+1)(2r+1),\qquad f=(r-1)(r+2).
\]
$r=\infty$ では $(a,b)=(0,1)$ を代入する。これにより以下を得る。
右列には各値で実際に交わる成分を一つ記した。
\[
\begin{array}{c|c|l}
 r &\dim T_\pi\operatorname{Pois}(X)& \text{所属を与える表示}\\\hline
 1,-2&15&C_{(1,2)(1,0)(0,1)}\quad (F L)/(U M^2)\ \text{または }(F M)/(U L^2)\\
 -1/2&13&C_{(1,1)(1,0)(0,2)}\quad F^2/(U^2 L M)\\
 0,-1&13&C_{(1,1)^2(0,1)}\quad F/(U M)\ \text{または }F/(U L)\\
 \infty&14&C_{\mathrm{fib}}\quad U F (M\,dL-L\,dM)\\
 \text{上記以外}&12& C_4\text{ の滑らかな点}
\end{array}
\]
例えば $r=1$ では定理\ref{p12:th:threecomponents} 第1族で
$(1,2)$ 型の分子を $FL$ に特殊化した形になる。
最後の行では12次元成分上にあることと接空間次元から滑らかさが従う。

同じ一般配置で $a\ne0$ とし、$x=L/M$, $y=F/(UM)$, $z=w$ と置くと、
$\mathbb C(X)=\mathbb C(x,y,z)$ で、葉層を $d\log y+r\,d\log x$ が定める。
$r\notin\mathbb Q$ なら有理第一積分は存在しない。実際、第一積分は
$\partial_z$ で消えるから $\mathbb C(x,y)$ に属し、
$x\partial_x-r y\partial_y$ でも消える。互いに素な分子・分母は
それぞれこの導分の固有多項式で、固有値は等しい
（整除と全次数から従う）。異なる単項式は異なる重みを持つので、
比は定数である。


\subsubsection{成分の相異性と混合階数1}
\begin{proposition}\label{p12:prop:distinct}
ここまでに構成した $C_{\mathrm{Ric}}$ 以外の八つの既約成分は互いに異なる。
\end{proposition}
\begin{proof}
次元が異なれば成分は異なる。12次元の二成分は第\ref{p12:sec:four}節で区別した。11次元の四成分は次の幾何的不変量で区別する。
$C_{\mathrm{fib}}$ のすべての点では
$C(v,w)\equiv0$ である。他の三つの11次元成分では一般点で $C\ne0$。
$C_{\mathrm{tor}}$ の一般点では $C$ は相異なる三直線の積で、
この性質の閉包は完全分解可能な三次式の錐の中にある。
一方 $C_{\mathrm{line\text{-}conic}}$ の一般点の $C$ は直線と
既約で滑らかな二次曲線の積であるから、両者は一致しない。
$C_{(1,1)^2(0,1)}$ は $\alpha=L\alpha'$ と書けるため、
$\mathcal O(0,1)$ 型の余次元1零因子をもつ切断の閉集合に含まれる。
対して $C_{\mathrm{line\text{-}conic}}$ の一般点は
定理\ref{p12:th:lineconic} の通り余次元1の零因子を持たない。
従って後者と pencil 成分も異なる。
一般のトーラス双ベクトルの零点集合は余次元2以上である。実際、開トーラス軌道と各境界因子の開軌道では接方向が少なくとも二つ残り、一般の係数なら外積は零にならない。従ってトーラス成分と pencil 成分も異なる。
\end{proof}

\begin{proposition}\label{p12:prop:rankone}
$X$ 上の大域双ベクトル
\[
 \pi=\operatorname{pr}_2^*\beta_C+h(u)\,\partial_u\wedge V,
 \qquad \beta_C=C(v,w)\partial_v\wedge\partial_w,
\]
を考える。$h(u)\partial_u$ は $\PP^1$ 上の非零大域ベクトル場、
$V=P(v,w)\partial_v+Q(v,w)\partial_w$ は $\PP^2$ 上の大域ベクトル場とする。
このとき $\pi$ が Poisson である必要十分条件は
\[
\mathcal L_V\beta_C=0,
\quad\text{すなわち}\quad
 P\,C_v+Q\,C_w-C(P_v+Q_w)=0
\]
である。
\end{proposition}
\begin{proof}
\eqref{p12:eq:jacobi} に $A=hP,\ B=hQ$ を代入すると、
$h'$ に関する二項が打ち消され、
$J=h(P C_v+Q C_w-C(P_v+Q_w))$ となる。
$h\ne0$ ならこの恒等式が示す条件に同値である。
\end{proof}

\subsection{直線と滑らかな二次曲線を射影三次曲線とする場合}
\label{p12:sec:conic-stratum}
$Y=\PP^2$ とし、その Poisson 双ベクトル $G$ の零因子が直線 $L$ と滑らかな
二次曲線 $Q$ の和である場合を考える。$L$ が $Q$ に接する場合も許す。
平面に沿う射影を $G$ に固定する。大域ベクトル場 $V_i$ により
\[
 \pi=G+V(t)\wedge\partial_t,\qquad
 V(t)=V_0+tV_1+t^2V_2
\]
と表す。この符号では \eqref{p12:eq:jacobi} は
$\mathcal L_{V(t)}G=V(t)\wedge V'(t)$ である。$V_i$ がすべて $Q$ に
接することを「$Q$ に沿って対数的」と呼ぶ。これは対応する平坦接続の
すべての成分が $Q$ の一般点で対数極を持つという条件（正則なものも含む）
に一致する。$L$ に沿う対数性も同様に定める。

\begin{proposition}\label{p12:prop:one-bad-conic}
すべての $\pi$ は $L,Q$ の少なくとも一方に沿って対数的である。
\end{proposition}
\begin{proof}
$Y\setminus(L\cup Q)$ の基本群は二次曲線の正のメリディアン $m_Q$ で生成される
無限巡回群で、直線のメリディアンは $m_L=m_Q^{-2}$ である。
実際、$L$ を無限遠とすると、接線の場合はアフィン二次曲線を $y=x^2$ と書け、
補集合は $\mathbb C\times\mathbb C^*$ である。割線の場合は二次曲線を $xy=1$
と書ける。補集合からさらに $x=0$ を除くと座標 $(x,xy)$ により
$\mathbb C^*\times(\mathbb C\setminus\{1\})$ を得る。
この開集合は元の補集合の基本群へ全射し、第1因子のループは $y=0$ 上で縮む。
残る生成元の位数は $xy-1$ の巻き数により無限である。いずれの場合も
有理関数 $Q/L^2$ の巻き数がメリディアンの関係を与える。

$Q$ に沿って非対数的なら、対応するトレース零の階数2平坦接続で
$m_Q$ のモノドロミー固有値はともに $-1$ である。実際、一般点の
局所座標 $z,s$（$Q=(z=0)$）で接続行列を $(A\,dz+B\,ds)/z$ と書くと、
平坦性から $[A|_{z=0},B|_{z=0}]=B|_{z=0}\ne0$ を得る。
従って $A|_{z=0}$ の固有値は $\pm1/2$ であり、横断円板上の正則特異接続の
モノドロミーの固有値はともに $-1$ である。
このとき $m_L=m_Q^{-2}$ の固有値はともに $+1$ であるから、
同じ理由で $L$ に沿う非対数極は存在しない。
\end{proof}

\begin{theorem}\label{p12:th:conic-log-classification}
$G$ を固定し、$Q$ に沿って対数的な $\pi$ を考える。混合テンソルの階数が
$1$ 以下の場合は $C_{\mathrm{line\text{-}conic}}$ の表示
\eqref{p12:eq:lineconic} から得られる。階数が $2$ 以上の場合はすべて
定理\ref{p12:th:threecomponents} の第2族 $C_{(1,1)(1,0)(0,2)}$ の表示から得られ、
階数はちょうど $2$ である。
\end{theorem}
\begin{proof}
$W=\operatorname{Lie}\operatorname{Aut}(Y,Q)\cong\mathfrak{sl}_2$ とする。
$W$ の外積は $Q$ 上で零であるから
\[
 j:\textstyle\bigwedge^2 W\longrightarrow Q\cdot H^0(Y,\mathcal O(1))
\]
を定める。$\operatorname{Aut}(Y,Q)$ は $Q$ 外で開軌道を持つので
$j$ は零写像でない。両辺とも $\operatorname{PGL}_2$ の3次元既約表現である。
実際 $Y$ を二変数二次形式の射影平面と同一視すれば $Q$ は判別式二次曲線であり、
$H^0(Y,\mathcal O(1))$ は二次対称冪である。従って $j$ は同型である。
$k=[\ ,\ ]\circ j^{-1}$ とおき $h=k(G)\ne0$ とする。

仮定は $V_i\in W$ である。$k$ の同変性と $k(V\wedge V')=[V,V']$ から
Jacobi 方程式は
\[
 [V_0,h]=[V_0,V_1],\quad [V_1,h]=2[V_0,V_2],\quad
 [V_2,h]=[V_1,V_2]
\]
となる。階数1以下なら $V_i$ はすべて $h$ と可換である。
$\mathfrak{sl}_2$ の非零元の中心化代数は一次元なので $V(t)=f(t)h$ であり、
これは \eqref{p12:eq:lineconic} に一致する。実際 $\mathcal L_hG=0$ で、
$L\,dQ-2Q\,dL$ に対応する非零ベクトル場は $W$ に属し $G$ を保つ
（$G$ で割ると $d\log(Q/L^2)$ に比例する）。

階数2以上なら、二つの独立な値を $0,\infty$ で取るように
$\PP^1$ の座標を選べるから、$V_0,V_2$ は独立としてよい。第1、第3式より
$h-V_1=aV_0,\ h+V_1=bV_2$ であり、第2式より $ab=-4$ である。
従って $V(t)=(1-at/2)V_0+t(t+b/2)V_2$ は $t=2/a$ で零となる。
この点を無限遠に移せば、新座標の二次係数は零で、第2式は
$[V_1,h]=0$ である。$V_0,V_1$ は独立なので第1式から $V_1=h$ を得る。
こうして標準形は $V(t)=Z+th$, $Z\in W\setminus\mathbb Ch$ となる。

最後にこの標準形を対数族と照合する。$Q$ の方程式を固定し、
$G$ の式が $2QL$ となるように直線の方程式 $L$ を選ぶ。
定理\ref{p12:th:threecomponents} 第2族で $U$ を無限遠、
$F=tL+M$（$M$ は任意の平面一次形式）とすると、アフィン表示は
\[
 2QL\,dt+(2Q\,dM-M\,dQ)+t(2Q\,dL-L\,dQ)
\]
となる。$M\mapsto 2Q\,dM-M\,dQ$ に対応するベクトル場は $Q$ に接する。
この線形写像は単射である。非零の $M$ が核にあれば
$d(M^2/Q)=0$ となり、$Q$ が一次形式の平方となるためである。
従ってこの写像は $W$ 上への同型である。$M$ を $L$ と独立に選べば階数は2で、
既に示した方程式から一次係数は $h$ である。この係数は $M$ と独立であり、
定数係数は $W$ の任意の元を取るので、すべての標準形が得られる。
\end{proof}

\begin{theorem}\label{p12:th:secant-classification}
$L$ が $Q$ と横断的に交わるとする。射影双ベクトル $G$ を固定した
集合論的ファイバーは、\eqref{p12:eq:lineconic} および
定理\ref{p12:th:threecomponents} の三つの表示で $G$ を誘導するものの和集合である。
特に、三次曲線が一般の「滑らかな二次曲線とその割線」であるような
Poisson 既約成分は、これらの四成分に尽くされる。
\end{theorem}
\begin{proof}
命題\ref{p12:prop:one-bad-conic}と定理\ref{p12:th:conic-log-classification}により、
$L$ に沿って対数的で $Q$ に沿って非対数的な場合を調べればよい。
$L$ を無限遠、$Q=(q=xy-1=0)$ ととる。$L$ に接する射影ベクトル場は
$x,y$ に関しアフィン一次式を係数にもつ。従って $G$ を係数とする接続行列は
\[
 \Omega=\frac{(a+bx+cy)\,dx+(d+ex+fy)\,dy}{q},
 \qquad a,b,c,d,e,f\in\mathfrak{sl}_2.
\]
ここで分子の定数倍に $G$ の尺度を吸収した。平坦性
$d\Omega+\Omega\wedge\Omega=0$ の六つの係数は
\begin{align*}
 [e,b]=b,\qquad [c,f]=f,\qquad [b,f]+[c,e]=0,\\
 a+[a,e]+[b,d]=0,\qquad -d+[a,f]+[c,d]=0,\qquad [a,d]=e-c .
\end{align*}
$Q$ の一般点で $x$ を一定にした横断円板上の留数は
$e+d/x+f/x^2$ である。前命題の局所計算からその固有値は
$\pm1/2$ である。行列式の恒等式で $x\to\infty$ とすると
$e$ も同じ固有値を持つ。$x,y$ を交換すれば $c$ についても同様である。

$H=\operatorname{diag}(1/2,-1/2), J=E_{12}, K=E_{21}$ とする。
$b,f$ の非零なものはそれぞれ $\operatorname{ad}e,\operatorname{ad}c$ の
重み1のベクトルなので冪零である。両方が非零で比例しないなら、
それぞれの核の直線を基底として $b=BJ,f=DK$ と書ける。
最初の二式から $e=H+rJ,c=-H+sK$ であるが、第3式の
$J,K,H$ 成分から順に $r=s=0,BD=0$ となり矛盾する。
両方が非零で比例する場合は $b=BJ,f=DJ$ とできる。このとき
$e=H+rJ,c=H+sJ$ で、第3式から $r=s$ である。
定数共役により $c=e=H$ としてよい。

一方だけ非零なら $x,y$ の交換で $b=BJ\ne0,f=0$ とできる。
$e=H$ としたとき第3式と $c$ の固有値から $c=\pm H$ となる。
$c=-H$ なら第5式より $d=d_-K$、第6式より $d_-\ne0$ かつ
$a$ の $H$ 成分は零である。第4式の $H$ 成分 $2Bd_-=0$ はこれと矛盾する。
$b=f=0$ の場合も $[c,e]=0$ より $c=\pm H$, $e=H$ とできる。
このうち $c=-H$ のとき、第4--6式から
$a=AJ$, $d=A^{-1}K$ ($A\ne0$) を得る。

残った $c=e=H$ の場合は $b=BJ,f=DJ$（零も可）であり、第4式の
$K$ 成分から $a_K=0$、第5式の $K$ 成分から $d_K=0$ を得る。
両式の $H$ 成分から $a_H=d_H=0$ となるので $a,d\in\mathbb CJ$ である。
以上により定数共役の後に接続行列は次の二形に限られる。
\begin{align*}
 &\Omega=(H\,dq+J\,dP)/q,
    &&P=p_1x+p_2y+p_3 x^2+p_4 y^2,\\
 &\Omega=\big[(AJ-Hy)\,dx+(A^{-1}K+Hx)\,dy\big]/q,
    &&A\ne0.
\end{align*}
第1形の射影水平方程式は $d(qt+P)=0$ である。
$F=qt+P$ を斉次化すると $(1,2)$ 型となり、得られる双ベクトルは
定理\ref{p12:th:threecomponents} の第1族（$U,L$ が無限遠）に属する。
第2形は
\[
 M=\begin{pmatrix}y&-A\\-A^{-1}&x\end{pmatrix},\qquad
 \det M=q,\qquad
 \Omega=M^{-1}dM-\tfrac12\operatorname{tr}(M^{-1}dM)I
\]
である。従って $M(t,1)^{\mathsf T}$ の成分比が第一積分である。
この二成分の斉次化はともに $(1,1)$ 型であり、
分母を払った1形式に無限遠の $L$ を掛けたものが
定理\ref{p12:th:threecomponents} の第3族である。
実際、無限遠を除くチャートでは両辺の $dt$ 係数は
$q$ の定数倍で等しく、正規化した水平方程式も同じであるから、
$G$ の尺度を合わせれば双ベクトルとして一致する。
\end{proof}

\subsubsection{直線が二次曲線に接するときの非対数的な極}
\label{p12:sec:tangentgoodline}
\begin{proposition}\label{p12:th:tangent-line}
平面三次曲線 $L+Q$ が滑らかな二次曲線 $Q$ とその接線 $L$ からなり、
平坦接続が $L$ の一般点で対数的、$Q$ の一般点で非対数的であるとする。
対応する Poisson 構造は $C_{(1,2)(1,0)(0,1)}$ 型の族または
$C_{(1,1)^2(0,1)}$ 型の族に属する。
\end{proposition}
\begin{proof}
平面の射影変換と全体のスカラー倍を施し、
$L$ を無限遠直線、$Q$ のアフィン方程式を $q=y-x^2$ とする。
$L$ で対数的であることは混合項の各射影ベクトル場が $L$ に接することと同値であり、
それらはアフィン線形な係数を持つ。従って水平持ち上げを
$\partial_i+\Omega_i$ と書く規約では、
\[
 \Omega=(X\,dx+Y\,dy)/q,\qquad
 X=Mx+Ny+A,\quad Y=Px+Ry+B,
\]
ここで係数は $\mathfrak{sl}_2$ に属する。
$\mathfrak{sl}_2$ を射影直線のベクトル場として用いるときもその Lie 括弧を保つ。
平坦性を
$q(\partial_x Y-\partial_y X)-q_xY+q_yX+[X,Y]=0$
と展開すると、次の六式を得る。
\begin{align*}
 [M,P]+P+N&=0, & [M,R]+[N,P]+2R&=0, & [N,R]&=0,\\
 (1-\operatorname{ad}B)M+2B+[A,P]&=0, &
 P+[N,B]+[A,R]&=0, & [B,A]&=A.
\end{align*}
$(x,q)$ 座標では $dq/q$ の係数を $Q$ に制限した行列は
$Y(x,x^2)=B+xP+x^2R$ である。$dx/q$ の係数の $Q$ 上の制限が
一般点で非零なので、平坦性の二位の極の係数から
$\operatorname{ad}Y(x,x^2)$ は固有値1を持つ。従って全ての $x$ に対し
$\operatorname{tr}(Y(x,x^2)^2)=1/2$ である。
特に定数共役で $B=h=\operatorname{diag}(1/2,-1/2)$ とできる。
上・下三角の標準行列をそれぞれ $e,f$ とし、
$[h,e]=e,\ [h,f]=-f,\ [e,f]=2h$ とする。
先の六式の最後から $A=ae$、トレース恒等式の一次係数から
$P=be+cf$ と書ける。$R=E e+K h+D f$ とおくと、その二・三・四次係数は
\[
 K=-2bc,\qquad bD+cE=0,\qquad K^2+4ED=0
\]
となる。六式の下段の最初の二式、ついで上段の第一式を比較すると
\begin{align*}
 M&=m e+H h, & H&=-2(1+ac),\\
 N&=(b-aK)e+n h-c f, & aD&=0,\qquad Hc=0,\qquad n=-2mc.
\end{align*}

$c=0$ なら $K=0,\ H=-2,\ N=P=be$ である。
$2R+[m e-2h,R]=0$ の $f$ 成分から $D=0$、従って $R=Ee$ となる。このとき
\[
 \Omega=h\,dq/q+e\,dg/q,\qquad
 g=m x^2/2+bxy+ax+E y^2/2 .
\]
基底 $h,e,f$ を $t\partial_t,t^2\partial_t,-\partial_t$ で実現すると、
水平な方向では $d(q/t+g)=0$ となる。$q$ と $g$ を
平面上の二次式に斉次化すれば、この有理第一積分は
$F/(U L^2)$ の形をとる。ここでアフィン座標では $F=q+tg$、$U=t$ であり、
その双次数は $(1,2),(1,0)$ である。
実際 $U L\,dF-FL\,dU-2FU\,dL$ のアフィン式は
$-q\,dt+t\,dq+t^2\,dg$ で、水平な1形式の定数倍に一致する。
従って主張の第一の族に入る。

$c\ne0$ なら $H=0,\ ac=-1$ である。$aD=0$ から $D=0$、
トレース条件から順に $K=b=E=0$、従って $R=0$ となる。
$[N,P]=0$ から $m=0$ が従い、$M=0$, $N=-P=-cf$ を得る。
対角共役で $a=1,c=-1$ とする。同じベクトル場の実現では
水平な1形式は
\[
 q\,dt+(y-t^2)\,dx-(t+x)\,dy
 = G\,dF-F\,dG,\qquad F=t+x,\quad G=tx+y
\]
である。$F,G$ を双次数 $(1,1)$ に斉次化し $G\,dF-F\,dG$ に
$L$ を掛けることで元の大域双ベクトルを得る。
\end{proof}


\begin{theorem}\label{p12:th:wholeconic}
$C=LQ$ が直線と滑らかな二次曲線の積なら、$q^{-1}(C)$ の閉点を尽くすのは
$C_{\mathrm{line\text{-}conic}}$、$C_{(1,1)(1,0)(0,2)}$、
$C_{(1,2)(1,0)(0,1)}$、$C_{(1,1)^2(0,1)}$ の四族である。
\end{theorem}
\begin{proof}
命題\ref{p12:prop:one-bad-conic} による二分のうち、
二次曲線に沿って対数的な場合は定理\ref{p12:th:conic-log-classification} に従う。
残りは割線の場合は定理\ref{p12:th:secant-classification}、接線の場合は命題\ref{p12:th:tangent-line} に従う。
\end{proof}

\subsection{三本の非共点直線の上のファイバー}\label{p12:sec:triangle}
\begin{theorem}\label{p12:th:triangle}
射影三次曲線が三本の相異なる非共点直線の和である Poisson 双ベクトルは、
$C_{\mathrm{tor}}$, $C_4$, $C_{(1,2)(1,0)(0,1)}$ のいずれかに入る。
より正確には、三直線のうち平坦接続が一般点で非対数的なものの本数は
$0,1,2$ のいずれかであり、各場合に順に上記の成分に属する。
\end{theorem}
\begin{proof}
補集合は $(\mathbb C^*)^2$ で、三本のメリディアンは可換で積が1である。
非対数極のメリディアンのモノドロミー固有値はともに $-1$ である。
可換な行列は同時三角化できるので、この固有値条件が三本すべてでは成立しない。

まず少なくとも二本に沿って対数的とし、それらを $x=0,y=0$、残りを無限遠とする。
両軸に接する射影ベクトル場は
$x(a+\ell)\partial_x+y(b+\ell)\partial_y$ と書ける。
ここで $a,b$ は定数、$\ell$ は斉次一次式である（射影ベクトル場の
二次部分は共通の一次式と Euler 場の積である）。
$G$ のアフィン方程式は $xy$ の定数倍なので、接続行列は定数行列
$A,B,S,T\in\mathfrak{sl}_2$ を用いて
\[
 \Omega=(A+Sx+Ty)\,d\log x+(B-Sx-Ty)\,d\log y
\]
となる。この表示において無限遠でも対数的となる条件は $S=T=0$ である。
実際 $d\log(x/y)$ は無限遠直線に接する一般の方向への制限が非零で、
$Sx+Ty$ がその一般点に極を持たないためには $S=T=0$ が必要である。
平坦性は係数比較により
\[
 [A,B]=0,\qquad [S,A+B]=S,\qquad [T,A+B]=T
\]
と同値である。従って $S=T=0$ なら $A,B$ は共通の
$\mathfrak{sl}_2$ 元のスカラー倍である。この元が半単純の場合は
トーラス上の対数双ベクトルであり、冪零の場合も共通の元を半単純元から
特殊化できるので $C_{\mathrm{tor}}$ に属する。

$(S,T)\ne(0,0)$ なら $A+B$ の随伴表現は固有値 $-1$ をもつから、
定数共役で $A+B=-H$, $H=\operatorname{diag}(1/2,-1/2)$ とできる。
$A,B$ は $A+B$ と可換であり $A=aH,B=bH$, $a+b=-1$ である。
また $S=sJ,T=rJ$ ($J=E_{12}$) である。$\eta=a\,d\log x+b\,d\log y$
と書くと射影水平方程式は
\[
 dt+t\eta+(sx+ry)(d\log x-d\log y)=0 .
\]
一般の $a,b$ には $a(1+a)\ne0$ が成立する。このとき
$P=p x+k y$ を $(1+a)p=s,\ ak=r$ で定めると
$dP+P\eta=(sx+ry)(d\log x-d\log y)$ であり、
方程式は $d(t+P)+(t+P)\eta=0$ に一致する。
これは $t+P$ の斉次化を $(1,1)$ 因子とする四因子対数形式
\eqref{p12:eq:four} である。$a=0,-1$ についても、
固定した $s,r$ のもとで上の行列表示の $a$ を特殊化すれば $C_4$ に属する。

次に非対数極が二本の場合、それらを座標軸 $x=0,y=0$ に取る。
無限遠で対数的なので、$q=xy$ として
\[
 \Omega=\big[(a+bx+cy)\,dx+(d+ex+fy)\,dy\big]/q
\]
と書く。定理\ref{p12:th:secant-classification} の $xy-1$ の計算と比較すると、
平坦性の六式のうち定数式だけが $[a,d]=0$ に変わる。
$y=0$ と $x=0$ を横切る円板上の留数はそれぞれ $e+d/x$, $c+a/y$ であり、
その固有値が $\pm1/2$ であることから、$e,c$ にも同じ固有値がある。
$H,J,K$ を先の証明と同じ行列とする。
$b,f$ が非比例の場合の排除と、$b= B J, f=D J, c=e=H$ への帰着
（両方非零かつ比例する場合）は最初の三式だけを用いたので、そのまま適用できる。
一方だけ非零なら対称性から $b=BJ\ne0,f=0$ とでき、$e=H,c=\pm H$ である。
$b=f=0$ でも $e=H,c=\pm H$ とできる。
$c=-H$, $b=BJ\ne0,f=0$ のとき $d=d_-K$,
$a_K=0,\ a_H=-2Bd_-$ である。
$[a,d]=0$ のもとで $d_-\ne0$ なら $a\in\mathbb CK$ となり
$a_H=0$ と矛盾する。従って $d=0$ となるが、
このとき $x=0$ 上で $dy$ の分子 $d+fy$ が零となり対数的である。
$c=-H$, $b=f=0$ のときは $a=AJ,d=DK$ となり $[a,d]=0$ から
$AD=0$ を得るので、やはりどちらかの軸で対数的になる。

残る $c=e=H$ では、先の計算の定数式以外から
$a,b,d,f\in\mathbb CJ$ を得る。
従って $\Omega=(H\,dq+J\,dP)/q$（$P$ は二次以下の多項式）であり、
$q t+P$ が有理第一積分である。$qt+P$ を斉次化した $(1,2)$ 因子と
無限遠の二直線を用いた定理\ref{p12:th:threecomponents} 第1族の表示を得る。
\end{proof}

\subsection{共点三直線の Riccati 成分}
中心 $p$ で交わる三本の相異なる直線の和を $D$ とする。
$p$ を $y$ 軸方向の無限遠点に置き、三直線の方程式を
$f(x)=\prod_{i=1}^3(x-r_i)=0$ と取る。平面双ベクトルのスカラー倍は
以下の族に含める。ファイバー座標 $t$ に関する二次以下の多項式
$h_0,h_1$ を選び、
\begin{equation}\label{p12:eq:concurrent-family}
 \pi= f(x)\,\partial_x\wedge\partial_y+
 h(t,x)\,\partial_t\wedge\partial_y,\qquad h=h_0(t)+x h_1(t)
\end{equation}
とする。これは式~\eqref{p12:eq:jacobi} を満たす。この族を座標変換で動かした閉包を
$C_{\mathrm{Ric}}$ とする。

\begin{theorem}\label{p12:th:concurrent-riccati}
$C_{\mathrm{Ric}}$ は12次元の既約成分である。その一般閉点で
$\operatorname{Pois}(X)$ の接空間は13次元であり、この成分において $\operatorname{Pois}(X)$ は
生成点で非被約である。
\end{theorem}
\begin{proof}
共点三直線を与える非零の平面双ベクトルの層は、中心 $p$ の2次元と、
$p$ を消す二つの一次式についての binary cubic の4次元からなる。
中心は一意である。$p$ を固定し、同じ中心の直線束の各葉に接する
ベクトル場は $W_p$（命題~\ref{p12:prop:zero-cubic}）をなし、
今の座標では $(a+bx+cy)\partial_y$ と書ける。このような混合項では
Jacobi 式は $-f\partial_y(a+bx+cy)=0$（$a,b,c$ は $t$ の二次以下の式）なので
$c=0$ である。従って~\eqref{p12:eq:concurrent-family} は、この層上で
各 $V_i$ が三直線に接する族の全体であり、ファイバーは6次元線形空間である。
実際、三直線を保存する射影変換の単位成分は、交点 $p$ を固定し、
直線束上で三点を固定するので束に自明に作用する。
したがって族は既約な12次元である。

この族が既約成分となることをモノドロミーで示す。
式~\eqref{p12:eq:concurrent-family} の接続は、直線束の底上の
\[
 d+\lambda(M_0+x M_1)\,dx/f(x)
\]
の引き戻しである。ここで $M_0,M_1\in\mathfrak{sl}_2$ は各 $h_i\partial_t$
に対応し、$\lambda$ は混合項を動かすスカラーである。
$[M_0,M_1]\ne0$ とする。$r_i$ の周りの基点付きメリディアンのホロノミーは
$I-2\pi i\lambda R_i+O(\lambda^2)$、$R_i=(M_0+r_iM_1)/f'(r_i)$ である。
異なる二つの $R_i$ の括弧は非零なので、十分小さな $\lambda\ne0$ で
対応するホロノミー行列は可換でない。一方、三つのメリディアンのトレースは
いずれも $2$ に近く、$-2$ とは異なる。

上の一つの点の近傍で、三次曲線の被約性は保たれる。
有限個の基点付き閉路を固定すると、その像を避ける近傍でホロノミーが
パラメータに連続に依存するので、二つの行列の非可換性も保たれる。
被約平面三次曲線のうち、補集合の基本群が非可換なのは共点三直線だけである
\cite[\S2.3.2, (3a)--(3e)]{p12-Cogolludo}。
よってこの近傍のすべての Poisson 構造の平面成分は共点三直線である。
さらに近傍を縮めれば三本を局所的に番号付けでき、上で選んだ各メリディアンは
動いた直線の対応するメリディアンに自由ホモトピックである。
実際、各直線の滑らかな点に置いた小さい横断円板を使えばよい。
定理~\ref{p12:th:irreducible-cubic} の曲率計算により、ある直線に沿って
接続が非対数的なら、その一般のメリディアンのホロノミーのトレースは $-2$ である。
この局所的結論には基本群に関する仮定を必要としない。
トレースの連続性により三本とも対数的でなければならない。
従って $V_i$ はすべてこの三直線に接し、近傍の閉点は先に記述した
12次元族に属する。これより $C_{\mathrm{Ric}}$ は既約成分である。

次に接空間を調べる。固定 $f$ に対する六つの混合方向があるほか、
射影 $q$ の微分の像は $x$ の三次以下の多項式全体を含む。
さらに、任意の二次以下の $g(x)$ について $\dot F=y g(x)$ もこの像に入る。
実際、$x$ に関して二次以下の $P(x,t)$ を
\[
 P(r_i,t)=-h(t,r_i)g(r_i)/f'(r_i)\qquad (1\le i\le3)
\]
で補間し、その $x^2$ 係数を $c(t)$ とする。
$\dot A=P$, $\dot B=y(d(t)+c(t)x)$ は正則な混合項であり、
この変形の一次 Jacobi 条件は
\[
 (f'+h_t)P-hP_t-f(P_x+c x)+h g=f d
\]
となる。左辺は各 $r_i$ で消える。$t$ の微分を含む二項は
二次以下の多項式ベクトル場の括弧であり $t$ の次数は高々2である。
左辺の $x^4$ の係数も $f'P-f(P_x+c x)$ において打ち消し合う。
したがって左辺を $f$ で割れば $t$ の二次以下の多項式 $d$ が得られる。
以上から接空間の次元は少なくとも $6+4+3=13$ である。

最後に一点で階数の逆向きの評価を確認する。
$f=(x-1)(x+2)(x-3)$、
$h=1+2t+3t^2+x(5+7t+11t^2)$ と置く。
Jacobi 式の微分は
\[
 (f'+h_t)\dot A-h\dot A_t-f(\dot A_x+\dot B_y)+h\dot F_y .
\]
$t=0,1,2$、$x,y=0,1,2,3$ の組でこれを評価した行列は
$\mathbb F_{101}$ 上で階数21である
（付属の \path{p1p2/verify_concurrent.py}）。よって標数零で
一般点の接空間は高々 $34-21=13$ 次元である。
台の既約成分の次元は12であり、他の成分との交わりを除けば
この次元差は生成点の非被約性を意味する。
\end{proof}

\begin{theorem}\label{p12:th:double}
$C_{\mathrm{Ric}}$ の生成点 $\eta$ における $\operatorname{Pois}(X)$ の局所環の長さは2である。
\end{theorem}
\begin{proof}
$J$ を30個の斉次二次式の組として用いる。
一次変形 $\delta$ が2次まで延長する条件は、
$J(\delta)\in\operatorname{im} dJ_\pi$ である。実際、
$J(\pi+\epsilon\delta+\epsilon^2\delta_2)$ の $\epsilon^2$ 係数は
$J(\delta)+dJ_\pi(\delta_2)$ である。
前の証明の補間で $g=1$ を取り、
\[
 f=x(x-1)(x-2),\qquad h=1+2t+3t^2+x(4+7t+11t^2)
\]
に特殊化する。ここで $\pi=(0,h,f)$ とし、補間から得る一次変形を明記すると
\[
 \delta=(P,\ y(d+cx),\ y),
\]
\begin{align*}
 P={}&-\tfrac12\bigl[(42t^2+27t+15)x^2
              -(73t^2+47t+26)x+(3t^2+2t+1)\bigr],\\
 c={}&-\tfrac12(42t^2+27t+15),\qquad
 d=\tfrac12(17t^2+19t+10)
\end{align*}
である。$P$ の $x^2$ 係数が $c$ なので、$\delta$ は大域双ベクトルを表す。
\eqref{p12:eq:jacobi} への代入により $dJ_\pi(\delta)=0$ となり、
\[
 J(\delta)=\frac y4\bigl[
 (336t^2+336t-12)x^2+(-417t^2-474t+81)x-123t^2-68t-51\bigr]
\]
を得る。多項式 $R$ の単項式 $m$ の係数を $[m]R$ と書き、線形汎関数
\[
 \ell(R)=\sum_{k,i=0}^2\ell_{ki}[t^kx^iy]R,\qquad
 (\ell_{ki})=
 \begin{pmatrix}
 886&-739&61\\
 -1521&279&-171\\
 756&81&81
 \end{pmatrix}
\]
を定める。線形化の式への代入は
\[
 \ell\circ dJ_\pi=0,\qquad \ell(J(\delta))=-72900\ne0
\]
を与える。従ってこの一次変形には2次障害がある。

階数21の証明書も記す。順序付き単項式列
\[
 \mathcal M=(1,y,y^2,y^3,x,xy,xy^2,x^2,x^2y,x^3)
\]
を用い、$dJ_\pi$ の行を $t^k m$ の順
（$k=0,1,2$、各 $k$ で $m\in\mathcal M$ の順）に並べる。
列は先に $(0,0,m)$ を $\mathcal M$ の順に置き、次に
$(t^kP_j,t^kQ_j,0)$ を並べる。ここで $k=0,1,2$、
$(P_j,Q_j)$ はアフィン表示の節に掲げた八つの平面ベクトル場の係数を
$(v,w)=(x,y)$ としたものである。添字を1から数えると、行・列集合
\begin{align*}
 I={}&(1,2,3,5,6,8,9,10,11,12,15,16,18,19,20,21,22,25,26,28,30),\\
 K={}&(2,3,4,6,7,9,11,13,14,16,17,18,19,21,22,24,25,26,27,30,34)
\end{align*}
で切り出した小行列の行列式は
\[
 \det(dJ_\pi)_{I,K}=-2^{19}3^{13}5^2\cdot11\cdot19\ne0
\]
である。前の証明で構成した13個の独立な接方向と合わせて
$\operatorname{rank}dJ_\pi=21$ が従う。したがって
$\operatorname{rank}[dJ_\pi\mid J(\delta)]=22$ である。
以上の代入と行列式は \path{p1p2/verify_double.py} でも有理数上で検証できる。

補間は $f$ の判別式を逆にすれば係数の有理式である。
固定した座標内の族でこの非零小行列式は開条件であり、
座標変換でその族が $C_{\mathrm{Ric}}$ を稠密に覆う。
従って成分の一般閉点に、2次に延長しない一次変形がある。

一般点の近くで微分が独立な21本の Jacobi 方程式を選ぶ。
その共通零点は滑らかな13次元多様体 $Z$ であり、
$C_{\mathrm{Ric}}$ の被約構造は一般点の近くで $Z$ の滑らかな因子である。
従って $\mathcal O_{Z,\eta}$ は離散付値環で、残余の方程式のイデアルは
一様化元 $s$ のある冪 $(s^m)$ である。他の成分を除く開集合をさらに縮めれば、
そこで同じ等式が成り立つ。定理\ref{p12:th:concurrent-riccati} の接空間次元より $m\ge2$。
$m\ge3$ なら、滑らかな $Z$ 上で任意の接ベクトルを2次まで延長すると
$s^m$ は $\epsilon^3$ の剰余で自動的に零になる。一般点での上の障害に反する。
従って $m=2$ である。
\end{proof}

\subsubsection{三本の相異なる共点直線の上の全解}
\label{p12:sec:concurrentfibers}
\begin{theorem}\label{p12:th:all-concurrent}
平面成分の零曲線が相異なる三本の共点直線の和であるとする。
少なくとも一本の直線の一般点で接続は対数的である。
三本とも対数的なら解は直線束からの引き戻し \eqref{p12:eq:concurrent-family} である。
ちょうど二本が非対数的なら解は $C_{(1,2)(1,0)(0,1)}$ に属する。
ちょうど一本が非対数的なら解は $C_4$ に属する。
\end{theorem}
\begin{proof}
三本を $x=0,1$ と無限遠直線に取る。平面成分を定数倍で
$q=x(x-1)$ に正規化する。射影ベクトル場の二次部分は動径方向にあるから、
水平持ち上げ $\partial_i+\Omega_i$ を表す行列1形式は
\[
 \Omega=(X\,dx+Y\,dy)/q,\quad
 X=A+Mx+Ny+Dxy+Ey^2,\quad
 Y=B+Px+Ry-Dx^2-Exy
\]
と書ける。各係数は $\mathfrak{sl}_2$ の元である。

初めに三本とも非対数的と仮定する。各直線の近傍で法線座標を $s$、
接方向の座標を $z$ として $\Omega=(S\,ds+T\,dz)/s$ と書く。
平坦性から $[S|_{s=0},T|_{s=0}]=T|_{s=0}$ であり、
非対数的なら一般点で $T|_{s=0}$ は非零べき零である。
三直線に沿うこの行列は、零でない有理関数を掛ける自由度の下で順に
\[
 B+Ry,\qquad B+P-D+(R-E)y,\qquad D+Ey'
\]
である。最後は $x'=1/x,\ y'=y/x$ と変換して得る。
各式は線形で、$\mathfrak{sl}_2$ のべき零元の射影化は滑らかな二次曲線であるから、
それぞれ一定のべき零方向 $e_i$ を張る。

この三方向は互いに異なる。実際、共点を通らない一般の直線 $y=ax+b$ に制限すれば、
自明な階数2束の行列1形式で三つの単純極だけをもつ。その留数 $S_i$ は
$\sum_i S_i=0$ を満たし、接方向項の寄与は $e_i$ の倍数なので
$[S_i,e_i]=e_i$ となる。トレース零より $S_i$ は $\ker e_i$ 上に $1/2$ で作用し、
固有値は $\pm1/2$ である。二方向が一致すれば第三の留数は同じ直線上で
$-1$ として作用するので矛盾する。

ところが、三つの相異なるべき零方向は滑らかな射影二次曲線上の三点なので
線形独立である。$R,R-E,E$ はそれぞれこの三方向上にあって
$R-(R-E)-E=0$ だから $R=E=0$ である。
$Y=Y_0(x)=B+Px-Dx^2$ とおくと、平坦性の $y$ の係数は $[N+Dx,Y_0(x)]=0$ となる。
$Y_0(0),Y_0(1)$ は非零べき零で、その $\mathfrak{sl}_2$ における中心化代数は
それぞれの張る直線である。従って $N,N+D,D$ はそれぞれ $e_0,e_1,e_\infty$
の方向にあり、同じ線形独立性から $D=0$ となる。
これは無限遠直線で非対数的であることに反する。

対数的な直線を無限遠に選び直すと $D=E=0$ となる。
$X=A+Mx+Ny,\ Y=B+Px+Ry$ を平坦性の式に代入して
\begin{align*}
 [M,P]&=P+N, & [M,R]+[N,P]&=2R,&[N,R]&=0,\\
 N-2B+[M,B]+[A,P]&=0, & R+[N,B]+[A,R]&=0,&[A,B]&=-B
\end{align*}
を得る。$x=0,1$ の接方向の係数は対数的なら零であり、いずれの場合にも
$B+Ry$ と $B+P+Ry$ は全ての $y$ に対してべき零である。

$R\ne0$ なら $R,B,P$ は同じべき零方向 $e$ を持ち、
両直線とも非対数的である。$[N,R]=0$ から $N$ も $e$ の倍数となる。
$[M,R]=2R,\ [A,R]=-R$ から $M=2h+m e,\ A=-h+a e$、ついで $N=P$ を得る。
ここで $[h,e]=e$ とし、$h,e$ は前節と同じ標準形に取る。
よって $P=p e,\ R=r e,\ B=b e$ と書くと
\[
  \Omega=h\,dq/q+e\,dg/q,\qquad
  g=mx^2/2+p xy+r y^2/2+ax+by.
\]
これは有理第一積分 $q/t+g$ の $(1,2)$ 型の族である。

$R=0$ とする。べき零元 $B,B+P$ が一次独立なら $[N,B]=[N,P]=0$ から $N=0$ である。
このとき $P$ は独立な二つのべき零元の差なので $\operatorname{tr}(P^2)\ne0$ だが、
$[M,P]=P$ に反する。従って $B=b e,\ P=p e$ と置ける。
$p=0$ なら $N=0$ である。$b\ne0$ なら $[M,B]=2B,\ [A,B]=-B$
から再び上の $(1,2)$ 型を得る。$p=b=0$ なら三本とも対数的で、
$Y=0,\ X=A+Mx$ が全解である。

$p\ne0$ なら $[N,P]=0$ なのである $m,u,v,a\in\mathbb C$ を用いて
\[
 M=m h+u e,\qquad N=(m-1)p e,\qquad A=a h+v e
\]
と書ける。実際、最初の式は $[M,P]=P+N$ を比較し、$A$ の式は
$N-2B+[M,B]+[A,P]=0$ を比較して得る。
残る条件は $b(a+1)=0,\ p(a+m-1)=b(2-m)$ である。
$b(b+p)\ne0$ なら両直線とも非対数的で、$a=-1,m=2$ となり上の $(1,2)$ 型である。
残る場合には $0$ と $1$ を交換すれば $b=0$ とでき、このとき $a=1-m$ である。
$k=m-1$ と置くと、接続の係数は
\[
 h\left(k\frac{dx}{x}+\frac{dx}{x-1}\right)
 +e\left(\frac{ux+kpy+v}{x(x-1)}\,dx+\frac{p}{x-1}\,dy\right).
\]
$k\ne0,-1$ とし、$\ell(x)=u x/(k+1)+v/k$ と置く。
$h=t\partial_t,\ e=t^2\partial_t$ で水平方程式を書けば、
$H=(x-1)/t+p y+\ell(x)$ は $dH+kH\,dx/x=0$ を満たす。
$H$ を $F/(U L_\infty)$ と斉次表示すると $U=t$、$F$ は $(1,1)$ 型で、
$x$ の零直線と $L_\infty$ を合わせて \eqref{p12:eq:four} の四因子を得る。
分母を払った1形式の $dt$ 係数は $q$ の非零定数倍なので、
この表示は大域双ベクトル自身の定数倍までの一致を与える。
$k=0,-1$ でも、$p,u,v$ を固定して $k$ を動かした元の接続は $k$ の多項式族であり、
解は閉包 $C_4$ に入る。
\end{proof}


\subsection{二重直線と単純直線}
\subsubsection{二重直線と単純直線のうち単純直線で対数的な場合}
\begin{proposition}\label{p12:th:double-good}
平面成分の零因子を $2L+M_0$ とし、$L\ne M_0$ とする。
単純直線 $M_0$ の一般点で接続が対数的なら、解は
$C_{(1,2)(1,0)(0,1)}, C_{\mathrm{Ric}}, C_4, C_{\mathrm{tor}}$
のいずれかに属する。
\end{proposition}
\begin{proof}
$M_0$ を無限遠直線、$L=(x=0)$ として平面成分を $q=x^2$ に正規化する。
$M_0$ における対数性より、水平接続を定める行列1形式は
\[
 \Omega=((A+Bx+Cy)\,dx+(D+Ex+Gy)\,dy)/x^2
\]
である。各係数は $\mathfrak{sl}_2$ の元であり、平坦性は
\begin{align*}
 [B,E]&=E+C,&[B,G]+[C,E]&=2G,&[C,G]&=0,\\
 [B,D]+[A,E]&=2D,&[C,D]+[A,G]&=0,&[A,D]&=0
\end{align*}
に等しい。とくに全ての $y$ で $[A+Cy,D+Gy]=0$ である。

$D,G$ が一次独立なら、$[A,D]=[C,G]=0$ および中心化代数の一次元性から
$A=aD,C=cG$ と書ける。混合係数の消滅と $[D,G]\ne0$ より $a=c$ となる。
$y'=y+ax$ と変換すると、新しい表示の $A,C$ は零になり、$D,G$ は変わらない。
ところがこの表示で $[B,D]=2D,\ [B,G]=2G$ を満たすのは不可能である。
実際 $\mathfrak{sl}_2$ の随伴作用の非零固有値に対する固有空間は高々一次元である。

$D,G$ は一次従属で少なくとも一方は非零とする。
$D=dK,\ G=gK$ と置くと、先の交換関係から
$A=aK,\ C=cK$ と書ける。従って
\[
 g[B,K]+c[K,E]=2gK,\qquad d[B,K]+a[K,E]=2dK.
\]
$(d,g),(a,c)$ が一次独立なら $[B,K]=2K,\ [K,E]=0$ である。
よって $K$ はべき零、$E$ は $K$ の定数倍であり、$[B,E]=E+C$ から $C=E$ を得る。
$[h,e]=e,\ K=e$ の標準表示を取ると $B=2h+b e$ となる。
すると $h$ 成分は $dq/q$ で、$e$ 成分の分子は $C=E$ によって閉じたアフィン線形1形式である。
従ってある二次以下の多項式 $g_0$ について
$\Omega=h\,dq/q+e\,dg_0/q$ となり、第一積分 $q/t+g_0$ をもつ $(1,2)$ 型の族に入る。
二組が比例すれば再び $y'=y+\lambda x$ により $A=C=0$ とできる。
$[B,K]=2K,\ [B,E]=E$ が同時に成り立つことから $E=0$ であり、同じ形の解を得る。

以下 $D=G=0$ とする。$E=0$ なら $C=0$ で、$\Omega=(A+Bx)dx/x^2$ は
直線束からの引き戻しである。射影直線上の binary cubic の重根を動かして
\eqref{p12:eq:concurrent-family} からの特殊化と書ける。
$E\ne0$ なら $[A,E]=[C,E]=0$ から $A=aE,C=cE$ と書け、
$[B,E]=(1+c)E$ となる。

$c\ne-1$ なら $E$ はべき零で、$E=e,\ B=(1+c)h+b e$ と置ける。
接続は
\[
 (1+c)h\,\frac{dx}{x}+
 e\left(\frac{a+bx+cy}{x^2}\,dx+\frac{dy}{x}\right).
\]
さらに $c\ne0$ とする。$h=t\partial_t,\ e=t^2\partial_t$ とすれば
\[
 W=\frac1t+\frac{y+a/c}{x}+\frac{b}{1+c}
\]
は水平方向で $dW+(1+c)W\,dx/x=0$ を満たす。
$W=F/(U L)$ と斉次表示でき、アフィン式は $U=t$, $L=x$,
$F=x+t(y+a/c+bx/(1+c))$ である。
分母を払うと $(1,1)$ 型の $F$、$U$、$L$、無限遠直線を用いた
四因子形式 \eqref{p12:eq:four} が得られる。実際その $dt$ 係数は
$x^2$ の非零定数倍である。$c=0$ でも $a,b$ を固定して $c$ を動かせば、
元の行列表現が $c$ に多項式的に依存するので $C_4$ への所属が従う。

最後に $c=-1$ なら $B=bE$ であり、
\[
 \Omega=E\left(b\,d\log x+d\frac{y-a}{x}\right)
\]
となる。任意の $E$ についてこれがトーラス成分に属することは、三角形からの
次の明示的特殊化で分かる。$\delta\ne0$ とし、
\[
 q_\delta=x(x+\delta(y-a)),\qquad
 \Omega_\delta=E\left(b\,d\log x+
 \delta^{-1}d\log\frac{x+\delta(y-a)}{x}\right)
\]
とする。極の三直線は非共点である。平面のトーラス座標を $x,x+\delta(y-a)$
と取れば、この接続は同じ $E$ の定数倍に値を取る不変対数形式である。
$E$ が半単純なら両因子のトーラス作用から生じ、一般の $E$ も半単純元の極限である。
$q_\delta\Omega_\delta$ の分母を払うと $\delta^{-1}$ は消え、
$\delta\to0$ で $q_\delta\to q$, $q_\delta\Omega_\delta\to q\Omega$
となるので、大域双ベクトルの特殊化が得られる。
\end{proof}

\subsubsection{二重直線と単純直線のうち単純直線で非対数的な場合}
\begin{proposition}\label{p12:prop:doublebad}
射影三次曲線が $2L+M$（$L\ne M$ は直線）であり、単直線 $M$ の一般点で
接続が非対数的であるとする。この解は $C_4$ または
$C_{(1,2)(1,0)(0,1)}$ に属する。
\end{proposition}
\begin{proof}
$L$ を無限遠、$M$ を $x=0$ に取り、平面成分を $q=x$ とする。
接続行列の分子の二次部分は射影ベクトル場の動径成分から生じ、
行列を
\[
 \Omega=\frac{(a+bx+cy+mxy+ny^2)\,dx+
                (d+ex+fy-mx^2-nxy)\,dy}{x}
\]
と書ける。各係数は $\mathfrak{sl}_2$ の元である。
$x=0$ が非対数的なので $d+fy$ は一般点で非零であり、常にべき零である。
従って $d=DJ$, $f=FJ$ となる一定のべき零方向
$J=E_{12}$ があり $(D,F)\ne(0,0)$ である。平坦性の二位の極から
$[a+cy+ny^2,J]=J$ を得る。従って $c=CJ$, $n=NJ$ であり、
$J$ を固定する定数共役により $a=H=\operatorname{diag}(1/2,-1/2)$ とできる。

二次成分の消滅を示す。$S=H+(Cy+Ny^2)J$ とおくと
\[
 \Omega_x=S/x+b+my,\qquad
 \Omega_y=(D+Fy)J/x+e-mx-NJy.
\]
平坦性 $\partial_x\Omega_y-\partial_y\Omega_x+
[\Omega_x,\Omega_y]=0$ で $x$ の係数および定数項を比較すると
\[
 [b,m]=0,\qquad -2m-[H,m]+[b,e]=0,\qquad
 -C[J,m]-N[b,J]+[m,e]=0 .
\]
$m\ne0$ なら $b=\lambda m$ と書け、これらは
$[H+\lambda e,m]=-2m$ および $[m,e+(C-N\lambda)J]=0$ を与える。
後者から $e=-(C-N\lambda)J+\mu m$ である。
従ってある定数 $s$ について $[H+sJ,m]=-2m$ を得るが、
$\operatorname{ad}(H+sJ)$ の固有値は $0,\pm1$ なので矛盾する。
ゆえに $m=0$ である。上記の定数項と $x^{-1}$ の係数の
$y$ に関する各係数は
\begin{align*}
 [b,e]=0,\qquad N[b,J]=0,\qquad [H,e]+D[b,J]=CJ,\\
 C[J,e]+F[b,J]=3N J,\qquad N[J,e]=0
\end{align*}
となる。$N\ne0$ なら左辺の二つの括弧がともに零となり $3NJ=0$ に反する。
従って $n=0$ であり、すべての混合ベクトル場は $L$ に接する。

次に全行列が同時に上三角であることを示す。$b,e$ は可換なので
$b=\lambda Z,e=\mu Z$ と書ける。$Z$ の下三角成分が非零とすると
$[H,e]+D[b,J]=CJ$ の下三角成分から $\mu=0$ が従う
（任意の $Z$ で $[Z,J]$ の下三角成分は零）。
$[Z,J]$ の対角成分は非零なので同じ等式から $D\lambda=0$ かつ $C=0$、
ついで $C[J,e]+F[b,J]=0$ より $F\lambda=0$ である。
$(D,F)\ne(0,0)$ なので $b=e=0$ を得る。従っていずれにせよ全係数は上三角である。

従ってあるアフィン一次式 $\ell$、次数1以下の多項式を係数とする
平面1形式 $\xi$ を使い
\[
 \Omega=H\,(d\log x+d\ell)+J\,\xi/x
\]
と書ける。水平な射影座標を $(t,1)^{\mathsf T}$ とすれば、
$s=xt$ で水平方程式は $ds+s\,d\ell+\xi=0$ となり、
$d\xi+d\ell\wedge\xi=0$ である。
$d\ell=0$ なら $\xi=dP$ をみたす次数2以下の多項式 $P$ がある。
このとき $xt+P$ の斉次化が $(1,2)$ 型で、定理\ref{p12:th:threecomponents} 第1族
（$U,L$ を無限遠に選ぶ）の水平方程式を得る。

$d\ell\ne0$ なら次数1以下の $P$ で $\xi=dP+P\,d\ell$ と書ける。
実際に $d_\ell=d+d\ell\wedge$、$d\ell(Z)=1$ となる定数ベクトル場
$Z$ をとると $d_\ell\iota_Z+\iota_Zd_\ell=1+\mathcal L_Z$ である。
定数方向微分 $\mathcal L_Z$ は多項式係数上局所べき零なので、
$P=(1+\mathcal L_Z)^{-1}\iota_Z\xi$ は同じ次数の上限を満たす。
$F=xt+P$ の斉次化は $(1,1)$ 型で、水平方程式は $dF+F\,d\ell=0$ となる。
$\ell=L'/L$ とし、$U$ をファイバーの無限遠とすると
これは四因子対数族の特殊化である。すなわち
\[
 FUL(L+\epsilon L')\left(
 d\log\frac{F}{UL}+\epsilon^{-1}d\log\frac{L+\epsilon L'}{L}\right)
\]
は $\epsilon\ne0$ で \eqref{p12:eq:four} に属し、
$U=L=1$ のチャートでの極限は $dF+F\,d\ell$ である。
いずれの場合も $dt$ 成分が $x$ であるから、同じ平面成分をもつ
双ベクトルとしても一致する。
\end{proof}

\subsection{三重直線}
\begin{proposition}\label{p12:prop:triple-cubic}
平面成分の零因子が三重直線の場合、対応する点は
$C_{\mathrm{Ric}}$、$C_{\mathrm{line\text{-}conic}}$、
$(1,0)+(1,1)+(0,1)+(0,1)$ 型の対数形式の閉包の和集合に入る。
\end{proposition}
\begin{proof}
三重線を無限遠直線とし、スカラー倍で $\sigma=\partial_x\wedge\partial_y$ とする。
射影ベクトル場の二次の部分が一次式と Euler ベクトル場の積であることから
$\Omega=P\,dx+Q\,dy$ の斉次二次部分は
$L(y\,dx-x\,dy)$、$L=xH+yI$ となる。行列はすべてトレース零とする。
$P,Q$ の一次部分を $P_1,Q_1$ とし、$M=xP_1+yQ_1$ とおく。
曲率の三次部分は $-[M,L]$ である。$L\ne0$ と仮定する。

まず $H,I$ が独立なら、$\mathbb C(x,y)$ 上での非零トレース零行列の中心化環は
$\mathfrak{sl}_2$ 内で1次元なので、$M$ は $L$ の有理関数倍である。
$L$ の成分に公約因子がないので $M=(a x+b y)L$ と書ける。
平行移動 $x\mapsto x+r,\ y\mapsto y+s$ によって $M$ は
$(s x-r y)L$ だけ変化する。実際、二次部分に由来する一次部分の増分は
$(rH+sI)(y\,dx-x\,dy)+L(s\,dx-r\,dy)$ で、Euler ベクトル場に代入すればよい。
従って $M=0$ にでき、$\Omega=(L+C)(y\,dx-x\,dy)+A\,dx+D\,dy$ となる。
曲率二次項より $[xA+yD,L]=0$、したがって $xA+yD=kL$（$k$ は定数）。
一次項は $-3L+k[C,L]=0$ であり、$k\ne0$ かつ $H,I$ が
$\operatorname{ad}C$ の同じ非零固有値に属することになる。
この固有空間は高々1次元なので矛盾する。

次に $H,I$ が従属で $L\ne0$ なら、線形座標変換と定数の取り替えで $L=xH$ とできる。
$[M,L]=0$ により $M$ は二次式と $H$ の積であり、前と同じ平行移動により
$M=g y^2 H$ とできる。従って
\[
 \Omega=(xH+C)(y\,dx-x\,dy)+gHy\,dy+A\,dx+D\,dy .
\]
曲率二次項の $x^2,xy$ の係数は $-[A,H],[H,D]$ だから $A,D$ は $H$ の倍数である。
定数項 $-2C+[A,D]=0$ より $C=0$ となり、一次項 $-3xH$ が消えない。
よってすべての場合に接続の二次部分は零である。

今や $P=A+P_1$, $Q=D+Q_1$ であり、曲率二次項は $[P_1,Q_1]=0$ を与える。
$P_1$ の二つの行列係数が独立なら、先の中心化環の議論から $Q_1=kP_1$ である。
底の線形座標変換で一次部分を $P'_1\,dx$ にでき、その二つの係数はなお独立である。
平坦性より $[P'_1,D']=0$ なので $D'=0$、さらに $\partial_yP'_1=0$ となり矛盾する。
$Q_1$ についても同じである。従って一次部分が非零なら、その全行列係数は共通の
行列 $H$ の倍数である。$P_1=lH,Q_1=mH$ で一次式 $l,m$ が独立なら
一次の曲率条件より $A,D$ も $H$ の倍数となる。従属な場合は座標変換により
\[
 \Omega=(A+(bx+cy)H)dx+D\,dy,\qquad
 D=dH,\qquad d[A,H]=cH
\]
と書ける。一次部分が零の場合は平坦性が $[A,D]=0$ を意味する。

全係数が可換な場合には $\Omega=H\,dK$, $\deg K\le2$ である。
$1+\epsilon K$ をアフィン二次式、$L$ を無限遠直線として
$\mathrm{line\text{-}conic}$ 構成を適用すると、
ファイバーの多項式を $1/\epsilon$ 倍することでこの形を極限に得る。
ここで $K$ が退化していても一般の二次式からの特殊化で従う。
係数が可換でなく $c=0$ の場合は $D=0$ で、
$(A+b xH)dx$ は相異なる三根を持つ $f_\epsilon\to1$ による
$(A+b xH)dx/f_\epsilon$ の極限だから $C_{\mathrm{Ric}}$ に属する。

最後に $c\ne0$ なら $\alpha=c/d\ne0$ とおく。
$[A,H]=\alpha H$ によりファイバーの座標変換で
$H=E_{12}$, $A=\alpha\,\mathrm{diag}(1/2,-1/2)$ とできる。
水平式は $dt+\alpha t\,dx+(bx+cy)\,dx+d\,dy=0$ である。
$T=t+d y+(b/\alpha)x-b/\alpha^2$ とおくとこの式は $dT+\alpha T\,dx=0$ となる。
これは $(1+\epsilon x)dT+\alpha Tdx$ の極限である。
後者は $d\log T+(\alpha/\epsilon)d\log(1+\epsilon x)$ の分母を払った形式で、
$T$ の分子は $(1,1)$ 型、分母はファイバーと平面の無限遠切断の積である。
もう一つの直線 $1+\epsilon x=0$ と合わせて主張の四因子型となる。
\end{proof}


\subsection{既約成分の網羅性}\label{p12:sec:conclusion}
\begin{proof}[定理\ref{p12:th:complete}の証明]
平面三次式 $C$ を固定する。$C=0$ のときは命題\ref{p12:prop:zero-cubic}、
被約既約の場合は定理\ref{p12:th:irreducible-cubic} がファイバーを尽くす。
$C=LQ$ で $Q$ が滑らかな二次曲線なら定理\ref{p12:th:wholeconic}、
相異なる三直線の積なら定理\ref{p12:th:triangle}, \ref{p12:th:all-concurrent} が適用できる。
$C=L^2M$ で $L\ne M$ なら命題\ref{p12:th:double-good}, \ref{p12:prop:doublebad} が単純成分に沿う二つの場合を尽くし、
$C=L^3$ は命題\ref{p12:prop:triple-cubic} で尽きる。
これらは複素数上の全三次式を覆う。各場合、すべての点は上表の九つの既約閉集合の
いずれかに属する。他方、各閉集合が実際に成分であることは、 $C_{\mathrm{Ric}}$ 以外の八つについては
像の次元と接空間による証明から、$C_{\mathrm{Ric}}$ については
定理\ref{p12:th:concurrent-riccati} から従う。
$C_{\mathrm{Ric}}$ は前の二つの12次元成分と異なる。実際、その一般点で平面三次式は
共点三直線を定め、他では非共点三直線または直線と滑らかな二次曲線を定める。
台がそのような共点三直線またはその特殊化となる条件は閉である
（中心を動かす二変数三次式の射影パラメータ空間の像）。
他の成分同士の相異性は命題\ref{p12:prop:distinct} で示した。
前の八成分の各々には全 Poisson 概型の滑らかな点があり、従って生成点での長さは1である。
$C_{\mathrm{Ric}}$ の長さは定理\ref{p12:th:double} による。
\end{proof}

この分類は固定した $X$ 上のアフィン概型についてのものである。
異なる成分の全交叉イデアル、埋め込まれた素因子の有無、
および一般の自己同型軌道の分類は未解決である。
各成分の階数証明書、固定因子の制限、四因子型の Smith 不変因子、
共点三直線型の2次障害は同梱のコードで再現できる。

\section{三重積の Poisson 概型の分類}\label{cube:sec:classification}
\subsection{分類と成分の幾何}
$X=\PP^1_1\times\PP^1_2\times\PP^1_3$、$L=\mathcal O_X(1,1,1)$ とし、
$\mathfrak g_i=H^0(\PP^1_i,T_{\PP^1_i})\simeq\mathfrak{sl}_2$ とおく。
因子には順序を付ける。$T_X$ を分解すると
\[
 V=H^0(X,\wedge^2T_X)=\bigoplus_{i<j}\mathfrak g_i\otimes\mathfrak g_j,\qquad
 H^0(X,\wedge^3T_X)=\mathfrak g_1\otimes\mathfrak g_2\otimes\mathfrak g_3.
\]
いずれも次元27である。$V$ における $[\pi,\pi]=0$ の零点概型を
$\operatorname{Pois}(X)$ と書く。各対のテンソルをその二因子上の双ベクトル場と同一視する。
$-K_X\simeq L^2$ の同型を一つ選ぶと、
$\wedge^2 T_X\simeq\Omega_X^1\otimes L^2$ により双ベクトルを捩れた一形式で表せる。
同型を変える自由度は全体のスカラーである。

\begin{theorem}\label{cube:th:classification}
$\operatorname{Pois}(X)_{\mathrm{red}}$ の既約成分は以下の八つで尽きる。
\[
\begin{array}{c|c|c|l}
\text{成分}&\text{個数}&\text{アフィン次元}&\text{一般元あるいは構成}\\ \hline
\mathcal R&1&13& (1,1,1)\text{ 型の pencil}\\
\mathcal L_i&3&10& (1,1)\text{ 型の二つの曲線と第 }i\text{ 因子の二点による対数形式}\\
\mathcal T&1&9& \text{三つの因子のベクトル場によるトーラス型}\\
\mathcal S_{ij}&3&9& \text{第 }i,j\text{ 因子の曲面上の任意の双ベクトル}
\end{array}
\]
$\PP\mathcal R\simeq\operatorname{Gr}(2,H^0(X,L))$ であり、全八成分は $\operatorname{Pois}(X)$ の生成点で被約である。
射影化は成分の個数を保ち、各次元を一つ下げる。
全自己同型群による商スタックでは四つの既約成分になる。
\end{theorem}

ここで各族を具体的に定義する。$\mathcal R$ は $s,t\in H^0(X,L)$ に対する
$\omega=s\,dt-t\,ds$ からなる。$s/t$ は一般元の有理第一積分である。
一般の pencil の基底曲線は二つの $L$ の完全交叉である。
Bertini と随伴公式からこれは滑らかな連結な種数1曲線で、各因子への射影は次数2である。
滑らかな一般の葉の閉包 $D\in |L|$ では $-K_D=L|_D$、$K_D^2=L^3=6$ である。

$\mathcal S_{ij}=\mathfrak g_i\otimes\mathfrak g_j$ は線形部分空間である。$\mathcal T$ は $u_i\in\mathfrak g_i$ を選び
$\sum_{i<j}c_{ij}u_i\wedge u_j$ と書けるものの閉包である。
これは $(\PP^2)^3$ 上の階数3のベクトル束の像の閉包であり、
一般点から三つの直線 $\langle u_i\rangle$ を回収できる。よって既約で次元9である。
半単純な $u_i$ をとれば、座標変換で $u_i$ は $x_i\partial_{x_i}$ の定数倍になる。

$\{i,j,k\}=\{1,2,3\}$ とする。$\mathcal L_i$ は
\begin{equation}\label{cube:eq:log-family}
 \pi=D+u_i\wedge(h_j+h_k),\quad D\in\mathfrak g_j\otimes\mathfrak g_k,\quad
 [h_j+h_k,D]=0,\quad \delta(h_j)=\delta(h_k)\ne0
\end{equation}
と書ける族の閉包である。$\delta$ はベクトル場の係数の二次式としての判別式である。
$u_i$ も一般には半単純である。例えば $i=2$ とし、
三因子の座標を $x,y,z$ とすると、開稠密な族は座標変換後
\begin{equation}\label{cube:eq:normal-log}
 \lambda y\partial_y\wedge(x\partial_x+z\partial_z)
 +(a x^2+b xz+c z^2)\,\partial_z\wedge\partial_x
\end{equation}
の形をとる。ここで $\lambda\ne0$ は二つの半単純ベクトル場の規格化から生じる係数である。
実際、縮約した一形式は $\lambda yz\,dx+(a x^2+b xz+c z^2)\,dy-\lambda xy\,dz$ で、
$t=x/z$, $q(t)=a t^2+bt+c$ とすると $yz^2 q(t)$ で割って
\[
 \frac{dy}{y}+\lambda\,\frac{dt}{q(t)}
\]
を得る。一般の $q$ は相異なる二根を持ち、その部分分数分解は二つの対数微分の差である。
従って極は $y=0,\infty$ の二因子と、二つの $t=\text{定数}$ の閉包である。
後者は第 $j,k$ 因子の積における bidegree $(1,1)$ の曲線の逆像である。
逆に、一般の二つの $(1,1)$ 曲線の交点二つを $(0,0),(\infty,\infty)$ に移せば
この形になる。従って式~\eqref{cube:eq:normal-log} は表の対数型を与える。

\paragraph{先行研究}
Loray--Pereira--Touzet \cite[定理9.1]{p12-LPT} は Picard 数1の Fano 三次元多様体の
Poisson 構造の成分を分類した。ここで扱う $X$ の Picard 数は3である。
Belmans--Fatighenti--Tanturri \cite[\S2.2]{p12-BFT} は Fano 三次元多様体の多重ベクトル場の
コホモロジーを決定し、自己括弧を調べる段階を明示した。$X$ は彼らの
$\mathrm{MM}_{3.27}$ であり、対応する Fanography の項目 \cite{cube-Fano} もこの分類問題を掲げる。
pencil から Poisson 構造を作る一般構成は \cite[\S3.2.2]{p12-Pym} を参照されたい。
本稿の主眼はこの構成の像の判別式による特徴付け、および非一致軌跡も含む成分の網羅である。
八成分の分類について、文献上の新規性と優先関係は未確定である。

\subsection{三つの判別式条件と pencil}
アフィン座標 $x,y,z$ で捩れた一形式を
\[
 \omega=A(y,z)\,dx+B(x,z)\,dy+C(x,y)\,dz
\]
と書く。各係数は二変数のそれぞれについて次数が高々2である。
$J=\omega\wedge d\omega/(dx\wedge dy\wedge dz)$ は三重次数が高々 $(2,2,2)$ で、
その27係数が Jacobi 方程式である。例えば $x$ 頂点では
\[
 D_x=\operatorname{disc}_y C(x,y)-\operatorname{disc}_z B(x,z)
 \in H^0(\PP^1_x,\mathcal O(4)).
\]
他の二頂点でも巡回的に $D_y,D_z$ を定める。各 $\mathfrak g_i$ において
$\delta(p_0+p_1t+p_2t^2)=p_1^2-4p_0p_2$ は随伴不変な二次形式であるため、
この定義は座標変換と可換である。

\begin{proposition}\label{cube:prop:res}
$J=D_x=D_y=D_z=0$ の零点集合は $\mathcal R$ である。
\end{proposition}
\begin{proof}
$W=H^0(X,L)$ は三つの二次元空間のテンソル積である。
Gaussian 写像 $\gamma:\wedge^2W\to V$ を $s\wedge t\mapsto s\,dt-t\,ds$ とする。
$\mathrm{SL}_2^3$ の表現として $\wedge^2W$ は、三つの $\mathfrak g_i\otimes\mathfrak g_j$ と
一次元自明表現の直和である。実際、二つのテンソルを入れ替えて反対称になる項は、
三因子のうち一つまたは三つで外積を取る項である。$\gamma$ は三つの9次元直和因子
それぞれに非零である。例えば $s,t$ が一つの変数でのみ異なる分解テンソルを
取れば、該当する微分の係数だけが非零になる。
核は各二次元空間の交代テンソルの積 $\CC\Omega$ であり、
$\Omega\in\wedge^2 W$ は $W^*$ 上の非退化交代形式である。
特に $\PP\gamma$ の中心は階数8で、Grassmann 多様体の割線および接線の張る
直線上の点の階数は高々4である。中心からの射影は $\operatorname{Gr}(2,W)$ 上で閉埋め込みとなる。

$\gamma(s\wedge t)$ で $J=0$ は $s\,dt-t\,ds$ の可積分性から従う。
$x$ を固定すると $s,t$ は $y,z$ の同次座標の双一次形式である。
これらを $Y^{\mathsf t}FZ,Y^{\mathsf t}GZ$ とすると、$dy,dz$ の係数は
それぞれ列または行の二本組の行列式に等しい（符号は判別式に影響しない）。
この二つの二次形式の判別式は等しい。実際 $F$ が可逆な稠密開集合で
基底変換により $F=I$ に帰着できる。両辺は同じ行列式指標で変換し、
$F=I$ ではともに $(\mathrm{tr}\,G)^2-4\det G$ になる。
多項式恒等式によりすべての $F,G$ で等しい。従って $D_x=D_y=D_z=0$ である。

$V$ 上でこの射影された Grassmann 多様体の錐に消える二次方程式の空間を計算する。
$\wedge^2W$ 上で分解テンソルに消える二次式は Plücker 二次式
$q_\ell(\eta)=\ell(\eta\wedge\eta)$、$\ell\in(\wedge^4 W)^*$ の全体である。
$q_\ell$ が $V$ に降りる条件は $\ell(\Omega\wedge\wedge^2W)=0$ である。
ここで $\Omega\wedge-:\wedge^2W\to\wedge^4W$ は単射である。
実際、$L_\Omega=\Omega\wedge-$ とし、Darboux 基底での対応する縮約を
$\Lambda_\Omega$ とすると、次数二では
$\Lambda_\Omega L_\Omega-L_\Omega\Lambda_\Omega=2\operatorname{id}$ である。
$L_\Omega\beta=0$ なら
$\beta=-\frac12(\Lambda_\Omega\beta)\Omega$ となり、
$\Omega^2\ne0$ から $\beta=0$ を得る。従って二次式の空間は
$\binom84-\binom82=42$ 次元である。
一方、$J$ の係数は $\mathfrak g_1\otimes\mathfrak g_2\otimes\mathfrak g_3$ の双対、各 $D_i$ の係数は
第 $i$ 因子の二項四次式の双対を与える。これら四表現は互いに非同型かつ既約であり、
いずれの二次写像も非零である。例えば $A=z,B=1,C=0$ なら $J=1$、
$B=0,C=y$ なら $D_x=1$ である。従って $27+3\cdot5=42$ 個の係数が
この空間の基底をなす。よって三つの判別式条件と Jacobi 方程式を満たすことは、
$\gamma(\eta)=\pi$ に対して
\begin{equation}\label{cube:eq:primitive}
 \eta\wedge\eta\;\in\;\Omega\wedge\wedge^2W
\end{equation}
と同値である。

この二次条件の十分性を、重固有値および冪零部分を含めて示す。
$\eta,\Omega$ を $W^*$ 上の交代形式とみなし、
$\eta(v,w)=\Omega(Tv,w)$ で $T$ を定める。$T$ は $\Omega$ に関して自己随伴である。
Darboux 分解 $W^*=P\oplus P^*$、$\dim P=4$ をとって
$T=M\oplus M^{\mathsf t}$ と書けることを確認する。
これはシンプレクティック・スペクトル定理 \cite[命題8]{cube-SM} の場合であるが、
必要な線形代数をここに示す。異なる固有値の一般固有空間は
$\Omega$ に関して直交するので、それぞれに帰着する。
その上で $N=T-\lambda I$ の最大 Jordan 長を $m$ とし、
$N^{m-1}v\ne0$ となる $v$ と
$\Omega(N^{m-1}v,w)\ne0$ となる $w$ を選ぶ。
$N$ の自己随伴性と交代性から
$\Omega(N^iv,N^jv)=\Omega(v,N^{i+j}v)=0$ であり、$w$ にも同様である。
二つの巡回空間の間の対は
$\Omega(N^iv,N^jw)=\Omega(v,N^{i+j}w)$ で与えられ、
$i+j\ge m$ で零、$i+j=m-1$ で非零である。
従って両巡回空間の直和は非退化な $2m$ 次元の $T$-不変部分空間であり、
その直交補空間も $T$-不変である。帰納的にこの分解を繰り返し、
$v$ の巡回空間の直和を $P$、$w$ の巡回空間の直和をその補空間とする。
後者の基底を $\Omega$ による双対基底に取れば、所望の行列表現を得る。
任意の二次元部分空間 $S\subset P$ について $S\oplus S^\perp$ は
Lagrange 部分空間である。ここで $S^\perp\subset P^*$ は零化空間を表す。
式~\eqref{cube:eq:primitive} から $\eta^2$ はそこで消える。
制限された $\eta$ の交差項は $S\to P/S$, $v\mapsto Mv\bmod S$
の表現行列であり、二つの対角項は零である。従って任意の $S$ について
$S+MS\ne P$ である。

これよりある $\lambda$ について $\operatorname{rank}(M-\lambda I)\le1$ を得る。
実際、すべての固有空間の次元が2以下なら、$M$ の有理標準形の巡回直和因子は高々二つである。
不変因子が整除列をなし、最小の非定数不変因子の根が全因子に現れるためである。
次元の分割が $4$ なら巡回ベクトル $u$ に対し $S=\langle u,M^2u\rangle$ を取る。
$2+2$ なら二つの巡回ベクトルで $S$ を張る。
$3+1$ なら三次元ブロックの巡回ベクトル $u$ と一次元ブロックの固有ベクトル $v$ を用い、
$S=\langle u,Mu+v\rangle$ とする。
いずれも $S+MS=P$ で矛盾する。従って $\eta-\lambda\Omega$ の階数は高々2であり、
それは $s\wedge t$ と書ける。$\pi=\gamma(s\wedge t)$ となる。
\end{proof}
\subsection{判別式が異なる場合と既約分解}
\subsubsection{頂点の四次式が一致しない場合}
各因子のベクトル場の空間を $\mathfrak g_i=H^0(\mathbb P^1,T_{\mathbb P^1})$
とする。アフィン座標ではベクトル場をその二次以下の係数で表し、
$[p,q]=pq'-qp'$ と書く。$p=p_0+p_1t+p_2t^2$ に対して
$\delta(p)=p_1^2-4p_0p_2$ と置く。$p\ne0$ のとき
$\operatorname{ad}_p$ の $\mathfrak g_i/\langle p\rangle$ 上の固有値は
$\pm\sqrt{\delta(p)}$ である。これは $p$ の零点が二つなら
$p=ct$、重根を持つなら $p=1$ に座標を変えて確認できる。

座標を $x,y,z$ とし、体積形式との縮約から得る一形式を
\[
 \omega=A(y,z)\,dx+B(x,z)\,dy+C(x,y)\,dz
\]
と書く。すなわち二重ベクトル場の表示
$\pi=T+\partial_x\wedge(U+V)$ では
$T=A\partial_y\wedge\partial_z$, $U=C\partial_y$, $V=-B\partial_z$ である。
$x$ を固定して $B=B(x)\in\mathfrak g_z$、$C=C(x)\in\mathfrak g_y$
と扱い、$\dot B,\dot C$ は $x$ 微分とする。
Jacobi 方程式は
\begin{equation}\label{cube:eq:nonres-frob}
 (-\operatorname{ad}_C\otimes1+1\otimes\operatorname{ad}_B)A
  + C\otimes\dot B-\dot C\otimes B=0.
\end{equation}
実際、その係数は $A(C_y-B_z)+B(A_z-C_x)+C(B_x-A_y)$ である。

三辺がすべて非零で、$x$ に付随する二つの四次式
$\delta(C(x))$ と $\delta(B(x))$ が異なると仮定する。一般の $x$ で
$B,C\ne0$ である。$\mathfrak g_y/\langle C\rangle\otimes
\mathfrak g_z/\langle B\rangle$ 上で
$-\operatorname{ad}_C\otimes1+1\otimes\operatorname{ad}_B$ は可逆である。
実際、一方が冪零の場合も二つの固有値集合は交わらない。
従って \eqref{cube:eq:nonres-frob} より
\begin{equation}\label{cube:eq:nonres-inc}
 A\in C(x)\otimes\mathfrak g_z+\mathfrak g_y\otimes B(x).
\end{equation}
特に $\operatorname{rank}A\le2$ である。

\paragraph{階数二の排除}
$\operatorname{rank}A=2$ ならば \eqref{cube:eq:nonres-inc} によって
$C(x),B(x)$ はそれぞれ $A$ の左右の二次元の像に属し、射影化した
二点は固定した射影同型で対応する。実際、右側の直線 $\langle B(x)\rangle$ は
左側の二次元空間の双対における $C(x)$ の零化直線の $A$ による像である。
従って適当な基底 $u,v$ および $r,s$ を用いると
\[
 A=a(u\otimes s-v\otimes r),\quad
 C=hF,\quad B=kG,\quad F=fu+gv,\quad G=fr+gs,\quad a\ne0.
\]
ここで $f,g$ は射影写像の共通因子を除いた同次数の二項形式で、
$h,k$ は残りの因子である。以下の微分はアフィン座標で行う。
この射影写像が定数でない場合、
$W=f\dot g-g\dot f\ne0$ とおくと \eqref{cube:eq:nonres-frob} は
\[
 ah[u,v]\otimes G+akF\otimes[r,s]
 +(h\dot k-k\dot h)F\otimes G
 +hkW(u\otimes s-v\otimes r)=0
\]
となる。各二次元の像で割ると $[u,v]=l u+m v$、$[r,s]=n r+p s$
を得る。特にこれらの平面は $\mathfrak{sl}_2$ の Borel 部分代数である。
左または右に共変ベクトル $(-g,f)$ を適用すれば
\[
 kW=a(fm-gl),\qquad hW=a(gn-fp).
\]
$[F,[u,v]]=(fm-gl)[u,v]$、$[G,[r,s]]=(fp-gn)[r,s]$ であり、
Borel の導来直線上の固有値の二乗は判別式に等しい。
よって $\delta(C)=h^2(fm-gl)^2=k^2(fp-gn)^2=\delta(B)$ となり矛盾する。

射影写像が定数ならば $C=hu$, $B=kr$, $A=u\otimes s+v\otimes r$ と書ける。
\eqref{cube:eq:nonres-frob} より $[u,v]\in\langle u\rangle$,
$[r,s]\in\langle r\rangle$ を得る。$u,v$ および $r,s$ はそれぞれ独立で
あるから、$u,r$ は各 Borel の導来直線に属し、いずれも冪零である。
再び $\delta(C)=\delta(B)=0$ となる。

\paragraph{階数一と残りの辺}
従って $A=u\otimes r$ と書ける。\eqref{cube:eq:nonres-inc} より
$C(x)\in\langle u\rangle$ または $B(x)\in\langle r\rangle$
が一般の $x$ で成り立つ。これらは $x$ に関する閉条件なので、
いずれかが恒等的に成り立つ。対称性により $C=h\otimes u$ としてよい。
このとき \eqref{cube:eq:nonres-frob} は
\begin{equation}\label{cube:eq:axial-centralizer}
 (\operatorname{ad}_h\otimes1-1\otimes\operatorname{ad}_r)B=0
 \quad\text{in }\mathfrak g_x\otimes\mathfrak g_z
\end{equation}
に帰着する。ここで $h,r\ne0$ である。
$\delta(h)\ne\delta(r)$ ならばこの核は $\langle h\otimes r\rangle$ である。
実際、双方半単純なら共通固有値は $0$ のみで、一方が冪零なら
他方の零固有空間上でその冪零作用の核を取ればよい。この場合、三辺は
各頂点に一つずつ選んだ直線から作られる。

$\delta(h)=\delta(r)$ の場合は \eqref{cube:eq:axial-centralizer} の核が
三次元である。非零の判別式の場合は三つの固有値から各一次元を得る。
両方が非零冪零なら両方の随伴作用が大きさ三の単一 Jordan ブロックであり、
その間の intertwiner の空間は三次元である（不変双線形形式による双対化で
符号が変わっても Jordan 型は同じである）。
$\delta(h)=\delta(r)$ は六次元アフィン空間内の非退化二次錐であり、
$h,r\ne0$ の開集合上で核は階数三のベクトル束になる。
$\langle u\rangle\in\mathbb P(\mathfrak g_y)$ を動かすと、
得られる族の閉包は既約で次元は $2+5+3=10$ である。実際、
$u$ の代表を局所的に固定すると $A,C$ から $h,r$ が一意に定まる。
$\delta(h)=\delta(r)\ne0$ 上で一般の $B$ は階数三である。
従って三通りの共通頂点は一般点の階数から区別できる。

最後に、辺がちょうど二本非零ならば共有する頂点における二つの像の任意の
元が互いに可換である。$\mathfrak{sl}_2$ の任意の非零元の中心化代数は
一次元なので、両方の像は同じ直線になる。この場合も各頂点における
直線の選択から得られる。辺が一本だけ非零ならその辺のテンソルは任意である。
\subsubsection{成分の極大性と一般点での被約性}
命題~\ref{cube:prop:res} と前節の議論から八つの既約閉集合で全点が覆われる。
三つの $\mathcal S_{ij}$ は支持する辺で区別され、三つの $\mathcal L_i$ は
一般元で階数3をもつ辺で区別される。$\mathcal T$ と各 $\mathcal S_{ij}$ は同次元だが、
$\mathcal T$ の一般元では全ての辺が非零である。
各閉集合が実際に成分であることと一般点の被約性を同時に証明するため、
各型の一点で接空間を計算する。

以下、$w=(A,B,C)$ を用い、三つの辺の変数を順に $(y,z),(z,x),(x,y)$ とする。
三重添字 $(i,a,b)$ は第 $i$ 辺のこの順序の単項式を表す。
具体的な四点を次で選ぶ。
\begin{itemize}
\item $\mathcal R$ では $s,t$ の係数を $(x,y,z)$ の指数の辞書式昇順
$000,001,010,011,100,101,110,111$ でそれぞれ
$(2,-1,4,2,-3,5,1,3)$ および $(1,3,-3,4,1,2,-5,-1)$ とする。
\item $\mathcal S_{23}$ では $A$ の係数を $(y,z)$ の指数の辞書式昇順で
$(0,1,3,-1,1,5,5,0,2)$ とし、$B=C=0$ とする。
\item $\mathcal T$ では $w=(yz,3zx,5xy)$ とする。
\item $\mathcal L_2$ では $w=(2yz,3x^2+5xz+7z^2,-2xy)$ とする。
\end{itemize}
接空間の余次元は、$v$ を27個の単項式方向に動かしたときの
\[
 v\longmapsto v\cdot\mathrm{curl}\,w+w\cdot\mathrm{curl}\,v
\]
の階数である。像は各変数について次数2以下なので、
$\{0,1,2\}^3$ の27点で評価すれば係数を可逆な行変換で置き換えたことになる。
この $27\times27$ 行列を $M_\pi$ とする。
行は点 $(x,y,z)\in\{0,1,2\}^3$、列は $(i,a,b)$ の辞書式昇順に並べ、
いずれも1から番号を付ける。以下は非零小行列式の証明書である。
$I,J$ はそれぞれ行と列の順序付き添字列であり、最後の欄は
$\det M_\pi[I,J]$ の $101$ による剰余である。
\[
\begin{array}{c|l|r}
 & (I,J)&\det\pmod {101}\\\hline
\mathcal R&
\begin{aligned}
 I&=(1,2,3,4,5,6,10,11,12,7,8,13,14,19),\\
 J&=(1,2,\ldots,14)
\end{aligned}&96\\[4pt]
\mathcal S_{23}&
\begin{aligned}
 I&=(1,10,19,4,11,20,8,13,22,2,5,12,7,14,15,21,23,24),\\
 J&=(10,11,\ldots,27)
\end{aligned}&57\\[4pt]
\mathcal T&
\begin{aligned}
 I&=(10,11,12,13,14,16,15,17,9,22,4,18,5,23,6,3,20,24),\\
 J&=(1,2,3,4,6,7,8,9,10,12,13,15,16,18,19,21,25,27)
\end{aligned}&74\\[4pt]
\mathcal L_2&
\begin{aligned}
 I&=(2,3,11,5,6,14,8,9,15,13,4,17,23,16,12,18,19),\\
 J&=(1,2,3,4,5,6,7,8,9,10,11,15,18,19,21,25,27)
\end{aligned}&52
\end{array}
\]
したがって階数は順に $14,18,18,17$ 以上である。
各点を含む既約族の次元が既に $13,9,9,10$ と分かっているので、
接空間の次元による逆向きの評価と合わせて等号を得る。
付属の \texttt{check\_jacobi.py} も同じ行列を整数で組み立て、
有理数による行消去で次の値を再現する。

\[
\begin{array}{c|rrrr}
 &\mathcal R&\mathcal S_{23}&\mathcal T&\mathcal L_2\\\hline
 \operatorname{rank} dJ&14&18&18&17\\
 \dim T_\pi\operatorname{Pois}(X)&13&9&9&10\\
 \dim G\pi&9&6&6&7
\end{array}
\]
である。最下段は $G=(\mathrm{PGL}_2)^3$ の九つの基本ベクトル場の
Lie 微分の階数を示す。同じコードがこれも計算する。
この $9\times27$ 行列を $N_\pi$ とし、行を
$(t,d)\in\{1,2,3\}\times\{0,1,2\}$ に対応する
$\mathcal L_{x_t^d\partial_{x_t}}\pi$ の辞書式昇順、列を上と同じ順序とすると、
軌道次元の下界も次の非零小行列式で確認できる。
\[
\begin{array}{c|l|l|r}
 &I&J&\det N_\pi[I,J]\pmod {101}\\\hline
\mathcal R&(4,5,7,6,8,9,2,3,1)&(1,2,3,4,5,6,10,11,13)&52\\
\mathcal S_{23}&(4,5,7,8,6,9)&(1,2,3,4,5,6)&82\\
\mathcal T&(4,7,9,6,1,3)&(2,4,6,8,13,15)&92\\
\mathcal L_2&(4,7,9,6,1,2,3)&(2,4,6,8,11,12,15)&75
\end{array}
\]
軌道次元の上界は後の商の証明で説明する。
上の四点で接空間の次元はその点を含む既約閉集合の次元に等しい。
局所次元は両者の間にあるので、各点の局所環は正則で、対応する閉集合は
その点を通る唯一の成分である。因子の置換により残りの四成分にも従う。
これで定理~\ref{cube:th:classification} の固定 $X$ に関する主張が示された。
Jacobi イデアル自体が根基的かどうかは、この一般点での計算からは決まらない。

\subsection{自己同型群による商}
$G=\operatorname{Aut}^0(X)=(\mathrm{PGL}_2)^3$、$\operatorname{Aut}(X)=G\rtimes\mathfrak S_3$ である。
これは nef 錐の三本の端射線に対応する射影が自己同型により置換されることから従う。
連結群 $G$ は各成分を保ち、$\mathfrak S_3$ はそれぞれ三つの $\mathcal L_i$、
三つの $\mathcal S_{ij}$ を推移的に置換する。従って商スタックの成分は
$G$ で八つ、$\operatorname{Aut}(X)$ で四つである。

\begin{proposition}\label{cube:prop:git}
アフィン GIT 商および $\mathcal O(1)$ による射影 GIT 商でもこの個数は保たれる。
アフィン商の各成分の次元は $\mathcal R$ で4、他の型で3であり、
射影商ではそれぞれ3と2である。特に、因子の順序を保つ射影商では
\[
 (\PP\mathcal R)\mathbin{/\mkern-6mu/}G\simeq\PP(1,2,2,3).
\]
\end{proposition}
\begin{proof}
$\mathfrak g_i$ をコンパクト実形 $\mathfrak{su}(2)$ の複素化とみなし、$V$ にその誘導実構造を入れる。
各成分で実点は Zariski 稠密である。$\mathcal S_{ij}$ と $\mathcal T$ では実直線および実係数を用いる。
$\mathcal L_i$ ではコンパクト実形で長さが等しい二つの非零元を選べるので、
非退化二次錐 $\delta(h_j)=\delta(h_k)$ は滑らかな実点をもつ。
その実点は稠密で、実構造を保つ階数3の核束でも同様である。
$\mathcal R$ では $W$ の三因子に標準の四元数構造を入れると、
そのテンソル積 $j$ は $j^2=-1$ を満たす。$\operatorname{Gr}(2,W)$ の $j$-不変な
複素二平面は四元数直線の空間 $\mathbb H\PP^3$ となる。これは非特異な実形なので
Zariski 稠密である。$s\wedge js$ は $\wedge^2j$ で不変であり、その実数倍が
錐で稠密になる。$X$ の各因子の反足実構造がこの構造とベクトル場のコンパクト
実構造を誘導するため、$L^2\simeq -K_X$ の同型を実構造と両立するように
スカラー調整すれば $\gamma$ は実となる。

コンパクト群に関して不変な実正定値内積を使うと、実ベクトルにおいて
モーメント写像は零である。従って Kempf--Ness の定理 \cite{cube-KN} により各実点の
$G$ 軌道は閉じている。各成分で他成分の和を避ける実点を取る。
この点の商像が他成分の像にも属するなら、ファイバーの唯一の閉軌道が他成分にも
含まれ、矛盾する。不変閉部分集合の像は閉じているので、商像はそれぞれ別の既約成分となる。
内積を複素双線形に延長した不変二次式は非零実点で正であり、これらの点は
射影 GIT でも半安定である。有限群 $\mathfrak S_3$ の商で成分はその軌道ごとにまとまる。

各型の最大軌道次元は上表の値である。上限は $\mathcal R$ では $\dim G=9$、
$\mathcal S_{ij}$ では残りの因子の自己同型、$\mathcal T$ では三つの中心化群から従う。
$\mathcal L_i$ の一般元には $u_i$ の中心化トーラスと $h_j+h_k$ の対角トーラスが作用する。
表の計算はこれらの上限に達している。
最大軌道次元の開集合で選んだ実点の閉軌道は、その成分内で商のファイバーに一致する。
真に大きな軌道の境界に入るなら軌道次元の最大性に反するからである。
ファイバー次元定理によりアフィン商の次元は主張通りである。射影商の次元は
正次数のスカラー倍を除くため一つ小さくなる。

最後に $\PP\mathcal R$ の偏極は $\operatorname{Gr}(2,W)$ の Plücker 偏極である。
射影像の斉次座標環はこの偏極の全切断環の部分環であり、
一次部分の次元はそれぞれ27と28である。十分高い次数では両者は一致するので、
$G$ 不変部分環の $\operatorname{Proj}$ も一致する。
二つの切断の基底変換 $\mathrm{SL}_2$ の不変式は二行行列の最大 minors で生成され、
その環は Grassmann 多様体の Plücker 切断環である。従ってさらに $G$ 不変式を取ると、
$(\CC^2)^{\otimes4}$ の $\mathrm{SL}_2^4$ 不変式環が得られ、テンソル次数は
Plücker 次数の二倍になる。Luque--Thibon \cite[\S\S II--III]{cube-LT} による
自由生成元のテンソル次数は $2,4,4,6$ なので、この環の $\operatorname{Proj}$ は
$\PP(1,2,2,3)$ である。
\end{proof}

\subsection{六次 del Pezzo 曲面上の括弧の延長}
\begin{proposition}\label{cube:prop:delpezzo}
滑らかな $S=(s=0)\in|L|$ と $\sigma\in H^0(S,-K_S)$ に対し、
$S$ を Poisson 部分多様体にして $\pi|_S=\sigma$ を満たす
$\pi\in\operatorname{Pois}(X)$ はただ一つ存在する。この延長は $\mathcal R$ に属し、
$S,\sigma$ を動かした像の閉包は $\mathcal R$ である。
\end{proposition}
\begin{proof}
随伴公式から $-K_S\simeq L|_S$ である。
$0\to\mathcal O_X\xrightarrow{s}L\to L|_S\to0$ と $H^1(X,\mathcal O_X)=0$ により、
対応する切断を $t\in H^0(X,L)$ に持ち上げる。$s\,dt-t\,ds$ は
$S$ 上で $-t\,ds$ である。局所座標の一つに $s$ を取ると、
これが $S$ に接する双ベクトルで、その誘導する切断は随伴同型の符号を選べば
$\sigma$ である。持ち上げの不定性 $t\mapsto t+cs$ は消える。
二つの延長の差は $H^0(X,\wedge^2 T_X\otimes L^{-1})$ に属するが、
各直和因子は次数 $(1,1,-1)$ の置換であるから、この空間は零である。
一般の pencil は滑らかなメンバーを持つので閉包についても従う。
\end{proof}



\part{可換模型の Poisson 双有理幾何}
\section{古典的 Poisson 双有理写像の blow-up 表示}
\label{sec:classical-factorizations}

標数零の滑らかな完備多様体間の双有理写像には、滑らかな中心の
blow-up と blow-down による弱分解がある
\cite[定理0.1.1]{AKMW}。
Poisson 構造を各中間模型上で正則に保つには、各中心に追加条件が必要になる。
本節では使用する中心を具体的に指定し、共通の滑らかな Poisson 模型を構成する。

\subsection{曲線中心と点中心の持ち上げ条件}

\begin{proposition}\label{prop:poisson-centres}
滑らかな三次元 Poisson 多様体 $(X,\pi)$ と滑らかな中心 $Z$ を考える。
\begin{enumerate}
\item $Z$ が曲線ならば、$\pi$ が $\operatorname{Bl}_Z X$ に正則に持ち上がる
必要十分条件は $\pi|_Z=0$ である。
\item $Z=\{p\}$ が点ならば、必要十分条件は $\pi(p)=0$ と、$p$ を原点とする
局所座標での一次項が
\[
 \pi_1=\bigl(x\partial_x+y\partial_y+z\partial_z\bigr)\wedge V
\]
と書けることである。ここで $V$ は定係数ベクトル場である。
\end{enumerate}
いずれの場合も持ち上げは一意的で、blow-down 射は Poisson 射となる。
\end{proposition}

\begin{proof}
曲線の場合、$Z=\{v=w=0\}$ として
\[
 \pi=A\partial_v\wedge\partial_w+
 B\partial_w\wedge\partial_x+C\partial_x\wedge\partial_v
\]
と書く。$w=v\tau$ のチャートでは
\[
 \widetilde\pi=\frac{A}{v}\partial_v\wedge\partial_\tau+
 C\partial_x\wedge\partial_v-
 \frac{B+\tau C}{v}\partial_x\wedge\partial_\tau.
\]
例外因子上の任意の $\tau$ で正則となる条件は
$A|_Z=B|_Z=C|_Z=0$ である。他方のチャートも同じ条件を与える。

点の場合、$A=\{y,z\}$、$B=\{z,x\}$、$C=\{x,y\}$ と置く。
$u=y/x$、$v=z/x$ のチャートでは
\[
 \{x,u\}=C/x,\qquad \{x,v\}=-B/x,\qquad
 \{u,v\}=\frac{xA+yB+zC}{x^3}.
\]
三つのチャートの正則性から定数項は全て零となる。
$A_1,B_1,C_1$ を一次項とすると、残る条件は
\[
 xA_1+yB_1+zC_1=0.
\]
一次式のこの関係は、$(A_1,B_1,C_1)$ の係数行列が歪対称であることと同値であり、
ちょうど命題の $\pi_1$ を与える。二次以上の項については、上の最後の分子は
三次以上なので、各 blow-up チャートへの代入後に分母が消える。
これで十分性も従う。
一意性と Jacobi 恒等式は中心の補集合の稠密性による。
\end{proof}

この条件は Polishchuk \cite[定理8.2、命題8.3]{Polishchuk1997} の
退化した余法 Lie 代数という条件の三次元での表示である。
例えば $\pi=x\partial_y\wedge\partial_z$ は原点で零であるが、
$\{y/x,z/x\}=1/x$ なので原点の blow-up には持ち上がらない。
一方、以下の二次曲面 pencil の点中心では一次項の条件も満たされる。

\subsection{二次曲面 pencil と垂直型の Poisson 双有理同値}
\label{sec:birational-rat22-vertical}

この節では基礎体を $\mathbf C$ とし、双有理写像が誘導する関数体の同型が
Poisson 括弧を保存するとき、二つの Poisson 多様体は Poisson 双有理同値であるという。
ここで用いる双有理同値は、滑らかな共通モデルからの正則 Poisson 射によって実現される。

\begin{theorem}
\label{thm:rat22-vertical}
$\mathbf P^3$ 上の二次式 $q_0,q_1$ の共通零点
$E=\{q_0=q_1=0\}$ が滑らかな完全交叉曲線であるとし、
\[
 \alpha=q_0\,dq_1-q_1\,dq_0
 \in H^0\bigl(\mathbf P^3,\Omega^1_{\mathbf P^3}(4)\bigr)
\]
に対応する Poisson 双ベクトルを $\pi$ とする。
任意の $p\in E$ に対し、滑らかな平面三次曲線
$C\subset\mathbf P^2$ と、滑らかな射影三次元多様体 $Z$ 上の正則 Poisson 構造
$\widetilde\pi$ が存在し、
\[
 (\mathbf P^3,\pi)
 \xleftarrow{\ \beta\ }
 (Z,\widetilde\pi)
 \xrightarrow{\ \psi\ }
 (\mathbf P^1\times\mathbf P^2,\pi_C)
\]
の両射は双有理かつ正則な Poisson 射となる。
ここで $\pi_C$ は $\mathbf P^2$ 上の三次式 $C$ に対応する Poisson 構造の
$\mathbf P^1\times\mathbf P^2$ への垂直な延長である。
また $C$ は $E$ と同型であり、具体的に
\[
 Z=\operatorname{Bl}_{E'}\bigl(\operatorname{Bl}_p\mathbf P^3\bigr)
 \cong\operatorname{Bl}_{\Gamma}(\mathbf P^1\times\mathbf P^2),
\]
と取れる。ここで $E'$ は $E$ の狭義変換であり、
$\Gamma$ は $(1,1)$ 型および $(1,2)$ 型因子の滑らかな完全交叉である。
第二射影は $\Gamma\cong C$ を与え、第一射影 $\Gamma\to\mathbf P^1$ の次数は $2$ である。
逆に、任意の非零三次式で定義される滑らかな平面三次曲線と、その垂直型 Poisson 構造は、
この構成から得られる。
\end{theorem}

\begin{proof}
$p=[1:0:0:0]$ とする斉次座標 $[x_0:x_1:x_2:x_3]$ を取り、
$y=[x_1:x_2:x_3]$ と書く。
$q_i(p)=0$ より、一意的に
\[
 q_i=x_0\ell_i(y)+r_i(y)\qquad(i=0,1)
\]
と書ける。ここで $\ell_i$ は一次式、$r_i$ は二次式である。
$E$ の $p$ における滑らかさから $\ell_0,\ell_1$ は一次独立である。
\[
 C(y)=\ell_1(y)r_0(y)-\ell_0(y)r_1(y)
\]
と置き、次の双有理写像を考える。
\[
 \varphi\colon\mathbf P^3\dashrightarrow\mathbf P^1\times\mathbf P^2,
 \qquad x\longmapsto\bigl([q_0(x):q_1(x)],y\bigr).
 \tag{\ref{sec:birational-rat22-vertical}.1}
\]
逆写像を斉次座標で表すため、$\mathbf P^1$ の座標を $[a:b]$ とし、
\[
 D=a\ell_1-b\ell_0,\qquad N=br_0-ar_1
\]
と置く。このとき逆写像は
\[
 \chi\colon\mathbf P^1\times\mathbf P^2\dashrightarrow\mathbf P^3,
 \qquad ([a:b],[y_1:y_2:y_3])
 \longmapsto[N:Dy_1:Dy_2:Dy_3]
\]
である。実際、$q_0(N,Dy)=aDC$, $q_1(N,Dy)=bDC$ が成り立つ。
$a\ne0$ の開集合で $u=b/a$ とすれば、これは
$x_0=(u r_0-r_1)/(\ell_1-u\ell_0)$ と同値である。

まず Poisson 双ベクトルの変換を計算する。
$x_3=1$ とし、$v=x_1/x_3$, $w=x_2/x_3$ を使う。
このアフィン座標で $\iota_\pi(dx_0\wedge dv\wedge dw)=\alpha$ とすると、
\[
 \alpha=q_0^2\,du,
 \qquad
 \frac{\partial u}{\partial x_0}
 =\frac{\ell_1q_0-\ell_0q_1}{q_0^2}
 =\frac{C(v,w,1)}{q_0^2}.
\]
$u$ は Casimir であり、$\{v,w\}_\pi$ は $\alpha$ の $dx_0$ 係数だから、
\[
 \varphi_*\pi
 =C(v,w,1)\,\partial_v\wedge\partial_w.
 \tag{\ref{sec:birational-rat22-vertical}.2}
\]
これは $\mathbf P^1\times\mathbf P^2$ 全体で正則である。
この等式は双ベクトルそのものの変換公式である。

次に $C=0$ の滑らかさを確認する。
$E$ は楕円曲線であり、$H=\mathcal O_E(1)$ は次数 $4$ の直線束である。
$p$ からの射影が $E$ 上に誘導する写像は完全線形系 $|H(-p)|$ で与えられる。
次数 $3$ の直線束は楕円曲線上で非常に豊富なので、これは $E$ を滑らかな平面三次曲線に埋め込む。
$q_0=q_1=0$ から $x_0$ を消去すると像の方程式は $C=0$ となる。
したがって、この三次式は滑らかな曲線を定義する。

$B=\operatorname{Bl}_p\mathbf P^3$ と置く。
$p$ からの射影は $B\to\mathbf P^2$ に延長する。
$p$ の近傍で $s=q_0$, $t=q_1$ を適切な局所自明化で表し、
$E=\{s=t=0\}$ となるエタール局所座標 $(s,t,e)$ を取る。
$\pi$ は、正則可逆関数 $h$ を用いて
\[
 \pi=h(s,t,e)\,(s\partial_s+t\partial_t)\wedge\partial_e
 \tag{\ref{sec:birational-rat22-vertical}.3}
\]
と書ける。符号は $h$ に含めている。
点の blowup の $s$ チャート
$(s,t,e)=(a,ab,ac)$ では
\[
 (s\partial_s+t\partial_t)\wedge\partial_e
 =\partial_a\wedge\partial_c,
\]
$e$ チャート $(s,t,e)=(ea,eb,e)$ では
\[
 (s\partial_s+t\partial_t)\wedge\partial_e
 =(a\partial_a+b\partial_b)\wedge\partial_e.
\]
$t$ チャートも同様なので、$\pi$ は $B$ 上の正則 Poisson 構造に持ち上がる。
$E'$ は $e$ チャートの $a=b=0$ であり、持ち上げた双ベクトルは $E'$ 上で消える。
さらに $E'$ を blowup し、例えば $b=ac$ と置くと
\[
 (a\partial_a+b\partial_b)\wedge\partial_e
 =a\partial_a\wedge\partial_e.
\]
よって $Z=\operatorname{Bl}_{E'}B$ にも正則に持ち上がる。
各段階の Jacobi 恒等式は稠密開集合で成立するため全体で成立する。

$B$ 上の pencil の基底イデアルは、点の blowup の例外因子を共通因子として除くと
$E'$ のイデアルになる。実際、$s$ チャートでは残余イデアルは $(1,b)$ であり、
$e$ チャートでは $(a,b)$ である。
したがって $Z$ 上で $[q_0:q_1]$ も射に延長し、
射影 $Z\to\mathbf P^2$ と合わせて双有理射
$\psi\colon Z\to\mathbf P^1\times\mathbf P^2$ を得る。
式~(\ref{sec:birational-rat22-vertical}.2) より、$\psi$ は稠密開集合で Poisson 射である。
双方の構造と射が正則であるから、$\psi$ は全体で Poisson 射となる。
$\beta$ についても同様である。

$\psi$ を blowup として同定する。
$\chi$ の四つの座標は $\mathcal O(1,2)$ の切断であり、その基底イデアルは
\[
 (N,Dy_1,Dy_2,Dy_3)=(N,D)
\]
である。ここで右辺との等式はイデアル層の等式であり、各点で少なくとも一つの $y_j$ が
可逆であることによる。
したがって $\chi$ のグラフの閉包は
$\operatorname{Bl}_{\Gamma}(\mathbf P^1\times\mathbf P^2)$ となる。ただし
\[
 \Gamma=\{D=N=0\}.
\]
一方、$B$ は $p$ からの射影のグラフの閉包であり、
$Z=\operatorname{Bl}_{E'}B$ はその上の残余 pencil のグラフの閉包である。
従って $Z$ は $\varphi$ のグラフの閉包そのものである。
互いに逆な双有理写像のグラフの閉包は因子の順序を交換すると一致するため、
\[
 Z\cong\operatorname{Bl}_{\Gamma}(\mathbf P^1\times\mathbf P^2)
\]
を得る。この同型のもとで $\psi$ は blowup の射である。

$\Gamma$ の滑らかさも直接確認できる。
$p_0=\{\ell_0=\ell_1=0\}$ と置くと、$C(p_0)=0$ である。
$y\ne p_0$ では $D=0$ が $[a:b]=[\ell_0(y):\ell_1(y)]$ を定め、
$N=0$ は $C(y)=0$ と同値になる。
$C$ の滑らかさから $(r_0(p_0),r_1(p_0))\ne(0,0)$ である。
例えば $r_0(p_0)\ne0$ なら、$a=1$ とするチャートで
$N=0$ から $b=r_1/r_0$ を得て、残りの方程式は
$D=C/r_0=0$ となる。
従って第二射影はスキームの同型 $\Gamma\cong C$ を与える。
特に $\Gamma$ は滑らかな楕円曲線であり、表示通り $(1,1),(1,2)$ 型の完全交叉である。
第一射影は $C$ 上の線形系 $|\mathcal O_C(1)(-p_0)|$ で与えられる。
この直線束の次数は $2$ で、完全線形系は基点を持たないので、
$\Gamma\to\mathbf P^1$ の次数は $2$ である。

逆に、滑らかな平面三次式 $C(y)$ と点 $p_0\in\{C=0\}$ を取る。
$p_0$ のイデアルを生成する独立な一次式 $\ell_0,\ell_1$ を選ぶと、
$C\in(\ell_0,\ell_1)$ より二次式 $r_0,r_1$ を用いて
\[
 C=\ell_1r_0-\ell_0r_1
\]
と書ける。$q_i=x_0\ell_i+r_i$ と置く。
$p=[1:0:0:0]$ では二つの微分 $dq_0,dq_1$ が一次独立である。
$y\ne p_0$ では $q_0=q_1=0$ は $C(y)=0$ と $x_0$ を一意に定める一次方程式に同値であり、
この部分は $C\setminus\{p_0\}$ と同型である。
$C$ の $p_0$ における滑らかさから
$(r_0(p_0),r_1(p_0))\ne(0,0)$ であり、$y=p_0$ の上に $p$ 以外の点はない。
以上により $\{q_0=q_1=0\}$ は滑らかな完全交叉曲線となる。
先の変換公式は、指定した三次式 $C$ の係数を含めてそのまま与える。
\end{proof}

\begin{remark}
二次曲面 pencil 成分 $R(2,2)$ は
$\operatorname{Gr}(2,H^0(\mathbf P^3,\mathcal O(2)))$ に対応し、
$\dim H^0(\mathbf P^3,\mathcal O(2))=10$ なので、アフィン Poisson 空間内での次元は
$2(10-2)+1=17$ である。
垂直成分 $V_{\mathrm{vert}}$ は
$H^0(\mathbf P^2,\mathcal O(3))$ であり、アフィン次元は $10$ である。
定理~\ref{thm:rat22-vertical} は両成分の滑らかな楕円曲線を伴う開集合の間の
Poisson 双有理同値を構成する。写像の選択には、基底楕円曲線上の点の選択が含まれる。
\end{remark}

\subsection{楕円曲線に沿う基本変換と三重積のペンシル}
$S=\mathbb P^1_y\times\mathbb P^1_z$, $X=\mathbb P^1_x\times S$ とおく。
$s,t\in H^0(X,\mathcal O(1,1,1))$ を
\[
 s=xa(y,z)+b(y,z),\qquad t=xc(y,z)+d(y,z),\qquad
 a,b,c,d\in H^0(S,\mathcal O(1,1))
\]
と書き、$\alpha=s\,dt-t\,ds$ に対応する Poisson 構造を $\pi$ とする。
\[
 D=bc-ad\in H^0(S,\mathcal O(2,2)),\qquad u=t/s
\]
とおけば、関数体上で $u$ は Casimir であり
\[
 \{u,y\}=\{u,z\}=0,\qquad
 \{y,z\}=s\partial_xt-t\partial_xs=D(y,z).
\]
従って有理写像
\[
 \Phi\colon X\dashrightarrow X'=\mathbb P^1_u\times S,
 \qquad (x,y,z)\longmapsto(t/s,y,z)
\]
は $\pi$ を実際の双ベクトル $D(y,z)\partial_y\wedge\partial_z$ に送る。
変換式は双ベクトル自体の等式である。
$D\ne0$ なら逆写像は
\[
 x=\frac{d-ub}{ua-c}
\]
なので $\Phi$ は双有理である。

一般の $s,t$ に対して $E=\{s=t=0\}\subset X$ は滑らかな楕円曲線である。
また $\Gamma=\{D=0\}\subset S$ は滑らかな $(2,2)$ 曲線である。
$a,b,c,d$ は同時には零にならず、$\Gamma$ 上で行列
$\left(\begin{smallmatrix}a&b\\c&d\end{smallmatrix}\right)$ は階数一なので、
射影 $E\to\Gamma$ は同型になる。
$[x_0:x_1]$, $[u_0:u_1]$ を $x=x_1/x_0$, $u=u_1/u_0$ に対応する座標とすれば、
$\Phi$ の閉グラフは
\[
 B=\{x_1(u_0c-u_1a)+x_0(u_0d-u_1b)=0\}
 \subset\mathbb P^1_x\times\mathbb P^1_u\times S.
\]
これは同時に
\[
 B\cong\operatorname{Bl}_{E}X
 \cong\operatorname{Bl}_{E'}X',\qquad
 E'=\{u_0c-u_1a=u_0d-u_1b=0\}
\]
である。実際、どちらの射影に関しても二つの切断の比のグラフであり、
それらの共通零点は滑らかな余次元二の完全交叉である。
同様に $E'\to\Gamma$ は同型で、$B$ は滑らかである。

命題\ref{prop:poisson-centres}の曲線中心の条件を適用する。

$E$ 上では $s=t=0$ より $\pi=0$ であり、$E'$ 上では $D=0$ である。
従って $B$ 上には正則 Poisson 構造があり、上の二つの射影はいずれも
Poisson 射である。二つの持ち上げは稠密開集合で等しいので一致する。
一般の $u$ に対する葉の閉包は、元の側では $(1,1,1)$ 型の
次数六の del Pezzo 曲面、変換先では $\mathbb P^1\times\mathbb P^1$ である。
$E'\to\mathbb P^1_u$ の次数は二であり、共通モデルの上では後者を二点で
ブローアップして前者を得る。

この構成は一般の $(2,2)$ 型の Poisson 双ベクトルを与える。
実際、$V=H^0(S,\mathcal O(1,1))$, $W=H^0(S,\mathcal O(2,2))$ とおくと、
行列式写像 $V^4\to W$, $(a,b,c,d)\mapsto bc-ad$ の微分は
\[
 (\dot a,\dot b,\dot c,\dot d)\longmapsto
 c\dot b+b\dot c-d\dot a-a\dot d.
\]
$a,b,c,d$ が $V$ を張る点では、その像は $V\cdot V=W$ である。
従って行列式写像は優越的であり、$D$ が滑らかである開条件と
$a,b,c,d$ が同時には消えない開条件を満たす一般点で、上の構成が適用できる。

\subsection{平面と二重積の曲面リンク}
\label{subsec:classical-surface-link}
相異なる $p,q\in\mathbb P^2$ と、$r\in\mathbb P^1\times\mathbb P^1$ に対して
標準的な曲面の同型
\[
 S_0=\operatorname{Bl}_{p,q}\mathbb P^2
 \cong\operatorname{Bl}_{r}(\mathbb P^1\times\mathbb P^1)
\]
がある。左側では $p,q$ を結ぶ直線の狭義変換を右側の $r$ へ収縮する。
座標を
\[
 p=[0:1:0],\quad q=[0:0:1],\quad
 [T:Y:Z]=[1:y:z],\quad r=(\infty,\infty)
\]
と取れば、有理写像は共通アフィン平面上で恒等写像である。
このリンクに $\mathbb P^1_x$ を掛けた三次元のリンクを考える。

三重積上の双ベクトルを
\[
 \rho=A(y,z)\partial_y\wedge\partial_z+
       B(x,z)\partial_z\wedge\partial_x+
       C(x,y)\partial_x\wedge\partial_y
\]
と書く。$A=\sum_{0\le i,j\le2}a_{ij}y^iz^j$ とし、
$B=b_0(x)+b_1(x)z+b_2(x)z^2$、
$C=c_0(x)+c_1(x)y+c_2(x)y^2$ とする。
この同じ有理双ベクトルが $\mathbb P^1\times\mathbb P^2$ に正則に延びる条件は
\begin{equation}\label{eq:surface-link-conditions}
 a_{22}=0,\qquad b_2(x)\equiv0,\qquad c_2(x)\equiv0.
\end{equation}
これは $\rho|_{\mathbb P^1\times\{r\}}=0$ と同値である。
直接には、$a_{22}=0$ は $A$ の全次数が三以下となる条件であり、
残り二条件は混合ベクトル場が平面の射影ベクトル場となる条件である。
一般の射影ベクトル場の二次部分は
$\ell(y,z)(y\partial_y+z\partial_z)$ なので、
各係数が一方の変数だけに依存する本表示では二次部分は消えなければならない。

逆向きには、$\mathbb P^1_x\times\mathbb P^2$ 上の双ベクトルを
\[
 \rho=F\,\partial_y\wedge\partial_z+
       \sum_{i=0}^2 x^i\partial_x\wedge V_i
\]
と書く。ここで $F$ は斉次平面三次式のアフィン表示であり、$V_i$ は
平面の射影ベクトル場である。三重積への正則延長の必要十分条件は
\[
 F(p)=F(q)=0,\qquad V_i(p)=V_i(q)=0\quad(0\le i\le2).
\]
すなわち $\mathbb P^1\times\{p,q\}$ 上で双ベクトルが零になる条件である。
固定した座標の下で両側の条件を満たす双ベクトルの空間はともに二十次元で、
共通アフィンチャートでの恒等写像がそれらを線形に同一視する。
Jacobi 方程式も同一視される。
これらの条件が成り立つとき $\mathbb P^1\times S_0$ の上への持ち上げは正則で、
両射影が Poisson 射になる。

特に $F=0$ が滑らかな平面三次曲線 $\Gamma_2$ で、混合項が零ならば、
$p,q\in\Gamma_2$ を相異なる一般の二点に選べる。
その狭義変換は $S_0$ 上でも滑らかな楕円曲線であり、
$p,q$ を結ぶ直線の狭義変換と一点で横断的に交わる。
従って収縮後の曲線 $\Gamma_{11}$ は滑らかな $(2,2)$ 曲線である。
クラスを用いると、$H$ を平面の直線、$E_p,E_q$ を例外曲線として
\[
 F_1=H-E_p,\quad F_2=H-E_q,\quad E_r=H-E_p-E_q,
 \qquad 3H-E_p-E_q=2F_1+2F_2-E_r
\]
となる。これが $(2,2)$ 型を与える。
逆に一般の滑らかな $(2,2)$ 曲線上の一般点を $r$ に選べば、
二本の座標線を収縮して滑らかな平面三次曲線を得る。
従って純粋な曲面 Poisson 構造については、平面三次曲線の一般点と
二重積の $(2,2)$ 曲線の一般点が、正則 Poisson 共通モデルを介して対応する。

上の曲面リンクを三重積の一般のペンシルに直接適用するには、
\eqref{eq:surface-link-conditions} を満たす特殊化が必要である。
実際、一般のペンシルでは三つの辺の $3\times3$ テンソルはいずれも階数三である。
例えば $s=1+xyz$, $t=x+y+z$ では各辺の係数は
$1-yz(y+z)$ とその置換であり、係数行列の行列式は非零である。
$\mathbb P^1_x\times\{(y_0,z_0)\}$ 上で混合項が恒等的に消える条件は、
階数三のテンソルを $y_0$ または $z_0$ における非零評価写像で縮約して零にする条件なので、
成立しない。ペンシルから純粋な曲面型へ先に楕円曲線に沿う基本変換を行うことで、
その後に曲面リンクを適用できる。

\subsection{三つの標準的なコンパクト化のトーラス型}
$T=(\mathbb C^*)^3$ の座標を $x_1,x_2,x_3$ とする。
$\lambda=(\lambda_{ij})$ を定数歪対称行列として
\[
 \pi_\lambda=\sum_{i<j}\lambda_{ij}
 (x_i\partial_{x_i})\wedge(x_j\partial_{x_j})
\]
と置く。三つのベクトル場は可換なので Jacobi 条件が成り立つ。
任意の滑らかなトーリックコンパクト化でこれらは大域正則ベクトル場に
延長する。従って $\pi_\lambda$ は $\mathbb P^3$、
$\mathbb P^1\times\mathbb P^2$、$(\mathbb P^1)^3$ のそれぞれで正則であり、
共通の稠密トーラス上で恒等写像によって同じ Poisson 双ベクトルを与える。

この対応には滑らかな共通 Poisson 解消が存在する。
$N=\mathbb Z^3$ において $v_i=e_i$ $(1\le i\le3)$、
$v_0=-e_1-e_2-e_3$ とする。非空真部分集合 $I\subset\{0,1,2,3\}$
について $v_I=\sum_{i\in I}v_i$ を取り、部分集合の鎖に対応する錐を
すべて取ると $A_3$ 型 permutohedral fan を得る。
$N_{\mathbb R}$ の実座標を $\xi_1,\xi_2,\xi_3$ とすると、
その極大錐は $0,\xi_1,\xi_2,\xi_3$ の大小関係を一つに固定する錐であり、
その各原始生成元は $N$ の基底をなす。この扇は
$\mathbb P^3$ の扇、座標直交象限からなる $(\mathbb P^1)^3$ の扇、
および、$\mathbb P^1$ の方向を $e_3$、$\mathbb P^2$ の三つの半直線を
$e_1,e_2,-e_1-e_2$ とする積の扇を細分する。
対応する滑らかな射影トーリック三次元多様体 $Z_{\rm per}$ から三つの
コンパクト化への固有双有理トーリック射が得られる。
$\pi_\lambda$ は $Z_{\rm per}$ 上でも正則であり、各射が稠密トーラス上で
Poisson であることから全体で Poisson である。
三つの射はいずれも滑らかなトーラス不変中心の blow-up の合成である。
その中心を以下に明示する。滑らかな扇の錐
$\tau=\langle r_1,\ldots,r_c\rangle$ に対応する軌道閉包を blow-up すると、
$r_\tau=r_1+\cdots+r_c$ を加える star subdivision になる。
三次元錐を中心にすれば点、二次元錐を中心にすれば曲線の blow-up である。
中心とそれに対応する錐は、各段階の模型上で取る。

\begin{proposition}\label{prop:toric-explicit-blowups}
共通模型 $Z_{\rm per}$ は次の三つの列のいずれによっても得られる。
\begin{enumerate}
\item $\mathbb P^3$ から四つの座標点を blow-up し、続いて六つの座標直線の
狭義変換を blow-up する。対応する中心錐は、まず
$\langle v_i,v_j,v_k\rangle$ $(i<j<k)$ の四つ、続いて
$\langle v_i,v_j\rangle$ $(i<j)$ の六つである。
\item $(\mathbb P^1)^3$ から、二つの点
$\langle e_1,e_2,e_3\rangle$、$\langle-e_1,-e_2,-e_3\rangle$
を blow-up する。続いて六つの曲線
$\langle e_i,e_j\rangle$ および $\langle-e_i,-e_j\rangle$ $(i<j)$
の狭義変換を blow-up する。
\item $\mathbb P^1\times\mathbb P^2$ の扇を上と同じ規約で取り、
$c=-e_1-e_2$、$w=-e_1-e_2-e_3$ と置く。
次の表の順序で、一点と八曲線を blow-up する。
\[
\begin{array}{c|c|c}
 \text{段階}&\text{中心錐}&\text{新しい半直線の原始生成元}\\\hline
 1&\langle e_1,e_2,e_3\rangle&e_1+e_2+e_3\\
 2&\langle c,-e_3\rangle&w\\
 3&\langle c,e_2\rangle&-e_1\\
 4&\langle w,e_2\rangle&-e_1-e_3\\
 5&\langle c,e_1\rangle&-e_2\\
 6&\langle w,e_1\rangle&-e_2-e_3\\
 7&\langle e_2,e_3\rangle&e_2+e_3\\
 8&\langle e_1,e_3\rangle&e_1+e_3\\
 9&\langle e_1,e_2\rangle&e_1+e_2
\end{array}
\]
\end{enumerate}
各列の全中間模型で $\pi_\lambda$ は正則であり、各 blow-down は Poisson 射である。
従って三つの端点の間のトーラス型の双有理写像は、この共通模型を通る
Poisson blow-up と blow-down の列として表せる。
\end{proposition}
\begin{proof}
各列について star subdivision の規則を適用する。
以下では $\sigma(1)$ を錐 $\sigma$ の半直線の原始生成元の集合とする。
中心錐 $\tau$ を含む最大錐 $\sigma$ は、$r\in\tau(1)$ ごとに
\[
 \bigl\langle r_\tau,\ \sigma(1)\setminus\{r\}\bigr\rangle
\]
で置き換え、それ以外の最大錐は保つ。
上の各中心は指定された段階で実際に錐をなし、新しく得られる最大錐の
三本の生成元の行列式は $\pm1$ である。
また各段階の二次元面の線形包は
$\xi_i=0$ または $\xi_i=\xi_j$ の形になる。
これらは、表および二つの対称な列の生成元について外積を取れば確認できる。
従って $0,\xi_1,\xi_2,\xi_3$ の大小関係による扇は、全中間扇を細分する。

三つの初期扇の最大錐の個数はそれぞれ $4,8,6$ である。
三次元錐の star subdivision は一つの最大錐を三つに置き換える。
二次元錐は滑らかな完備三次元扇のちょうど二つの最大錐に含まれ、
それぞれを二分する。従っていずれの操作も最大錐の個数を二つ増やす。
操作の回数はそれぞれ $10,8,9$ なので、最終扇の最大錐は全て24個である。
これらは24個の大小関係の錐によって細分されるため、各最大錐が一つずつ一致する。
同じ結論は、最終最大錐が
\[
 \bigl\langle
 v_{i_1},\ v_{i_1}+v_{i_2},\ v_{i_1}+v_{i_2}+v_{i_3}
 \bigr\rangle,
 \qquad (i_1,i_2,i_3,i_4)\in\mathfrak S_4
\]
となることを上の置換規則で列挙しても得られる。
付属の \texttt{check\_birational.py} は中心の存在、全段階の滑らかさと面の条件、
最終24錐の集合の一致を整数演算で検証する。
Poisson 構造については、各中間模型が同じトーラスの滑らかなコンパクト化なので、
三つのトーラス作用のベクトル場の外積がそのまま正則な延長を与える。
\end{proof}

\subsection{\texorpdfstring{$(\mathbb P^1)^3$}{P1 x P1 x P1} の対数型からトーラス型への変換}
$S=\mathbb P^1_x\times\mathbb P^1_z$ とし、
$F,G\in H^0(S,\mathcal O_S(1,1))$ を一般の切断とする。
第三因子の座標を $y$ とし、対応する対数型の双ベクトルを局所体積形式
$dx\wedge dy\wedge dz$ との縮約で
\[
 \alpha=yFG\left(a\frac{dy}{y}+b\,d\log\frac FG\right)
\]
と表す。これは $y=0,\infty,F=0,G=0$ に対応する対数型である。
各係数は
\[
 \alpha_x=by(GF_x-FG_x),\quad
 \alpha_y=aFG,\quad
 \alpha_z=by(GF_z-FG_z)
\]
であり、それぞれの多重次数は高々 $(0,1,2)$、$(2,0,2)$、$(2,1,0)$
となる。従って対応する双ベクトルは三重積全体で正則である。
$F=xA(z)+B(z)$、$G=xC(z)+D(z)$ と書き、
$\Delta=AD-BC\in H^0(\mathbb P^1_z,\mathcal O(2))$ と置く。
有理写像
\[
 \Phi(x,y,z)=\left(u=\frac FG,\,y,\,z\right)
\]
は、$\Delta\ne0$ の開集合上で
$x=(B-uD)/(uC-A)$ を逆写像に持つ。
$\partial u/\partial x=\Delta/G^2$ であるから、体積形式の変換も含めて
\[
 \iota_{\Phi_*\pi}(du\wedge dy\wedge dz)
   =\frac{\Delta}{G^2}\alpha
   =\Delta\bigl(a u\,dy+b y\,du\bigr).
\]
従って双ベクトルそのものの像は
\[
 \Phi_*\pi
   =\Delta(z)\left(b y\partial_y\wedge\partial_z
                         +a u\partial_z\wedge\partial_u\right).
 \tag{T}\label{eq:birational-log-toric}
\]
直接に Poisson 括弧を計算しても
\[
 \{u,y\}=0,\qquad \{y,z\}=by\Delta,\qquad
 \{z,u\}=au\Delta
\]
を得る。像の各項も三つの因子上の大域正則ベクトル場の外積である。
一般の $F,G$ に対して $\Delta\partial_z$ は相異なる二つの零点を持ち、
適当な射影座標 $w$ で $\Delta\partial_z=\kappa w\partial_w$ と書ける。
式 \eqref{eq:birational-log-toric} はこの座標でトーラス型となる。
その係数行列は $\lambda_{uy}=0$ を満たすので、この構成の像は
トーラス型のうち一つの係数が零となる部分族である。

この双有理写像も Poisson blow-up の組合せで実現できる。
一般の $F,G$ の共通零点は $S$ 上の相異なる二点 $p_1,p_2$ であり、
その $z$ 座標は $\Delta$ の二つの零点である。
$X=S\times\mathbb P^1_y$ 上の不定点集合は
\[
 E=(F=G=0)=\bigl(\{p_1\}\sqcup\{p_2\}\bigr)\times\mathbb P^1_y.
\]
$E$ は滑らかな余次元二の零 Poisson 部分多様体である。
グラフの閉包は $\mathrm{Bl}_E X$ で、方程式は
$u_0F-u_1G=0$ である。反対側から見ると、これは
\[
 E'=\{u_0A-u_1C=u_0B-u_1D=0\}
       \subset\mathbb P^1_u\times\mathbb P^1_z\times\mathbb P^1_y
\]
の blow-up である。$E'$ も二本の相異なる $\mathbb P^1_y$ の和であり、
式 \eqref{eq:birational-log-toric} はその上で零となる。
滑らかな余次元二の Poisson 中心では conormal Lie algebra の退化条件は
自動的に成り立つので、Polishchuk の blow-up 定理によって両側の
blow-up に正則 Poisson 構造が持ち上がる。両構造は共通の稠密開集合で
一致するため、グラフ全体で一致する。これにより、両コンパクト化上の正則な双ベクトルと共通 Poisson 解消が得られる。


\subsection{\texorpdfstring{$X_{12}$}{X12} の有理 pencil 成分}
分類節の $C_{(1,1)^2(0,1)}$ を考える。
$F,G\in H^0(X_{12},\mathcal O(1,1))$、
$L\in H^0(\PP^2,\mathcal O(1))$ に対し
\[
 \alpha=L(F\,dG-G\,dF)
\]
と置く。$F=xa(v,w)+b(v,w)$、$G=xc(v,w)+d(v,w)$ と書くと、
$u=G/F$ は Casimir であり
\[
 \{v,w\}=L(bc-ad),\qquad \{u,v\}=\{u,w\}=0.
\]
従って $(x,v,w)\mapsto(u,v,w)$ は、零三次曲線が直線と二次曲線の和である
垂直型への Poisson 双有理写像である。
一般の係数では、行列式 $bc-ad=0$ は滑らかな二次曲線であり、
$a,b,c,d$ は平面上で同時には零にならない。
$F=G=0$ および逆写像の基底曲線は、この二次曲線と同型な滑らかな有理曲線である。
両曲線上でそれぞれの双ベクトルは零になる。
前述の比の閉グラフの議論と命題\ref{prop:poisson-centres}により、
そのグラフは両側の Poisson blow-up である。

より一般に $\CC(S)(x)$ の括弧が $\CC(S)$ 上の括弧 $\sigma$ を保ち、
$u=(ax+b)/(cx+d)$ が非定数の Casimir ならば、
$(x,s)\mapsto(u,s)$ は同じ $\sigma$ を垂直に持つ積への Poisson 双有理写像を与える。
正則な共通模型の構成では、基底イデアルと各中心上の双ベクトルを別に確認する。

\subsection{共通模型上の古典的 SOD}

滑らかな中心 $Y\subset X$ の余次元を $c$、$f:\widetilde X\to X$ をその blow-up、
$j:E\hookrightarrow\widetilde X$、$p:E\to Y$ を例外因子の射とする。
Orlov \cite[定理4.3]{OrlovBlowup} の式は
\[
 \operatorname{Perf}(\widetilde X)=
 \left\langle
 j_*\bigl(p^*\operatorname{Perf}(Y)\otimes\mathcal O_E(1-c)\bigr),\ldots,
 j_*\bigl(p^*\operatorname{Perf}(Y)\otimes\mathcal O_E(-1)\bigr),
 Lf^*\operatorname{Perf}(X)
 \right\rangle.
\]
三次元では曲線中心につき $\operatorname{Perf}(Y)$ が一個、
点中心につき $\operatorname{Perf}(\CC)$ が二個現れる。
例えば定理\ref{thm:rat22-vertical}の共通模型では、埋込み関手を省略すると
\[
\begin{aligned}
 \operatorname{Perf}(Z)
 &=\left\langle\operatorname{Perf}(E'),\operatorname{Perf}(\CC),
                  \operatorname{Perf}(\CC),\operatorname{Perf}(\PP^3)\right\rangle,\\
 \operatorname{Perf}(Z)
 &=\left\langle\operatorname{Perf}(\Gamma),
                  \operatorname{Perf}(\PP^1\times\PP^2)\right\rangle.
\end{aligned}
\]
$E'\cong\Gamma$ は楕円曲線である。
三重積 pencil の共通模型でも、両側にそれぞれ楕円曲線の圏が一個加わる。

命題\ref{prop:toric-explicit-blowups}の $Z_{\rm per}$ は長さ24の full exceptional collection を持つ。
実際、三つの端点の標準的な積の例外集の長さは順に $4,6,8$ である。
各点 blow-up には二つの例外対象が加わり、各曲線中心は滑らかな完備トーリック曲線
$\PP^1$ なので、その二項例外集が加わる。従って三つの構成から得られる長さは
\[
 4+2\cdot4+2\cdot6
 =6+2\cdot1+2\cdot8
 =8+2\cdot2+2\cdot6=24
\]
と一致する。この記述は上の可換な共通模型の導来圏についてのものである。
これらの式の非可換化には、中心とその例外部分を含む blow-up 圏の構成が必要になる。
その構成に適用できる先行結果を第\ref{sec:nc-blowups}節で整理する。


\part{非可換射影模型と関数斜体}
\section{分数体と半古典極限}
\label{sec:nc-fields}

基礎体を標数零の体 $k$ とする。連結次数付き noetherian domain $A$ に対して
\[
 k(\operatorname{qgr}A)=Q_0(A):=Q_{\mathrm{gr}}(A)_0
\]
と書く。$Q_{\mathrm{gr}}$ は非零斉次元を可逆にした次数付き分数環である。
本稿での非可換双有理同値は、この division ring の $k$-同型を意味する。
必要な場合には、半古典極限との整合性と、共通の非可換 blowup による実現も指定する。
一般の smooth proper dg 圏の変形から、この分数体を取り出すには
追加の射影表示またはアフィン表示が必要である。

\begin{lemma}\label{lem:segre-field}
$A$ を $A_1\ne0$ を満たし次数1で生成される連結次数付き noetherian domain とし、
$R=A\mathbin{\#}k[s,t]=\bigoplus_{n\geq0}A_n\otimes k[s,t]_n$ と置く。
このとき $Q_0(R)\cong Q_0(A)(u)$、ただし $u=t/s$ は中心変数である。
\end{lemma}
\begin{proof}
同次数の非零斉次元 $a,b\in A_n$ の比は
$(a s^n)(b s^n)^{-1}$ によって $Q_0(R)$ に入る。
非零 $a_1\in A_1$ を取れば、
$(a_1t)(a_1s)^{-1}=u$ も入る。
逆に $a_ns^{n-i}t^i=a_ns^n u^i$ なので、
任意の斉次元の比は $Q_0(A)(u)$ に入る。
$u$ は中心で超越的であり、二つの包含は逆になる。
\end{proof}

\begin{lemma}\label{lem:ore-projective-bundle}
$G$ を次数1で生成される連結局所有限 noetherian 次数付き代数とし、
$\mathcal C=\operatorname{Perf}(\operatorname{qgr}G)$ を smooth proper とする。
$\rho$ を $G$ の次数付き自己同型とする。
$H=G[Y_0][Y_1;\rho]$ とし、$Y_0$ は中央、
$\deg G_1=(1,0)$、$\deg Y_i=(0,1)$ とする。
右加群の twist $M\mapsto M_\rho$ が誘導する自己同値を $\Phi$ と書く。
$H$ の二重次数付き射影圏には
\[
 \operatorname{Perf}(\operatorname{Qgr}_{\mathbb Z^2}H)
 =\langle p^*\mathcal C,\ p^*\mathcal C(0,1)\rangle
\]
という半直交分解があり、これは smooth proper な dg 圏である。
対角部分環の商圏を取っても同じ結論を得る。
\end{lemma}
\begin{proof}
$p^*$ は $M\mapsto M\otimes_GH$ から誘導される。
第二次数についての相対 cohomology は、PBW 基底と二変数 Koszul 複体から
\[
 Rp_*p^*M=M,\qquad Rp_*(p^*M(0,-1))=0,\qquad
 Rp_*(p^*M(0,1))=M\oplus\Phi M
\]
となる。
従って二つの functor は fully faithful であり、指定した順序で半直交である。
また $Y_0,Y_1$ による Koszul 複体を射影圏へ移すと
\[
 0\longrightarrow p^*(\Phi M)(0,-2)
 \longrightarrow p^*M(0,-1)\oplus p^*(\Phi M)(0,-1)
 \longrightarrow p^*M\longrightarrow0
\]
を得る。これは全ての第二次数の twist を二つの隣接する次数から生成する。
誘導加群の全 twist は二重次数付き射影圏の perfect 部分を生成するため、
上の半直交分解は full である。
二成分の間の bimodule は $\mathrm{id}_{\mathcal C}\oplus\Phi$ の graph bimodule である。
これは $\mathcal C$ の smoothness により perfect であり、properness により
各 Hom 複体は有限次元 cohomology を持つ。
二つの diagonal bimodule とこの graph bimodule を結ぶ三角形から、
gluing の diagonal も perfect となる。
これで smoothness と properness が従う。
最後に対角部分環 $S=\bigoplus_nH_{n,n}$ を考える。
$H$ は noetherian Ore 拡大である。
斉次右イデアル $I\subset S$ について $IH$ の有限個の斉次生成元を
$i_1,\ldots,i_m\in I$ から選ぶと、対角部分を取って
$I=\sum_j i_jS$ を得る。従って $S$ は graded noetherian であり、
非負次数の最高次項を用いる標準的な議論から noetherian でもある。
左イデアルについても同様である。
二重次数付き商圏では、各元が十分大きい矩形次数の全ての元で
零になる加群を torsion とする。
PBW 基底と $G$ の次数1生成性から、非負添字について
$H_{a,b}H_{c,d}=H_{a+c,b+d}$ である。
対角抽出 $M\mapsto M_\Delta$ と $N\mapsto N\otimes_SH$ を取ると、
後者の対角抽出は $N$ 自身になる。
counit $M_\Delta\otimes_SH\to M$ の余核は、各斉次元を
十分大きい対角次数へ移すと像に入るため torsion である。
核の対角部分は零なので、同じ積の全射性から核も torsion である。
従って両 functor は商圏で互いに逆の同値を与える。
通常の tensor product の場合の定式化には
\cite[\S3.2、定理3.32]{BallardFarman} も参照されたい。
\end{proof}

\subsection{半古典極限を備えた division ring}

滑らかな点付き曲線の局所環を $R$、その一様化元を $\hbar$、分数体を $k$ とする。
$R$ 上平坦な Ore domain $\mathcal A$ について、
$\mathcal A/\hbar\mathcal A$ が可換 domain で、
$\hbar$-進フィルターが分離的であると仮定する。
PBW 基底を持つ以下の族はこの条件を満たす。
\[
 v(a)=\max\{n:a\in\hbar^n\mathcal A\},\qquad
 v(ab)=v(a)+v(b)
\]
だから、$v$ は $\Frac\mathcal A$ に離散付値として延長する。
その付値環の剰余体は $\Frac(\mathcal A/\hbar\mathcal A)$ であり、
\[
 \{\bar a,\bar b\}=\overline{\hbar^{-1}[a,b]}
\]
で Poisson 括弧が定まる。
実際、分数の分子と分母から同じ $\hbar$ の冪を除いて剰余を取れば
剰余体の記述を得る。積と交換子についての主張は Ore 分数表示によって延長する。
この構造と剰余体 $\mathbb C(X)$ の同定を、分数体の Poisson 標識と呼ぶ。
次数付き族の $Q_0$ に対しても同じ構成を用いる。

\begin{proposition}\label{prop:toric-common-field}
乗法的交代行列 $q=(q_{ij})$ に対して
\[
 \mathbb T_q=k\langle X_1^{\pm1},X_2^{\pm1},X_3^{\pm1}\rangle/
 (X_iX_j=q_{ij}X_jX_i)
\]
と置く。三模型 $\PP^3,\PP^1\times\PP^2,(\PP^1)^3$ の Cox 環を
同じ三次元トーラスの双指標で twist した斉次環は、
いずれも $\Frac\mathbb T_q$ を次数零分数体に持つ。
$q_{ij}=1+\hbar\lambda_{ij}+O(\hbar^2)$ なら、半古典極限は
$\{x_i,x_j\}=\lambda_{ij}x_ix_j$ である。
\end{proposition}
\begin{proof}
各射影因子の一つの斉次座標にトーラス重み零を与え、残る三つに
$\mathbb Z^3$ の標準基底を与える。
各因子の重み零の座標を可逆にした次数零部分は量子アフィン三空間であり、
さらに三つの座標を可逆にすると $\mathbb T_q$ を得る。
通常の斉次環または対角 Segre 部分環から同じ比を生成できるので、
全ての次数零分数体は一致する。
単項式基底が PBW 基底であることから次数ごとの平坦性が従い、
交換関係を $\hbar$ の一次まで取れば括弧を得る。
射影模型の smoothness と矩形 exceptional collection は
同梱の量子化原稿第3節で証明した。
\end{proof}

$M=(m_{ia})\in\mathrm{GL}_3(\mathbb Z)$ に対し、順序を固定した単項式
$Y_i=X_1^{m_{i1}}X_2^{m_{i2}}X_3^{m_{i3}}$ を取ると
\[
 q'_{ij}=\prod_{a,b}q_{ab}^{m_{ia}m_{jb}},\qquad
 \lambda'=M\lambda M^{\mathsf T}
\]
となる。逆行列の単項式は可逆なスカラー因子を除いて $X_i$ を復元するので、
これは実際の division ring の同型である。
この公式を用いて高次の量子化パラメーターも対応させる。

\begin{proposition}\label{prop:log-coefficient-invariant}
log-canonical Poisson 関数体
$\mathbb C(x_1,x_2,x_3)$、$\{x_i,x_j\}=\lambda_{ij}x_ix_j$
について、加法群
\[
 \Gamma_\lambda=\sum_{i<j}\mathbb Z\lambda_{ij}\subset\mathbb C
\]
は Poisson 双有理不変量である。
\end{proposition}
\begin{proof}
反復 Laurent 級数体への埋め込みで単項式付値を取る。
有理関数 $f,g$ の先頭指数を $m,n\in\mathbb Z^3$ とする。
$\{f,g\}=c fg$ が定数 $c$ について成立する場合、
先頭項の比較により $c=m\lambda n^{\mathsf T}$ となる。
従って全ての定数 log-bracket は $\Gamma_\lambda$ に属する。
逆に座標対は全ての $\lambda_{ij}$ を与えるので、
このような定数の生成する加法群はちょうど $\Gamma_\lambda$ である。
これは関数体そのものから定義される。
同じ議論の一般次元の定式化は
\cite[命題5.2(b)、定理5.5]{GoodearlLaunois} にある。
\end{proof}

従って係数に $\mathbb Q$-線形独立性を課すと、
一般の三係数型、二係数型、一係数型はそれぞれ
$\operatorname{rank}_{\mathbb Z}\Gamma_\lambda=3,2,1$ で区別できる。
以下の成分間の対応では係数を照合している。

\subsection{pencil の分数線形変換}

\begin{proposition}\label{prop:quantum-pencil}
$D$ を曲面の Poisson 関数体の量子化とし、$U$ を中心変数とする。
四つの有理関数 $a,b,c,d$ の持ち上げ $A,B,C,D_1\in D$ を選び、
$bc-ad\ne0$ と仮定する。このとき
\[
 X=(UA-C)^{-1}(D_1-UB),\qquad
 U=(CX+D_1)(AX+B)^{-1}
\]
は互いに逆な非可換分数式であり、$D(U)$ は $D$ と $X$ で division ring
として生成される。半古典極限は
\[
 x=\frac{d-ub}{ua-c},\qquad u=\frac{cx+d}{ax+b}
\]
である。
\end{proposition}
\begin{proof}
$UA-C$ の剰余は非零である。
$AX+B$ の剰余は $(ad-bc)/(ua-c)$ なので、これも非零である。
定義式の左から $UA-C$ を掛けて整理すると
$U(AX+B)=CX+D_1$ を得る。$U$ の中心性だけを用いており、
$A,B,C,D_1$ の相互の可換性は仮定していない。
逆式から $U$ を復元でき、剰余を取ると所望の有理式になる。
Poisson 括弧は付値環の交換子から定まるため、この剰余体の同型と整合する。
\end{proof}

三重積の一般の pencil を
\[
 s=xa(y,z)+b(y,z),\quad t=xc(y,z)+d(y,z),\quad
 \Delta=bc-ad
\]
と書くと、古典的な変換は $u=t/s$ であり、変換後の括弧は
$\{y,z\}=\Delta(y,z)$、$u$ は Casimir である。
従って、双二次式 $\Delta$ の非可換 quadric の関数体 $D$ に
命題\ref{prop:quantum-pencil}を適用すれば、pencil と二因子型を
同時に標識した量子化分数体 $D(U)$ が得られる。

$\PP^3$ の二次曲面 pencil でも、
$q_i=x_0\ell_i+r_i$ として
$(a,b,c,d)=(\ell_0,r_0,\ell_1,r_1)$ を用いれば同じ命題を適用できる。
曲面側の三次式は $\ell_1r_0-\ell_0r_1$ である。
これらは分数体とその Poisson 標識の構成である。
指定した射影模型の斉次環への延長と、その量子化の選択は
第\ref{sec:elliptic-links}節および第\ref{sec:p3-models}節で扱う。


\section{\texorpdfstring{$\PP^1\times\PP^2$}{P1 x P2} の一般元の有理模型}
\label{sec:p1p2-rational-models}\label{sec:p1p2-models}

この節の成分名と正規形は、第\ref{p12:classification}節の分類による。
以下の議論は各正規形の式から直接行う。
一般元の有理 Poisson 模型と、その模型上で選ぶ代数的量子化は次の通りである。
表の $c$ は Poisson 中心に属し、記していない括弧は零である。
係数 $\kappa,\lambda,\mu,a,b$ は変換後も保持する。
\[
\begin{array}{c|l|l}
\text{成分}&\text{有理 Poisson 模型}&\text{選ぶ量子化}\\\hline
V_{\mathrm{vert}}&\{v,w\}=f(v,w),\ c\text{ 中心}
 &\text{非可換平面と中心変数}\\
C_{\mathrm{tor}}&\{x_i,x_j\}=\lambda_{ij}x_ix_j
 &\text{三変数量子トーラス}\\
C_{\mathrm{fib}}&\{r,s\}=\kappa rs,\ c\text{ 中心}
 &\text{量子平面と中心変数}\\
C_{\mathrm{line\text{-}conic}}&\{r,s\}=\mu rs,\ \{s,t\}=\lambda st
 &\text{三変数量子トーラス}\\
C_{(1,1)^2(0,1)}&\{r,s\}=\kappa rs,\ c\text{ 中心}
 &\text{量子平面と中心変数}\\
C_{(1,1)(1,0)(0,2)}&\{r,s\}=\kappa rs,\ c\text{ 中心}
 &\text{量子平面と中心変数}\\
C_4&\{r,t\}=-b rt,\ \{s,t\}=a st
 &\text{三変数量子トーラス}\\
C_{(1,2)(1,0)(0,1)}&\{r,s\}=\kappa rs,\ c\text{ 中心}
 &\text{量子平面と中心変数}\\
C_{\mathrm{Ric}}&
 (f(x)\partial_x+h(t,x)\partial_t)\wedge\partial_y
 &\text{微分型 Ore 拡大}
\end{array}
\]
最初の行の一般の $f$ は滑らかな平面三次曲線を定める。
その射影量子化は 第\ref{sec:surface-quantization}節 にある。
表の量子トーラスは有限型基底上の PBW 族であり、
有理座標の古典極限を通して指定された Poisson 関数体を量子化する。
この表示を元の射影コンパクト化の斉次環へ延長する問題は、
各模型ごとの条件を伴う。次節では直線・二次曲線型について
両側の射影斉次環まで構成する。

\begin{proposition}\label{prop:p1p2-seven-toric}
上表の $V_{\mathrm{vert}}$ と $C_{\mathrm{Ric}}$ を除く七成分の一般元は、
表示した log-canonical Poisson 関数体を持つ。
従って係数行列が一致する二つの表示には、
同じ量子トーラスの分数体を用いる有理量子化を選べる。
\end{proposition}
\begin{proof}
体積形式 $du\wedge dv\wedge dw$ と
$\alpha=\iota_\pi(du\wedge dv\wedge dw)$ の規約を用いる。
トーラス成分については定義から従う。

\emph{直線束型。}
$\alpha=H\,dw$ とし、
$H=A(u)+vB(u)+wC(u)$ と書く。$A,B,C$ は二次以下である。
$s=H$ と置くと
\[
 \{u,s\}=B(u)s,\qquad w\text{ は Casimir}
\]
となる。一般の $B$ について
$B(u)=\beta(u-r_1)(u-r_2)$、$r_1\ne r_2$ と書き、
\[
 r=\frac{u-r_1}{u-r_2},\qquad c=w,\qquad
 \kappa=\beta(r_1-r_2)
\]
と置けば $\{r,s\}=\kappa rs$ である。
逆変換は $u=(rr_2-r_1)/(r-1)$ と
$v=(s-A(u)-cC(u))/B(u)$ で与えられる。
族全体で根を選ぶ場合は、判別式非零の開集合の有限 \'etale 被覆を取る。

\emph{$(1,2)$ 型。}
$U=L=1$ のチャートで
$\alpha=dF$、$F=uQ(v,w)+R(v,w)$ と書ける。
$c=F$ は Casimir で、$\{v,w\}=Q$ である。
一般の二次式をアフィン座標変換によって
$Q=\kappa(vw+1)$ とする。
この座標変換の行列式による定数も $\kappa$ に含め、
$F,R$ を同じ定数で取り直す。
$r=v,s=vw+1$ と置けば $\{r,s\}=\kappa rs$、
\[
 w=(s-1)/r,\qquad u=(c-R(v,w))/(\kappa s)
\]
である。

\emph{$(1,1)$-pencil 型。}
$L=1$ とし、$\alpha=G\,dF-F\,dG$、$F=ua+b,G=uc+d$ と書く。
$a,b,c,d$ は平面上の一次式である。
ここだけ Casimir を $z=F/G$ と書くと
\[
 \{v,w\}=ad-bc,\qquad
 u=\frac{zd-b}{a-zc}.
\]
一般の二次式 $ad-bc$ を $\kappa(vw+1)$ に直し、
前段と同じ二変数変換を適用する。

\emph{$(1,1)(1,0)(0,2)$ 型。}
$U=1$ とし、$F$ の $u$ 係数を平面のアフィンチャートの分母に取ると
$F=u+\ell(v,w)$ とできる。$z=F$ と置けば
\[
 \alpha=2Q\,dz-z\,dQ.
\]
$Q=\kappa_0(vw+1)$ と正規化すると
\[
 \{z,v\}=-\kappa_0zv,\quad
 \{z,w\}=\kappa_0zw,\quad
 \{v,w\}=2\kappa_0(vw+1).
\]
従って $c=z^2/(vw+1)$ は Casimir であり、
$(r,s)=(z,v)$ と $\kappa=-\kappa_0$ を用いれば表の形となる。
逆変換は $w=(r^2/c-1)/s$、$u=r-\ell(v,w)$ である。

\emph{直線・二次曲線型。}
$L=1$ とし、第一因子上の半単純ベクトル場を対角化する。
二次式も正規化し、係数を吸収すると
\[
 \alpha=\lambda(vw+1)\,du+\mu u\,d(vw+1)
\]
となる。$(r,s,t)=(u,v,vw+1)$ に対して
\[
 \{r,s\}=\mu rs,\qquad \{s,t\}=\lambda st,
 \qquad \{r,t\}=0,
\]
逆式は $w=(t-1)/s$ である。

\emph{四因子型。}
一般配置を $U=1,L=v,M=w,F=u+v+w$ とし、$z=F$ と置く。
このとき
\[
 \alpha=avw\,dz+bzw\,dv-(a+b)zv\,dw.
\]
$r=z/w,s=v/w,t=w$ と取れば
\[
 \{r,s\}=0,\quad \{r,t\}=-brt,\quad \{s,t\}=ast.
\]
逆変換は $w=t,v=st,u=(r-s-1)t$ である。
以上は各場合の有理逆写像を伴うので関数体の同型である。

量子化には $q_{ij}=1+\hbar\lambda_{ij}$ を基本パラメーターとして
対応する量子トーラスを取る。
整数単項式変換を戻す場合の新しい $q$ は、基本の $q_{ij}$ の整数べきの積で定める。
これにより、古典的な係数の整数関係と対応する量子交換関係を全次数で保てる。
\end{proof}

\section{直線・二次曲線型と \texorpdfstring{$\PP^3$}{P3} の対数型の量子双有理変換}
\label{sec:line-conic-quantum-link}\label{sec:line-conic-quantum}

$p=1+\hbar\mu,q=1+\hbar\lambda$、
$R=\CC[\hbar,(pq)^{-1}]$ とする。
$g(U,Z)$ を二次斉次式とし、
\[
 \sigma(U)=p^{-1}U,\quad \sigma(Z)=Z,\qquad
 F_g=qg-\sigma(g)
\]
と置く。四生成元を全て次数1として
\begin{equation}\label{eq:ore-log112}
 A_g=\frac{R\langle U,Z,V,W\rangle}
 {\bigl(Z\text{ は中心},\ UV-pVU,\ UW-p^{-1}WU,
       VW-qWV-F_g(U,Z)\bigr)}
\end{equation}
を定める。

\begin{proposition}\label{prop:ore-log112}
$A_g$ は次数ごとに自由な $R$-代数で、各ファイバーは
Hilbert 級数 $(1-s)^{-4}$ の四次元 AS-regular 代数である。
古典的なファイバーは $\CC[U,Z,V,W]$ であり、
$Z=1$ 上の半古典極限は、$Q=vw+g(u,1)$ として
\[
 \alpha=\lambda Q\,du+\mu u\,dQ
\]
に対応する。$(\lambda,\mu,g)$ を一般に選ぶと、これは $\PP^3$ の
$\operatorname{Log}(1,1,2)$ 成分の一般元を与える。
\end{proposition}
\begin{proof}
$B=R[U,Z][V;\sigma]$ 上で
\[
 \tau|_{R[U,Z]}=\sigma^{-1},\quad \tau(V)=q^{-1}V,
 \qquad \delta|_{R[U,Z]}=0,\quad \delta(V)=-q^{-1}F_g
\]
と置く。$\tau$ は自己同型である。
$b\in R[U,Z]$ に対して、関係式 $Vb=\sigma(b)V$ に
$\tau$-導分を適用した両辺は、ともに $-q^{-1}F_gb$ になる。
従って $\delta$ は well-defined で、
$A_g=B[W;\tau,\delta]$ である。
PBW 基底は $U^iZ^jV^kW^\ell$ であり、平坦性、Hilbert 級数、
全ファイバーの noetherian domain 性を与える。
AS 正則性と大域次元は、graded Ore extension が
twisted Calabi--Yau 次元を1増やす定理から従う
\cite[定理0.2、\S1.2--1.3]{LiuWangWu}。

$g=g(u,1)$ と略記すれば
\[
 \left.\frac{F_g}{\hbar}\right|_{\hbar=0}
   =\lambda g+\mu u g'(u).
\]
従ってアフィンチャートでの括弧は
\[
 \{u,v\}=\mu uv,\qquad \{u,w\}=-\mu uw,\qquad
 \{v,w\}=\lambda(vw+g)+\mu u g'.
\]
これが主張した1形式を与える。
その同次表示は
\[
 UZQ\left(\lambda\frac{dU}{U}
    +(-\lambda-2\mu)\frac{dZ}{Z}
    +\mu\frac{dQ}{Q}\right),\qquad Q=VW+g(U,Z)
\]
である。
一般の二平面を $U=0,Z=0$ に移し、一般の二次式の
$V,W$ に関する非退化部分を $VW$ に直す。
さらに $V,W$ に $U,Z$ の一次式を加えると混合項が消え、
二次式は $VW+g(U,Z)$ になる。これにより一般の
$\operatorname{Log}(1,1,2)$ を覆う。
この斉次 lift では $Z$ を中心として選んでいる。
\end{proof}

$g=Z^2$ の場合に、
\[
 B_q=R\langle Z,V,W\rangle/
       (Z\text{ は中心},\ VW-qWV-(q-1)Z^2)
\]
と置く。これは三次元 AS-regular 代数であり、
$\rho(V)=pV,\rho(W)=p^{-1}W,\rho(Z)=Z$ は自己同型である。
$Y_0$ を中心として
\[
 T=B_q[Y_0][Y_1;\rho],\qquad
 \deg B_{q,1}=(0,1),\quad\deg Y_i=(1,0),\qquad
 C_{p,q}=\bigoplus_{n\ge0}T_{(n,n)}
\]
を取る。ここで $\rho(Y_0)=Y_0$ とする。

\begin{proposition}\label{prop:projective-line-conic-link}
$C_{p,q}$ は $\PP^1\times\PP^2$ の Segre 環の次数ごとに自由な
noetherian 変形であり、
\[
 \operatorname{rank}_R(C_{p,q})_n=(n+1)\binom{n+2}{2}
\]
を満たす。$\lambda,\mu$ と座標変換を動かした族の半古典極限は
直線・二次曲線成分の一般元を覆い、
各幾何学的ファイバーの perfect 圏は smooth proper である。
$A_g$ と $C_{p,q}$ の非可換関数体は、任意の $g$ に対して同型である。
この同型の古典極限は Poisson 双有理変換である。
\end{proposition}
\begin{proof}
PBW 基底より $T_{(a,b)}$ の階数は $(a+1)\binom{b+2}{2}$ である。
対角環は六つの $Y_iZ,Y_iV,Y_iW$ で生成される。
$V,W$ の補助重みを1、$Z,Y_i$ の補助重みを0とする濾過を入れると、
対応する次数環は量子 Segre 環である。
これは有限生成可換半群環の双指標 twist なので noetherian である。
従って $C_{p,q}$ も noetherian である。
$B_q$ の AS 正則性と補題\ref{lem:ore-projective-bundle}により、
射影模型の perfect 圏は smooth proper である。

中心的な次数1元 $Y_0Z$ を局所化した次数零部分は
\[
 u=Y_1/Y_0,\qquad v=V/Z,\qquad w=W/Z
\]
で生成され、関係式は
\[
 uv=pvu,\qquad uw=p^{-1}wu,\qquad
 vw-qwv=q-1.
\]
この半古典極限は前節の直線・二次曲線型の正規形である。

一方、$A_g$ の $Z=1$ のチャートで
\[
 \Omega=wv+g(u,1)
\]
と置くと、関係式から
\[
 u\Omega=\Omega u,\qquad v\Omega=q\Omega v,
 \qquad uv=pvu
\]
を得る。PBW 基底で $\Omega$ のべきの最高 $w$ 次数を比較すると、
$u,v,\Omega$ の量子単項式は一次独立である。
さらに
\[
 w=(\Omega-g(u,1))v^{-1}
\]
なので、分数体は三変数量子トーラスの分数体
\[
 D_{p,q}=
 \Frac\bigl(\operatorname{Frac}(R)
 \langle u^{\pm1},v^{\pm1},\Omega^{\pm1}\rangle/
 (uv-pvu,u\Omega-\Omega u,v\Omega-q\Omega v)\bigr)
\]
に一致する。$C_{p,q}$ のチャートでは $g=1$ なので同じ分数体を得る。
具体的には $u,v$ を固定して
\[
 w_g=w_1+(1-g(u,1))v^{-1}
\]
とする写像が同型である。逆写像も同じ式から得られる。
古典極限は $vw_g+g(u,1)=vw_1+1$ を保ち、
$\{u,v\}=\mu uv,\{v,Q\}=\lambda vQ,\{u,Q\}=0$ と整合する。
\end{proof}

\section{Riccati 成分の微分型量子化と双有理障害}
\label{sec:riccati-quantum}

共点三直線のアフィン正規形を
\[
 \pi=f(x)\partial_x\wedge\partial_y
       +h(t,x)\partial_t\wedge\partial_y,
 \qquad h(t,x)=h_0(t)+xh_1(t)
\]
とする。$f$ は三次式、$h_0,h_1$ は二次以下である。
\[
 \partial=f(x)\partial_x+h(t,x)\partial_t
\]
を $\CC[x,t]$ の導分とする。

\begin{proposition}\label{prop:riccati-ore}
代数
\[
 \mathcal R=\CC[\hbar,x,t][y;-\hbar\partial]
\]
は有限型の平坦な代数族であり、関係式は
\[
 [x,t]=0,\qquad [x,y]=\hbar f(x),\qquad
 [t,y]=\hbar h(t,x)
\]
である。その半古典極限は指定した Riccati 構造に一致する。
各ファイバーは noetherian domain である。
\end{proposition}
\begin{proof}
Ore 拡大の PBW 基底は $x^it^jy^k$ であり、
導分の係数は $\hbar$ と元のパラメーターの多項式である。
これが任意の係数族に対する平坦性と半古典式を与える。
noetherian 性と domain 性も導分型 Ore 拡大から従う
\cite[\S1.3]{LiuWangWu}。
\end{proof}

第\ref{sec:riccati-p3-link}節では、一般元について通常の $\PP^3$ を
古典ファイバーとする AS-regular コンパクト化を構成する。
元の $\PP^1\times\PP^2$ の射影量子化との比較は追加の問題である。

\begin{lemma}\label{lem:differential-poisson-obstruction}
$K$ を標数零の可換体、$\partial$ をその導分とする。
$F=K(y)$ に
\[
 \{a,b\}=0\ (a,b\in K),\qquad \{a,y\}=\partial(a)
\]
という Poisson 構造を入れる。
非零定数 $\lambda\in\ker\partial$ と $A,B\in F^*$ に対して
$\{A,B\}=\lambda AB$ は成立しない。
\end{lemma}
\begin{proof}
$y$ に関する有理関数の次数を、分子の次数と分母の次数の差で定める。
$\partial$ はこの次数を増やさず、$\partial_y$ は少なくとも1下げるので
\[
 \deg_y\{A,B\}\le\deg_y A+\deg_y B-1.
\]
一方、$\lambda AB$ の次数は $\deg_y A+\deg_y B$ である。
\end{proof}

\begin{corollary}\label{cor:riccati-not-log-poisson}
Riccati 正規形の Poisson 関数体は、非零の log-canonical
Poisson 関数体と Poisson 双有理同値でない。
従って、上表の残る七つのトーラス型に帰着する成分の一般元とは区別される。
\end{corollary}
\begin{proof}
Riccati 型には補題を $K=\CC(x,t)$ として適用する。
非零の log-canonical 型では、ある二つの座標が
$\{x_i,x_j\}=\lambda_{ij}x_ix_j$、$\lambda_{ij}\ne0$ を満たす。
その逆像は補題に反する。
\end{proof}

この分離は既知の Poisson 双有理不変量からも従う。
実際、Riccati 型では $\{x,y/f(x)\}=1$ である。
Goodearl--Launois \cite[命題5.2(a)、系5.3]{GoodearlLaunois}
は、log-canonical Poisson 関数体には括弧が1となる二元が存在しないことを示した。
上の次数による証明は、逆に Riccati 関数体における
非零 log-canonical 対の不存在を直接示している。

量子側では Alev--Dumas 不変量
\[
 \AD_k(D)=[D^*,D^*]\cap k^*
\]
を用いる。定義と量子トーラスでの計算は
\cite[定義3.1、命題3.3]{GoodearlYakimov} にまとめられている。

\begin{lemma}\label{lem:differential-ad}
$k\subset\ker\partial\subset K$ とし、
$D=\Frac(K[y;\partial])$ と置くと、$\AD_k(D)=\{1\}$ である。
\end{lemma}
\begin{proof}
非零の Ore 多項式 $P=\sum a_i y^i$ の最高次係数を
$\operatorname{lc}(P)$ と書く。
導分の項は次数を下げるため
\[
 \operatorname{lc}(PQ)=\operatorname{lc}(P)\operatorname{lc}(Q)
\]
となる。右分数で
$\operatorname{lc}(PQ^{-1})=\operatorname{lc}(P)/\operatorname{lc}(Q)$
と定めると、共通右分母に通分することで表示に依らないことが分かる。
実際、$Qa=Sb$ と $Pa=Rb$ が同時に成立する表示について、
最高次係数の等式から $\operatorname{lc}(P)/\operatorname{lc}(Q)
=\operatorname{lc}(R)/\operatorname{lc}(S)$ が従う。
同じ通分により乗法性も従い、$\operatorname{lc}:D^*\to K^*$ は群準同型である。
$K^*$ は可換なので全ての乗法交換子の像は1である。
$\lambda\in k^*$ の像は $\lambda$ 自身なので、主張を得る。
\end{proof}

\begin{corollary}\label{cor:riccati-not-log-quantum}
上の Riccati Ore 模型の量子分数体は、交換係数の少なくとも一つが
1と異なる量子トーラスの分数体と $k$-同型でない。
ここで $k=\CC(\hbar)$、または非零パラメーターに特殊化した後の $k=\CC$
と取れる。
\end{corollary}
\begin{proof}
量子トーラスでは $X_iX_jX_i^{-1}X_j^{-1}=q_{ij}$ なので、
$q_{ij}\ne1$ は非自明な $\AD_k$ の元である。
補題が Riccati Ore 分数体でのその存在を排除する。
\end{proof}

最後の系はここで指定した量子化同士の比較である。
同じ一次 Poisson 構造を持つ量子化の選択を変えた場合には、
新しい代数について不変量を計算する必要がある。

\section{Riccati 型の四次元 AS-regular 射影模型}
\label{sec:riccati-p3-link}

Riccati 成分には、通常の $\PP^3$ を古典ファイバーとする
AS-regular コンパクト化も構成できる。
ここで得られる $\PP^3$ 上の Poisson 構造は、
平面の次数2葉層の線形引き戻し成分 $S(2,3)$ に属する特別な族である。

\begin{proposition}\label{prop:riccati-p3-quantum-link}
Riccati 成分の一般元に対して、その微分型 Ore 量子化と
非可換双有理同値な四次元 quadratic AS-regular 代数の
一変数代数族が存在する。
この族の Hilbert 級数は $(1-s)^{-4}$、
古典的なファイバーは $\CC[Z,X,T,V]$ である。
分数体の同型の古典極限は、Riccati 成分から
$S(2,3)$ 内の特別な族への Poisson 双有理変換を与える。
\end{proposition}
\begin{proof}
元の座標を $(x,t,y)$ とし、
\[
 \partial=f(x)\partial_x+(h_0(t)+xh_1(t))\partial_t,
 \qquad f(x)=a(x-r_1)(x-r_2)(x-r_3)
\]
と書く。三つの根は相異なるとし、$d=r_3-r_1$、$k=x-r_3$ と置く。
有理座標
\[
 \xi=\frac{x-r_1}{x-r_3},\qquad \eta=\frac{y}{x-r_3}
\]
に移ると、Poisson 構造は
\[
 \left(F(\xi)\partial_\xi
     +(H_0(t)+\xi H_1(t))\partial_t\right)\wedge\partial_\eta
\]
となる。ただし
\begin{align*}
 F(\xi)&=-a\xi\bigl((r_3-r_2)\xi+r_2-r_1\bigr),\\
 H_0(t)&=-\frac{h_0(t)+r_1h_1(t)}d,
 &H_1(t)&=\frac{h_0(t)+r_3h_1(t)}d.
\end{align*}
実際、$\{\xi,\eta\}=f(x)\xi'(x)/k$、
$\{t,\eta\}=(h_0+xh_1)/k$ であり、
$1/k=(\xi-1)/d$ を代入すればこれらの式を得る。

一般の $h_0,h_1$ に対して $H_1(t)\partial_t$ は半単純である。
これは $H_1$ の二つの零点が異なるという Zariski 開条件であり、
$h_0+r_3h_1$ の係数は任意の二次式の係数を動けるので非空である。
その二つの零点を $s_1,s_2$ として
$H_1(t)=b(t-s_1)(t-s_2)$ と書き、
$\tau=(t-s_1)/(t-s_2)$ に移ると
\[
 H_1(t)\partial_t=\kappa\tau\partial_\tau,
 \qquad \kappa=b(s_1-s_2).
\]
$H_0(t)\partial_t$ は $\PP^1$ の大域ベクトル場なので、
同じ座標では二次以下の式 $A(\tau)\partial_\tau$ になる。
従って、座標を再び $(x,t,y)$ と書き直すと、
\begin{equation}\label{eq:riccati-quadratic-normal}
 \pi=\left(F(x)\partial_x+H(x,t)\partial_t\right)\wedge\partial_y,
 \qquad H(x,t)=A(t)+\kappa xt,
\end{equation}
を得る。$F,H$ は全次数が高々2である。
これらの正規化を族で行うときは、根を標識する有限 \'etale 被覆を用いる。

量子座標の変換も厳密に定義できる。
元の微分型 Ore 分数体では、左除算を用いて
\[
 \xi=\frac{x-r_1}{x-r_3},\qquad
 \eta=(x-r_3)^{-1}y
\]
と置く。係数体は可換なので、任意の $a\in\CC(x,t)$ に対して
\[
 [a,\eta]=\hbar k^{-1}\partial(a)
\]
である。続く $t$ の M\"obius 変換も係数体の同型なので、
\eqref{eq:riccati-quadratic-normal} の導分を用いた
Ore 分数体と元の分数体は $\CC(\hbar)$ 上で同型である。
逆写像は $x=(r_3\xi-r_1)/(\xi-1)$、$y=k\eta$ と、
$t$ の逆 M\"obius 変換によって与えられる。

$F,H$ の二次斉次化を
\[
 F^{\mathrm h}=Z^2F(X/Z),\qquad
 H^{\mathrm h}=Z^2A(T/Z)+\kappa XT
\]
とし、
\[
 \ell=-\frac{\partial_XF^{\mathrm h}+\partial_TH^{\mathrm h}}4,
 \qquad
 \Delta=(F^{\mathrm h}+X\ell)\partial_X
          +(H^{\mathrm h}+T\ell)\partial_T+Z\ell\partial_Z
\]
と置く。$\ell$ は一次式なので、Euler の恒等式から
\[
 \operatorname{div}(\Delta)
  =\partial_XF^{\mathrm h}+\partial_TH^{\mathrm h}+4\ell=0.
\]
全生成元の次数を1として
\[
 \mathcal A=\CC[\hbar,Z,X,T][V;-\hbar\Delta]
\]
を定める。これは homogeneous differential Ore extension であり、
PBW 基底 $Z^iX^jT^kV^\ell$ を持つ。
従って次数ごとに自由で、Hilbert 級数は $(1-s)^{-4}$ である。
各ファイバーは noetherian domain であり、
Ore 拡大の定理 \cite[定理0.2、\S1.2--1.3]{LiuWangWu}
によって大域次元4の AS-regular 代数となる。

$Z$ は正規元であり、
\[
 VZ=Z(V-\hbar\ell)
\]
となる。従って $Z$ のべきを Ore 局所化できる。
局所化の次数零部分において
\[
 x=X/Z,\qquad t=T/Z,\qquad y=Z^{-1}V
\]
と取ると、
\[
 \Delta(x)=ZF(x),\qquad \Delta(t)=ZH(x,t)
\]
なので
\[
 [x,t]=0,\qquad [x,y]=\hbar F(x),\qquad
 [t,y]=\hbar H(x,t).
\]
これらは次数零部分を生成する。
PBW 基底の $V$ 次数を比較すると追加の関係はなく、
次数零部分は正規化後の Riccati Ore 代数そのものである。

最後に $S(2,3)$ への所属を確認する。
アフィンでの1形式は $H\,dx-F\,dt$ であり、同次化は
\[
 \omega=ZH^{\mathrm h}\,dX-ZF^{\mathrm h}\,dT
            +(TF^{\mathrm h}-XH^{\mathrm h})\,dZ.
\]
これは $V$ を含まず、$\PP^2_{[Z:X:T]}$ の次数2葉層の
線形射影による引き戻しである。
一般の係数について $\gcd(F^{\mathrm h},H^{\mathrm h})=1$ であり、
$F=f_2x^2+f_1x$、$A=a_2t^2+a_1t+a_0$ と書くと
\[
 (TF^{\mathrm h}-XH^{\mathrm h})|_{Z=0}
    =XT\bigl((f_2-\kappa)X-a_2T\bigr)\ne0.
\]
従って $\omega$ の係数に共通因子がなく、葉層の次数は正確に2である。
これが主張した成分内の特別な族を与える。
\end{proof}

\section{三重積の対数型の非可換双有理変換}
\label{sec:cube-log}

三重積の三つの対数型成分の一般元は、因子を並べ替えた後、
\[
 \pi=(a x^2+b xz+c z^2)\partial_x\wedge\partial_z
       +\kappa y\partial_y\wedge(x\partial_x+z\partial_z)
\]
と正規化される。ここで $b^2-4ac\ne0$ を仮定する。
同梱の量子化原稿第4節では同変な cubic $\mathbb Z$-代数を用いた。
この節では分数体が直接計算できる別の斉次表示を構成する。

$B_R=R[Y,Z]$ とし、$R$ は必要な非零多項式を逆にした
$\mathbb C[\hbar]$ の局所化とする。
\[
 L(Y)=\tfrac b2Y+cZ,\qquad
 L(Z)=-aY-\tfrac b2Z,\qquad
 \sigma=1+\hbar L\quad\text{on }\langle Y,Z\rangle.
\]
$\det\sigma$ を逆にして $\sigma$ を $B_R$ の自己同型にし、$t=YZ$ と置く。
generalized Weyl algebra
\[
 G_R=B_R\langle X,W\rangle/
 \bigl(Xf-\sigma(f)X,\;Wf-\sigma^{-1}(f)W,\;
       WX-t,\;XW-\sigma(t)\bigr)
\]
の全生成元に次数1を与える。

\begin{proposition}\label{prop:gwa-quadric}
$G_R$ は次数ごとに自由な noetherian domain で、Hilbert 級数は
$(1+s)/(1-s)^3$ である。$\hbar=0$ では
$\mathbb C[X,Y,Z,W]/(XW-YZ)$ になる。
基底を $0$ の近傍に制限すると、各幾何学的ファイバーの射影圏は
smooth な非可換 quadric であり、その半古典極限は
$(a x^2+b xz+c z^2)\partial_x\wedge\partial_z$ である。
\end{proposition}
\begin{proof}
$G_R$ は $B_R[X^{\pm1};\sigma]$ に $W=tX^{-1}$ として埋め込まれ、
\[
 B_R\ \oplus\!\bigoplus_{n\ge1} B_RX^n
       \ \oplus\!\bigoplus_{n\ge1} B_RW^n
\]
が正規形となる。これが平坦性、domain 性、Hilbert 級数を与える。
さらに
\[
 S_R=B_R[X;\sigma][W;\tau,\delta],\quad
 \tau|_{B_R}=\sigma^{-1},\quad\tau(X)=X,\quad
 \delta|_{B_R}=0,\quad\delta(X)=t-\sigma(t)
\]
と置くと、$\tau,\delta$ は最初の Ore 関係を保存する。
$g=WX-t$ は中心的な正則二次元で、$G_R=S_R/(g)$ である。
各ファイバーの $S_R$ は四次元の homogeneous Ore AS-regular 代数である。
これは Ore 拡大の twisted Calabi--Yau 性と正の次数付けからも従う
\cite[定理0.2]{LiuWangWu}。
$S_b$ は $Y,Z,X,W$ の順序に関する quadratic PBW 代数なので Koszul であり、
中心的な正則二次元による商 $G_b=S_b/(g)$ も Koszul である。

smoothness を有限次元の条件で確認する。
Smith--Van den Bergh \cite[補題5.1、命題5.2、定理5.6]{SmithVdB}
により、$G_R$ の quadratic dual に付随する有限次元代数
\[
 C(G_R)=G_R^![w^{-1}]_0
\]
が半単純なら $\operatorname{qgr}G_R$ は smooth である。
ここで $w$ は $G_R^!$ の中心的な正則次数2の元であり、
$G_R^!/(w)=S_R^!$ である。
各ファイバーで $\dim(G_R^!)_4=\dim(G_R^!)_8=8$ で、
$w^2:(G_R^!)_4\to(G_R^!)_8$ は同型である。
従って基底を縮めれば、$C(G_R)$ は $(G_R^!)_4$ 上にこの同型で
積を戻した階数8の代数族として構成できる。
古典点では通常の smooth quadric の even Clifford 代数
$M_2(\mathbb C)\oplus M_2(\mathbb C)$ である。
正則表現の trace pairing の行列式が非零である開集合を取れば、
全ファイバーで半単純性と smoothness が従う。

半古典括弧は
$\{Y,Z\}=0$、$\{X,f\}=L(f)X$、$\{W,f\}=-L(f)W$ である。
$x=Y/X,z=Z/X$ とすると
\[
 \{x,z\}=z\bigl(\tfrac b2x+cz\bigr)
             -x\bigl(-ax-\tfrac b2z\bigr)
          =ax^2+bxz+cz^2.
\]
\end{proof}

重み $(X,Y,Z,W)=(0,1,1,2)$ の作用 $\rho$ は上の関係を保つ。
$p=1+\hbar\kappa$ をさらに $R$ 上で可逆にして、
\[
 H_R=G_R[Y_0][Y_1;\rho_p],\qquad
 \mathcal R_R=\bigoplus_{n\ge0}(H_R)_{n,n}
\]
と置く。$Y_0$ は中心的であり、$Y_i$ は第二次数1を持つ。
$\mathcal R_R$ は階数 $(n+1)^3$ の次数ごとに平坦な斉次環で、
古典点は $(\PP^1)^3$ の Segre 環である。
同変 Ore 拡大と Segre 積には補題\ref{lem:ore-projective-bundle}を適用できる。
各幾何学的ファイバー $G_b$ の dg smoothness と properness は、Orlov の分解
\cite[定理2.5]{OrlovSing} と Smith--Van den Bergh の
MCM 圏と $D^b(\operatorname{mod}C(G_b))$ の同値から従う。
実際、Gorenstein parameter は2なので二つの構造層 twist と
$C(G_b)$ の二つの単純加群が full exceptional collection を与える。
従って得られる三次元模型は次数ごとに平坦で、
各幾何学的ファイバーが dg smooth proper である。
アフィン座標 $y=Y_1/Y_0$ の交換子は混合項
$\kappa y\partial_y\wedge(x\partial_x+z\partial_z)$ を与える。

\begin{theorem}\label{thm:cube-log-torus}
上の射影量子化の division ring は、
一つの交換係数が1である三次元量子トーラスの division ring と同型である。
従って三重積の対数型の一般元には、トーラス型と非可換双有理な
射影量子化を選べる。
\end{theorem}
\begin{proof}
$L$ の固有一次式を $U,V$、固有値を $\nu,-\nu$ とする。
$\nu^2=(b^2-4ac)/4$ である。
$\sigma(U)=\lambda U,\sigma(V)=\mu V$、
$\lambda=1+\hbar\nu,\mu=1-\hbar\nu$ と置く。
次数零の分数体は $r=UX^{-1},s=VX^{-1}$ で生成され、
\[
 rs=\frac\lambda\mu sr,
 \qquad yr=p\,ry,\qquad ys=p\,sy.
\]
実際、$W=YZ X^{-1}$ なので $W$ は追加の分数体生成元を与えない。
$a_0=rs^{-1}$ と置けば
\[
 ya_0=a_0y,\qquad
 a_0s=\frac\lambda\mu sa_0,\qquad ys=p\,sy.
\]
これが所望の量子トーラス表示である。
命題\ref{prop:toric-common-field}により、同じ分数体を持つ
三模型のトーラス型 compactification が得られる。
\end{proof}

パラメーターを $q=\mu/\lambda$ と合わせると、
この分数体は第\ref{sec:line-conic-quantum}節の line--conic 型とも一致する。
この対応では量子化パラメーターの高次項も上式によって選んでいる。


\section{楕円型の曲面リンクと両立する量子化}
\label{sec:nc-surface-links}\label{sec:elliptic-links}

この節では、平面三次曲線と二重積の双二次曲線を結ぶ標準的な曲面リンクを扱う。
量子化の高次項および二周期的な表示を選択し、両側が同じ非可換関数体を持つ族を構成する。
本節の既知の曲面双有理定理は Presotto--Van den Bergh
\cite{PVdB2016}、Presotto \cite{Presotto2018}、
Rogalski--Sierra--Stafford \cite{RSS2023} による。
古典的ファイバーと一次括弧を指定する相対的な構成は、以下でこれらの定理から導く。

\subsection{楕円曲線上のパラメーターの照合}

滑らかな楕円曲線 $E$ 上の次数 $3$ の直線束 $L$ と平行移動 $\psi$ に付随する
三次元 quadratic Sklyanin 代数を $S(E,L,\psi)$ と書く。
$\tau=\psi^3\ne\mathrm{id}$ とし、$p,q\in E$ を選ぶ。
Presotto \cite[定理4.4、\S4.2、式(4.7)--(4.10)]{Presotto2018} の構成では
\[
 d_{2j}=\tau^{-j}p,\qquad d_{2j+1}=\tau^{-j}q,\qquad
 G_i=\psi^{i*}L\otimes\mathcal O_E(-d_i)
 \tag{S1}
\]
に付随する cubic AS-regular $\mathbb Z$-代数 $D$ と包含
\[
 D\hookrightarrow\check S
 \tag{S2}
\]
が得られ、この包含は非可換関数体の同型を誘導する。
ここで $\check S_{i,j}=S_{j-i}$ である。
使用する開条件には、隣接する $G_i,G_{i+1}$ が非同型であることを含める。
三次関係式は
\[
 R_i=\ker\!\left(
 H^0(G_i)\otimes H^0(G_{i+1})\otimes H^0(G_{i+2})
 \longrightarrow H^0(G_iG_{i+1}G_{i+2})\right)
 \tag{S3}
\]
で与えられる。
式 (S2) は点イデアルの積を用いて構成される。
その関数体上の全射性は、同論文 \S4.2 の Step 7 で証明されている。

平行移動 $\alpha$ を $\alpha^2=\tau$ となるように選べば
\[
 G_{i+2}\cong\alpha^*G_i
 \tag{S4}
\]
であり、$D$ は二周期的な表示を持つ。
従って even Veronese は通常の斉次代数 $Q$ として表せる。
正確には周期同型 $\phi:D_{i,j}\to D_{i+2,j+2}$ を選び、
$Q_n=D_{0,2n}$ に $a*b=a\phi^m(b)$、$a\in Q_m$、という積を入れる。
これにより $\check Q\cong D^{\mathrm{even}}$ である。
また
\[
 Q_0(Q)\cong\operatorname{Frac}(D^{\mathrm{even}})_{0,0}
       \cong\operatorname{Frac}(D)_{0,0}.
\]
最後の同型では、奇数終点 $N$ の分数 $ca^{-1}$ の分子分母に
同じ非零元 $b\in D_{N,N+1}$ を右から掛け、偶数終点の分数に直す。
以下の比較にはこの corner division ring の同定を用いる。
滑らかな楕円曲線と $\alpha^2\ne\mathrm{id}$ を持つ
一般の非可換 quadric から始める逆向きの構成は
\cite[定理4.1、\S4.1、式(4.4)]{Presotto2018} にあり、
平面側の平行移動は
\[
                 \psi^3=\alpha^2
                 \tag{S5}
\]
によって照合される。
特に、一周期的な cubic Sklyanin 代数の平行移動を $\sigma$ とする場合、
$\alpha=\sigma^2$ と取れるため、式 (S5) は
$\psi^3=\sigma^4$ となる。

\subsection{古典的な曲面リンクを持つ相対族}

\begin{proposition}\label{prop:relative-elliptic-surface-link}
$\sigma_f$ を、零曲線 $E=\{f=0\}\subset\mathbb P^2$ が滑らかな
平面 Poisson 構造とし、相異なる $p,q\in E$ を取る。
標準的なリンク
\[
 \operatorname{Bl}_{p,q}\mathbb P^2
 \cong \operatorname{Bl}_{r}(\mathbb P^1\times\mathbb P^1)
\]
で $\sigma_f$ から得られる Poisson 構造を $\sigma_D$ と書く。
滑らかな点付き代数曲線 $(B,0)$ 上で、両側を量子化する
次数ごとに平坦な斉次代数の族 $\mathcal S,\mathcal Q$ を選べる。
これらは
\[
 \mathcal S_0=\mathbb C[X,Y,Z],\qquad
 \mathcal Q_0=\bigoplus_{n\ge0}
 H^0(\mathbb P^1\times\mathbb P^1,\mathcal O(n,n))
\]
を満たし、一般の幾何学的ファイバーで
\[
 \operatorname{Frac}_0(\mathcal S_b)
 \cong \operatorname{Frac}_0(\mathcal Q_b)
 \tag{S6}
\]
となる。この同型は古典的な関数体のリンクを持ち上げる。
$\mathcal Q$ は cubic AS-regular $\mathbb Z$-代数の even Veronese として構成される。
\end{proposition}

\begin{proof}
まず $E$ を固定した Sklyanin 代数の族で $\sigma_f$ を量子化する。
この族の古典的な点と規格化は Hesse 表示で直接確認できる。
\[
 F=X^3+Y^3+Z^3-3\lambda XYZ,\qquad \lambda^3\ne1
\]
とし、点 $(b,h)=(-1,0)$ を持つ滑らかな代数曲線
\[
 1+b^3-h^3+3\lambda bh=0
 \tag{S7}
\]
のこの点の近傍上で
\[
 \mathcal S=
 \mathcal O_B\langle X,Y,Z\rangle/
 (YZ+bZY-hX^2,\ ZX+bXZ-hY^2,\ XY+bYX-hZ^2)
 \tag{S8}
\]
を取る。式 (S7) の $b$ による偏微分は古典点で $3$ なので、$h$ は局所パラメーターであり、
$b=-1+\lambda h+O(h^2)$ である。
従ってアフィン錐上の半古典括弧は
\[
 \{Y,Z\}=X^2-\lambda YZ=\tfrac13F_X
\]
およびその巡回置換である。射影化した括弧は $\sigma_F$ となる。
座標変換と $h$ の一次係数の調整で $\sigma_f$ を得る。

式 (S8) の point-module 行列の行列式は $bhF$ である。
従って $h\ne0$ の一般のファイバーの楕円曲線は同じ $E=\{F=0\}$ である。
$E$ 上でこの行列は古典点でも階数 $2$ であり、その核が定める自己同型
$\psi_b$ は代数的に $\psi_0=\mathrm{id}_E$ へ延びる。
基底を連結な近傍に取ると $\psi_b$ は平行移動である。
この Sklyanin 表示に対する AS 正則性の非退化条件を用いて基底を縮めると、$\mathcal S$ は
各次数の階数が $\binom{n+2}{2}$ の族となる。

$p,q$ をこの固定された楕円曲線上の定数切断として、式 (S1) を相対的に定義する。
$\tau_b=\psi_b^3$ と置く。
乗法 $[2]\colon E\to E$ は \'etale なので、基底の \'etale 近傍への置換により
\[
 \alpha_b^2=\tau_b,\qquad\alpha_0=\mathrm{id}_E
\]
を満たす平行移動の族を選べる。
式 (S4) は相対直線束について成立する。
必要な直線束の同型とその基底方向の自明化を選び、
二周期的な生成元・関係式の族として表示する。

次にこの族の平坦性を有限の条件で確認する。
$G_i$ の次数は $2$ なので $H^0(G_i)$ は階数 $2$ のベクトル束である。
古典点では
\[
 G_{2j}=L(-p),\qquad G_{2j+1}=L(-q),\qquad L=\mathcal O_E(1).
 \tag{S9}
\]
$p\ne q$ よりこれら二つは非同型である。
隣接する二直線束の非同型性を保つ開集合では、base-point-free pencil trick により
式 (S3) の写像は階数 $6$ の全射となる。
従って $R_i$ は階数 $2$ のベクトル束であり、基底変換と可換である。

四重テンソル中の
\[
 (R_i\otimes H^0(G_{i+3}))\cap(H^0(G_i)\otimes R_{i+1})
 \tag{S10}
\]
は古典点で一次元である。
一般ファイバーでは Presotto の cubic AS-regular 代数の四次 syzygy なので
同じく一次元である。従って交わりの次元の上半連続性から、
古典点の近傍で次元が一に等しく、基底変換と可換な直線束となる。
古典点のテンソルは linear quadric の geometric quintuple である。
geometric 条件は隣接二因子による縮約の非消滅で表される開条件なので、
\cite[\S5.2、定義5.3.1、定理5.5.10]{VdBQuadrics} によって、
基底を一度縮めれば相対 cubic AS-regular $\mathbb Z$-代数 $\mathcal D$ が得られる。
二周期性により確認すべきテンソルは有限個である。
このため全次数で同じ基底開集合を用いることができ、
$\mathcal D_{i,i+2n}$ の階数は $(n+1)^2$、
$\mathcal D_{i,i+2n+1}$ の階数は $(n+1)(n+2)$ となる。

古典点の関係式も明示する。
$p=[0:1:0]$, $q=[0:0:1]$ とする座標を取ると、
偶数位置と奇数位置の生成元はそれぞれ
\[
 V_p=\langle X,Z\rangle,\qquad V_q=\langle X,Y\rangle
\]
である。三次関係式は同じ偶奇の二因子を交換する関係であり、linear quadric を与える。
even Veronese の次数 $n$ は
\begin{align*}
 \mathcal D_{0,2n}|_0
 &=V_p^nV_q^n\\
 &=\operatorname{span}\{X^{2n-i-j}Y^iZ^j\mid0\le i,j\le n\}\\
 &=H^0(\mathbb P^2,I_p^nI_q^n(2n)).
 \tag{S11}
\end{align*}
従って古典的な even Veronese は、$p,q$ を結ぶ直線を収縮する
線形系 $|2H-E_p-E_q|$ による滑らかな quadric の切断環である。

式 (S2) の包含を相対化する。
$H^0(G_i)$ は、$\psi_b^{i*}L$ の切断として、
$\check{\mathcal S}_{i,i+1}$ の階数 $2$ の部分束に入る。
生成元のこの写像が三次関係式を零に送ることは、一般のファイバーでは
\cite[\S4.2]{Presotto2018} の包含から従う。
これは有限階数自由加群間の写像の恒等式なので、古典点を含めて成立する。
得られる $\mathcal D\to\check{\mathcal S}$ は、各斉次部分について
一般ファイバー上で単射である。始域は基底上平坦なので、全体でも単射である。
一般の幾何学的ファイバーの関数体の同型は同論文の定理4.4による。
この同型は基底の関数体 $K=\mathbb C(B)$ 上でも成立する。
実際、包含 $\mathcal D_K\hookrightarrow\check{\mathcal S}_K$ は $K$ 上定義されている。
元 $v$ を $\check{\mathcal S}_K$ の次数零関数体に取る。
$\overline K$ 上の全射性と Ore 分数表示から、ある $N$ と
$a,c\in (\mathcal D_{\overline K})_{0,N}$ が存在して
\[
                  va=c,\qquad a\ne0
\]
となる。この方程式は有限次元 $K$-空間
$\mathcal D_{0,N}\oplus\mathcal D_{0,N}$ 上の線形方程式である。
標的を、$v\mathcal D_{0,N}$ と $\mathcal D_{0,N}$ が張る有限次元空間に制限すればよい。
幾何学的に domain であるため、分数環はこの体拡大に関して単射を保つ。
従って解空間の $a$ 成分への射影が非零であることは $K$ に下降し、
$K$ 上でも $a\ne0$ なる解が得られる。
よって $v=ca^{-1}$ であり、$K$ 上の全射性が従う。
二周期性を用いた even Veronese を $\mathcal Q$ と置けば式 (S6) を得る。

古典点での包含は式 (S11) の可換な包含そのものである。
周期同型 $\phi$ は基底上線形で、古典点でも同型なので $h$-付値を保つ。
また式 (S11) の各斉次部分の包含は古典点でも単射であるため、
原始元を原始元へ送り、corner 分数体の包含も $h$-付値を保つ。
従って $Q_0(\mathcal Q)$ と $Q_0(\mathcal S)$ の同型は
Poisson 標識を保ち、剰余体上で古典的なリンクを回復する。
交換子を $h$ で割って得る一次括弧は、この剰余体の同型によって一致する。
従って $\mathcal Q$ の半古典極限は指定された $\sigma_D$ である。
\end{proof}

\begin{corollary}\label{cor:threefold-elliptic-surface-link}
$\mathbb P^1\times\mathbb P^2$ の垂直型成分と、
$(\mathbb P^1)^3$ の各二因子型成分の一般元について、
第\ref{subsec:classical-surface-link}節の曲面リンクと両立する射影量子化を選べる。
非可換関数体の同型は
\[
 \operatorname{Frac}_0(\mathcal S_b)(u)
 \cong \operatorname{Frac}_0(\mathcal Q_b)(u)
 \tag{S12}
\]
で与えられ、$u$ は中央な超越元である。
\end{corollary}

\begin{proof}
前命題の二つの族と $\mathcal O_B[U_0,U_1]$ の Segre 積を取る。
各次数の階数はそれぞれ $(n+1)\binom{n+2}{2}$ と $(n+1)^3$ となる。
次数零関数体では $u=U_1/U_0$ が中央であり、式 (S6) が式 (S12) を誘導する。
古典的なリンクの両側で第三因子の括弧は零なので、一次括弧も所望のものとなる。
\end{proof}

\subsection{共通ブローアップについての既知の定理}

非可換曲面が、四次元 Sklyanin 代数の中央二次式による商
$Q=Q_{\mathrm{VdB}}(E,\rho,\Omega)$ として与えられる場合には、
共通ブローアップの環まで既知である。
$\rho$ が無限位数、$\operatorname{qgr}Q$ が滑らかであると仮定し、
$T'=Q^{(2)}$、$\tau=\rho^2$ と置く。
点 $r\in E$ を通る二つの ruling の line module が楕円曲線上に定める
残余点を $p,q$ とする。
Rogalski--Sierra--Stafford \cite[Basic Facts 7.2、定理7.3]{RSS2023} は
\[
 T'(r)\cong S(E,\psi)^{(3)}(\tau^{-1}p+\tau^{-1}q),
 \qquad \psi^3=\tau
 \tag{S13}
\]
を証明している。
括弧内は elliptic algebra の点ブローアップを表す。
従って、この表示を持つ模型については一方の一点ブローアップと
他方の二点ブローアップが同じ環になり、標準的な曲面リンクの
非可換版が得られる。

式 (S13) は上記の仮定を持つ幾何学的ファイバーについての既知の定理である。
命題\ref{prop:relative-elliptic-surface-link} は別途、古典点を含む
両立した量子化と関数体の同型を構成している。
第\ref{sec:relative-common-blowup}節では、二点 evaluation と例外対象の持上げから、
古典点の上まで延びる共通 smooth proper 圏と二つの収縮を構成する。
そこから通常の斉次 quadric 環も復元し、同じ一次 Poisson 括弧を得る。
第\ref{sec:anticanonical-common-ring}節では、共通圏の反標準 orbit 環を
式 (S13) の平面側の部分環へ同定する。
その section functor からの共通圏の復元と、quadric 側の特定の反標準表示との
標識付き比較には、coherent heart と周期性の同型をさらに照合する。
固定された一般の Sklyanin quadric に対する式 (S13) の環と
categorical blow-up の比較は第\ref{blowup:quadric-comparison}節で証明する。

\section{\texorpdfstring{$\mathbb P^3$}{P3} の量子化と関数体}
\label{sec:p3-quantum-models}\label{sec:p3-models}

Pym \cite[定理1.1、表1、第3節]{Pym2015} は、四変数多項式環の
次数付き Calabi--Yau 変形の六成分と一般元の関係式を記述した。
対応する射影 Poisson 成分は
$L(1,1,1,1)$、$L(1,1,2)$、$R(2,2)$、$R(1,3)$、$S(2,3)$、$E(3)$ である。
それぞれ skew-polynomial 環、二次の Ore 拡大、四次元 Sklyanin 代数、
三次元 Sklyanin 代数の中心拡大、可換多項式環の微分 Ore 拡大、
および二次元非可換 Lie 代数に対する CGG 積によって量子化される。
以下では $Q_0(A)$ を次数付き分数環の次数零部分とする。

\subsection{引き戻し成分の微分 Ore 表示}

$F_i\in\mathbb C[x_1,x_2,x_3]_2$ とし、
$\delta=\sum_{i=1}^3F_i\partial_{x_i}$ と置く。
Pym の $S(2,3)$ 型は $\operatorname{div}\delta=0$ の場合に対応し、
量子化は
\[
 A_h=\mathbb C[h,x_1,x_2,x_3][x_0;h\delta]
\]
で与えられる。$x_0,x_1,x_2,x_3$ の次数を1とする。
可換係数環の非零元を Ore 局所化し、
\[
 v=x_1/x_3,\qquad w=x_2/x_3,\qquad T=x_3^{-1}x_0,
 \qquad f_i(v,w)=F_i(v,w,1)
\]
と置けば
\[
 Q_0(A_h)=\operatorname{Frac}\bigl(\mathbb C(h)(v,w)[T;hD]\bigr),
 \qquad
 D=(f_1-vf_3)\partial_v+(f_2-wf_3)\partial_w.
\]
実際、$[T,v]=h x_3^{-1}\delta(v)=h(f_1-vf_3)$ であり、$w$ も同様である。
局所化した次数付き環は、この次数零の環に $x_3^{\pm1}$ を加えて得られる。
共役は $x_3Tx_3^{-1}=T-hf_3$、$x_3v=vx_3$、$x_3w=wx_3$ なので、
その次数零の分数環は表示した微分 Ore 環の分数環に一致する。

\subsection{例外成分の Weyl 関数体}
\label{subsec:exceptional-weyl}

Pym \cite[第3.6節、補題3.3]{Pym2015} の二つの導分を
\[
 X=-\frac54x_0\partial_{x_0}-\frac14x_1\partial_{x_1}
       +\frac34x_2\partial_{x_2}+\frac74x_3\partial_{x_3},
 \qquad
 Y=4(x_0\partial_{x_1}+x_1\partial_{x_2}+x_2\partial_{x_3})
\]
とする。$[Y,X]=Y$ であり、積を
\begin{equation}\label{eq:exceptional-cgg}
 f\star g=\sum_{k\ge0}h^kY^k(f)\binom{X}{k}(g),
 \qquad \binom Xk=\frac{X(X-1)\cdots(X-k+1)}{k!}
\end{equation}
で定める。この積の一次の交換子を半古典括弧の符号規約とする。
$Y$ は局所冪零だから、この式は各 $f,g$ に対して有限和となる。

\begin{proposition}\label{prop:exceptional-weyl}
式\eqref{eq:exceptional-cgg}による例外型量子化の次数零の関数体は、
$h\ne0$ の各ファイバーにおいて
\[
 \operatorname{Frac}\bigl(\mathbb C(c)\langle r,P\rangle/(Pr-rP-h)\bigr)
\]
と同型である。従ってその中心は $\mathbb C(c)$ である。
\end{proposition}
\begin{proof}
$Y(x_0)=0$ であり、$x_0$ は $X$ の固有ベクトルだから、$x_0$ は正規正則元である。
具体的に $a=-5/4$ と置くと
$f\star x_0=x_0\star(1+hY)^a f$ であり、$(1+hY)^a$ は局所冪零性により
可逆な線形写像である。従って $x_0$ を Ore 局所化できる。
同じ式は可換環 $\mathbb C[x_0^{\pm1},x_1,x_2,x_3]$ 上で有限和のまま定義され、
その次数零部分の可換座標 $z_i=x_i/x_0$ 上の導分は
\[
 \overline X=z_1\partial_{z_1}+2z_2\partial_{z_2}+3z_3\partial_{z_3},
 \qquad
 \overline Y=4(\partial_{z_1}+z_1\partial_{z_2}+z_2\partial_{z_3})
\]
となる。ここで $z_i$ は積を変形する前のベクトル空間の座標である。
\[
 s=\frac{z_1}{4},\qquad
 p=z_2-\frac{z_1^2}{2},\qquad
 q=z_3-z_1z_2+\frac{z_1^3}{3}
\]
と置くと、可換環として $\mathbb C[z_1,z_2,z_3]=\mathbb C[s,p,q]$ であり、
\[
 \overline Y=\partial_s,\qquad
 \overline X=s\partial_s+2p\partial_p+3q\partial_q.
\]
従って式\eqref{eq:exceptional-cgg}から厳密に
\[
 [s,p]_\star=2hp,\qquad [s,q]_\star=3hq,\qquad [p,q]_\star=0
\]
を得る。$\star$ 冪 $s^{\star k}$ は通常の $s^k$ と、$s$ に関する低次項の和である。
従って $p^iq^js^{\star k}$ は基底をなし、チャートの代数は
\[
 \mathbb C[h,p,q][s;h(2p\partial_p+3q\partial_q)]
\]
に同型である。

分数環で
\[
 r=q/p,\qquad c=q^2/p^3,\qquad P=r^{-1}s
\]
と置く。$p,q$ は可換であり、$\mathbb C(p,q)=\mathbb C(c,r)$、
逆変換は $p=r^2/c$、$q=r^3/c$ である。
導分 $2p\partial_p+3q\partial_q$ は $r\partial_r$ になるから、
$[P,r]=h$ かつ $c$ は中心的である。逆に $s=rP$ により全分数環が復元される。
標数零の Weyl 分数環の中心はその係数体に一致するので、最後の主張も従う。
\end{proof}

この構成の半古典極限は $\{P,r\}=1$、$c$ は Casimir である。
従って例外成分の一般元は、$\mathbb P^1_c\times\mathbb P^2$ の
三次式 $X_0^3$ に対応する垂直型の量子化と非可換双有理同値になる。
これは垂直成分内の特別な退化との対応である。

\subsection{Sklyanin pencil の中心部分体を拡大した比較}
\label{subsec:sklyanin-central-extension}

$S_4=S_4(E,\alpha)$ を四次元 Sklyanin 代数とし、
$\alpha$ は無限位数の平行移動とする。
以下の中央二次式と quadric pencil の性質は
\cite[\S10.1--10.4、定理10.2]{SmithVdB} による。
$\Omega_1,\Omega_2\in Z(S_4)_2$ を一次独立に取り、
\[
 u=\Omega_1\Omega_2^{-1},\qquad K=\mathbb C(u),\qquad
 Q_u=(S_4\otimes_{\mathbb C}K)/(\Omega_1-u\Omega_2)
\]
と置く。$u$ は超越的であり、$K$ は $Q_0(S_4)$ の中心部分体である。
実際、二つの中央二次式の同次な代数関係があれば、
複素数上の二変数斉次式の分解と domain 性から一次従属を生じる。
$Q_0(Q_u)\simeq Q_0(S_4)$ が成り立つ。
実際、$\Omega_2$ を逆にすると $u=\Omega_1\Omega_2^{-1}$ による同一視ができ、
残りは中心的な非零多項式の局所化になる。

\begin{proposition}\label{prop:sklyanin-finite-base-extension}
上の設定で、有限拡大 $K'/K$ と、$\mathbb C$ 上で定義された平行移動
$\sigma$ で $\sigma^3=\alpha^2$ を満たすものが存在して
\[
 Q_0(Q_u\otimes_KK')
 \simeq
 Q_0(S_3(E,\sigma)\otimes_{\mathbb C}K')
\]
が成り立つ。
\end{proposition}
\begin{proof}
$\overline K$ 上で $Q_u$ は滑らかな Van den Bergh quadric である。
これは pencil の特異な quadric が四つであることから従う
\cite[定理10.2]{SmithVdB}。
Rogalski--Sierra--Stafford \cite[定理1.2、定理7.3]{RSS2023} により、
$Q_u^{(2)}$ の一点ブローアップは、ある $S_3(E,\sigma)^{(3)}$ の二点ブローアップと同型であり、
$\sigma^3=\alpha^2$ である。従って $\overline K$ 上で表示した二つの関数体は同型になる。
$\alpha$ は定数で、$[3]\colon E\to E$ のファイバーは $\mathbb C$ 上有限であるから、
$\sigma$ は $\mathbb C$ 上で選べる。次数3の直線束の選択も、$\overline K$ 上の平行移動により
固定したものに置き換えられる。

両側の $2$ 次および $3$ 次 Veronese は次数1で生成され、次数1の中心元を持つ。
$Q_u$ 側ではその中心元を $z=\Omega_2$ と取り、
$(Q_u^{(2)}[z^{-1}])_0$ は $Q_u^{(2)}$ の次数1の生成元を $z$ で割った元で生成される。
$S_3^{(3)}$ 側も次数3の中心元で同様である。
従って両方の次数零部分は有限生成であり、それぞれの分数環が対象の関数体となる。
上の同型と逆同型による有限個の生成元の像を分数で書くと、現れる係数は有限個である。
それらは $K$ 上代数的なので、全てを含む有限拡大 $K'/K$ を取れる。
同型と逆同型の等式もその拡大上で成立し、結論を得る。
\end{proof}

命題\ref{prop:sklyanin-finite-base-extension}は中心部分体の有限拡大後の主張である。
$K$ 上の同型を得るための下降は残る課題である。
特に二つの ruling の交換と、得られる三次元 Sklyanin 代数の形式を同時に調べる必要がある。
固定した $\Omega$ の各ファイバーの比較を $K$ 上で同時化することが、ここでの下降問題である。


\part{楕円型の共通圏と反標準環}
\section{二つの収縮を持つ相対的な非可換共通模型}
\label{sec:relative-common-blowup}

本節では有限次元代数の変形を用いて、平面の二点 blow-up と
quadric の一点 blow-up を同時に変形する。
基底の \'etale 近傍への置換を許す。
古典的ファイバーと一次の Poisson 方向を固定し、
高次の量子化係数を比較に適合するように選ぶという本稿の方針に沿った構成である。
得られる共通模型は有限型基底上の smooth proper dg 圏であり、
二つの contraction は相対的な fully faithful functor と SOD として与える。

\subsection{有限次元 directed 代数と例外集の変形}

以下で $R$ は滑らかな複素アフィン曲線 $U$ の座標環、$0\in U$ は閉点とする。
$R$-線形な有限対象圏 $\mathsf A$ が directed であるとは、
対象を $T_1,\ldots,T_N$ と順序付けて
\[
 \operatorname{Hom}_{\mathsf A}(T_j,T_i)=0\quad(j>i),\qquad
 \operatorname{End}_{\mathsf A}(T_i)=R
\]
と書けることをいう。
全ての Hom は有限局所自由 $R$-加群と仮定する。
同じ記号をその有限次元代数にも用い、右加群の perfect 圏を取る。

\begin{lemma}\label{lem:directed-relative-smooth}
上の $\operatorname{Perf}(\mathsf A)$ は $R$ 上 smooth proper である。
この性質は任意の基底変換と両立する。
\end{lemma}

\begin{proof}
$S=\bigoplus_iRe_i$、$J=\bigoplus_{i<j}\operatorname{Hom}(T_i,T_j)$ と置く。
$\mathsf A=S\oplus J$ は分裂した $S$-代数である。
正規化された相対 bar 分解は
\[
 \cdots\longrightarrow
 \mathsf A\otimes_SJ^{\otimes_S2}\otimes_S\mathsf A
 \longrightarrow\mathsf A\otimes_SJ\otimes_S\mathsf A
 \longrightarrow\mathsf A\otimes_S\mathsf A
 \longrightarrow\mathsf A
\]
となる。$J^{\otimes_SN}=0$ なので有限長である。
各項は $\mathsf A$-両側加群として有限生成射影加群の直和因子であり、
対角両側加群が perfect であることが従う。
また $\mathsf A$ 自体が $R$ 上有限局所自由なので proper である。
分解は $R$ 上分裂する正規化 bar 分解から得られ、基底変換後も同じ議論が成り立つ。
\end{proof}

\begin{lemma}\label{lem:exceptional-etale-lift}
$\mathsf A$ を上の directed $R$-代数とする。
$\operatorname{Perf}(\mathsf A_0)$ の有限個の exceptional objects は、
$0$ の共通 \'etale 近傍上に同時に持ち上がる。
指定した対象に既に持上げがある場合、その持上げを固定してよい。

さらに $(F_{1,0},\ldots,F_{m,0})$ が full strong exceptional collection ならば、
基底を縮めた後、その持上げ $(F_1,\ldots,F_m)$ は
相対的な full strong exceptional collection となる。
対応する SOD と Morita 同値は任意の基底変換と両立する。
\end{lemma}

\begin{proof}
各対象を、$e_i\mathsf A_0$ の有限直和からなる有界射影複体 $P_0$ で表す。
その次数付き加群の部分は、$e_i\mathsf A$ の同じ直和により持ち上げる。
局所パラメーターを $h$ とし、微分 $d$ を $h$-進的に構成する。
$h^{n+1}$ を法として $d^2=0$ が成り立つ場合、次の段階の障害は
\[
 H^2\operatorname{Hom}^{\bullet}_{\mathsf A_0}(P_0,P_0)
 =\operatorname{Ext}^2(F_0,F_0)
\]
に属する。これは零である。
具体的には、$h^{n+1}$ で割った $d^2$ の剰余が $[d_0,-]$ の二次 cocycle であり、
その原始元を一次の微分に加えて消せばよい。
従って有限長の同じ次数付き射影加群上に形式的な微分が得られる。

有限個の Hom 加群を局所自明化すると、$d^2=0$ は有限個の行列係数に関する
有限個の多項式方程式である。
Artin 近似 \cite[定理16.13.2]{RelativeArtin} により、
古典的な微分を保つ解が $0$ の \'etale 近傍上に存在する。
有限個の対象について同じ操作を同時に行う。
この障害計算は一般の対象の変形理論 \cite{RelativeLowen} の、
有限射影複体に対する場合でもある。

持ち上げた二対象の $\operatorname{RHom}$ は $R$ 上 perfect である。
古典的ファイバーで次数零にのみコホモロジーを持つならば、
近傍を縮めると次数零の有限局所自由加群になる。
特殊ファイバーで零となる Hom 複体は、同じく近傍で零となる。
従って strongness、半直交性、および相対的な exceptional 性が保たれる。
特に
\[
 \mathsf C=\operatorname{End}\!\left(\bigoplus_{i=1}^mF_i\right)
\]
は finite locally free directed $R$-代数になる。

生成性も確認する。
$T=\bigoplus_iF_i$ とし、$T$ が定める fully faithful functor
$\operatorname{Perf}(\mathsf C)\to\operatorname{Perf}(\mathsf A)$ を取る。
補題\ref{lem:directed-relative-smooth} により $\mathsf C$ は相対的に smooth である。
このため右随伴 $\operatorname{RHom}_{\mathsf A}(T,-)$ は perfect 圏に制限できる。
自由生成元 $\mathsf A$ における counit の cone は $R$ 上 perfect であり、
古典的ファイバーでは $T_0$ の生成性によって零である。
そのコホモロジーの台を除けば cone は零となる。
従って $\mathsf A$ は $T$ の生成する thick 部分圏に入り、fullness が従う。
用いた Hom 同型、bar 分解、unit と counit の同型はすべて基底変換を保つ。
\end{proof}

\begin{remark}
smooth proper なスキームの族における SOD の \'etale 局所変形は
Belmans--Okawa--Ricolfi \cite{RelativeBOR} により一般的に扱われている。
上の補題では directed 代数と full strong exceptional collection に対象を絞り、
必要な変形と生成性を有限個の行列で証明した。
\end{remark}

\subsection{二点の evaluation による共通圏}

$\mathcal S$ を命題\ref{prop:relative-elliptic-surface-link} の平面量子化の族とする。
その Beilinson 代数を
\[
 A=\operatorname{End}\bigl(\mathcal O\oplus\mathcal O(1)\oplus\mathcal O(2)\bigr)
\]
と書く。
有限表示としては三対象 $T_0,T_1,T_2$ と
$\operatorname{Hom}(T_i,T_j)=\mathcal S_{j-i}$ $(i\le j)$ を用いる。
$R$ 上の階数は $3+3+3+6=15$ である。

固定された滑らかな楕円曲線を $E$ とし、その相異なる二点を $p,q$ とする。
点に対応する右 $A$-加群を $M_p,M_q$ と書く。
各頂点成分は階数一であり、古典点では
\[
 (M_p)_i=\operatorname{Hom}_{\mathbb P^2}(\mathcal O(i),k_p),
 \qquad
 (M_q)_i=\operatorname{Hom}_{\mathbb P^2}(\mathcal O(i),k_q)
\]
である。
後述する $E$ への restriction による evaluation を使えば、
これらを基底全体で代数的に定義できる。

五対象 $T_0,T_1,T_2,P_p,P_q$ の $R$-線形圏 $B$ を
\[
 \begin{gathered}
 \operatorname{Hom}_B(T_i,P_\nu)=(M_\nu)_i,
 \qquad \operatorname{End}_B(P_\nu)=R,\\
 \operatorname{Hom}_B(P_\nu,T_i)=0,
 \qquad \operatorname{Hom}_B(P_p,P_q)=\operatorname{Hom}_B(P_q,P_p)=0
 \end{gathered}
\]
によって定める。
最初の三対象の間の Hom と積は $A$ のものであり、
$T_i\to T_j\to P_\nu$ の合成は $M_\nu$ の右作用による。
$B$ は階数 $15+6+2=23$ の directed 代数である。
その perfect 圏を $\mathcal H=\operatorname{Perf}(B)$ と置く。

\begin{theorem}\label{thm:relative-two-contractions}
上の族について、$0$ の \'etale 近傍上で次が成り立つ。
\begin{enumerate}
\item
$\mathcal H$ は $R$ 上 smooth proper であり、
\[
 \mathcal H=
 \langle\operatorname{Perf}(A),\operatorname{Perf}(R),\operatorname{Perf}(R)\rangle.
 \label{eq:relative-plane-sod}\tag{RB1}
\]
古典的ファイバーは
$Y=\operatorname{Bl}_{p,q}\mathbb P^2$ の perfect 圏である。
\item
$Y\to\mathbb P^1\times\mathbb P^1$ を、$p,q$ を結ぶ直線の狭義変換を収縮する射とする。
その例外曲線を $E_r$ と書く。
$\mathcal H$ には、この側の pullback と $\mathcal O_{E_r}$ を持ち上げる
相対的な full strong exceptional collection
\[
                 (Q_0,Q_1,Q_2,Q_3,\mathcal E_r)
\]
が存在する。$Q_0$ は最初の模型の $T_0=\mathcal O$ として固定できる。
$C=\operatorname{End}(Q_0\oplus\cdots\oplus Q_3)$、
$\mathcal C_{\mathrm{qd}}=\langle Q_0,Q_1,Q_2,Q_3\rangle$ とすると
\[
 \mathcal H=
 \langle\mathcal C_{\mathrm{qd}},\operatorname{Perf}(R)\rangle,
 \qquad \mathcal C_{\mathrm{qd}}\simeq\operatorname{Perf}(C).
 \label{eq:relative-quadric-sod}\tag{RB2}
\]
これらの SOD は任意の基底変換と両立する。
\item
$C$ の階数は24であり、四対象間の Hom の階数は
\[
 \begin{pmatrix}
 1&2&4&6\\
 0&1&2&4\\
 0&0&1&2\\
 0&0&0&1
 \end{pmatrix}.
 \label{eq:relative-quadric-ranks}\tag{RB3}
\]
各 $\operatorname{Hom}(Q_i,\mathcal E_r)$ の階数は一である。
\end{enumerate}
\end{theorem}

\begin{proof}
古典的ファイバーで $f:Y\to\mathbb P^2$ とすると、
Orlov の右側表示は
\[
 \operatorname{Perf}(Y)
 =\langle f^*\operatorname{Perf}(\mathbb P^2),
          \mathcal O_{E_p},\mathcal O_{E_q}\rangle
\]
である。
$\operatorname{RHom}(f^*\mathcal O(i),\mathcal O_{E_\nu})$
は点 $\nu$ での evaluation の一次元空間であり、
逆向きの Hom は $H^*(\mathbb P^1,\mathcal O(-1))=0$ により零となる。
従って
\[
 (f^*\mathcal O,f^*\mathcal O(1),f^*\mathcal O(2),
       \mathcal O_{E_p},\mathcal O_{E_q})
\]
は full strong exceptional collection で、その End は $B_0$ である。
一般の基底上の式 (RB1) は directed gluing の構成から従い、
smooth properness は補題\ref{lem:directed-relative-smooth} による。

次に $H=f^*\mathcal O_{\mathbb P^2}(1)$、$a=H-E_p$、$b=H-E_q$ と置く。
quadric 側の pullback と例外対象は
\[
 \mathcal O_Y,\quad\mathcal O_Y(a),\quad\mathcal O_Y(a+b),\quad
 \mathcal O_Y(2a+b),\quad\mathcal O_{E_r}
 \label{eq:relative-classical-collection}\tag{RB4}
\]
である。
quadric 上の最初の四対象は strong であり、
Euler 完全列
\[
 0\longrightarrow\mathcal O(0,1)
 \longrightarrow\mathcal O(1,1)^{\oplus2}
 \longrightarrow\mathcal O(2,1)\longrightarrow0
\]
を用いると矩形 Beilinson collection も生成するため full である。
右側 Orlov 表示により、式 (RB4) 全体が $\operatorname{Perf}(Y)$ の
full strong exceptional collection となる。
補題\ref{lem:exceptional-etale-lift} によってこれを同時に持ち上げる。
古典的な Hom は quadric 上の切断の次元であり、式 (RB3) と
$\operatorname{Hom}(Q_i,\mathcal E_r)$ の階数一を与える。
持上げ後も同じ階数となり、式 (RB2) が従う。
\end{proof}

\subsection{四対象の代数から射影 quadric を再構成する}

\begin{proposition}\label{prop:relative-quadric-reconstruction}
定理\ref{thm:relative-two-contractions} の $C$ は、基底を縮めると
cubic AS-regular $\mathbb Z$-代数の族 $\mathcal D$ の四頂点部分代数として実現できる。
各斉次部分は $R$ 上有限局所自由であり、全ての幾何学的ファイバーが
三次元 cubic AS-regular $\mathbb Z$-代数となる。
各ファイバーでは
\[
 \operatorname{Perf}(C_{\bar b})
 \simeq D^b(\operatorname{qgr}\mathcal D_{\bar b}).
\]
古典的ファイバーは linear quadric である。
\end{proposition}

\begin{proof}
右加群の向きを明記する。
$0\le i\le j\le3$ に対し
\[
 D_{i,j}=\operatorname{Hom}(Q_{3-j},Q_{3-i})
\]
と置けば、通常の End の積は $D_{i,j}D_{j,k}\subset D_{i,k}$ を満たす。
この四頂点代数は、冪等元の順序を逆にした $C$ である。
$V_i=D_{i,i+1}$ $(i=0,1,2)$ と置く。
長さ二の乗法は古典点で同型なので近傍でも同型であり、
長さ三の乗法は全射である。
従って
\[
 R_0=\ker(V_0\otimes V_1\otimes V_2\longrightarrow D_{0,3})
\]
は階数二の部分束である。
$V_3=R_0^\vee$ とし、包含 $R_0\hookrightarrow V_0\otimes V_1\otimes V_2$
が定める coevaluation tensor を
\[
 w\in V_0\otimes V_1\otimes V_2\otimes V_3
\]
と書く。

古典点では $w$ は linear quadric の geometric quintuple である。
geometric 条件の失敗は、隣接二因子の射影双対束の積上で
縮約が零となる閉集合の像として表せる。
この射は proper なので、geometric 条件は基底上で開である。
従って一度基底を縮めると全てのファイバーでこの条件が成り立つ。
Van den Bergh \cite[定理5.2.2、系5.2.3、定理5.5.10]{VdBQuadrics} により
この quintuple は quadric を定める。
具体的には $V_{i+4}=V_i$、$w_i$ を $w$ の巡回置換とし、
$w_i$ の最後の因子での縮約像を三次関係式とする。
従って構成は有限個の代数的データからなり、相対的に定義できる。

全ての幾何学的ファイバーで
\[
 \dim D_{i,i+2n}=(n+1)^2,\qquad
 \dim D_{i,i+2n+1}=(n+1)(n+2)
\]
となる。
各部分は有限表示 $R$-加群であり、この一定次元性から局所自由性が従う。
基底を次数ごとに取り替える必要はない。
四個の標準対象の Hom は四頂点部分代数であり、
cubic AS-regular 分解 \cite[定理5.5.7]{VdBQuadrics} により全ての twist を生成する。
従ってその collection は full strong で、上の導来同値を得る。
右加群の標準射影加群を $P_i=e_iD$ と書くと、
$(\pi P_3,\pi P_2,\pi P_1,\pi P_0)$ が $(Q_0,Q_1,Q_2,Q_3)$ に対応する。
\end{proof}

\begin{lemma}\label{lem:relative-helix-basechange}
四対象 $Q_0,\ldots,Q_3$ の geometric helix $(Q_i)_{i\in\mathbb Z}$ は、
同じ基底上の perfect objects として取れる。
$j\ge i$ に対して $\operatorname{RHom}(Q_i,Q_j)$ は次数零の
有限局所自由 $R$-加群であり、任意の基底変換と可換する。
$j-i=2n$ のとき階数は $(n+1)^2$、$j-i=2n+1$ のとき階数は
$(n+1)(n+2)$ である。
従って各ファイバーでの全ての正向き Hom は、同一の代数的 helix の特殊化から得られる。
\end{lemma}

\begin{proof}
$\operatorname{Perf}(C)$ の相対 Serre functor を $S_{C/R}$ とする。
$C$ は相対的に smooth proper なので、その Serre 両側加群と逆両側加群は perfect であり、
構成は基底変換と両立する。
$Q_{i+4}=S_{C/R}^{-1}(Q_i)[2]$ によって全ての整数へ延長する。
各幾何学的ファイバーでは前命題の cubic AS-regular 模型と同定できる。
その AS-Gorenstein 次数は4であり、最小分解とその双対
\cite[定理5.5.7]{VdBQuadrics} から、この延長は標準的な
$(\pi P_{3-i})_{i\in\mathbb Z}$ に一致する。
同じ分解で計算すると、正向きの Hom の高次 Ext は零で、
次数零の Hom は対応する $D_{3-j,3-i}$ である。
従って全てのファイバーにおいて、主張した階数とコホモロジーの消滅が成立する。

相対 Hom は $R$ 上 perfect である。
その全ての幾何学的ファイバーが、同じ階数を持つ次数零のベクトル空間であることから、
相対 Hom 自体が有限局所自由加群であり、任意の基底変換と可換する。
この議論は既に全ファイバーで AS 正則性を確保した同じ基底上で成り立つため、
添字ごとの追加の縮小を必要としない。
\end{proof}

\subsection{固定楕円曲線への restriction と相対的な even Veronese}

固定楕円曲線を保持する平面の族では、上の $\mathbb Z$-代数を
通常の次数付き環へ相対的に移すこともできる。
まず evaluation 加群の規約を固定する。
twisted homogeneous coordinate ring の積を
$u*v=u\,\psi^{m*}v$ $(\deg u=m)$ と書く。
$L=\mathcal O_E(1)$ とし
\[
 L_n=\bigotimes_{a=0}^{n-1}(\psi^a)^*L,\qquad
 N_i^{\mathrm{pl}}=(\psi^{-i})^*L_i
 =\bigotimes_{a=1}^{i}(\psi^{-a})^*L
 \quad (0\le i\le2)
\]
と置く。
平面の矢 $u:T_i\to T_j$ を
$(\psi^{-j})^*u\in H^0(E,N_j^{\mathrm{pl}}(N_i^{\mathrm{pl}})^{-1})$
に送る。
実際、合成 $v\circ u$ に対応する $\mathcal S$ の積は $vu$ であり、
\[
 (\psi^{-k})^*(v\,\psi^{(k-j)*}u)
 =(\psi^{-k})^*v\,(\psi^{-j})^*u
\]
なので、これは合成を保つ。
この restriction の evaluation により
$M_\nu=\operatorname{Hom}_E(\bigoplus_iN_i^{\mathrm{pl}},k_\nu)$
$(\nu=p,q)$ と定義する。

\begin{proposition}\label{prop:relative-even-rollup}
上の evaluation により $B$ を構成する。
定理\ref{thm:relative-two-contractions} の対象は、
$0$ の \'etale 近傍上で、次の性質を持つように取れる。
\begin{enumerate}
\item
$R$-線形な dg functor
\[
 \mathcal R_E:\mathcal H\longrightarrow\operatorname{Perf}(E\times U)
\]
があり、古典点では $E\hookrightarrow Y$ への導来引戻しである。
$N_i=\mathcal R_E(Q_i)$ は直線束となる。
\item
$G_i=N_i^{-1}N_{i+1}$ $(i=0,1,2)$ は次数二の直線束であり、
\[
 (G_0,G_1,G_2)|_0=(L(-p),L(-q),L(-p)).
\]
四対象の全ての正向き Hom は
\[
 \operatorname{Hom}(Q_i,Q_j)
 \xrightarrow{\ \sim\ }H^0(E,N_i^{-1}N_j)
 \qquad(0\le i<j\le3)
 \label{eq:relative-restriction-hom}\tag{RB5}
\]
によって記述される。
\item
同じ四頂点代数を持つ $\mathcal D$ を相対的に二周期とすることができる。
従って $\mathcal Q_n=\mathcal D_{0,2n}$ は通常の次数付き環の族となり、
\[
 \operatorname{rank}_R\mathcal Q_n=(n+1)^2,\qquad
 \mathcal Q\otimes_R k(0)
 =\bigoplus_{n\ge0}H^0(\mathbb P^1\times\mathbb P^1,\mathcal O(n,n))
\]
を満たす。
\end{enumerate}
\end{proposition}

\begin{proof}
$T_i$ を $N_i^{\mathrm{pl}}$、$P_p,P_q$ をそれぞれ $k_p,k_q$ へ送り、
矢を上の restriction と evaluation に送る。
これは有限対象圏 $B$ から perfect complexes への厳密な作用を定める。
対応する両側加群による導来 tensor が $\mathcal R_E$ である。
古典点では $E$ と $E_p,E_q$ は横断的に交わり、
$L i^*\mathcal O_{E_\nu}=k_\nu$ となるため、$\mathcal R_E$ は $L i^*$ に一致する。

各 $\mathcal R_E(Q_i)$ の古典的ファイバーは直線束である。
perfect complex が階数一の局所自由加群になる条件は開であり、
$E\times U\to U$ は proper なので、基底を縮めると $N_i$ はすべて直線束になる。
その次数から $\deg G_i=2$ を得る。
古典的な値は式 (RB4) を $E$ に制限すればよい。

式 (RB5) を古典点で確認する。
quadric 上で必要な直線束は $\mathcal O(1,0)$、$\mathcal O(0,1)$、
$\mathcal O(1,1)$、$\mathcal O(2,1)$ である。
$E\in|\mathcal O(2,2)|$ なので、restriction の核を与える差の直線束は
少なくとも一方の次数が $-1$ である。
その $H^0,H^1$ は零であり、restriction は同型となる。
相対的には両辺が同じ階数の局所自由加群なので、同型性は近傍で保たれる。
合成との両立性により、四頂点の三次関係式は
$G_0,G_1,G_2$ の切断の乗法の核である。

右加群の座標順序に合わせ
$H_0=G_2,H_1=G_1,H_2=G_0$ と書く。
$H_0\not\cong H_1$ と $H_1\not\cong H_2$ は古典点で成り立ち、近傍でも保たれる。
$\Delta=H_2H_0^{-1}\in\operatorname{Pic}^0(E\times U/U)$ と置くと $\Delta_0=\mathcal O_E$ である。
楕円曲線の原点を固定し、平行移動 $\alpha$ による次数二の直線束の変化を用いる。
対応する写像は $[-2]:E\to\operatorname{Pic}^0(E)$ であり、標数零では \'etale である。
従って基底を \'etale 局所化して
\[
 \alpha_0=\mathrm{id},\qquad
 \alpha^*H_0\cong H_2,\qquad
 \alpha^*H_1\cong H_1\otimes\Delta
\]
を得る。
relative Picard 上の等式と実際の直線束の同型の差は基底からの直線束であり、
有限個のこれらの因子は基底上で局所自明化できる。

全ての整数 $m$ に対し
\[
 H_{2m}=(\alpha^m)^*H_0,\qquad
 H_{2m+1}=(\alpha^m)^*H_1
\]
と定義し、$V_i=H^0(E,H_i)$ および
\[
 R_i=\ker\bigl(V_i\otimes V_{i+1}\otimes V_{i+2}
                   \longrightarrow H^0(E,H_iH_{i+1}H_{i+2})\bigr)
\]
から $\mathcal D$ を構成する。
隣接する二直線束は非同型なので、pencil trick により各 $R_i$ は階数二である。
$H_0\not\cong H_2$ のファイバーでは、これは Van den Bergh の admissible quadruple による構成である。
$H_0\cong H_2$ のファイバーでは、三次関係式は同じ偶奇の二因子を交換する
linear quadric の関係式となる。
従って全てのファイバーで AS 正則であり、前命題と同じ四頂点部分代数を持つ。

$\alpha^*$ は全ての乗法写像と可換なので、
相対的な二周期同型 $\phi:\mathcal D_{i,j}\to\mathcal D_{i+2,j+2}$ を与える。
これは無限個の同型を独立に選ぶ構成ではなく、単一の $\alpha$ の反復による。
$a\in\mathcal Q_m$、$b\in\mathcal Q_n$ に
\[
                         a*b=a\phi^m(b)
\]
と積を定める。
その Hilbert 階数と古典的ファイバーは linear quadric の計算から従う。
この相対的な議論は \cite[命題5.6.1、系5.6.2]{VdBQuadrics} の
楕円曲線の二分点を選ぶ議論を、古典点での恒等平行移動を指定して行ったものである。
\end{proof}

\begin{proposition}\label{prop:relative-elliptic-chart}
命題\ref{prop:relative-even-rollup} の族では、基底を縮めると
二次の正規 regular 元 $g\in\mathcal Q_2$ と全射
\[
 \mathcal Q\longrightarrow T(E,N,\alpha),\qquad N=H_0\otimes H_1,
 \qquad \mathcal Q/(g)\cong T(E,N,\alpha)
\]
が存在する。ここで $\deg N=4$ であり、右辺は相対 twisted homogeneous coordinate ring である。
$g|_0$ は quadric 上の $E\in|\mathcal O(2,2)|$ の方程式となる。

$K=\operatorname{Frac}(R)$ とすると、$\mathcal Q_K$ は左右 noetherian domain であり、
非零斉次元は左右 Ore 集合をなす。
固定楕円曲線の補集合に対応する affine chart は
\[
 \mathfrak U_K=(\mathcal Q_K[g^{-1}])_0
\]
である。その関数斜体は
\[
 \operatorname{Frac}(\mathfrak U_K)=Q_0(\mathcal Q_K)
\]
となる。
\end{proposition}

\begin{proof}
切断の乗法から、各次数において
\[
 \mathcal Q_n\longrightarrow
 H^0\!\left(E,N\otimes\alpha^*N\otimes\cdots
                    \otimes(\alpha^{n-1})^*N\right)
\]
を得る。
楕円型ファイバーでは \cite[定理5.5.6]{VdBQuadrics} により全射であり、
その核は $\mathbb Z$-次数4の正規 regular 元により左右から生成される。
線形ファイバーでは $H_0\cong H_2$ である。
このとき $\alpha$ は両方の次数二の直線束を保ち、
二つの線形系を通じて quadric の自己同型に延びる。
even rollup は Segre 環のこの自己同型による twist として表せる。
restriction の核は $E$ の方程式によって生成され、同じ正規性と regular 性を満たす。
従って相対的な写像も全ての次数で全射となる。

二次の核は階数 $9-8=1$ の直線束である。
基底上でこれを自明化して生成元 $g$ を取る。
各次数の核も有限局所自由であり、
$g$ による左乗法と右乗法は、対応する核への写像として
全ての幾何学的ファイバーで同型である。
従って基底全体で同型となり、
\[
 \ker(\mathcal Q\longrightarrow T)=g\mathcal Q=\mathcal Qg
\]
および $g$ の正規 regular 性が得られる。
ここでも次数ごとの基底の縮小は用いない。

各幾何学的ファイバーの noetherian 性は
\cite[系5.5.9、補題3.5]{VdBQuadrics} の noetherian 性と Veronese 同値から従う。
これを generic fiber に適用し、体拡大から降下させると
$\mathcal Q_K$ は左右 noetherian である。
$T(E_K,N,\alpha)$ は整な楕円曲線の twisted section ring なので domain である。
$g$ の正規 regular 性を用い、非零斉次元から割り切れる限り $g$ を取り除くと、
次数が有限であるため有限回で止まる。
$g$ を法とした積が非零であることから、元の積も非零となる。
従って $\mathcal Q_K$ も domain である。
左右 noetherian domain に対する Goldie の定理を斉次元に適用すれば Ore 性を得る。

$g$ の冪による局所化は、その正規性から定義できる。
第二 Veronese に移って $g$ の次数を一とすると、局所化した環は strongly graded であり、
その次数零環が $\mathfrak U_K$ となる。
最後の斜体の等式は直接確認できる。
同次数の非零斉次元による分数 $ab^{-1}$ を取り、
必要なら分子と分母に同じ非零一次斉次元を右から掛け、共通次数を $2n$ とする。
すると
\[
 ab^{-1}=(ag^{-n})(bg^{-n})^{-1},\qquad
 ag^{-n},bg^{-n}\in\mathfrak U_K
\]
である。
反対向きの包含は局所化の定義から従う。
\end{proof}

\begin{remark}
式 (RB1)--(RB2) は、古典点を含む代数的な共通 smooth proper 圏と二つの SOD を与える。
この構成で得た共通圏と RSS の共通 blow-up 環の qgr との同定、
および選択済みの特定の高次量子化係数との一致は、それぞれ比較すべき追加の問題である。
本節で構成した二つの射影模型の一次 Poisson 括弧と generic localization の同定は、
以下で扱う。
\end{remark}

\subsection{共通関数斜体と半古典極限}
\label{subsec:relative-common-generic}

定理\ref{thm:relative-two-contractions}で得られる二つの contraction は、
generic な局所化とも両立する。この点を、古典点における有限複体の計算から証明する。

\begin{lemma}\label{lem:valued-complex-cancellation}
$D$ を離散付値 $v:D^\times\to\mathbb Z$ を持つ斜体とし、
$\mathcal V=\{a:v(a)\geq0\}$、$\mathfrak m=\{a:v(a)>0\}$、
$\kappa=\mathcal V/\mathfrak m$ と置く。
$P^\bullet$ を有限階数自由右 $\mathcal V$-加群の有界複体とする。
その $\kappa$ 上の還元が acyclic なら、$P^\bullet\otimes_{\mathcal V}D$ も acyclic である。
還元の cohomology が次数零の一次元空間だけなら、
$P^\bullet\otimes_{\mathcal V}D\simeq D[0]$ である。
後者の複体二つの間の射が還元の cohomology 上で非零なら、
その $D$ 上の射も非零である。
\end{lemma}
\begin{proof}
還元された微分の非零成分を一つ選ぶと、その lift は $\mathcal V$ の単元である。
行と列の基本変形により、この成分を1とし、同じ行と列の残りを零にする。
複体条件 $d^2=0$ から、前後の微分の対応成分も零となる。
従って $\mathcal V\xrightarrow{1}\mathcal V$ という可縮な直和因子を除ける。
これを有限回繰り返す。還元が acyclic なら全ての項が消え、
cohomology が一つの一次元空間なら次数零の $\mathcal V$ 一個が残る。
これらの基本変形を射にも施すと、最後の主張は還元が非零のスカラーの可逆性から従う。
この証明に $\mathcal V$ の完備性は不要である。
\end{proof}

\begin{lemma}\label{lem:normal-projective-localization}
$G$ を本節の平面または quadric の次数付き環とし、
$g$ を楕円曲線を定める斉次正規非零因子とする。
$U=(G[g^{-1}])_0$ と置くと、制限 functor は
\[
 \operatorname{Perf}(\operatorname{qgr}G)/\ker(-|_U)
 \longrightarrow\operatorname{Perf}(U)
\]
の冪等完備化として同値を与える。
\end{lemma}
\begin{proof}
$g$ の次数の Veronese を取ってよいので、$\deg g=1$ とする。
$G[g^{-1}]$ は強次数付きであり、その次数付き加群圏は $U$ の加群圏と同値である。
正規元の Ore 局所化は平坦かつ完全であり、核は $g$-冪 torsion である。
有限生成加群と有界複体に必要な分母払いを確認する。
$G[g^{-1}]$ 上の有界な有限生成次数付き複体で、各項の有限個の斉次生成元と
その微分の像を取ると、微分で保たれる有限生成 $G$-格子が得られる。
複体条件により像をもう一度追加する必要はない。
有限個の chain maps と homotopies についても、その像を格子に加えれば同時に持ち上がる。
格子の拡大は局所化後に同型になる。
従って分母を可逆にする roof を用いれば、有界導来圏の局所化は
$g$-冪 torsion cohomology を持つ複体による Verdier 商となる。
有限次元 graded torsion はこの核に含まれ、十分高次の truncation により
同じ議論を $\operatorname{qgr}G$ に移せる。

本節の射影圏は full strong exceptional collection を持つ smooth な圏なので、
その有界 coherent 導来圏は perfect 圏である。
Veronese 後の ample 列 $\{\pi G(n)\}_{n\in\mathbb Z}$ はこの圏を thick に生成する。
実際、coherent 圏の有限 homological dimension を $d$ とし、ample 性で
$0\to K\to L_d\to\cdots\to L_0\to F\to0$ を取る。
各 $L_i$ は有限個の $\pi G(n)$ の直和とできる。
対応する三角 $K[d]\to[L_d\to\cdots\to L_0]\to F\to K[d+1]$ は
$\operatorname{Ext}^{d+1}(F,K)=0$ により分裂するので、$F$ はこの有限複体の直和因子である。
全ての twist $\pi G(n)$ は perfect であり、局所化後には自由階数一となる。
従って上の quotient の対象は $U$ 上 perfect である。
逆に $U$ 自身が像に入り、その有限複体と直和因子は全ての
$\operatorname{Perf}(U)$ を生成する。
これで主張の冪等完備化を得る。
\end{proof}

\begin{proposition}\label{prop:relative-common-field}
定理\ref{thm:relative-two-contractions}の基底曲線の関数体を $K$ とし、
平面の量子化を $\mathcal S_K$、再構成した quadric の even roll-up を $\mathcal Q_K$ とする。
二点 gluing 代数を $B$、その圏を $\mathcal H=\operatorname{Perf}(B)$ と書く。
このとき
\[
 D=Q_{\rm gr}(\mathcal S_K)_0
 \cong Q_{\rm gr}(\mathcal Q_K)_0
 \tag{RG1}
\]
であり、$\mathcal H_K$ と両端点に共通の generic localization
$\mathcal H_K\to\operatorname{Perf}(D)$ を指定できる。
二つの SOD の全ての例外対象は、この局所化で零となる。
同型 (RG1) の古典的剰余は、指定した平面と quadric の曲面リンクである。
\end{proposition}
\begin{proof}
基底の古典点での局所環を $R$、その一様化元を $h$ とする。
$\mathcal S_R$ は各次数で自由であり、$\mathcal S_R/h\mathcal S_R$
は可換多項式環である。従って $h$ の位数は積に関して加法的であり、
Ore 分数環へ延長する。次数零部分 $D$ の剰余斜体は
$\CC(\mathbb P^2)$ である。各有理関数は斉次分数で lift できるため、
この剰余の同定は全射である。

平面の Beilinson 代数を $A$ とする。
三角代数の射 $B\to A$ と平面の generic localization の合成を
\[
 F:\operatorname{Perf}(B_K)\longrightarrow\operatorname{Perf}(D)
 \tag{RG2}
\]
と置く。$B\to A$ による導来拡大は二つの追加点の例外成分を核とする商である。
実際、三角代数の標準的な射影加群の三角により
$\operatorname{Perf}(B)/\langle E_p,E_q\rangle\simeq\operatorname{Perf}(A)$
を直接得る。
平面側では中央三次式を可逆化した次数零環が noetherian Ore domain となる。
補題\ref{lem:normal-projective-localization}と命題\ref{prop:generic-perf}から、
$F$ も核による Verdier 商の冪等完備化を与える。

平面の三つの line bundle を、一つの非零な一次式を使って generic point 上で自明化する。
古典点でこの一次式が零でないように選ぶと、$A_R$ の射の像は全て $D$ の付値環に入る。
$B_R$ の二つの追加頂点は零へ送る。
したがって $B_R$ 上の有限射影複体の像は、付値環上の有限自由複体として計算できる。

第三の例外対象 $E_r$ の古典的ファイバーは、
$Y=\operatorname{Bl}_{p,q}\mathbb P^2$ の曲線 $H-E_p-E_q$ の構造層である。
その generic restriction は零なので、補題\ref{lem:valued-complex-cancellation}より
$F(E_r)=0$ となる。
同じ補題により、quadric の四つの対象 $Q_i$ とその geometric helix の各対象は
$D[0]$ へ送られる。これらの古典的ファイバーは全て line bundle だからである。

第二の SOD を $\mathcal H_K=\langle\mathcal C_{\rm qd},E_r\rangle$ と書く。
$E_r\in\ker F$ なので、SOD の射影三角を $F$ に適用すると、
その制限 $F|_{\mathcal C_{\rm qd}}$ も同じ $\operatorname{Perf}(D)$ への Verdier 局所化を与える。
ここでは各対象の $E_r$ 成分を零とする商を先に取ってよい。

この商を quadric の内在的な関数斜体と照合する。
helix の二対象間の非零射を $h$ のべきで規格化し、その還元を非零にする。
全 helix の各 Hom は有限自由 $R$-加群で、その古典点への基底変換は
古典的な Hom と同型である。従って primitive な射の還元は非零である。
古典的な line bundle 間の非零射は
有理関数体上でも非零である。
補題\ref{lem:valued-complex-cancellation}より、その像も $D$ 上非零となる。
従って全ての非零斉次射が $D$ 上で可逆となり、
quadric の Ore corner division ring $D_{\rm qd}$ から $D$ への埋込みを得る。
even roll-up を用いれば $D_{\rm qd}=Q_{\rm gr}(\mathcal Q_K)_0$ である。

quadric 側にも、楕円曲線を定める正規元の局所化による affine Ore chart がある。
従ってその canonical generic functor は、補題\ref{lem:normal-projective-localization}と
命題\ref{prop:generic-perf}により $\operatorname{Perf}(D_{\rm qd})$ への Verdier 局所化である。
上の制限 functor は
\[
 \mathcal C_{\rm qd}\longrightarrow\operatorname{Perf}(D_{\rm qd})
 \xrightarrow{-\otimes_{D_{\rm qd}}D}\operatorname{Perf}(D)
 \tag{RG3}
\]
と分解する。斜体拡大による右の functor は conservative である。
従って左の functor と合成の核は一致する。
両者がその核による Verdier 商の冪等完備化であることから、右の functor は同値である。
自由一次元加群の自己準同型環を比較して $D_{\rm qd}\cong D$ を得る。

各 helix 対象の像を付値環上で階数一に簡約すると、primitive な斉次射の像は
付値零となり、その剰余は古典的な切断である。
斉次分数に対して付値は分子と分母の付値の差なので、
同型 (RG1) は $h$-付値を保つ。
全ての自明化と有限複体の還元は、古典的な共通開集合への制限と両立する。
従って (RG1) の剰余は $k(Y)=k(\mathbb P^2)=k(\mathbb P^1\times\mathbb P^1)$
の指定した同一視となる。
\end{proof}

\begin{proposition}\label{prop:relative-common-poisson}
上で再構成した quadric の一次 Poisson 構造は、
$\sigma_f$ を古典的な曲面リンクで移した $\sigma_D$ である。
\end{proposition}
\begin{proof}
命題\ref{prop:relative-common-field}の同型 $\iota:D_{\rm qd}\to D$ は
$K$-線形であり、$h$-付値と古典的な曲面リンクを保つ。
付値零の元 $a,b$ に対して
\[
 \overline{h^{-1}[\iota(a),\iota(b)]}
 =\bar\iota\!\left(\overline{h^{-1}[a,b]}\right)
 \tag{RG4}
\]
となる。左辺は選んだ平面量子化の Poisson 括弧である。
右辺は quadric の一次括弧を同じ関数体へ移したものである。
従って二つの括弧は共通の稠密開集合上で一致する。
quadric の一次変形は
$HH^2(\mathbb P^1\times\mathbb P^1)=H^0(\mathcal O(2,2))$
によって正則 bivector を定める。
正則 bivector は稠密開集合への制限で一意に定まるので、主張を得る。

共通圏の一次変形にも同じ記述が成り立つ。
$Y=\operatorname{Bl}_{p,q}\mathbb P^2$ では $H^2(\mathcal O_Y)=0$ であり、
$\operatorname{PGL}_3$ の相異なる順序付き二点への推移的作用と
blow-up の接束完全列から $H^1(T_Y)=0$ となる。
従って $HH^2(Y)=H^0(Y,\wedge^2T_Y)$ である。
有限代数の包含 $A\subset B$ と古典的な共通開集合 $W$ への制限について、
\[
 \begin{array}{ccc}
 HH^2(Y)&\longrightarrow&HH^2(\mathbb P^2)\\
 \downarrow&&\downarrow\\
 HH^2(W)&=&HH^2(W)
 \end{array}
\]
は可換となる。実際、上の矢は full subcategory への Hochschild cochain の制限であり、
下への矢は同じ tilting bimodule を $W$ へ制限して得られる。
$B\to A$ の二つの追加点の成分は $W$ 上で零なので、両経路の bimodule は一致する。
従って $B_R$ の一次クラスは、$\sigma_f$ の $Y$ への Poisson lift である。
\end{proof}

\begin{corollary}\label{cor:relative-common-threefold}
$X_{12}$ の垂直楕円型と $X_{111}$ の二因子楕円型について、
古典的な曲面リンクを持つ射影量子化の代数的な族と、
その共通の smooth proper dg 圏 $\mathcal Y$ を選べる。
基底の \'etale 近傍上で
\[
 \begin{aligned}
 \mathcal Y
 &=\left\langle
 \mathcal X_{12},\operatorname{Perf}(\mathbb P^1_U),
 \operatorname{Perf}(\mathbb P^1_U)\right\rangle,\\
 \mathcal Y
 &=\left\langle
 \mathcal X_{111},\operatorname{Perf}(\mathbb P^1_U)\right\rangle
 \end{aligned}
 \tag{RG5}
\]
という相対 admissible SOD があり、基底変換と両立する。
generic な共通関数斜体は $D(u)$ であり、$u$ は中央な超越元である。
古典的ファイバーは
$\operatorname{Perf}(Y\times\mathbb P^1)$ で、full exceptional collection の長さは10である。
\end{corollary}
\begin{proof}
$\mathcal Y=\mathcal H\otimes\operatorname{Perf}(\mathbb P^1_U)$ とする。
定理\ref{thm:relative-two-contractions}の二つの SOD に
$\operatorname{Perf}(\mathbb P^1_U)$ を掛けると (RG5) を得る。
各射影模型は平面側および quadric 側の斉次環と
$R[U_0,U_1]$ の対角 Segre 積で与えられる。
その次数ごとの階数は $(n+1)\binom{n+2}{2}$ と $(n+1)^3$ である。
補題\ref{lem:segre-field}により、関数斜体は両方とも $D(u)$ となる。
一次 Poisson 構造は命題\ref{prop:relative-common-poisson}で照合される。
共通圏は $B_R$ と $\mathbb P^1$ の Beilinson 代数のテンソル積でも表示でき、
この代数は階数 $23\cdot4=92$、頂点数10を持つ。
\end{proof}

この構成では二点 $p,q$、基底の \'etale 近傍、周期同型を選んでいる。
得られる通常の斉次 quadric 環と第\ref{sec:nc-surface-links}節の特定の表示の
高次係数の一致は、別途そのパラメーターを照合して判断する。
定理\ref{thm:relative-two-contractions}と上の系は、
指定した古典的リンクと一次括弧を持つ量子化を両側で選ぶという存在主張である。

\section{共通圏の反標準環と楕円 blow-up 環}
\label{sec:anticanonical-common-ring}

本節では、階数23の directed 代数 $B$ から構成した共通圏に、
全次数で平坦な反標準切断環を付随させる。
楕円曲線への restriction が与える Serre 三角を用いるため、
次数を増やすたびに基底を縮める操作を必要としない。
この構成は、非可換 del Pezzo 曲面上の rigid sheaf から中心拡大を作る
Rains の構成 \cite[命題5.1]{PencilRainsFiltered} の、
有限次元代数による相対的な実現を与える。
ここでの相対的な証明と明示的な二点 blow-up 環との比較については、
先行研究との新規性の比較を継続する。

\subsection{restriction の Serre 三角}

$R=\mathcal O(U)$、$\mathcal H=\operatorname{Perf}(B)$、
$O=T_0$ とする。$U$ は古典点 $0$ を含む滑らかな連結アフィン曲線である。
本節で基底を縮める操作は、$0$ を含む Zariski 開近傍への置換を意味する。
固定楕円曲線 $E$ 上の次数3の線束 $L$ と平行移動 $\psi$ を用い、
\[
 N_0=\mathcal O_E,\qquad
 N_1=(\psi^{-1})^*L,\qquad
 N_2=(\psi^{-1})^*L\otimes(\psi^{-2})^*L
\]
と置く。相異なる定数切断 $p,q$ と合わせて
\[
 \mathcal P=N_0\oplus N_1\oplus N_2\oplus k_p\oplus k_q
 \quad\text{on }E_U=E\times U
\]
を考える。既に定義した evaluation による $B$-作用を用いて
\[
 r=-\otimes_B^{\mathbf L}\mathcal P\colon
 \mathcal H\longrightarrow\operatorname{Perf}(E_U)
\]
とする。$\operatorname{Perf}(E_U)$ の Serre 関手は $[1]$ である。
この同一視には $H^1(E,\mathcal O_E)\simeq\mathbb C$ の自明化を固定する。

\begin{lemma}\label{lem:ac-serre-triangle}
$\mathcal H$ の相対 Serre 関手を $\mathbb S$ とし、
$\theta=\mathbb S[-2]$ と置く。
自然な完全三角
\begin{equation}\label{eq:ac-serre-triangle}
 \theta\xrightarrow{\alpha}\mathrm{id}_{\mathcal H}
 \longrightarrow r^Rr\longrightarrow\theta[1]
\end{equation}
が存在する。この三角と $\theta$ は基底変換と両立する。
基底を一度縮めると、$r$ は相対 spherical functor となる。
その twist を
\[
 T=\operatorname{Cone}(rr^R\longrightarrow\mathrm{id}),\qquad
 F=\theta^{-1}=\mathbb S^{-1}[2],\qquad
 \Psi=T^{-1}[2]
\]
と書けば、自然同型
\begin{equation}\label{eq:ac-intertwining}
                 rF\simeq\Psi r
\end{equation}
が得られる。
\end{lemma}

\begin{proof}
$B$-両側加群の perfect 複体
\[
 D=\mathbf R\Gamma(E_U,
       \mathbf R\mathcal Hom_{E_U}(\mathcal P,\mathcal P))
\]
を考える。$B\to D$ は $B$ の作用から得られる。
線束の次数は $0,3,6$ であり、$p,q$ は相異なるので、
\[
 H^0(D)=B,\qquad H^1(D)=B^\vee,
 \qquad H^i(D)=0\ (i\ne0,1).
\]
最初の等式は、最初の三対象間では
$\mathcal S_d\simeq H^0(E,L_d)$ が $d=0,1,2$ で成立することと、
線束から点への evaluation による。
二番目の等式は楕円曲線上の相対 Serre 双対性であり、
積に関する trace の巡回性により両側加群の同型である。
従って標準的な truncation により
\[
 \operatorname{Cone}(B\longrightarrow D)\simeq B^\vee[-1]
\]
を得る。$\mathbb S$ は $-\otimes_B^{\mathbf L}B^\vee$ であるから、
これを一つ回転した三角が \eqref{eq:ac-serre-triangle} である。

spherical 性について用いる開性を具体的に述べる。
$r$ とその左右随伴は perfect な $B$--$E_U$ 両側加群で表される。
左随伴は $r^L=\mathbb S^{-1}r^R[1]$ である。
Anno--Logvinenko の spherical 判定
\cite[定理1.1および定理5.1]{ACAnnoLogvinenko} に現れる比較射
\[
                  r^R\longrightarrow\theta r^L[1]
\]
の cone も、この種の perfect 両側加群で表される。
古典点では $\mathcal H_0=\operatorname{Perf}(Y)$、
$Y=\operatorname{Bl}_{p,q}\mathbb P^2$ であり、
$r_0=\mathbf L i^*$ は反標準楕円曲線 $i\colon E\hookrightarrow Y$ への restriction
である。このとき比較射は同型である。
cone を $E_U$ 上の perfect 複体として見ると、そのコホモロジーの台は閉集合であり、
$E_U\to U$ は proper である。その像を除けば比較射は同型になる。
$\theta$ は既に可逆なので、同判定を適用できる。

twist の kernel も明示できる。
$p_1,p_2\colon E_U\times_U E_U\to E_U$ とすると、$rr^R$ は
\[
 p_1^*(\mathcal P^\vee)\otimes_B^{\mathbf L}p_2^*\mathcal P
\]
を kernel とする。$B$ の相対 smoothness によりこの kernel は perfect である。
評価射から対角の構造層 $\mathcal O_\Delta$ への射を取り、その cone が $T$ の kernel
となる。逆関手の kernel はその随伴 kernel で与えられる。
spherical 関手の標準的な同型 $Tr\simeq r\theta[2]$ から
\eqref{eq:ac-intertwining} が従う。
以上の構成は tensor、双対、proper pushforward と unit・counit によるため、
基底変換と両立する。
\end{proof}

\begin{lemma}\label{lem:ac-elliptic-automorphism}
基底をさらに一度縮めると、$\Psi$ は
\[
                 \Psi(G)=\rho_*(G\otimes M)
\]
と書ける。ここで $\rho$ は $E_U/U$ の平行移動、$M$ は各ファイバーで次数7の線束である。
さらに、基底からの線束の自明化を選ぶと
\begin{equation}\label{eq:ac-elliptic-data}
 \rho=\tau=\psi^3,\qquad
 M=L_3(-\tau^{-1}p-\tau^{-1}q),\qquad
 L_3=L\otimes\psi^*L\otimes(\psi^2)^*L
\end{equation}
となる。
\end{lemma}

\begin{proof}
古典点の Cartier divisor の計算により
\[
 T_0(G)=G\otimes\mathcal O_E(-E)[2],\qquad
 \Psi_0(G)=G\otimes\mathcal O_E(E).
\]
$E^2=K_Y^2=7$ である。
前補題の $\Psi$ の perfect kernel を $K$ とする。
その古典的ファイバーは $\Delta_*\mathcal O_E(E)$ である。
derived flatness の開性を $K$ の第一射影に適用すると、
$E\times\{0\}$ の近傍で、$K$ は第一射影に関して平坦な一つの coherent sheaf となり、
その各ファイバーは長さ1の零次元層となる。
この条件の失敗する閉集合を第一射影で押し出し、さらに $E_U\to U$ で押し出して除く。
どちらも proper な射の下の台を用いるので、$0$ を含む一つの開近傍を取れる。
長さ1の層の族は相対 Hilbert scheme $\operatorname{Hilb}^1(E_U/U)=E_U$ により、
ある射 $\rho\colon E_U\to E_U$ の graph 上の線束となる。
$\Psi$ が同値なので $\rho$ は自己同型である。
$\rho_0=\mathrm{id}$ と $U$ の連結性により、$\rho$ は平行移動である。
線束の次数は連結基底上で一定なので7となる。

平行移動と線束を $K_0$ で決定する。
五対象の Euler 行列は基底に依存せず、Serre 作用はその行列により決まる。
従って $F$ の $K_0$ 作用は、古典点での $-K_Y$ の tensor 作用と同じである。
\[
 \delta=[T_2]-2[T_1]+[T_0],\qquad
 \varepsilon=(\psi^{-2})^*L\otimes((\psi^{-1})^*L)^{-1}
\]
と置けば、$\delta$ は古典的な点の類であり、$r(\delta)=(0,\varepsilon)$ となる。
ここでは $K_0(E)\simeq\mathbb Z\oplus\operatorname{Pic}(E)$ を rank と determinant
により表示する。古典点で直接計算すると
\begin{align}
 F[P_p]&=[P_p]+\delta,\label{eq:ac-coxeter-point}\\
 F[O]&=[T_0]-3[T_1]+3[T_2]-[P_p]-[P_q].\label{eq:ac-coxeter-unit}
\end{align}
例えば二番目は $-K_Y=3H-E_p-E_q$ および
$\chi(\mathcal O_Y(-K_Y))=8$ から従う。

式 \eqref{eq:ac-coxeter-point} を restriction すると
\[
        \mathcal O_E(\rho(p)-p)\simeq\varepsilon.
\]
$\psi=t_s$ と書けば、次数3の線束に対して
$\varepsilon\simeq\mathcal O_E(3s-0)$ であるから $\rho=t_{3s}=\psi^3$ となる。
式 \eqref{eq:ac-coxeter-unit} からは
\[
 \Psi(\mathcal O_E)=\rho_*M
 \simeq ((\psi^{-2})^*L)^{\otimes3}(-p-q)
\]
を得る。$\tau^*$ を適用し、楕円曲線上の平行移動に関する恒等式
\[
             (\psi^*L)^{\otimes3}\simeq
             L\otimes\psi^*L\otimes(\psi^2)^*L
\]
を用いると \eqref{eq:ac-elliptic-data} が従う。
これらの線束の同型を各ファイバーで得た場合、seesaw により残る差は $U$ からの
線束である。$0$ の近傍でそれを自明化する。
\end{proof}

\subsection{全次数の平坦性と中心元}

\begin{theorem}\label{thm:relative-anticanonical-ring}
上の一つの基底上で
\[
 \mathscr A_n=\operatorname{Hom}_{\mathcal H}(O,F^nO),\qquad
 \mathscr A=\bigoplus_{n\ge0}\mathscr A_n
\]
と置く。乗法は、$a\in\mathscr A_m$、$b\in\mathscr A_n$ に対して
\begin{equation}\label{eq:ac-orbit-product}
                      a b=F^n(a)\circ b
\end{equation}
とする。すると次が成立する。
\begin{enumerate}
\item
全ての $n\ge0$ について
$\mathbf R\operatorname{Hom}_{\mathcal H}(O,F^nO)$ は次数零に集中し、
$\mathscr A_n$ は有限局所自由 $R$-加群である。その階数は
\[
         \operatorname{rk}_R\mathscr A_n=1+\frac{7n(n+1)}2.
\]
構成とこの同型は任意の基底変換と両立する。
\item
自然変換 $\alpha$ から得られる
$g=\alpha_{FO}\colon O\to FO$ は次数1の中心正則元であり、
\begin{equation}\label{eq:ac-central-extension}
 0\longrightarrow\mathscr A(-1)
 \xrightarrow{\,g\,}\mathscr A
 \longrightarrow B(E_U,M,\tau)\longrightarrow0
\end{equation}
は graded algebra の中心拡大である。
\item
$\mathscr A$ は次数1で生成され、局所的に8生成元、42個の二次関係式で表示される。
古典的ファイバーは
\[
          \mathscr A_0^{\mathrm{fiber}}
          \simeq\bigoplus_{n\ge0}H^0(Y,\mathcal O_Y(-nK_Y))
\]
である。各幾何学的ファイバーは domain である。
\end{enumerate}
\end{theorem}

\begin{proof}
$M_n=M\otimes\tau^*M\otimes\cdots\otimes(\tau^{n-1})^*M$ と置くと、
\[
                  \Psi^n\mathcal O_E=\tau_*^nM_n,
                  \qquad\deg M_n=7n.
\]
式 \eqref{eq:ac-serre-triangle} を $F^nO$ に適用し、$O$ から Hom を取ると
\begin{equation}\label{eq:ac-degree-triangle}
 \mathbf R\operatorname{Hom}(O,F^{n-1}O)
 \longrightarrow\mathbf R\operatorname{Hom}(O,F^nO)
 \longrightarrow\mathbf R\Gamma(E_U,\Psi^n\mathcal O_E)
 \longrightarrow
\end{equation}
となる。$n\ge1$ では右項は次数零に集中し、階数 $7n$ の局所自由加群である。
これは各楕円ファイバー上での正次数線束の $H^1$ 消滅と cohomology and base change
による。$\mathbf R\operatorname{Hom}(O,O)=R$ を出発点として帰納すれば
\begin{equation}\label{eq:ac-degree-short-exact}
 0\longrightarrow\mathscr A_{n-1}\longrightarrow\mathscr A_n
 \longrightarrow H^0(E_U,\Psi^n\mathcal O_E)\longrightarrow0
\end{equation}
および第一の主張を得る。全ての次数で同じ正次数消滅を用いるため、
この帰納の途中で $U$ を縮める必要はない。

中心性のための自然変換の整合性を確認する。
自然性から、任意の $Y$ について
\[
 \alpha_Y\circ(\alpha_{\theta Y}-\theta(\alpha_Y))=0,
 \qquad
 \alpha_{\theta Y}-\theta(\alpha_Y)\colon\theta^2Y\to\theta Y
\]
となる。$Y=F^{n+2}O$ と取る。
三角 \eqref{eq:ac-serre-triangle} の $Y$ での値に
$\operatorname{Hom}(F^nO,-)$ を適用する。
$rF^nO$ と $rY$ は線束であるため
\[
 \operatorname{Ext}^{-1}_{E_U}(rF^nO,rY)=0.
\]
従って $\alpha_Y$ による postcomposition は
$\operatorname{Hom}(F^nO,F^{n+1}O)$ 上単射であり、
\[
 \alpha_{F^{n+1}O}=F^{-1}(\alpha_{F^{n+2}O}),\qquad
 F(\alpha_{F^{n+1}O})=\alpha_{F^{n+2}O}
\]
が従う。$\eta_X=\alpha_{FX}\colon X\to FX$ と書くと、
$\eta_{F^nO}=F^n(\eta_O)$ である。
$a\colon O\to F^nO$ に対する $\eta$ の自然性から
\[
                  F(a)g=F^n(g)a
\]
となり、\eqref{eq:ac-orbit-product} の下で $ag=ga$ を得る。
ここでは必要な orbit 上で可換性を証明した。
式 \eqref{eq:ac-degree-short-exact} の左射はこの $g$ による乗法であるから、
$g$ は左右の正則元である。

restriction は乗法を保つ。
$H^0(E_U,\Psi^n\mathcal O_E)\simeq H^0(E_U,M_n)$ の下で、
式 \eqref{eq:ac-orbit-product} は
\[
        a b=a\,\tau^{m*}b\qquad(\deg a=m)
\]
となる。従って商は \eqref{eq:ac-central-extension} の THCR である。
次数7の線束による楕円 THCR は次数1で生成される。
各次数の乗法の全射性は、次数が少なくとも3の二線束についての
楕円曲線上の切断の乗法の全射性から従い、相対的にも base change により成立する。
中心元 $g$ を加えて帰納すれば $\mathscr A$ も次数1生成となる。

次数7の楕円 THCR は Koszul、特に二次関係式で定義される
\cite[定理3.19]{PencilRainsFiltered}。
中心正則元による拡大も二次関係式で表示されることを確認する。
商の二次関係式を $\mathscr A$ に持ち上げ、$g$ と残りの七生成元との可換関係を加えた
二次代数を $C$ とする。$C\twoheadrightarrow\mathscr A$ であり、
$C/gC\simeq\mathscr A/g\mathscr A$ である。
核の次数 $n$ の元は $g$ の倍数となる。
$g$ が $\mathscr A$ 上正則であることと次数に関する帰納により、その元は零である。
この議論は各ファイバーで成立する。
相対的には
$\ker(\mathscr A_1\otimes_R\mathscr A_1\to\mathscr A_2)$ を関係式空間に取る。
各次数の presentation の核は有限 $R$-加群であり、$\mathscr A_n$ の平坦性と
全ファイバーでの二次性から Nakayama により零となる。
従って同一基底上の相対的な二次表示を得る。
$\operatorname{rk}\mathscr A_1=8$、$\operatorname{rk}\mathscr A_2=22$ なので、
関係式空間の階数は $8^2-22=42$ である。

古典点では $F=-\otimes\mathcal O_Y(-K_Y)$ なので、主張した section ring となる。
各幾何学的ファイバーの商 THCR は domain である。
非零な斉次元を割り切る最大の $g$ の冪を取り除けば、
商でその積が非零となることから $\mathscr A$ の各ファイバーも domain となる。
\end{proof}

\begin{remark}\label{rem:ac-noetherian}
$B(E_U,M,\tau)$ は左右 noetherian である。
実際、$E$ 上の次数7の線束 $H$ を固定すると
$M_n\otimes H^{-n}$ は全ファイバーで次数零であり、相対的に numerically effective
である。従って Keeler の相対 Fujita 消滅
\cite[定理1.5]{ACKeelerAmple} により、$M$ は $\tau$-ample となる。
同論文の定理1.1は、可換 noetherian 基底上の proper scheme に適用でき、
左右の $\tau$-ampleness と THCR の noetherian 性を与える。
中心正則元による拡大について $g$-進 associated graded は
$B(E_U,M,\tau)[z]$ となる。
斉次イデアルの先頭項の有限個の生成元を持ち上げ、$g$ の次数が正であることを用いて
次数ごとに帰納すれば、$\mathscr A$ も左右 noetherian となる。
この議論は楕円 THCR の ample pair に関するものであり、
$\mathcal H$ 上の反標準関手 $F$ の ample 性は追加の問題である。
\end{remark}

\subsection{平面の二点 blow-up 環との同定}

平面側の圏を $\mathcal D=\operatorname{Perf}(A)$ とし、
$j\colon\mathcal D\hookrightarrow\mathcal H$ を包含とする。
$\mathcal E=\langle P_p,P_q\rangle$ と置けば
$\mathcal H=\langle\mathcal D,\mathcal E\rangle$ である。
右 mutation の規約は
\[
 R_{\mathcal E}(X)\longrightarrow X\longrightarrow
 P_p\otimes\mathbf R\operatorname{Hom}(X,P_p)^\vee
 \oplus P_q\otimes\mathbf R\operatorname{Hom}(X,P_q)^\vee
 \longrightarrow
\]
とする。従って $R_{\mathcal E}(X)\in{}^\perp\mathcal E$ である。

以下の Sklyanin 表示では平面側の反標準関手は
$F_{\mathrm{pl}}=\mathbb S_{\mathcal D}^{-1}[2]=(3)$ である。
この点は Nakayama automorphism の規約に関わるため確認しておく。
式 (S8) の三関係式は cyclic potential
\[
       W=XYZ+bXZY-\frac h3(X^3+Y^3+Z^3)
\]
の cyclic derivatives である。
AS-regular Koszul resolution に対応する三段の両側 Koszul complex は、
この cyclicity により内部次数3のずれを伴って自己双対になる。
従って Nakayama automorphism は恒等であり、
balanced dualizing complex の qgr への降下から
$\mathbb S_{\mathcal D}=(-3)[2]$ を得る。
自己双対両側分解による Calabi--Yau 判定は
\cite[定理4.2--4.3]{ACBocklandt} に述べられている。
この有限な Koszul complex と自己双対写像は相対的にも同じ式で定義される。

\begin{lemma}\label{lem:ac-mutation-map}
自然同型と自然射
\[
 F j\simeq R_{\mathcal E}jF_{\mathrm{pl}},\qquad
 c_1\colon Fj\longrightarrow jF_{\mathrm{pl}}
\]
が存在する。これを反復して得られる $c_n\colon F^nj\to jF_{\mathrm{pl}}^n$ は、
graded algebra homomorphism
\begin{equation}\label{eq:ac-plane-homomorphism}
                 \phi\colon\mathscr A\longrightarrow\mathcal S^{(3)}
\end{equation}
を定める。次数1の射は単射であり、その像は
\begin{equation}\label{eq:ac-degree-one-subspace}
 V=\bigl\{s\in\mathcal S_3\mid
     \bar s\in H^0(E_U,L_3(-\tau^{-1}p-\tau^{-1}q))\bigr\}.
\end{equation}
ここで $\bar s$ は楕円商への像である。
\end{lemma}

\begin{proof}
$M\in\mathcal D$ とし、$V_M=R_{\mathcal E}jF_{\mathrm{pl}}M$ と置く。
$Z\in\mathcal D$ に対しては半直交性により
\[
 \mathbf R\operatorname{Hom}_{\mathcal H}(V_M,jZ)
 \simeq\mathbf R\operatorname{Hom}_{\mathcal D}(F_{\mathrm{pl}}M,Z)
 \simeq\mathbf R\operatorname{Hom}_{\mathcal H}(jZ,jM[2])^\vee.
\]
$Z\in\mathcal E$ に対しては同じ両辺が零である。
二つの部分圏が $\mathcal H$ を生成するので、Serre 関手の特徴付けから
$V_M\simeq FjM$ を得る。これは $M$ に自然であり、mutation 三角から $c_1$ を得る。

$c_0=\mathrm{id}$ とし、$c_n$ を $c_1$ の反復合成で定める。
合成の定義から
\begin{equation}\label{eq:ac-c-coherence}
 c_{m+n}=(c_nF_{\mathrm{pl}}^m)\circ F^n(c_m)
\end{equation}
となる。$a\colon O\to F^mO$ に対し、
$\phi(a)$ を $c_m(O)\circ a\colon jO\to jF_{\mathrm{pl}}^mO$ に対応する
$\mathcal D$ の射とする。
$c_n$ の自然性を $\phi(a)$ に適用すると
\[
 c_{m+n}(O)\circ F^n(a)
   =jF_{\mathrm{pl}}^n(\phi(a))\circ c_n(O).
\]
この右に $b\colon O\to F^nO$ を合成すれば
$\phi(ab)=\phi(a)\phi(b)$ となる。
ここで両側の orbit 乗法に \eqref{eq:ac-orbit-product} を用いた。
平面側の orbit 環は $\mathcal S^{(3)}$ である。

$N_3=(\psi^{-3})^*L_3=\tau_*L_3$ と置く。
evaluation による gluing は任意の $M\in\mathcal D$ について
\[
 \mathbf R\operatorname{Hom}_{\mathcal H}(jM,P_\nu)
 \simeq\mathbf R\operatorname{Hom}_{E_U}(rjM,k_\nu)
\]
と両立する。従って $M=F_{\mathrm{pl}}O=O(3)$ の mutation 三角は
\[
 FO\longrightarrow jO(3)\longrightarrow
       P_p\otimes N_3|_p\oplus P_q\otimes N_3|_q\longrightarrow
\]
となる。$O$ から Hom を取り、定理\ref{thm:relative-anticanonical-ring} の消滅を用いると
\[
 0\longrightarrow\mathscr A_1\xrightarrow{\phi_1}\mathcal S_3
 \longrightarrow N_3|_p\oplus N_3|_q\longrightarrow0
\]
を得る。最後の射は $rj$ による evaluation である。
既存の右加群の規約では $s\in\mathcal S_3$ の restriction は
$(\psi^{-3})^*\bar s$ なので、これは $\bar s$ を
$\tau^{-1}p,\tau^{-1}q$ で評価する射である。
従って核は \eqref{eq:ac-degree-one-subspace} となる。
\end{proof}

\begin{theorem}\label{thm:ac-rss-identification}
関数体 $K=\mathbb C(U)$ 上で $\tau_K$ が無限位数であると仮定する。
幾何学的 generic fiber について RSS の意味で
\begin{equation}\label{eq:ac-rss-generic}
 \mathscr A_{\overline K}
 \simeq\mathcal S_{\overline K}^{(3)}
       (\tau^{-1}p+\tau^{-1}q)
\end{equation}
となる。相対的には、$\phi$ は単射であり、
\begin{equation}\label{eq:ac-rss-relative}
                 \mathscr A\simeq R\langle V\rangle\subset\mathcal S^{(3)}.
\end{equation}
基底をさらに一度縮め、中心元を単元倍で規格化すると、$\phi$ は
$g$ を $\mathcal S$ の次数三の中心元 $g_{\mathcal S}$ に移す。
各ファイバーでも \eqref{eq:ac-rss-relative} は単射な環の表示となる。
従ってこの式は古典点を含む、次数ごとに平坦な二点 blow-up 環の代数的な族を与える。
\end{theorem}

\begin{proof}
$\mathscr A$ は次数1で生成されるので、前補題により
$\phi(\mathscr A)=R\langle V\rangle$ である。
無限位数の場合、RSS の blow-up 環
$\mathcal S^{(3)}(d)$、$d=\tau^{-1}p+\tau^{-1}q$ は $V$ で生成される。
またその Hilbert 関数は
\[
             1+\frac{(9-\deg d)n(n+1)}2
             =1+\frac{7n(n+1)}2
\]
である \cite[定義5.1、定理5.3]{BlowupRSS2015}。
従って $\phi_{\overline K}$ は各次数で同じ有限次元空間の全射であり、同型となる。
忠実平坦な体拡大から $K$ 上でも単射である。
$\mathscr A_n$ は $R$ 上 torsion-free なので、$\phi_n$ の核は零である。
これで \eqref{eq:ac-rss-relative} を得る。

中心元と特殊化の整合性も確認する。
$\mathcal S_3\to H^0(E_U,L_3)$ の核は階数1の線束であり、
近傍で生成元 $g_{\mathcal S}$ を選べる。
generic fiber の Sklyanin の中心性により、これは全基底上で中心元となる。
古典点では $E$ の三次方程式を生成する。
無限位数の楕円 THCR の正次数の中心は零であるから、
RSS 環の次数1の中心は $\overline K g_{\mathcal S}$ である。
\eqref{eq:ac-rss-generic} と $g$ の中心性から
\[
                          \phi(g)=c g_{\mathcal S}
\]
を得る。両辺の局所自由性と楕円商への射により $c\in R$ である。
古典点では Serre 三角の射が反標準楕円曲線の方程式であるから、$c(0)\ne0$ となる。
基底を縮めて $c$ を可逆とし、$g$ を規格化する。

この規格化の下で $\phi$ は中心商の射
\[
 \mathscr A/g\mathscr A\longrightarrow
 \mathcal S^{(3)}/g_{\mathcal S}\mathcal S^{(3)}=B(E_U,L_3,\tau)
\]
を誘導する。その像は $H^0(E_U,L_3(-d))$ で生成される部分環である。
楕円曲線上の切断の乗法の全射性により、その次数 $n$ の部分は
\[
 H^0\!\left(E_U,(L_3)_n(-d_n)\right),\qquad
 d_n=d+\tau^{-1}d+\cdots+\tau^{-(n-1)}d
\]
であり、$n\ge1$ で階数 $7n$ である。
中心商の左辺も各次数で同じ階数なので、この像への射は同型である。
従って商の射は任意のファイバーで単射となる。
各ファイバーで $\phi(a)=0$ ならば、中心商の単射性により $a=g a'$ と書ける。
$g_{\mathcal S}$ の正則性から $\phi(a')=0$ である。
次数を一つずつ下げて帰納すると $a=0$ となる。
この証明でも全次数について新しい開近傍を取る操作は用いていない。
\end{proof}

\begin{remark}\label{rem:ac-coefficients-and-hearts}
式 \eqref{eq:ac-degree-one-subspace} と二次性により、係数の有限計算法を得る。
$\mathcal S_3$ の10次元表示で二点の evaluation 行列を作り、その核 $V$ の基底を選ぶ。
次に Sklyanin の二次関係式を用いて
$V\otimes_RV\to\mathcal S_6$ の積を計算する。
その核は階数42であり、これが \eqref{eq:ac-rss-relative} の全関係式を生成する。
この有限計算が十分である理由は、前の全次数の証明と二次性である。

RSS の quadric--plane 同型の仮定を満たすファイバーでは、
\eqref{eq:ac-rss-generic} の環は $Q^{(2)}(r)$ とも同型である
\cite[定理7.3]{RSS2023}。
一方、$\operatorname{Perf}(B)$ と $\operatorname{Perf}(\operatorname{qgr}\mathscr A)$
の同値を反標準 section functor そのもので構成するには、
選んだ coherent heart における $F$ の ample 性、または blow-up 圏との
derived comparison を確認する必要がある。
Rains の指定する非可換 del Pezzo heart では反標準類が nef となり、
特定の root の壁で ample 性が失われる
\cite[\S4]{PencilRainsFiltered}。
この現象と Van den Bergh の derived comparison
\cite[定理11.1.3]{BlowupVdB2001} を比較する際は、
heart、構造対象、blow-up 点の表示を揃える必要がある。
本節の環の同型と全次数平坦性は、上記の Serre 三角、楕円曲線上の消滅、
mutation、および RSS の Hilbert 関数によって証明されている。
\end{remark}


\part{中央 pencil、射影量子化と共通模型}
\section{最大 pencil 成分の Segre 射影量子化}
\label{sec:native-pencil}

三重積の最大成分 $\mathcal R$ の一般元に対し、
元の $\mathcal O(1,1,1)$ 偏極を持つ斉次環の代数的量子化を構成する。
Rains \cite[命題5.1、第6節]{PencilRainsFiltered} による
楕円代数の filtered deformation の族を用い、元の pencil に対応する
相対的な直線と、全次数の平坦性を具体的に記述する。
filtered deformation の直線から非可換 Fano 三次元多様体を作る方針は
同論文 p.14 に既に記されている。
本節で必要な曲面論は \cite{PencilRainsBirational} を基盤とする。

本節の $B^{-}(E,L,\tau)$ は Rains の規約に合わせ、
$L_n=L\otimes(\tau^{-1})^*L\otimes\cdots\otimes(\tau^{1-n})^*L$ と
積 $f*g=f(\tau^{-m})^*g$ を用いる。
本稿の他節の規約では $B^{-}(E,L,\tau)=B(E,L,\tau^{-1})$ である。

\subsection{次数六の filtered deformation の相対 ruling}
\label{pencil:subsec:ruling}

楕円曲線の Picard 群はこの小節では加法的に書く。
滑らかな楕円曲線 $E$ と $L_E\in\operatorname{Pic}^6(E)$ を固定する。
量子化を表す平行移動を $\eta\in\operatorname{Pic}^0(E)$ と書き、
Rains \cite[\S6]{PencilRainsFiltered} の次数六のパラメーターを
\[
 h\in\operatorname{Pic}^3(E),\qquad
 x_0,x_1,x_2\in\operatorname{Pic}^1(E)
\]
とする。これらはそれぞれ $\phi[\mathcal O_H(-1)]$ と
$\phi[\mathcal O_{e_i}(-1)]$ の行列式である。
例えば前者の楕円曲線への導来制限は交点に支持された長さ三の類であり、
その行列式の次数は三である。
$\phi[\mathcal O_X]=[\mathcal O_E]$ を指定した場合、
同論文の式(6.4)は
\begin{equation}\label{pencil:eq:period-torsor}
 3h-x_0-x_1-x_2=\ell_\eta,
 \qquad \ell_\eta=L_E-6\eta
\end{equation}
となる。この三次元の torsor を $T_\eta$ と書く。

\begin{lemma}\label{pencil:lem:relative-ruling}
Rains の次数六の族は、$T_\eta/W(A_2A_1)$ から
$B^{-}(E,L_E,\tau_\eta)$ の filtered deformation のモジュライへの代数的な射を与える。
$\eta$ を $0$ を含む滑らかな代数曲線 $U$ 上で動かすと、この商は
\[
 \bigl(T_0/W(A_2A_1)\bigr)\times U
 \simeq\mathbb P^1\times\mathbb P^2\times U
\]
と自明化できる。三つの $y_i$ が相異なる $\mathbb P^2$ の点を固定して得る
相対 $\mathbb P^1_U$ 上では、filtered deformation の自然な直線束の引戻しは
\[
 \mathcal O_{\mathbb P^1_U}(1)\otimes p_U^*N
\]
の形である。ここで $N$ は $U$ 上の直線束であり、$0$ の近傍で自明化できる。
\end{lemma}

\begin{proof}
使用する Rains の結果を指定する。
同論文の第4節では、rank と degree を保つ写像
$K_0(X)\to K_0(E)$ から、標識した楕円非可換曲面を構成する。
構造層に命題5.1を適用すると、反標準環とその filtered deformation を得る。
第6節の式(6.1)--(6.8)は、$L_E$ と平行移動を固定したこの族を
式\eqref{pencil:eq:period-torsor}で表示する。
同節の式(6.8)の後の Weyl 群に関する議論により、
filtered deformation への射は有限 Weyl 群による商に降りる。
ここで用いるのは構成された族からの射である。
すべての filtered deformation をこの族が尽くすという予想は用いない。

これらの構成は、Poincar\'e 直線束と楕円曲線上の点の universal family を用いて
代数的な族の上で実行される。
必要な正次数の Hom の消失と基底変換は、命題5.1の完全系列(5.5)から従う。
Weyl 群による不変性は、同節で説明される通り、根が収縮されない稠密部分での
同型を有限表示の係数の恒等式として延長することで得られる。
従って $\eta=0$ を含めて同じ射が定義される。

商と相対自明化を直接計算する。
$y_i=x_i-h\in\operatorname{Pic}^{-2}(E)$ とおけば
\[
 y_0+y_1+y_2=-\ell_\eta.
\]
$A_2$ の Weyl 群は三つの $y_i$ を置換する。
$A_1$ の Cremona 反射は
\[
 h\longmapsto2h-x_0-x_1-x_2=\ell_\eta-h,
 \qquad x_i\longmapsto h-x_j-x_k
\]
であり、三つの $y_i$ は固定する。
従って商は
\[
 \left(\operatorname{Pic}^3(E)/(h\sim\ell_\eta-h)\right)
 \times
 \left(\{(y_0,y_1,y_2)\mid\textstyle\sum_i y_i=-\ell_\eta\}/\mathfrak S_3\right)
\]
となる。第一因子は楕円曲線の二重商 $\mathbb P^1$、
第二因子は対称三乗の Abel 写像のファイバー $\mathbb P^2$ である。
さらに
\begin{equation}\label{pencil:eq:relative-period-trivialization}
 h=h^{(0)}-3\eta,\qquad
 x_i=x_i^{(0)}-\eta,\qquad
 y_i=y_i^{(0)}+2\eta
\end{equation}
は $T_0\times U\to T_U$ の同型であり、上の Weyl 作用と可換である。
これが必要な相対自明化を与える。

最後に直線束を確認する。
Rains の命題6.1では、$\phi[\mathcal O_X]=[\mathcal O_E]$ のとき、
自然な直線束の引戻しの偏極は
\[
 -\frac{h^2}{2}+\frac{x_0^2+x_1^2+x_2^2}{2}
 =h^2-\ell_\eta h+\frac{y_0^2+y_1^2+y_2^2}{2}
\]
である。式(6.13)ではこの直線束が商上の $\mathcal O(1,1)$ に対応する。
この下降は直線束そのものについて確認できる。
必要なら基底を \'etale 局所化して
$2h_* =\ell_\eta$、$3y_*=-\ell_\eta$ を選び、
$h=h_*+a$、$y_i=y_*+b_i$、$\sum_i b_i=0$ と書く。
命題6.1の直線束は、基底からの直線束を除いて、
\[
 \mathcal O_E(2[0])\boxtimes
 \left.\left(\bigotimes_{i=0}^2\operatorname{pr}_i^*
                \mathcal O_E([0])\right)\right|_{\sum b_i=0}
\]
となる。第一因子は $a\mapsto\{a,-a\}$ による
$\mathcal O_{\mathbb P^1}(1)$ の引戻しである。
第二因子は対称因子 $b_0+b_1+b_2\in|3[0]|$ による
$\mathcal O_{\mathbb P^2}(1)$ の引戻しである。
後者では $0$ を含む因子の超平面が、$b_i=0$ で定まる
三つの theta divisor の和に引き戻される。
従ってこの計算は偏極の数値類に加え、商上の直線束の表示を与える。
係数族の線形化についても下降を確認できる。
固定した相異なる $y_i$ の順序を選ぶと、残る商は
$h\mapsto\ell_\eta-h$ だけである。
式\eqref{pencil:eq:relative-period-trivialization}でこの反転と四つの固定点は一定になり、
古典点での安定化群は pencil の $\mathcal O(1)$ の引戻し上自明に作用する。
その文字は $\{\pm1\}$ 値で局所一定なので、相対族でも自明であり、
有限群商の直線束の下降条件を満たす。
固定した $\mathbb P^2$ の点上で見ると、$h$ の楕円曲線上での次数は二であり、
次数二の商写像により $\mathbb P^1$ 上の次数は一となる。
相対的には $\operatorname{Pic}(\mathbb P^1_U)
\simeq\operatorname{Pic}(U)\oplus\mathbb Z[\mathcal O(1)]$ から所望の表示を得る。
\end{proof}

\begin{lemma}\label{pencil:lem:classical-ruling}
$X=(\mathbb P^1)^3$ の一般の $(1,1,1)$ pencil の基底曲線を $E$ とし、
三つの因子の次数二の直線束の制限を $L_x,L_y,L_z$ とする。
その pencil の次数六 del Pezzo 曲面の族は、
補題\ref{pencil:lem:relative-ruling}の $\eta=0$ において、
第二因子の点
\[
 \{-L_x,-L_y,-L_z\}\in\mathbb P^2
\]
を固定した $\mathbb P^1$ ruling に対応する。
ここでは $L_E=L_x+L_y+L_z$ とする。
\end{lemma}

\begin{proof}
$s=xa+b$、$t=xc+d$ と書く。
$E$ 上の階数一の行列
$\left(\begin{smallmatrix}a&b\\c&d\end{smallmatrix}\right)$ は
\[
 \mathcal O_E^{\oplus2}\longrightarrow L_x
 \longrightarrow(L_y\otimes L_z)^{\oplus2}
\]
と因子化する。従って $u=t/s$ に対応する射 $E\to\mathbb P^1$ の直線束は
$M=L_y+L_z-L_x$ である。
一般の $u$ に対応する次数六の葉は
$Y_u=\operatorname{Bl}_{p,q}(\mathbb P^1_y\times\mathbb P^1_z)$ と書け、
$p+q\in|M|$ である。
$F_1,F_2$ を二つの ruling、$E_p,E_q$ を例外曲線とする。
平面への blowdown に対応する基底を
\[
 H=F_1+F_2-E_p,\quad e_0=F_1-E_p,\quad
 e_1=F_2-E_p,\quad e_2=E_q
\]
と取る。このとき $H-e_0-e_1-e_2=E_p-E_q$ であり、
\[
 h=L_y+L_z-p,\quad x_0=L_y-p,\quad
 x_1=L_z-p,\quad x_2=q.
\]
従って
\[
 (y_0,y_1,y_2)=(-L_z,-L_y,p+q-L_y-L_z)
 =(-L_z,-L_y,-L_x)
\]
となり、第二因子は一定である。
$p\mapsto L_y+L_z-p$ は $E\to\operatorname{Pic}^3(E)$ の同型で、
$p\leftrightarrow q$ を $h\leftrightarrow L_E-h$ に移す。
従って $|M|$ と第一因子 $\mathbb P^1$ の間の写像は同型である。
一般の $u$ で得た射の同定は、両辺の代数的な射の同一性により pencil 全体へ延びる。
分岐点 $p=q$ では、標識した滑らかな曲面は infinitely-near blow-up となり、
$E_p-E_q$ の $(-2)$-曲線を反標準環が収縮する。
従ってこの延長は、元の pencil の特異な次数六のメンバーとも整合する。
\end{proof}

\subsection{filtered deformation の直線からの二重中央拡大}
\label{pencil:subsec:central-extension}

以下の構成では、通常の斉次環を得る際の全次数の平坦性を一度に確認する。
整域 $R$ 上の次数付き代数 $B=\bigoplus_{j\geq0}B_j$ を取り、
$B_0=R$、各 $B_j$ は有限射影的とする。
有限射影的な $V=B_1$ と $P\subset V\otimes_RV$ によって
$B=T_R(V)/(P)$ と表示されると仮定する。
$S=R[s,t]$ とし、$s,t$ を中央元として
\begin{equation}\label{pencil:eq:two-parameter-presentation}
 A=
 T_R(V)[s,t]\big/
 \bigl(r-\ell(r)-q(r)\mid r\in P\bigr)
\end{equation}
を定める。ここで
\[
 \ell:P\longrightarrow V\otimes_R S_1,
 \qquad q:P\longrightarrow S_2
\]
は $R$-線形であり、$V,s,t$ はすべて全次数 $1$ とする。

\begin{lemma}\label{pencil:lem:double-pbw}
式\eqref{pencil:eq:two-parameter-presentation}に $V$-次数によるフィルトレーション
$F$ を入れる。このフィルトレーションでは $s,t$ の次数を $0$ とする。
$K=\operatorname{Frac}(S)$ 上で
\[
 \operatorname{gr}_F(A\otimes_SK)\simeq B\otimes_RK
\]
が成り立つと仮定する。このとき
\[
 \operatorname{gr}_F A\simeq B\otimes_RS,
 \qquad
 A\simeq\bigoplus_{j\geq0} B_j\otimes_RS
 \quad\text{（$S$-加群として）}
\]
となる。後の同型では $B_j$ の元に全次数 $j$ を与える。
特に各 $A_n$ は有限射影的な $R$-加群であり、
\begin{equation}\label{pencil:eq:central-hilbert}
 \operatorname{rank}_R A_n
 =\sum_{j=0}^n(n-j+1)\operatorname{rank}_R B_j.
\end{equation}
$s,t$ は中央正則列であり、$A/(s,t)\simeq B$ である。
これらの結論は任意の $R$-代数への基底変換と両立する。
\end{lemma}

\begin{proof}
各 $j$ について、有限射影性を用いて全射
$V^{\otimes j}\to B_j$ の $R$-線形切断を選ぶ。
その像を使うと $S$-線形写像
\[
 N:=\bigoplus_{j\geq0}B_j\otimes_RS\longrightarrow A
\]
を得る。この写像は全射である。実際、$V$-次数 $j$ の語は
$B_j$ の選んだ持上げと、$P$ の関係式を含む語の和に書ける。
後者で $r$ を $\ell(r)+q(r)$ に置き換えると $V$-次数が下がるので、
次数に関する帰納法を適用できる。
体 $K$ 上では仮定によりこの写像は同型である。
$N$ は整域 $S$ 上捩れを持たないので、その核も捩れを持たない。
核の $K$ へのテンソル積が零であることから核自身が零となる。
同じ帰納法はフィルトレーションを保つので、associated graded の同型も従う。

全次数 $n$ の部分を取ると、$A_n$ は
$\bigoplus_{j=0}^n B_j\otimes_R S_{n-j}$ と同型になり、
式\eqref{pencil:eq:central-hilbert}を得る。
$S$-加群としての表示から $s$ が単射的に作用し、
$A/sA\simeq\bigoplus_j B_j\otimes_RR[t]$ では $t$ が単射的に作用する。
$s=t=0$ と置くと元の二次関係式を得る。
選んだ切断を含む $S$-加群の同型をテンソルすることで、基底変換との両立性も分かる。
\end{proof}

この補題の generic な仮定は、filtered deformation のモジュライへの射から得られる。
Rains \cite[命題2.2--2.3、定理2.12]{PencilRainsFiltered} の規約では、
次数 $1$ の生成元に対する関係式の一次の補正係数の重みは $1$、
定数の補正係数の重みは $2$ である。
負次数の導分が存在しなければ、生成元の平行移動に関する線形補空間を選び、
filtered deformation を一意な正規形で表せる。
従って自然な直線束を $\mathcal O(1)$ と書いたとき、
モジュライへの射
\[
 \iota:\mathbb P^1_R\longrightarrow\mathfrak M_B,
 \qquad \iota^*\mathcal O(1)\simeq\mathcal O_{\mathbb P^1_R}(1)
\]
から、補正係数を二変数の同次形式として読める。
その全次数がそれぞれ $1,2$ であることが、
式\eqref{pencil:eq:two-parameter-presentation}の $\ell,q$ を与える。
基底から来る直線束の因子がある場合は、基底の近傍で自明化して同じ構成を行う。

相対的な正規形については、$V,P$ を基底の近傍で自由にし、
生成元の定数移動が一次補正に与える線形写像
\[
 V^*\longrightarrow\operatorname{Hom}_R(P,V)
\]
の補空間を選べばよい。古典点での単射性は
\cite[系2.5]{PencilRainsFiltered} による。
これは有限行列の階数条件であるため、一度の基底縮小で補空間を選べる。
各次数について個別の開集合を取り直す操作は必要ない。

\begin{proposition}\label{pencil:prop:degree-six-central-extension}
$E$ を滑らかな楕円曲線、$L_E$ を次数 $6$ の直線束とする。
滑らかな点付き代数曲線 $(U,0)$ 上で、$\tau_0=\mathrm{id}_E$ を満たす
平行移動 $\tau$ の族を取り、
$B_U=B^{-}(E,L_E,\tau)$ を相対 twisted homogeneous coordinate ring とする。
その filtered deformation のモジュライに、上記の自然な直線束を
$\mathcal O(1)$ に引き戻す相対 $\mathbb P^1$ の族が与えられていると仮定する。
すると基底の近傍で、通常の斉次代数 $A_U$ を構成でき、次を満たす。
\begin{enumerate}
 \item $A_U$ は次数 $1$ の八生成元と三十七個の二次関係式を持つ。
 \item 二つの中央元 $s,t\in(A_U)_1$ が正則列となり、
 $A_U/(s,t)\simeq B_U$ である。
 \item 全次数 $n\geq0$ について $(A_U)_n$ は局所自由であり、
 \[
 \operatorname{rank}(A_U)_n=(n+1)^3,
 \qquad
 H_{A_U}(z)=\frac{1+4z+z^2}{(1-z)^4}.
 \]
 \item 各幾何学的ファイバーは noetherian な domain であり、Koszul である。
\end{enumerate}
\end{proposition}

\begin{proof}
$B_U$ の各正次数成分は次数 $6n$ の楕円曲線上の直線束の大域切断であり、
その階数は $6n$ である。次数 $0$ の階数は $1$ である。
Rains \cite[定理3.19]{PencilRainsFiltered} により、
次数 $6$ の twisted homogeneous coordinate ring は二次 Koszul 代数である。
従って $V$ の階数は $6$、二次関係式の階数は $6^2-12=24$ となる。
与えられた射から上の $\ell,q$ を取り、
補題\ref{pencil:lem:double-pbw}を適用する。
二十四個の変形された関係式に加え、
$[s,x_i]=[t,x_i]=0$（$1\leq i\leq6$）と $[s,t]=0$ を入れるため、
二次関係式の総数は $24+12+1=37$ である。
Hilbert 級数は
\[
 H_{B_U}(z)=1+\sum_{n\geq1}6nz^n
 =\frac{1+4z+z^2}{(1-z)^2},
 \qquad
 H_{A_U}(z)=\frac{H_{B_U}(z)}{(1-z)^2}
\]
から計算できる。

任意の幾何学的ファイバーで、$F$ に関する associated graded は
$B^{-}(E,L_E,\tau)[s,t]$ である。
楕円曲線の twisted homogeneous coordinate ring は noetherian な domain であり、
その多項式拡大も同じ性質を持つ。
非負のフィルトレーションについての先頭項の議論から $A$ も domain となり、
associated graded から有限個の生成元を持ち上げる標準的な議論で
左右 noetherian 性が従う。
また $B[s,t]$ は全次数に関して Koszul である。
bar 複体に $F$ を入れた spectral sequence は、各全次数で有限の
フィルトレーションを持つ。associated graded の Tor が対角上に集中するので、
$\operatorname{Tor}^{A}_i(k,k)_j=0$（$i\ne j$）を得る。
\end{proof}

\subsection{中央 pencil による一次 Poisson 構造の同定}
\label{pencil:subsec:semiclassical}

$X=(\mathbb P^1)^3$、$L=\mathcal O_X(1,1,1)$ とする。
一般の $s,t\in H^0(X,L)$ に対し、$\pi_{s,t}$ は
$-K_X\simeq L^2$ と
$\bigwedge^2T_X\simeq\Omega_X^1\otimes L^2$ を用いて
$s\,dt-t\,ds$ に対応する双ベクトルとする。

\begin{lemma}\label{pencil:lem:central-pencil-poisson}
$s\,dt-t\,ds$ が余次元一の零点を持たないと仮定する。
正則双ベクトル $\rho\in H^0(X,\bigwedge^2T_X)$ が
$u=t/s$ を Casimir とするとき、ある $\lambda\in\mathbb C$ に対し
$\rho=\lambda\pi_{s,t}$ となる。
\end{lemma}

\begin{proof}
関数体上で $du\ne0$ であり、三次元の接空間において
$\rho^\sharp(du)=0$ を満たす双ベクトルの空間は一次元である。
従って $\rho=f\pi_{s,t}$、$f\in\mathbb C(X)$ と書ける。
$\pi_{s,t}$ が余次元一で消えないことから、$\rho$ の正則性は
$f$ がどの素因子上でも極を持たないことを意味する。
$X$ は正規で固有だから $f$ は定数である。
\end{proof}

一般の pencil はこの仮定を満たす。
例えば $s=xa+b$、$t=xc+d$ と書き、
$D=bc-ad$ が滑らかな $(2,2)$ 曲線を定める一般の場合を考える。
接束の直和分解に対応する一形式の第一成分は
$D\in H^0(X,\mathcal O(0,2,2))$ である。
したがって余次元一の共通零点があれば、その素因子は $D=0$ である。
第二成分は $H^0(X,\mathcal O(2,0,2))$ に属する。
これが $D=0$ 上で消えれば
$H^0(X,\mathcal O(2,-2,0))=0$ より恒等的に零となる。
一般の pencil では第二成分は非零なので、共通の素因子は存在しない。

\begin{proposition}\label{pencil:prop:semiclassical-identification}
命題\ref{pencil:prop:degree-six-central-extension}の族について、
古典的ファイバーと中央元が
\[
 (A_0,s,t)\simeq
 \left(\bigoplus_{n\geq0}H^0(X,L^n),s,t\right)
\]
と同定され、右辺の pencil は一般であると仮定する。
さらに $\dot\tau_0\ne0$ とする。
このとき $A_U$ の一次 Poisson 構造は
$\lambda\pi_{s,t}$（$\lambda\ne0$）である。
基底の局所パラメーターの一次係数を調整すると、
指定した $\pi_{s,t}$ を量子化する。
\end{proposition}

\begin{proof}
次数ごとの平坦性から、交換子を基底パラメーターで割ったものは
$A_0$ 上の全次数 $0$ の Poisson 括弧を定め、$X$ に降りる。
$s,t$ は族の上で中央なので、$u=t/s$ は降りた括弧の Casimir である。
補題\ref{pencil:lem:central-pencil-poisson}により
その括弧は $\lambda\pi_{s,t}$ の形を持つ。

$\lambda\ne0$ を確認する。
$B_U$ の斉次分数環は、適当な次数 $1$ の元 $w$ を用いて
\[
 \mathbb C(E)[w,w^{-1};\tau^{-1}]
\]
と書ける。従って $\dot\tau_0=v\ne0$ に対し、その一次括弧は
\[
 \{w,f\}=-v(f)w\qquad(f\in\mathbb C(E))
\]
となり非零である。
もし $X$ 上の括弧が零ならば、$s$ が中央であるため、
\[
 A_0[s^{-1}]
 =\bigl(A_0[s^{-1}]\bigr)_0[s,s^{-1}]
\]
上の括弧全体も零になる。
$A_0$ は domain だから元の括弧も零になり、
その商 $B_0=A_0/(s,t)$ の非零括弧に矛盾する。
従って $\lambda\ne0$ であり、最後の規格化が可能である。
\end{proof}

\begin{theorem}\label{pencil:thm:native-segre}
$X=(\mathbb P^1)^3$、$L=\mathcal O_X(1,1,1)$ とする。
一般の $s,t\in H^0(X,L)$ が定める Poisson 構造 $\pi_{s,t}$ に対し、
滑らかな点付き代数曲線 $(U,0)$ 上の通常の斉次代数の族 $A_U$ が存在し、
\[
 A_U\otimes_{\mathcal O_U}k(0)=\bigoplus_{n\geq0}H^0(X,L^n),
 \qquad \operatorname{rank}(A_U)_n=(n+1)^3
\]
を満たす。
この族は $\pi_{s,t}$ を量子化し、中央正則列をなす二つの次数一の元が
古典的な $s,t$ を持ち上げる。
各幾何学的ファイバーは八生成元・三十七二次関係式を持つ
noetherian な Koszul domain である。
\end{theorem}

\begin{proof}
基底楕円曲線 $E$ と $L_E=L|_E$ を固定する。
平行移動 $\eta$ が零から非零の一次速度で動く代数曲線 $U$ を取り、
補題\ref{pencil:lem:classical-ruling}の $\mathbb P^2$ の点を
式\eqref{pencil:eq:relative-period-trivialization}で動かす。
補題\ref{pencil:lem:relative-ruling}により、自然な直線束が
$\mathcal O_{\mathbb P^1_U}(1)$ となる filtered deformation の族を得る。
命題\ref{pencil:prop:degree-six-central-extension}を適用して $A_U$ を構成する。

古典的ファイバーの環を同定する。
$C_X=\bigoplus_{n\geq0}H^0(X,L^n)$ とおくと、
$s,t$ は $C_X$ の中央正則列であり、
$C_X/(s,t)\simeq B^{-}(E,L_E,\mathrm{id})$ である。
後の等式は、二つの切断の Koszul 完全系列と
$H^1(X,L^n)=H^2(X,L^n)=0$（$n\in\mathbb Z$）から従う。
$C_X$ は二次表示を持つ Segre 環なので、
$H^0(E,L_E)$ の六つの基底切断を持ち上げ、$s,t$ を加えると
式\eqref{pencil:eq:two-parameter-presentation}の形の表示を持つ。
中央正則列から
\[
 H_{C_X}(z)=\frac{H_{B^{-}(E,L_E,\mathrm{id})}(z)}{(1-z)^2}
\]
となり、graded な Nakayama の議論で選んだ持上げが
$\mathbb C[s,t]$-加群の基底を与えることが分かる。
従って $C_X$ 自身が、pencil に沿う filtered deformation の直線の二重中央拡大である。

その filtered deformation の射は、
補題\ref{pencil:lem:classical-ruling}によって $\eta=0$ の Rains の ruling と一致する。
生成元の平行移動に関する正規形を取る際の変更は
$x_i\mapsto x_i+a_i s+b_i t$ の形であり、中央 pencil を保つ。
正規形の一意性から、対応する同次係数が一致する。
自然な直線束の自明化を古典点で調整すると
$A_U|_0\simeq C_X$ を $s,t$ の標識を保って得る。
一次 Poisson 構造は命題\ref{pencil:prop:semiclassical-identification}で同定される。
\end{proof}

\begin{proposition}\label{pencil:prop:generic-center}
上で構成した斉次環の体 $k$ 上のファイバーを $A$ とし、
対応する楕円平行移動 $\tau$ が無限位数であるとする。
全分数斜体を $Q(A)$、斉次分数環の次数零部分を $\operatorname{Frac}_0(A)$ と書くと、
\[
 Z(A)=k[s,t],\qquad
 Z(Q(A))=k(s,t),\qquad
 Z(\operatorname{Frac}_0(A))=k(t/s)
\]
である。
\end{proposition}

\begin{proof}
$A$ は noetherian domain なので Ore 分数斜体を持つ。
$V$-次数の associated graded が domain であることから、
$\deg_F(ab)=\deg_F(a)+\deg_F(b)$ である。
従って $\deg_F(ab^{-1})=\deg_F(a)-\deg_F(b)$ は分数表示に依存せず、
$-\deg_F$ は $Q(A)$ 上の離散付値を定める。
すべての非零同次先頭項が可逆化されるので
\[
 \operatorname{gr}_F Q(A)
 \simeq k(s,t)(E)[w,w^{-1};\tau^{-1}]
\]
である。右辺では $w$ の次数を一、$k(s,t)(E)$ の次数を零とする。
$\tau$ は無限位数なので、この環の中心は $k(s,t)$ である。
実際、非零の次数 $n$ の中央元は $\tau^n=\mathrm{id}$ を強制する。
また $\tau$-不変有理関数の極の有限集合は平行移動で不変であるため空であり、
そのような有理関数は基礎体の元に限る。

$z\in Z(Q(A))$ の先頭項は $\operatorname{gr}_F Q(A)$ の中心に属する。
従って $\deg_Fz=0$ であり、その先頭項は $c\in k(s,t)$ である。
$z-c$ が非零ならば、これは中央元で $\deg_F(z-c)<0$ となり矛盾する。
従って $z=c$ である。
$A$ の場合も同じ議論を $\operatorname{gr}_FA=B[s,t]$ に適用して
$Z(A)=k[s,t]$ を得る。

$s$ は中央な次数一の元なので
$Q(A)=\operatorname{Frac}_0(A)(s)$ である。
従って $\operatorname{Frac}_0(A)$ の中央元は $Q(A)$ 全体でも中央である。
$k(s,t)$ の中の全次数零の元は $k(t/s)$ であり、最後の等式を得る。
\end{proof}

\begin{remark}\label{pencil:rem:scope}
この構成は通常の斉次環、元の Segre Hilbert 級数、中央 pencil、
指定された一次 Poisson 構造を同時に与える。
共通 blow-up から得た smooth proper 圏との導来同値、
およびその generic localization とこの斉次環の分数斜体との同定には、
両方の構成における関手と標識の比較をさらに行う。
\end{remark}

\section{楕円曲線を中心とする基本変換の相対共通圏}
\label{pcat:sec}

本節では三重積の最大 pencil 成分に対し、楕円曲線を例外成分に持つ
二つの相対 SOD を構成する。
二つの収縮は同じ楕円曲線の圏を用いる有限回の mutation によって得られる。
さらに共通の generic localization と一次 Poisson 括弧を指定する。

\subsection{古典的な基本変換と gluing のデータ}

$S=\mathbb P^1_y\times\mathbb P^1_z$、$L=\mathcal O_S(1,1)$ とする。
$a,b,c,d\in H^0(S,L)$ を取り、
\[
 s=xa+b,\qquad t=xc+d,\qquad D=bc-ad
 \tag{PC1}
\]
と置く。$D=0$ が滑らかな楕円曲線 $E\subset S$ を定め、
行列 $\left(\begin{smallmatrix}a&b\\c&d\end{smallmatrix}\right)$ の四成分が
同時に消えないと仮定する。
$E$ 上で $[u_0:u_1]=[a:c]=[b:d]$ により
\[
 \ell:E\longrightarrow\mathbb P^1_u,\qquad
 M=\ell^*\mathcal O_{\mathbb P^1}(1)
 \tag{PC2}
\]
を定める。一般の pencil では $\deg M=2$ である。
右辺の二つの表示は、それぞれが定義される開集合上で用いる。

$X=\mathbb P^1_x\times S$、$X'=\mathbb P^1_u\times S$ と置き、
$Z=\{s=t=0\}\subset X$、$Z'=(\ell,i)(E)\subset X'$ とする。
ここで $i:E\hookrightarrow S$ は包含である。
古典的な共通模型は
\[
 B=\operatorname{Bl}_{Z}X\cong\operatorname{Bl}_{Z'}X'.
 \tag{PC3}
\]
以下で $H_x,H_u,L$ はそれぞれ $B$ への引戻しも表す。
$f':B\to X'$ の例外因子を $F'$ とすると
\[
                 F'=H_u+L-H_x.
                 \tag{PC4}
\]
実際、$X'$ から $\mathbb P^1_x$ への有理写像は
$u_0c-u_1a,u_0d-u_1b$ の二切断によって定まり、
その線形系の固定因子が blow-up 後の $F'$ である。

滑らかな点付き代数曲線 $(U,0)$ の座標環を $R$ とする。
第\ref{sec:nc-surface-links}節の構成から、$S$ の Poisson 構造
$D(y,z)\partial_y\wedge\partial_z$ の量子化
$\mathcal C_R$ を取り、その固定楕円曲線への restriction と
Segre shift を
\[
 \rho:\mathcal C_R\longrightarrow\operatorname{Perf}(E_R),
 \qquad \mathsf s:\mathcal C_R\xrightarrow{\sim}\mathcal C_R
 \tag{PC5}
\]
と書く。$E_R=E\times U$ であり、古典点では
$\rho_0=Li^*$、$\mathsf s_0=(-)\otimes L$ である。
ここでは通常の斉次 quadric 環の次数一 shift を用いる。
有限 Beilinson 代数とその $E_R$ 上の加群によって、(PC5) は
$R$ 上の perfect bimodule を持つ代数的な quasi-functor になる。

$\mathcal C'_R=\operatorname{Perf}(\mathbb P^1_U)\otimes_R\mathcal C_R$ とし、
\[
 \theta(A\boxtimes G)=L\ell_R^*A\otimes^L\rho(G)
 \tag{PC6}
\]
で $\theta:\mathcal C'_R\to\operatorname{Perf}(E_R)$ を定める。
同じ記号で $\ell_R:E_R\to\mathbb P^1_U$ を表した。
二つの成分の間の Hom を
\[
 \operatorname{RHom}_{\mathcal H_R}(A,jF)
 =\operatorname{RHom}_{E_R}(\theta A,F),\qquad
 \operatorname{RHom}_{\mathcal H_R}(jF,A)=0
 \tag{PC7}
\]
とする upper triangular dg category の perfect hull を
$\mathcal H_R$ と定める。$j$ は楕円成分の埋込みを表す。

\begin{lemma}\label{pcat:gluing}
$\mathcal H_R$ は $R$ 上 smooth proper であり、
\[
 \mathcal H_R=\langle\mathcal C'_R,\mathcal E_R\rangle,
 \qquad \mathcal E_R=j\operatorname{Perf}(E_R).
 \tag{PC8}
\]
古典点のこの分解は、$f'$ の右側 Orlov 分解
\[
 \operatorname{Perf}(B)
 =\langle Lf'^*\operatorname{Perf}(X'),j'_*p'^*\operatorname{Perf}(E)\rangle
 \tag{PC9}
\]
と同一視される。構成と分解は任意の導来基底変換と両立する。
\end{lemma}

\begin{proof}
dg gluing とその SOD は \cite[定義3.4--3.5、命題3.7--3.8]{PencilOrlovGluing}
の構成である。各成分は $R$ 上 smooth proper であり、(PC7) の bimodule は
相対的に perfect である。
後者は、$\mathcal C'_R$ の有限個の生成元の像が $E_R$ 上 perfect であることと、
$E_R\to U$ の smooth properness から従う。
より具体的には、両成分の perfect な対角 bimodule を用いると、
$R$ 上 perfect な Hom を持つ bimodule は両成分のテンソル積上でも perfect になる。
gluing の対角は、二つの対角 bimodule とこの cross bimodule から有限個の cone で
作られる。これで相対 smoothness を得る。
properness は (PC7) の全 Hom が $R$ 上 perfect であることから従う。
体上の同じ判定は \cite[命題3.22、定理3.25]{PencilOrlovGluing} にある。
この証明は perfect な対角と有限個の cone だけを用いるため基底変換を保つ。

古典点では $\theta_0=L(\ell,i)^*$ である。
右側 Orlov 分解の二成分間の Hom は (PC7) に一致する。
従って gluing の再構成により (PC9) と一致する。
この codimension two の gluing 表示は
\cite[例3.10]{PencilOrlovGluing} にも明記されている。
随伴と射影関手は perfect bimodule の双対とテンソル積で作られるので、
SOD も任意の導来基底変換と両立する。
\end{proof}

\subsection{二回の mutation と三重積の成分}

$\mathcal C_0,\mathcal C_u\subset\mathcal C'_R$ を
$G\mapsto\mathcal O\boxtimes G$ および
$G\mapsto\mathcal O(1)\boxtimes G$ の像とする。
従って $\mathcal H_R=\langle\mathcal C_0,\mathcal C_u,\mathcal E_R\rangle$ である。
admissible な成分 $\mathcal D$ の包含を $a$、その左随伴を $a^*$ とし、
右 mutation を
\[
 \mathsf R_{\mathcal D}(T)
 =\operatorname{Cone}(T\longrightarrow aa^*T)[-1]
 \tag{PC10}
\]
という規約で用いる。部分圏としての mutation はシフトの規約に依存しない。

\begin{theorem}\label{pcat:two-contractions}
上のデータに対して
\[
 \mathcal W_R=\mathsf R_{\mathcal E_R}(\mathcal C_u),\qquad
 \mathcal E_R^{\rm new}=\mathsf R_{\mathcal W_R}(\mathcal E_R),\qquad
 \mathcal A_R=\langle\mathcal C_0,\mathcal W_R\rangle
 \tag{PC11}
\]
と定める。このとき
\[
 \mathcal H_R
 =\langle\operatorname{Perf}(\mathbb P^1_U)\otimes_R\mathcal C_R,
               \operatorname{Perf}(E_R)\rangle
 =\langle\mathcal A_R,\operatorname{Perf}(E_R)\rangle
 \tag{PC12}
\]
という相対 admissible SOD がある。
二番目の楕円成分の埋込みは $\mathsf R_{\mathcal W_R}\circ j$ である。
古典点で $\mathcal A_0$ は $Lf^*\operatorname{Perf}(X)$ に一致する。

さらに $0$ の近傍を縮めると、$\mathcal A_R$ は階数64・八頂点の
finite locally free directed $R$-代数 $A_R$ の perfect 圏となる。
その八つの対象は、古典点で $X=(\mathbb P^1)^3$ の矩形 Beilinson collection
$\mathcal O_X(\varepsilon_x,\varepsilon_y,\varepsilon_z)$
$(\varepsilon_i\in\{0,1\})$ の引戻しを与える。
全構成は任意の導来基底変換と両立する。
\end{theorem}

\begin{proof}
admissible な二成分の mutation を順に適用すると
\[
 \langle\mathcal C_0,\mathcal C_u,\mathcal E_R\rangle
 \longmapsto
 \langle\mathcal C_0,\mathcal E_R,\mathcal W_R\rangle
 \longmapsto
 \langle\mathcal C_0,\mathcal W_R,\mathcal E_R^{\rm new}\rangle.
 \tag{PC13}
\]
mutation は移動する成分上の同値であり、
$\mathcal W_R\simeq\mathcal C_R$ と
$\mathcal E_R^{\rm new}\simeq\operatorname{Perf}(E_R)$ を得る。
二つの楕円成分の同値はこの具体的な mutation によって指定される。

古典点で $T=\mathcal O_B(H_u)\otimes p_S^*G$ とする。
$\mathcal E_0$ への左射影の unit は例外因子への制限であり、
\[
 0\longrightarrow T(-F')\longrightarrow T
 \longrightarrow j'_*j'^*T\longrightarrow0
\]
が (PC10) の三角となる。
(PC4) により
\[
 \mathsf R_{\mathcal E_0}
       (\mathcal O(H_u)\otimes p_S^*G)
 =\mathcal O(H_x)\otimes p_S^*(G\otimes L^{-1}).
 \tag{PC14}
\]
$(-)\otimes L^{-1}$ は $\operatorname{Perf}(S)$ の自己同値なので、
$\mathcal W_0=\mathcal O(H_x)\otimes p_S^*\operatorname{Perf}(S)$ である。
従って $\langle\mathcal C_0,\mathcal W_0\rangle$ は
$Lf^*\operatorname{Perf}(X)$ であり、その補成分は $f$ の右側 Orlov 成分に一致する。

$\mathcal C_R$ に矩形例外集 $F_0,F_1,F_2,F_3$ を選ぶ。
古典点では順に $\mathcal O_S,\mathcal O_S(1,0),\mathcal O_S(0,1),\mathcal O_S(1,1)$
とする。これは第\ref{sec:relative-common-blowup}節の四対象からも作れる。
古典点が $(\mathcal O,\mathcal O(1,0),\mathcal O(1,1),\mathcal O(2,1))$ の表示では、
最後の二対象の評価三角から $\mathcal O(0,1)$ を得る。
基底を縮めると、この有限個の対象間の Hom は古典点と同じ次数零の階数を持ち、
相対 full strong exceptional collection となる。

八対象を
\[
 T_{0,i}=\mathcal O\boxtimes F_i,\qquad
 T_{1,i}=\mathsf R_{\mathcal E_R}
              (\mathcal O(1)\boxtimes\mathsf s F_i)
 \quad(0\le i\le3)
 \tag{PC15}
\]
と定める。(PC14) によってその古典点は所望の矩形例外集である。
各 Hom は $R$ 上 perfect であり、古典点で高次コホモロジーが零なので、
共通の近傍で次数零の有限局所自由加群になる。
半直交性、例外性および生成性は (PC13) と四対象の生成性から従う。
従って $A_R=\operatorname{End}(\bigoplus T_{\varepsilon,i})$ は directed である。
階数は古典点で計算でき、
\[
 \operatorname{rk}_R A_R
 =\prod_{r=1}^3\left(1+1+2\right)=64.
 \tag{PC16}
\]
基底変換に関する主張は、全ての Hom、随伴、mutation が perfect bimodule から
構成されることによる。
\end{proof}

\begin{proposition}\label{pcat:kernel-bimodule}
$C_R=\operatorname{End}(\bigoplus_{i=0}^3F_i)$ とすると $\operatorname{rk}_R C_R=16$ である。
$A_R$ の二つの四頂点 corner は $C_R$ に同型で、cross bimodule は
\[
 K_R(F,G)=\operatorname{Fib}\!\left(
 H^0(\mathbb P^1,\mathcal O(1))\otimes
       \operatorname{RHom}_{\mathcal C_R}(F,\mathsf sG)
 \longrightarrow
 \operatorname{RHom}_{E_R}(\rho F,M\otimes\rho(\mathsf sG))
 \right)
 \tag{PC17}
\]
で与えられる。写像は $\ell$ による切断の評価と $\rho$ から誘導される。
有限個の $F_i,F_j$ に対して $K_R$ は次数零に集中し、その全階数は32である。
\end{proposition}

\begin{proof}
(PC10) の三角に
$\operatorname{RHom}_{\mathcal H_R}(\mathcal O\boxtimes F,-)$ を適用する。
$j^*(\mathcal O(1)\boxtimes\mathsf sG)=M\otimes\rho(\mathsf sG)$ と
(PC7) により (PC17) が得られる。

古典点では determinant の行列から
\[
 0\longrightarrow (L^{-1})^{\oplus2}
 \xrightarrow{\left(\begin{smallmatrix}a&b\\c&d\end{smallmatrix}\right)}
 \mathcal O_S^{\oplus2}\longrightarrow i_*M\longrightarrow0
 \tag{PC18}
\]
を得る。$E$ 上で行列の階数は一であり、その像は
$\ell^*\mathcal O_{\mathbb P^1}(-1)\subset\mathcal O_E^{\oplus2}$ なので、
cokernel は (PC2) の $M$ である。
ここで (PC17) の評価写像の基底は $(-u_1,u_0)$ と取る。
これは通常の $(u_0,u_1)$ からの定数基底変換であり、
行列の像 $(u_0,u_1)$ を核とする quotient map を与える。
従って (PC17) の古典点は
\[
 K_0(F,G)=H^0(\mathbb P^1_x,\mathcal O(1))\otimes
                    \operatorname{RHom}_S(F,G)
 \tag{PC19}
\]
となる。
例えば $F=\mathcal O_S(\alpha_y,\alpha_z)$、
$G=\mathcal O_S(\beta_y,\beta_z)$ について
$m=1+\beta_y-\alpha_y$、$n=1+\beta_z-\alpha_z$ と置くと、
$m,n\in\{0,1,2\}$ である。(PC17) の古典点の写像は全射で、
その source と target の次元は $2(m+1)(n+1)$ と $2+2m+2n$、
kernel の次元は $2mn$ となる。
全射性は (PC18) を $\mathcal O_S(m,n)$ で捩り、
$H^1(\mathcal O_S(m-1,n-1))=0$ を用いれば従う。
これらの有限個の写像に相対的な perfectness と基底変換を適用すると、
一つの近傍で次数零への集中と同じ階数を得る。
従って $\operatorname{rk}K_R=2\operatorname{rk}C_R=32$ である。
\end{proof}

\subsection{共通 generic localization と一次括弧}

以下では $K=\mathbb C(U)$ とし、quadric の関数斜体を
$D_K=Q_{\rm gr}(\mathcal Q_K)_0$ と書く。
$D_K(u)$ は中央不定元 $u$ を持つ $D_K[u]$ の全分数斜体を意味する。

\begin{theorem}\label{pcat:generic}
定理\ref{pcat:two-contractions}の二つの収縮に対して、
同じ斜体 $D_K(u)$ への generic localization を指定できる。
これは $\mathcal H_K$ と両端点の適切な thick 部分圏による Verdier 商の
冪等完備化であり、二つの楕円成分はともに局所化で零となる。
古典点への剰余体の同一視は
\[
 u=\frac{xc+d}{xa+b},\qquad
 x=\frac{d-ub}{ua-c}
 \tag{PC20}
\]
である。$\mathcal A_R$ の一次 Hochschild 類は、
$X$ 上の $s\,dt-t\,ds$ に対応する Poisson 構造である。
\end{theorem}

\begin{proof}
最初の SOD の楕円成分を零に送る射影を用い、
\[
 q:\mathcal H_K\longrightarrow\mathcal C'_K
       \longrightarrow\operatorname{Perf}(D_K(u))
 \tag{PC21}
\]
と置く。後半は quadric の normal な元の projective localization と
affine Ore localization、および中央変数 $u$ の局所化である。
補題\ref{lem:normal-projective-localization}、命題\ref{prop:generic-perf}により、
(PC21) はその kernel による Verdier 商の冪等完備化を与える。

第二の楕円成分も kernel に入ることを、mutation の随伴から直接確認する。
$W(G)=\mathsf R_{\mathcal E_R}(\mathcal O(1)\boxtimes\mathsf sG)$ と書き、
$W^*$ をその左随伴とする。
$\operatorname{RHom}(jF,\mathcal O(1)\boxtimes\mathsf sG)=0$ なので、
mutation の三角から
\[
 \operatorname{RHom}_{\mathcal H_R}(jF,W(G))
 =\operatorname{RHom}_{E_R}(F,M\otimes\rho(\mathsf sG))[-1].
\]
$\rho^L$ を $\rho$ の左随伴とすると、これより
\[
          W^*jF=\mathsf s^{-1}\rho^L(M^{-1}\otimes F)[1]
          \tag{PC22}
\]
となる。
命題\ref{prop:relative-elliptic-chart}の楕円 divisor は正規 regular 元 $g$ で定義される。
その長さ一の graded resolution に双対を取ると、$\rho^L$ は $E_R$ からの
pushforward に normal twist とシフトを合成した関手である。
従って $\rho^L$ の像は $g$ の局所化で零となる。
mutation の三角と $qj=0$ より
\[
 qW(W^*jF)
 =q(\mathcal O(1)\boxtimes\rho^L(M^{-1}\otimes F))[1]=0.
\]
第二の mutation の三角からも
\[
             \mathcal E_K^{\rm new}\subset\ker q
 \tag{PC23}
\]
を得る。
二番目の SOD を先に商にすることにより、$q|_{\mathcal A_K}$ も
$\operatorname{Perf}(D_K(u))$ への Verdier localization となる。
(PC20) は古典点の二つの blow-down の関数体の同一視である。

最後に一次類を照合する。
$0$ における局所環の一様化元を $h$ とする。
quadric 族の次数ごとの平坦性と古典点での整性から、
$h$ の位数は $D_K(u)$ 上の離散付値に延びる。
その付値環 $V$ の剰余体は $\mathbb C(S)(u)=\mathbb C(B)$ である。
命題\ref{prop:relative-common-field}の証明と同様に、
quadric の有限 Beilinson 代数の generic functor は $V$ 上の有限自由複体で表せる。
$\mathbb P^1$ の矢には $1,u$ を用い、$\mathcal E_R$ を零へ送る。
これにより (PC21) の $h$-整な模型を得る。
その還元は $B$ の通常の generic restriction である。
$\mathcal A_0\simeq\operatorname{Perf}(X)$ と
$H^1(X,T_X)=H^2(X,\mathcal O_X)=0$ より
\[
 HH^2(\mathcal A_0)=H^0(X,\wedge^2T_X).
 \tag{PC24}
\]
上で作った $h$-整な局所化は、この一次類を $V/h^2V$ の一次類へ送る。
この両立性を代数で確認できる。
八対象の像を $V$ 上で自明化すると、対応する代数写像
$A_R\to\operatorname{Mat}_8(V)$ を得る。
具体的には、それぞれの像を有限自由 $V$-複体で表すと、
その還元は古典的な直線束の generic fiber であり、剰余体の次数零の一次元空間となる。
付値環上の有限複体の消去補題により各複体は $V[0]$ に簡約され、
次数零の射を取ると上の代数写像が得られる。
その一次項について乗法の保存を展開すると、source と target の
一次 Hochschild cocycle の差は写像の一次項の coboundary になる。
古典点で $T=\operatorname{Mat}_8(\mathbb C(X))$ と置けば、
この写像は tilting bundle の generic restriction である。
generic restriction の平坦性と tilting Morita 同値から
$T\otimes^L_{A_0}T\simeq T$ となる。
実際、generic point の包含 $j_\eta$ は flat monomorphism であり、
$j_\eta^*j_{\eta*}\simeq\mathrm{id}$ である。
tilting Morita 同値の下で、この fully faithful な右随伴は
上の代数写像の restriction of scalars に対応する。
従って導来随伴により
$HH^*(A_0,T)\simeq HH^*(T,T)\simeq HH^*(\mathbb C(X))$ を得る。
先の cocycle の等式は、この同定の下で (PC24) の通常の制限写像を与える。
従って $\mathcal A_R$ の bivector の関数体上の括弧は
\[
 \{u,y\}=\{u,z\}=0,\qquad \{y,z\}=D(y,z).
 \tag{PC25}
\]
(PC20) を用いて $X$ の座標に戻すと、これは (PC1) の
$s\,dt-t\,ds$ に対応する括弧である。
正則 bivector は稠密開集合への制限で一意に定まるので、主張が従う。
\end{proof}

\begin{remark}\label{pcat:scope}
構成には $E\to\mathbb P^1_u$ の全グラフを用いている。
次数二の写像の分岐点でも、$Z'\subset X'$ は滑らかな codimension two の中心であり、
二点への分裂やそのための基底拡大は不要である。
二つの楕円成分は (PC13) の mutation によって、固定された
$\operatorname{Perf}(E_R)$ と同値になる。

定理\ref{pcat:generic}は、選んだ有限代数 $A_R$ と generic localization の組を与える。
特定の Segre 偏極付き斉次環との同値を得るには、その斉次環の八対象と
命題\ref{pcat:kernel-bimodule}の cross bimodule を照合する必要がある。
また共通圏 $\mathcal H_R$ 自身の一次類を $HH^2(B)$ の全 HKR 成分に分解する問題は、
両端点の一次 Poisson 類と共通 generic localization の同定から独立に扱う。
\end{remark}


\section{中央な正則列の Rees 環と楕円中心の blow-up}
\label{sec:central-rees}
ここでは前節の中央 pencil から生じる blow-up を、導来圏の主張とは独立に代数的に定式化する。
\begin{proposition}\label{prop:rees-central-regular}
$A$ を結合的な $k$-代数、$s,t\in Z(A)$ を中央元の正則列とする。
すなわち $s$ は $A$ の左右非零因子で、$t$ は $A/sA$ の左右非零因子である。
$I=(s,t)A$ とし、$U,V$ を次数 $1$ の中央不定元とする。このとき自然な準同型
\[
 A[U,V]/(tU-sV)\longrightarrow\mathcal R_A(I)
 =\bigoplus_{n\ge0} I^n z^n\subset A[z],
 \qquad U\mapsto sz,\quad V\mapsto tz
\]
は Rees 次数付き代数の同型である。
\end{proposition}
\begin{proof}
各 $I^n$ は $s^it^{n-i}$（$0\le i\le n$）で生成されるので全射である。
単項式間の関係 $tU=sV$ は零へ写る。
逆に、次数 $n$ の核の元を $F=\sum_{i=0}^{n}a_iU^iV^{n-i}$ と書く。
$\sum_i a_is^it^{n-i}=0$ を $sA$ で割ると、
$\bar a_0\bar t^n=0$ が $A/sA$ で成立する。
$\bar t$ は非零因子なので $a_0=sb$ と書ける。
$tU-sV$ の $-bV^{n-1}$ 倍を加減して定数項を消すと、
新しい $F$ は $U$ で割れる。像も $0$ なので $s$ を消去し、
次数 $n-1$ の同じ主張に帰着する。$n=0$ で核は零であるから帰納法が閉じる。
\end{proof}
\begin{remark}
この命題は、中央元 $s,t$ の存在と正則性を前提として Rees 環を同定する。
$\operatorname{Proj}$ の適切な非可換導来圏に Orlov 型の SOD が存在することは、
この代数同型だけからは従わない。その主張には blow-up 圏、例外対象、
導来完全忠実性、局所化との両立性の別個の検証を要する。
また本命題から、二種類の pencil 量子化が導来同値であるとは結論できない。
\end{remark}

\section{中央 pencil の例外代数と圏論的 blow-up の構成}
\label{sec:restored-rees-sod}
前節の Rees 環の表示からは、例外因子の代数と、
\emph{比較対象としての} categorical blow-up をさらに構成できる。
しかし両者を同一視するには、別個の比較定理が必要である。

\begin{proposition}[随伴次数環]\label{prop:restored-associated-graded}
第\ref{prop:rees-central-regular}命題の仮定の下で $D=A/(s,t)$ とおく。
Rees 次数を $m$ とすると、自然な同型
\[
\operatorname{gr}_{I}A=\bigoplus_{m\ge0}I^m/I^{m+1}
 \simeq D[U,V],\qquad [s]\mapsto U,\ [t]\mapsto V
\]
がある。右辺で $U,V$ は $D$ と可換な中央不定元である。
特に $D$ が次数六の楕円 twisted homogeneous coordinate ring で、
$\operatorname{qgr}D\simeq\operatorname{coh}E$ なら、
Rees 方向の例外射影模型の局所チャートは
$D[w]$ と $D[w^{-1}]$ で記述できる。
偏極の共通 twist を除くと、例外モデルは
$E\times\mathbb P^1$ の非可換代数的表示になる。
\end{proposition}
\begin{proof}
Rees 代数 $\mathcal R_A(I)$ を、基底環 $A$ から来る
次数零の二側イデアル $I\mathcal R_A(I)$ で割ると
\[
\mathcal R_A(I)/I\mathcal R_A(I)
 =\bigoplus_{m\ge0}I^m/I^{m+1}
\]
を得る。
第\ref{prop:rees-central-regular}命題の表示ではこの商は
\[
 A[U,V]/(tU-sV,s,t)=(A/I)[U,V]
\]
である。$U$ と $V$ は中央なので、$U\ne0$ と $V\ne0$ の
二つの標準チャートを取ると、それぞれの相対次数零部分は
$D[V/U]$ と $D[U/V]$ になる。
交叉上で両座標が互いに逆になるから、最後の主張が従う。
\end{proof}

\begin{theorem}[制限関手からの categorical blow-up]\label{thm:restored-categorical-gluing}
$\mathcal C$ と $\mathcal D$ を smooth proper な冪等完備 dg 圏とし、
$\operatorname{Res}:\mathcal C\to\mathcal D$ を perfect dg bimodule によって
表される dg 関手とする。
すると smooth proper な dg 圏 $\mathcal B_{\mathrm{cat}}$ が構成でき、
\[
 \mathcal B_{\mathrm{cat}}
 =\langle\mathcal D,\mathcal C\rangle
\]
という admissible な半直交分解を持つ。
ここで混合 Hom は、$G\in\mathcal D,F\in\mathcal C$ に対して
\[
\operatorname{RHom}_{\mathcal B_{\mathrm{cat}}}(G,F)
 \simeq\operatorname{RHom}_{\mathcal D}(G,\operatorname{Res}F)[-1],
\qquad
\operatorname{RHom}_{\mathcal B_{\mathrm{cat}}}(F,G)=0
\]
と指定できる。
さらに $X$ を滑らかな射影多様体、$Z\subset X$ を滑らかな余次元二の中心とし、
$\mathcal C=\operatorname{Perf}(X)$、$\mathcal D=\operatorname{Perf}(Z)$、
$\operatorname{Res}=Li^*$ とすると、
この圏は標準的な dg enhancement の意味で
$\operatorname{Perf}(\operatorname{Bl}_Z X)$ と同値である。
\end{theorem}
\begin{proof}
まず $\mathcal D$ と $\mathcal C$ を二つの対角成分とし、
$\mathcal D$ から $\mathcal C$ への Hom を指定の perfect bimodule、
逆向きの Hom を零とする上三角 dg category を作る。
その perfect closure を $\mathcal B_{\mathrm{cat}}$ とおく。
上三角性から二つの fully faithful な包含と表示した半直交分解が得られる。
両成分の smooth properness と bimodule の perfect 性から、
Orlov の gluing 定理\cite{PencilOrlovGluing}によって
$\mathcal B_{\mathrm{cat}}$ も smooth proper である。

可換の場合の同定を確認する。
$f:Y=\operatorname{Bl}_Z X\to X$、$j:P(N_{Z/X})\hookrightarrow Y$、
$p:P(N_{Z/X})\to Z$ とおく。
$\Psi(G)=j_*(p^*G\otimes\mathcal O_{P(N)}(-1))$ とすれば、
Orlov の blow-up 公式は
$\operatorname{Perf}(Y)=\langle\Psi\operatorname{Perf}(Z),
Lf^*\operatorname{Perf}(X)\rangle$ を与える\cite{OrlovBlowup}。
$Rp_*\mathcal O_{P(N)}(-1)=0$ により
$\operatorname{RHom}_Y(f^*F,\Psi G)=0$ である。
また、例外因子の正規束が $\mathcal O_{P(N)}(-1)$ なので
$j^!f^*F\simeq p^*i^*F\otimes\mathcal O_{P(N)}(-1)[-1]$。
従って
\[
\operatorname{RHom}_Y(\Psi G,f^*F)
 \simeq\operatorname{RHom}_{Z}(G,Li^*F)[-1].
\]
この dg bimodule と両対角成分が一致するから、
gluing の普遍性により主張の同値が従う。
\end{proof}

\begin{remark}[中央 pencil への適用範囲]\label{rem:restored-sod-scope}
最大 pencil の代数 $A$ で $A/(s,t)=B(E,M,\tau)$ とすると、
$E$ は中心の候補である。
$\mathcal C=\operatorname{Perf}(\operatorname{qgr}A)$ が smooth proper で、
中央商に沿う導来制限 $\operatorname{Res}:\mathcal C\to\operatorname{Perf}(E)$ が
perfect bimodule で表されるなら、上の定理は
$\langle\operatorname{Perf}(E),\mathcal C\rangle$ を持つ
\emph{categorical} blow-up の存在を与える。
しかしこの圏が第\ref{prop:rees-central-regular}命題の
Rees 環の $\operatorname{qgr}$ から作る圏と同値であることは、
本書ではまだ証明していない。
また $\operatorname{Perf}(\operatorname{qgr}A)$ と
第\ref{pcat:sec}節の階数64の mutation 代数の同値も未解決である。
これらを混同すると、存在が未確認の SOD を Rees 模型へ誤って付与することになる。
\end{remark}

\begin{proposition}[generic localization の条件付き両立性]
上の $\mathcal B_{\mathrm{cat}}=\langle\mathcal D,\mathcal C\rangle$ において、
$L:\mathcal C\to\mathcal G$ を Verdier 局所化とする。
$\mathcal B_{\mathrm{cat}}/\mathcal D$ は冪等完備化を除き $\mathcal C$ と同値であり、
従って $L$ は $\mathcal B_{\mathrm{cat}}$ から $\mathcal G$ への
局所化へ合成され、$\mathcal D$ の対象はそこで零になる。
ただしこの事実だけから、別の Rees 圏の generic localization と同一視することはできない。
\end{proposition}
\begin{proof}
admissible SOD の Verdier 商の普遍性による。
\end{proof}

\section{中央 pencil の Rees 模型の基底変換とアフィン半直交分解}
\label{sec:central-basechange}
本節では第\ref{sec:restored-rees-sod}節に残された問題を、
中央二変数環に関する\emph{アフィン錐の段階}で解く。
次の議論は与えられた中央正則列について完全に代数的であり、
特定の量子化や形式的な変形理論には依存しない。
ただし最後に述べるように、射影圏への降下は別の問題である。

\begin{theorem}[中央正則列に関する次数付き自由性]
\label{thm:central-graded-freeness}
$k$ を体とし、$A=\bigoplus_{n\ge0}A_n$ を
$A_0=k$、$\dim_kA_n<\infty$ の連結次数付き結合代数とする。
$s,t\in Z(A)_1$ は中央元で、$s$ は $A$ の非零因子、
$t$ は $A/sA$ の非零因子であるとする。
$C=k[s,t]$、$D=A/(s,t)$ と置くと、
\begin{enumerate}
\item 自然な写像 $k[S,T]\to A$、$S\mapsto s,T\mapsto t$ は単射である。
\item $D$ の任意の斉次 $k$-基底を $A$ に斉次に持ち上げると、
それらの持ち上げは $C$-加群として $A$ の自由基底をなす。
\item 特に $A$ は $C$ 上忠実平坦で、Hilbert 級数は
\[
 H_A(z)=\frac{H_D(z)}{(1-z)^2}
\]
を満たす。
\end{enumerate}
\end{theorem}
\begin{proof}
$s,t$ の正則性から
$0\to A(-1)\xrightarrow{s}A\to A/sA\to0$ と
$0\to (A/sA)(-1)\xrightarrow{t}A/sA\to D\to0$
は完全である。従って第3項の Hilbert 級数の公式を得る。
$D$ の各次数の基底を $A$ に持ち上げた集合を $\{b_j\}$ とする。
非負次数について帰納すれば、$A_n$ の任意の元は $b_j$ の
$C$-線形結合で表せる。実際、$D_n$ の基底を引き去ると
$(s,t)A$ に属するので、一つ低い次数へ帰着する。
従って
$F=\bigoplus_jC(-\deg b_j)\twoheadrightarrow A$
は次数付き全射である。
$H_F=H_D/(1-z)^2=H_A$ なので、各次数で全射は同じ有限次元
$k$-ベクトル空間の間の同型である。従って $F\simeq A$。
もし $P(s,t)=0$ という非零の多項式関係があれば、
得られた自由 $k[S,T]$-加群構造に反する。
$D_0=k$ の持ち上げは $1$ と取れるので、自由加群は零でなく、
忠実平坦性も従う。
\end{proof}

\begin{theorem}[Rees 環の可換 blow-up からの基底変換]
\label{thm:central-rees-basechange}
定理\ref{thm:central-graded-freeness}の仮定の下で $I=(s,t)A$ とする。
$m=(s,t)\subset C$、$\mathcal R_C(m)=\bigoplus_{r\ge0}m^r u^r$ とおけば
\[
\boxed{\quad
 \mathcal R_A(I)\;\simeq\; A\otimes_C\mathcal R_C(m)
 \;\simeq\; A[U,V]/(tU-sV).
 \quad}
\]
ここで $U,V$ は $A$ と可換で、Rees 次数は $1$、元の次数も $1$ である。
$Y=\operatorname{Bl}_0\mathbb A^2=\operatorname{Proj}_{C}\mathcal R_C(m)$
とし、$f:Y\to\operatorname{Spec}C$ と置くと、
Rees 方向の非可換相対 $\operatorname{Proj}$ は
$Y$ 上の準連接な代数の層
\[
  \mathscr A=\mathcal O_Y\otimes_C A
\]
で記述できる。その二つの標準アフィンチャートは
\[
\mathscr A(U\ne0)\simeq A[w]/(t-sw),\qquad
\mathscr A(V\ne0)\simeq A[v]/(s-tv),
\]
ここで $w=V/U$、$v=U/V$ は中央元で、交叉では $wv=1$ である。
例外因子 $E=f^{-1}(0)\simeq\mathbb P^1$ に沿う制限は
\[
  \mathscr A|_E\simeq D\otimes_k\mathcal O_{\mathbb P^1}.
\]
さらに元の次数を $z$、Rees 次数を $w$ で数えた二変数 Hilbert 級数は
\[
 H_{\mathcal R_A(I)}(z,w)
   =H_A(z)\frac{1-z^2w}{(1-zw)^2}.
\]
\end{theorem}
\begin{proof}
$I^r=m^rA$ であり、$A$ は $C$-平坦だから
$\mathcal R_C(m)\hookrightarrow C[u]$ に
$A\otimes_C-$ を施した写像も単射である。その像は
$\bigoplus_{r\ge0}I^ru^r$ である。
可換二変数の blow-up に対する標準表示
$\mathcal R_C(m)\simeq C[U,V]/(tU-sV)$ を
基底変換すれば最初の二つの同型を得る。
二つのチャートは $U$ または $V$ で局所化し、
Rees 次数零部分を取ることで得られる。
$s=t=0$ を課すと $D[w]$ と $D[v]$ が生じる。
また二重次数 $(\deg_A,\deg_{\mathrm{Rees}})$ は
$\deg U=\deg V=(1,1)$、$\deg(tU-sV)=(2,1)$
なので、正則な関係式一つを除すことから Hilbert 級数の式が従う。
同じ式は $I^r=m^rA$ と $C$-自由性からも直接得られる。
\end{proof}

\begin{theorem}[アフィン非可換 blow-up の Orlov 分解]
\label{thm:central-affine-orlov}
さらに $A$ を左右 noetherian とする。
$\mathscr A=\mathcal O_Y\otimes_C A$ 上の perfect な右加群複体の
dg 圏を $\operatorname{Perf}(Y,\mathscr A)$ と書く。
$\iota:E\hookrightarrow Y$ を例外因子とし、
$\operatorname{Perf}(D)$ の対象 $M$ に
\[
 \Psi(M)=\iota_*\bigl((M\otimes_k\mathcal O_E)\otimes\mathcal O_E(-1)\bigr),
 \qquad
 \Phi(N)=N\otimes_A^{\mathbf L}\mathscr A
 \quad(N\in\operatorname{Perf}(A))
\]
を対応させる。両関手は dg fully faithful であり、
\[
\boxed{\quad
\operatorname{Perf}(Y,\mathscr A)
 =\langle\Psi\operatorname{Perf}(D),\,
               \Phi\operatorname{Perf}(A)\rangle.
\quad}
\]
混合 Hom は
\[
 \operatorname{RHom}_{Y,\mathscr A}(\Phi N,\Psi M)=0,
 \qquad
 \operatorname{RHom}_{Y,\mathscr A}(\Psi M,\Phi N)
 \simeq\operatorname{RHom}_{D}
       (M,N\otimes_A^{\mathbf L}D)[-1]
\]
で与えられる。この SOD は第\ref{thm:restored-categorical-gluing}定理の
上三角 bimodule の具体的な\emph{アフィン}実現である。
特にRees と gluing の比較は、
\emph{アフィン錐については}成立する。
\end{theorem}
\begin{proof}
可換な $f:Y\to\operatorname{Spec}C$ は滑らかな点の余次元二 blow-up
なので、Orlov の定理\cite{OrlovBlowup}は
\[
 \operatorname{Perf}(Y)
 =\langle\iota_*\mathcal O_E(-1)\otimes\operatorname{Perf}(k),
                 f^*\operatorname{Perf}(C)\rangle
\]
という $\operatorname{Perf}(C)$-線形の分解を与える。
$A$ の $C$-平坦性によって、これを $A$-係数の加群圏へ
基底変換できる。具体的には、$Y$ の二つのアフィンチャート上で
$\operatorname{Mod}(\mathcal O_Y)\otimes_{\operatorname{Mod}(C)}
\operatorname{Mod}(A)$ は
$\operatorname{Mod}(\mathcal O_Y\otimes_C A)$ と同値であり、
この同値はチャートの交わりと両立する。
従って準連接導来加群圏を貼り合わせ、compact objects を取れば、
$\operatorname{Perf}(Y)$ の上記分解を
$\operatorname{Perf}(Y,\mathscr A)$ へ移せる。
この基底変換の完備化・compact 性は
可換側の Orlov 分解の局所的構成からも確認できる
（半直交分解の基底変換について\cite{KuznetsovSODBaseChange}も参照）。

直接、完全忠実性と混合 Hom を確認しておく。
$Rf_*\mathcal O_Y=C$、$R^if_*\mathcal O_Y=0\;(i>0)$ と
$C$-平坦性により $Rf_*\mathscr A=A$ となり、$\Phi$ は fully faithful。
$E$ は Cartier 因子で、その法束は $\mathcal O_E(-1)$ である。
$\iota^!\Phi N\simeq
 ((N\otimes_A^{\mathbf L}D)\otimes_k\mathcal O_E(-1))[-1]$
から第二の混合 Hom を得る。
第一の混合 Hom は $R\Gamma(E,\mathcal O_E(-1))=0$ から消える。
同じ Cartier resolution を $\Psi M,\Psi M'$ に使うと
$\operatorname{RHom}(\Psi M,\Psi M')\simeq
 \operatorname{RHom}_D(M,M')$ となる。
全生成性は上記の $C$-線形 Orlov 分解の基底変換による。
\end{proof}

\begin{proposition}[射影圏への降下の際に残る非正則頂点]
\label{prop:central-vertex-obstruction}
最大 pencil の中央商が
$D=B(E,L,\tau)$、$E$ が楕円曲線で $\deg L=6$ であるとする。
このとき $A$ と $D$ の連結次数付き代数としての自明加群 $k$ は
いずれも無限射影次元を持つ。実際
\[
 H_D(z)=\frac{1+4z+z^2}{(1-z)^2},\qquad
 H_A(z)=\frac{1+4z+z^2}{(1-z)^4}
   =\sum_{n\ge0}(n+1)^3z^n.
\]
従って定理\ref{thm:central-affine-orlov}の
アフィン perfect 圏の SOD を、
単に「有限次元加群で Verdier 商を取る」ことによって
$\operatorname{Perf}(\operatorname{qgr}A)$ の SOD にすることはできない。
これはそのような射影 SOD が存在しないという主張ではなく、
\emph{導来 torsion 局所化を別途構成する必要がある}という障害である。
\end{proposition}
\begin{proof}
$\deg L=6$ の楕円曲線では $h^0(E,L_n)=6n$（$n>0$）なので
表示した $H_D$ を得る。中央正則列から $H_A=H_D/(1-z)^2$。
もし連結次数付き右 noetherian 代数 $R=A$ または $D$ 上で
$k=R/R_{\ge1}$ の射影次元が有限なら、有限個の有限生成次数付き
自由加群からなる $k$ の分解があり、Euler--Poincar\'e の公式から
$1/H_R(z)$ は Laurent 多項式となる。
しかし $1+4z+z^2$ は $(1-z)$ と互いに素であるため
\[
 H_D(z)^{-1}=\frac{(1-z)^2}{1+4z+z^2},\qquad
 H_A(z)^{-1}=\frac{(1-z)^4}{1+4z+z^2}
\]
はいずれも Laurent 多項式にならない。矛盾する。
従って irrelevant torsion の基本対象 $k$ は
$\operatorname{Perf}^{\mathrm{gr}}(A)$ および
$\operatorname{Perf}^{\mathrm{gr}}(D)$ に含まれない。
射影化ではこの非 perfect torsion も消すので、
アフィン perfect 圏の単純な商だけで射影圏を記述することはできない。
\end{proof}

\begin{remark}[次の具体的な比較]
上記の証明で Rees の\emph{アフィン} gluing は決定した。
射影三次元模型に対しては、元の次数と Rees 次数の\emph{二重次数付き}
導来加群圏を取り、まず Rees 方向の相対 Proj を構成し、
次に元の irrelevant torsion を消す操作が必要である。
そのとき投影関手が torsion 局所化と両立するか、
$\operatorname{qgr}D\simeq\operatorname{coh}E$ の例外成分が
正しい dg 埋込みとして残るか、を検証しなければならない。
次の目標は二重次数付き局所化後の
$\langle\operatorname{Perf}(E),\operatorname{Perf}(\operatorname{qgr}A)\rangle$
という SOD の成立を実際に証明することであり、
本節のアフィン結果から無条件に結論しない。
\end{remark}

\begin{proposition}[可換二点 blow-up の反標準性の境界]
\label{prop:anticanonical-commutative-boundary}
$Y$ を $\mathbb P^2$ の二点での逐次 blow-up とする。
二点が元の平面上で相異なれば $-K_Y$ は ample で、
$\operatorname{coh}Y\simeq\operatorname{qgr}R(Y,-K_Y)$ である。
二点目が最初の例外曲線上にあれば、$-K_Y$ は big, nef だが
ample ではなく、最初の例外曲線の狭義変換を収縮する。
両場合で反標準環の Hilbert 関数は同じく
\[
 h^0(Y,-nK_Y)=1+\frac{7n(n+1)}2\quad(n\ge0)
\]
である。従ってこの Hilbert 関数の一致は、
反標準 section functor の ample 性を保証しない。
\end{proposition}
\begin{proof}
元の平面上で二点が相異なる場合、Mori 錐は二つの例外曲線
$E_1,E_2$ と二点を通る直線の狭義変換
$L=H-E_1-E_2$ により生成される。
$-K_Y=3H-E_1-E_2$ はこれらと交点数 $1$ を持ち、
$(-K_Y)^2=7$ なので Nakai--Moishezon により ample。
従って通常の Serre の定理により反標準環から coherent 圏を復元できる。
二点目が infinitely near なら $C=E_1-E_2$ は
$C^2=-2$、$(-K_Y)\cdot C=0$ を満たす。
この二点配置の反標準類は big, nef で、$C$ を収縮するから
ample ではない（第\ref{ex:restored-weak-dp}注意も参照）。
両場合 $K_Y^2=7$, $\chi(\mathcal O_Y)=1$ なので
Riemann--Roch により $\chi(-nK_Y)=1+7n(n+1)/2$。
$n>0$ の高次コホモロジーは
$-nK_Y=K_Y+(n+1)(-K_Y)$ への Kawamata--Viehweg 消滅から零であり、
$n=0$ も有理曲面のコホモロジーから従う。
\end{proof}

\part{非可換 blow-up、SOD と局所化}
\section{双指標によるトーリック blow-up と半直交分解}
\label{sec:toric-nc-blowups}

命題\ref{prop:toric-explicit-blowups}の全ての blow-up を、
共通トーラスの同じ量子化と両立させる。
本節では $M=\mathbb Z^3$ とし、
\[
 B=\CC[q_{12}^{\pm1},q_{13}^{\pm1},q_{23}^{\pm1}],
 \qquad
 c(m,n)=\prod_{i<j}q_{ij}^{m_i n_j}
\]
を用いる。幾何学的ファイバーについて述べるときは $B\to k$ を取り、
$k$ を標数零の代数閉体、$q_{ij}\in k^\times$ とする。
$M$-斉次元の積を $a\star b=c(|a|,|b|)ab$ で変更する操作を
$c$-twist と呼ぶ。双指標恒等式によって結合性が成り立つ。
トーリック多様体の扇、正規単項式の Ore 局所化と層の貼合せを用いる構成は
Cirio--Landi--Szabo \cite[\S\S3--4]{ToricCLS} による。
以下では blow-up の Rees 環と導来圏の全ての射にも同じ twist を施す。

\subsection{射影模型と局所 Rees 環}

$X=X_\Sigma$ を滑らかな射影トーリック多様体とする。
錐 $\sigma\in\Sigma$ 上の局所環を
\[
 A_{\sigma,B}^c
 =\bigoplus_{m\in\sigma^\vee\cap M}B X^m,
 \qquad X^m\star X^n=c(m,n)X^{m+n}
\]
と定める。面への制限は単項式の Ore 局所化である。
従ってこれらの右加群圏を共通部分上で貼り合わせることができる。
得られる coherent 圏を $\operatorname{coh}X_c$、
その perfect dg 圏を $\mathcal C_c(X)$ と書く。
トーラス不変の閉部分多様体 $Z$ には、その二側単項式イデアルによる商を用いる。
軌道閉包では、これは有効商トーラスの指標部分格子への $c$ の制限に一致する。

この圏には有限表示の射影模型がある。
古典的な Cox 構成と coherent 加群の表示には
\cite[補題2.2、定理3.2、命題3.3]{ToricCox} を用いる。
Cox 環 $S=\CC[z_\rho:\rho\in\Sigma(1)]$ の
$\operatorname{Pic}(X)$ 次数を用い、完全列
\[
 0\longrightarrow M\longrightarrow
 \mathbb Z^{\Sigma(1)}
 \longrightarrow\operatorname{Pic}(X)\longrightarrow0
\]
の分裂を一つ選ぶ。$M$ への射影を Cox 変数の重みとし、
$S_B^c$ を $B\otimes S$ の $c$-twist とする。
これは可逆なスカラー交換関係を持つ量子多項式環である。
各最大錐 $\sigma$ について
$z^{\widehat\sigma}=\prod_{\rho\notin\sigma(1)}z_\rho$ を固定順に取る。
この正規元を可逆化した環は $\operatorname{Pic}(X)$ に関して強次数付きであり、
その次数零部分は $A_{\sigma,B}^c$ である。
実際、$\sigma$ が滑らかなので、$\sigma$ の外の Cox 変数の次数は
$\operatorname{Pic}(X)$ の基底になる。
その単項式は局所化後に各次数の可逆元を与える。
従って強次数付き環の次数零部分への制限により
\[
 \operatorname{coh}X_c
 \simeq
 \operatorname{gr}^{\operatorname{Pic}(X)}S_B^c/
 \{\text{全ての }z^{\widehat\sigma}\text{ 局所化で零になる加群}\}
 \tag{TNC1}
\]
を得る。有限個のアフィンチャート上の生成元と関係式の分母を払うと、
貼合せられた coherent 加群が右辺の有限生成次数付き加群から来ることも分かる。
分裂を変えても局所環 $A_{\sigma,B}^c$ の貼合せは同じである。

十分高いトーラス線形化された ample line bundle $H$ を取れば、
切断環
\[
 R_B(X,H)^c
 =\bigoplus_{n\geq0}H^0(X,H^{\otimes n})\otimes_\CC B
\]
にも重みによる同じ積を入れられる。
各次数は有限階数自由 $B$-加群であり、
環は有限生成量子多項式環の商になるので noetherian である。
頂点切断を可逆化した次数零環は $A_{\sigma,B}^c$ となるため、
その tails は式 (TNC1) の圏と同値である。
三つの端点では、$H=\mathcal O(1)$、$\mathcal O(1,1)$、
$\mathcal O(1,1,1)$ に対する環は
本書第\ref{sec:toric-quantization}節 のトーラス型量子化そのものである。

\begin{lemma}\label{lem:toric-twisted-rees}
$f:\widetilde X=\operatorname{Bl}_Z X\to X$ を、
滑らかなトーラス不変中心 $Z$ に沿う blow-up とする。
星状細分の量子模型 $\widetilde X_c$ は、各アフィンチャート上で
通常の積を用いる Rees 環
\[
 \mathcal R(I_c)=\bigoplus_{d\geq0}I_c^d t^d
 \subset A_\sigma^c[t],
 \qquad t\text{ は中央}
 \tag{TNC2}
\]
の相対 Proj と一致する。
その中心イデアル $I_c$ は $Z$ の単項式イデアルの $c$-twist である。
量子トーラスへの局所化後、$f_c$ は恒等写像となる。
\end{lemma}
\begin{proof}
中心と交わる滑らかなアフィンチャートで、
古典的な中心イデアルを $I=(x_1,\ldots,x_d)$ と書く。
ここで $d=\operatorname{codim}(Z,X)$ であり、
必要なら残りの座標の一部を可逆にする。
各 $x_i$ は $A_\sigma^c$ の正規元である。
単項式積の係数は全て $B$ の単元なので、
$I_c^n$ の単項式集合は古典的な $I^n$ と一致する。
従って $\mathcal R(I_c)=\mathcal R(I)^c$ が成り立つ。
$t$ の $M$ 重みは零とする。

$x_i t$ を可逆化した次数零環は
\[
 x_i,\qquad x_jx_i^{-1}\ (1\leq j\leq d,\ j\ne i),
 \qquad x_\ell\ (\ell>d)
 \tag{TNC3}
\]
で生成される。比の表示には非零スカラーを吸収する。
これらの重みは
$m_i,\ m_j-m_i,\ m_\ell$ であり、古典的な blow-up チャートの双対半群をなす。
その積も $X^m\star X^n=c(m,n)X^{m+n}$ である。
この計算はチャートの共通部分での単項式局所化と両立するので、
星状細分の貼合せと一致する。
全座標を可逆化すると、全てのチャートは同じ $B_c[M]$ に入る。
\end{proof}

\subsection{重みを持つ全射空間の twist}

次の補題で非同変導来圏の生成性を扱う。
トーラス線形化は変形する対象を指定するために用い、
射空間には全てのトーラス重みを含める。

\begin{lemma}\label{lem:toric-twist-homs-generation}
$X$ を滑らかな射影トーリック多様体とし、
$(E_1,\ldots,E_N)$ をトーラス線形化された perfect complexes からなる
full exceptional collection とする。
その $c$-twist $(E_1^c,\ldots,E_N^c)$ は
$D^b(\operatorname{coh}X_c)$ の full exceptional collection である。
各 $E_i^c$ は通常の非同変 coherent 圏の対象として取る。
トーラス線形化された perfect complexes $P,Q$ の間では
\[
 \operatorname{Ext}^r_{X_c}(P^c,Q^c)_m
 \simeq \operatorname{Ext}^r_X(P,Q)_m
 \tag{TNC4}
\]
であり、右辺における重み $m$ の射 $f$ と重み $n$ の射 $g$ の合成は、
左辺では $c(n,m)(g\circ f)$ に変わる。
\end{lemma}
\begin{proof}
まず右 $M$-次数付き $A$-加群 $P$ の作用を
$p\star a=c(|p|,|a|)pa$ とする。
重み $m$ の $A$-線形写像 $f:P\to Q$ に対し
\[
 f^c(p_\ell)=c(m,\ell)f(p_\ell)
 \tag{TNC5}
\]
と置けば、双指標性から $A^c$-線形である。
さらに
$(g^c\circ f^c)(p_\ell)=c(n,m)(g\circ f)^c(p_\ell)$
となる。有限生成加群間の射は有限個の重み成分の和なので、
この計算は非同変の全射空間を記述する。
重み零の微分を持つ自由分解とトーリック被覆の \v Cech 複体に適用すると
式 (TNC4) と合成則を得る。

同じ議論はトーラス線形化された perfect complexes に適用できる。
各対象は Cox 環上で
$\operatorname{Pic}(X)\oplus M$-次数付きの有限自由分解を持つ。
完全性は単項式局所化後にも保存される。
従って例外性と半直交性は式 (TNC4) から従う。

生成性には重みシフトも確認する。
$P\langle s\rangle_\ell=P_{\ell+s}$ と置くと、
重み $\ell$ の元 $p_\ell\in P$ を $c(s,\ell)p_\ell$ へ送る写像は
\[
 (P\langle s\rangle)^c\simeq P^c
 \tag{TNC6}
\]
という非同変加群同型を与える。
実際、その右作用に現れる係数の比は
$c(\ell-s,n)c(s,\ell+n)=c(s,\ell)c(\ell,n)$ である。

古典的なトーラス線形化された line bundle $L$ に対し、
full exceptional collection の評価三角を順に用いる。
その係数 $\operatorname{RHom}(E_i,L)$ は有限次元のトーラス表現であり、
評価写像を全重み成分に分けると、各写像は重み零に取れる。
これにより $L$ は $E_i$ の指標 twist から有限回の cone と直和因子で生成される。
この構成を $c$-twist し、式 (TNC6) を用いると、
$L^c$ は $E_i^c$ から生成される。

一方、量子 Cox 環 $S^c$ は有限大域次元を持つ。
$\dim X>0$ のとき、ample divisor $A$ との交点 $D\mapsto D\cdot A^{\dim X-1}$ は
各 Cox 変数に正の整数次数を与える。
この補助次数に関する graded Nakayama により、
有限生成の次数付き射影 $S^c$-加群は次数付き自由加群となる。
$\dim X=0$ では $S^c=k$ なので同じ結論が直ちに従う。
従って任意の有限生成 $\operatorname{Pic}(X)$-次数付き加群は、
有限個の Cox 自由加群の次数シフトから生成される。
式 (TNC1) を通すと、$D^b(\operatorname{coh}X_c)$ は全ての $L^c$ から生成される。
以上から全圏での生成性が従う。
\end{proof}

\begin{theorem}\label{thm:toric-nc-orlov}
$X$ と $Z$ を滑らかな射影トーリック多様体とし、
$Z\subset X$ を余次元 $d\geq2$ の軌道閉包とする。
$X$ と $Z$ がトーラス線形化された full exceptional collections を持つと仮定する。
共通双指標 $c$ の量子模型には
\[
 \mathcal C_c(\widetilde X)
 =
 \left\langle
 \Phi^c_{1-d}\mathcal C_c(Z),\ldots,
 \Phi^c_{-1}\mathcal C_c(Z),
 Lf_c^*\mathcal C_c(X)
 \right\rangle
 \tag{TNC7}
\]
という admissible SOD がある。
$Lf_c^*$ は補題\ref{lem:toric-twisted-rees}の射の導来 pullback である。
$\Phi^c_i$ は古典的な
$j_*(q^*(-)\otimes\mathcal O_E(i))$ の全重み Hom を twist した functor である。
各幾何学的ファイバーは smooth proper dg 圏となる。
\end{theorem}
\begin{proof}
古典的な Orlov の functor はトーラス作用と両立する。
例外因子の $\mathcal O_E(i)$ もトーラス線形化を持つ。
同変化については Elagin \cite[定理10.2]{ToricElagin} を参照できる。
各 functor の Hom 上の写像は重みを保つので、
式 (TNC5) による合成則と両立する。
従って full exceptional collection 上の dg functor を twist すると、
非同変 perfect 圏の間の functor が得られる。
局所の自由複体上では $Lf_c^*$ は環準同型
$A_\sigma^c\to A_{\sigma'}^c$ による導来拡大と一致する。
それゆえ全圏でも補題\ref{lem:toric-twisted-rees}の射の pullback である。

古典的な二つの full exceptional collections を Orlov の順序で並べると、
$\widetilde X$ のトーラス線形化された full exceptional collection を得る。
補題\ref{lem:toric-twist-homs-generation}をこの collection に適用すると、
全てのブロックの fully faithfulness、半直交性および生成性が従う。

properness は式 (TNC4) による Hom の有限次元性から従う。
smoothness は有限個の例外対象の directed な最小 $A_\infty$ 模型から分かる。
その augmentation ideal の因子を並べた reduced bar resolution は、
非零の射の列で対象の番号が狭義に増加するため有限長になる。
各 Hom は有限次元なので対角 bimodule は perfect である。
smooth proper 圏の間の上記 fully faithful functor は左右の随伴を持つ。
\end{proof}

この証明は $B$ 上でも行える。
古典的 Hom の有限次元な最小 $A_\infty$ 模型を重みごとに取り、
高次積の入力を $a_n,\ldots,a_1$、重みを $m_n,\ldots,m_1$ とすると、
その積を
\[
 \mu_n^c(a_n,\ldots,a_1)
 =
 \left(\prod_{i>j}c(m_i,m_j)\right)
 \mu_n(a_n,\ldots,a_1)
 \tag{TNC8}
\]
で定める。全 Stasheff 恒等式の各項には同じ係数が掛かる。
directed 性と厳密な単位性によって非零の高次積の arity は有限に限られる。
従って有限自由 $B$-加群上の代数的な smooth proper dg 圏の族となり、
上記 SOD は perfect dg 圏の任意の derived 基底変換と両立する。

\subsection{三つの模型の共通量子 blow-up}

\begin{corollary}\label{cor:toric-nc-common}
三つの標準模型のトーラス型量子化を $X_{3,c},X_{12,c},X_{111,c}$ とする。
命題\ref{prop:toric-explicit-blowups}の全列は、
共通の射影量子模型 $Z_{\mathrm{per},c}$ への blow-up 列に持ち上がる。
各段階は式 (TNC2) の Rees 構成を持ち、式 (TNC7) の SOD を持つ。
点中心は二つの $\operatorname{Perf}(k)$、
曲線中心は一つの $\operatorname{Perf}(\PP^1)$ を加える。
共通模型は三つの表示のいずれからも長さ24の full exceptional collection を持つ。
その共通の関数斜体は
\[
 D_q
 =
 \operatorname{Frac}\!
 \left(
 k\langle x_1^{\pm1},x_2^{\pm1},x_3^{\pm1}\rangle/
 (x_ix_j-q_{ij}x_jx_i)_{i<j}
 \right)
 \tag{TNC9}
\]
であり、全ての pullback がこの斜体上で恒等写像を誘導する。
\end{corollary}
\begin{proof}
三つの出発点の矩形 Beilinson collections はトーラス線形化されている。
各中心はトーラス固定点または滑らかな不変曲線 $\PP^1$ であり、
それぞれの full exceptional collection は一対象または二対象からなる。
$\PP^1$ の二対象 $(\mathcal O,\mathcal O(1))$ の間には
長さ二の非単位合成がないので、全 Hom の双指標 twist は
通常の Kronecker 圏と同じになる。
従って中心の圏はそれぞれ $\operatorname{Perf}(k)$、
$\operatorname{Perf}(\PP^1)$ である。
定理\ref{thm:toric-nc-orlov}を各段階に適用する。
最終扇は三つの列で同一なので、局所環の貼合せから共通模型も同一である。
例外対象の個数は
\[
 4+2\cdot4+2\cdot6
 =
 6+2\cdot1+2\cdot8
 =
 8+2\cdot2+2\cdot6
 =24
\]
となる。

全てのアフィンチャートで単項式を可逆化すると同じ量子トーラス $k_c[M]$ を得る。
各チャートは noetherian domain であり、その分数体は式 (TNC9) である。
Rees チャートの式 (TNC3) から各射はこの局所化と両立する。
例外ブロックの対象は量子トーラスへの制限で零となるので、
SOD の pullback 成分の generic quotient はそのまま保存される。
\end{proof}

基底を $q_{ij}=1+h\lambda_{ij}$ によって制限すると、
$h=0$ で古典的な共通 blow-up 列と Orlov の SOD を回復する。
各局所環の一次交換子は
$\{X^m,X^n\}
=\sum_{i<j}\lambda_{ij}(m_i n_j-m_j n_i)X^{m+n}$
であり、全中間模型で同じ $\pi_\lambda$ を持ち上げる。
この構成には $q_{ij}$ の無限位数の仮定はない。
トーラス型の三つの成分の一般元に対する、古典点を含む代数的な共通模型である。

% 2026-11-08 v7. Japanese LuaLaTeX fragment.
% Requires the theorem environments and macros of noncommutative_birational.tex.
\section{非可換 blowup と三次元の半直交分解}
\label{blowup:section}\label{sec:nc-blowups}

本節では基礎体を $k=\CC$ とする。
半直交分解の記法
$\mathcal D=\langle\mathcal A_1,\ldots,\mathcal A_m\rangle$
には $j>i$ のとき
$\operatorname{Hom}(\mathcal A_j,\mathcal A_i)=0$ となる順序を用いる。
三次元の対象には、射影圏 $\qgr A$、可換三次元多様体上の order、
smooth proper dg 圏としての三次元量子化という複数の表示が現れる。
以下では各結果の表示と仮定を指定する。
特に、三次元 Sklyanin 代数の射影圏は非可換曲面であり、
四次元 AS-regular 代数の射影圏が非可換射影三次元空間の候補となる。

\subsection{曲面の blowup と blowdown}
\label{blowup:surfaces}

Van den Bergh の構成は、可換 Cohen--Macaulay 曲線 $C$ を
両側閉な Cartier divisor として含む noetherian quasi-scheme $S$ を出発点とする。
$C$ に支えを持つ対象の $S$ 上での有限 injective dimension などの仮定の下で、
$p\in C$ における blowup $\alpha\colon\widetilde S\to S$ を構成する。
生成に関する仮定は、$C$ への制限が line bundle となる対象で $S$ が生成される場合に
満たされる。非可換射影平面と標準的な非可換 ruled surface はこの範囲に入る。
exceptional object $\mathcal E=\mathcal O_e(-1)$ が存在し、
有界 coherent derived category には
\begin{equation}\label{blowup:surface-sod}
 D^b(\operatorname{coh}\widetilde S)
 =\left\langle
     \mathcal E\otimes_k\operatorname{Perf}(k),
     L\alpha^*D^b(\operatorname{coh}S)
   \right\rangle
\end{equation}
が成り立つ。この結果は \cite[定理8.4.1]{BlowupVdB2001} による。
生成仮定を明示した定式化は \cite[定理11.12]{BlowupRains2025} にもある。
したがって、非可換曲面の点 blowup に対する Orlov 型の SOD は既知である。

blowdown の abelian category を回復する際には $t$-structure の指定も重要である。
Van den Bergh の \emph{Abstract blowing down} は、曲面の滑らかな点の blow-up において、
$\operatorname{coh}(\widetilde S)$ と $\mathcal O_e(-1)$ から、
その直交圏上の perverse $t$-structure を使って $\operatorname{coh}(S)$ を回復する。
\cite[定理1.1]{BlowupAbstract}

楕円代数 $T$ は、次数1の中心元 $g$ を持ち
$T/gT=B(E,\mathcal L,\tau)$ となる連結次数付き domain である。
Rogalski--Sierra--Stafford は $|\tau|=\infty$ と
$\deg D<\deg\mathcal L$ の下で有効因子 $D$ に沿う部分環 $T(D)\subset T$ を構成し、
分数環の保存と一重点 blowup の反復による表示を証明した。
また適切な自己交点 $-1$ の line module を収縮する環論的 blowdown を構成した。
前者は \cite[定理1.1、定理5.3、命題5.4]{BlowupRSS2015}、
後者は \cite[定理1.4、定理1.5]{BlowupRSS2017} による。
第\ref{sec:nc-surface-links}節の楕円リンクに用いた
共通 blowup の同型は \cite[定理1.2、定理7.3]{RSS2023} の具体的な因子分解である。
Presotto の $\mathbb Z^2$-代数による構成にも、
共通 blowup の射影圏との比較同値がある。
\cite[定理1.1]{BlowupPresotto2017}

これらの定理で用いる smoothness と比較同値には実質的な役割がある。
環論的 blowup の射影圏が有限 homological dimension を持つことは
各段階で確認を要する。実際、同じ点における反復操作によって
非 smooth な射影圏が生じる例が \cite[\S10]{BlowupRSS2017} にある。
環論的モデルに \eqref{blowup:surface-sod} を適用するときは、
Van den Bergh 型 blowup の圏との同一視を確認する。
第\ref{blowup:quadric-comparison}節では、quadric の一点 blow-up について
この比較を証明する。

Rains の近年の研究は、非可換 quasi-ruled surface の blowdown をさらに扱う。
その定理15.9は、$S$ のある反復 blowup が、曲線対 $(C_0,C_1)$ に付随する
quasi-ruled surface の反復 blowup と同一になると仮定する。
このとき $S$ 自身が同じ曲線対に付随する quasi-ruled surface の反復 blowup であるか、
あるいは $C_0\cong C_1\cong\PP^1$ かつ $S$ が非可換平面である。
この結果の出発点には共通 blowup の存在が置かれている。
\cite[定理15.9]{BlowupRains2025}

\subsection{曲面から得られる三次元の blowup と SOD}
\label{blowup:product}

dg 圏 $\mathcal C,\mathcal D$ の積
$\mathcal C\widehat\otimes_k\mathcal D$ は、
derived tensor product を pretriangulated かつ idempotent complete にした圏とする。
以下の命題を既知の曲面 SOD からの形式的帰結として記す。

\begin{proposition}[積による blowup の SOD]\label{blowup:product-proposition}
$\mathcal C$ と $\widetilde{\mathcal C}$ を smooth proper な $k$-線形 dg 圏とする。
dg-enhanced fully faithful functor
$F\colon\mathcal C\to\widetilde{\mathcal C}$ と exceptional object
$\mathcal E\in\widetilde{\mathcal C}$ があり、admissible な半直交分解
\[
 \widetilde{\mathcal C}
 =\langle\mathcal E\otimes_k\operatorname{Perf}(k),F\mathcal C\rangle
\]
を与えると仮定する。
$\mathcal P=\operatorname{Perf}(\PP^1)$ とおくと
\begin{equation}\label{blowup:product-sod}
 \widetilde{\mathcal C}\widehat\otimes_k\mathcal P
 =\left\langle
      \mathcal E\boxtimes\mathcal P,
      (F\widehat\otimes\mathrm{id})
      (\mathcal C\widehat\otimes_k\mathcal P)
   \right\rangle
\end{equation}
は smooth proper dg 圏の admissible な半直交分解であり、
第一成分は $\mathcal P$ と同値である。

さらに $\mathcal C$ と $\widetilde{\mathcal C}$ がそれぞれ
$\operatorname{Perf}(\qgr A)$ と $\operatorname{Perf}(\qgr\widetilde A)$ で表示され、
$A,\widetilde A$ が次数1生成の noetherian graded domain で
$Q_0(A)\cong Q_0(\widetilde A)=D$ を満たすと仮定する。
両 Segre 環が積圏を表示する場合、その関数斜体はともに
\begin{equation}\label{blowup:product-field}
 D(u),\qquad u\text{ は中心変数}
\end{equation}
となる。
\end{proposition}

\begin{proof}
外積対象間の Hom 複体には
\[
 R\operatorname{Hom}(M\boxtimes P,N\boxtimes Q)
 \cong
 R\operatorname{Hom}(M,N)\otimes_k^{\mathbb L}
 R\operatorname{Hom}(P,Q)
\]
が成り立つ。従って
$\mathcal E\boxtimes(-)$ と $F\widehat\otimes\mathrm{id}$ は fully faithful であり、
必要な半直交性も従う。
任意の $M\in\widetilde{\mathcal C}$ に元の分解三角形を適用してから
$P\in\mathcal P$ と外積を取ると、$M\boxtimes P$ が右辺の二成分から生成される。
これらの外積対象は積圏を生成するため、分解は full である。
元の射影 functor とその随伴を $\mathrm{id}_{\mathcal P}$ と tensor すれば
両成分の随伴を得るので、admissibility も保たれる。

properness は上の Hom 複体の式から、smoothness は
積圏の diagonal bimodule が両因子の diagonal bimodule の derived tensor product
であることから従う。
最後の主張は補題\ref{lem:segre-field}をそれぞれの次数付き表示に適用すればよい。
\end{proof}

\eqref{blowup:surface-sod} にこの命題を適用すると、
$S\times\PP^1$ の $p\times\PP^1$ に沿う非可換 blowup の積モデルが得られる。
可換特殊化が存在する場合、その特殊化は
\[
 \operatorname{Bl}_{p\times\PP^1}(S_0\times\PP^1)
 \cong(\operatorname{Bl}_{p}S_0)\times\PP^1
\]
であり、exceptional divisor は $\PP^1\times\PP^1$ である。
この同型は blowup の平坦基底変換性から従う。
非可換積モデルの SOD に加わる成分は $\operatorname{Perf}(\PP^1)$ となる。
従って、本研究の $X_{12}$ と $X_{111}$ のうち
曲面量子化と中央 $\PP^1$ の積で表されるモデルに適用できる。
代数的パラメータ族全体に同じ主張を行うには、曲面 functor と SOD の相対化を用いる。
系\ref{cor:relative-common-threefold}は、二つの収縮を同時に相対化した具体例である。

\begin{corollary}[共通曲面 blowup の積]\label{blowup:common-product}
$\mathcal C_{\mathrm{pl}}$ と $\mathcal C_{\mathrm{qd}}$ をそれぞれ
非可換平面と非可換 quadric surface の perfect dg 圏とする。
smooth proper dg 圏 $\mathcal B$ が両者の共通 blowup であり、
二つの Van den Bergh 型点 blowup と一つの点 blowup により
\[
 \mathcal B
 =\langle\mathcal E_2,\mathcal E_1,\mathcal C_{\mathrm{pl}}\rangle
 =\langle\mathcal E_r,\mathcal C_{\mathrm{qd}}\rangle
\]
と表示されると仮定する。
ここで各 $\mathcal E_i$ は exceptional object の生成する成分を表す。
$\mathcal Y=\mathcal B\widehat\otimes\operatorname{Perf}(\PP^1)$ とおけば
\begin{align}
 \mathcal Y
 &=\left\langle
    \operatorname{Perf}(\PP^1),\operatorname{Perf}(\PP^1),
    \mathcal C_{\mathrm{pl}}\widehat\otimes\operatorname{Perf}(\PP^1)
   \right\rangle,\label{blowup:common-plane}\\
 \mathcal Y
 &=\left\langle
    \operatorname{Perf}(\PP^1),
    \mathcal C_{\mathrm{qd}}\widehat\otimes\operatorname{Perf}(\PP^1)
   \right\rangle.\label{blowup:common-quadric}
\end{align}
\end{corollary}
\begin{proof}
命題\ref{blowup:product-proposition}を各 blowup に順次適用する。
\end{proof}

\cite[定理7.3]{RSS2023} の共通環に対しては、
各環論的 blowup と上記 categorical blowup の比較を確認したうえでこの系を適用する。
関数斜体 $D(u)$ の同一視は第\ref{sec:nc-surface-links}節で用いた Segre の計算から得られる。
一方、\eqref{blowup:common-plane}--\eqref{blowup:common-quadric} の
仮定は exceptional object と fully faithful functor も指定している。
第\ref{sec:relative-common-blowup}節は、選んだ量子化族に対してこれらのデータを構成する。

平面側の比較には、さらに具体的な先行結果がある。
Van den Bergh \cite[定理11.1.3]{BlowupVdB2001} は楕円量子平面の
反復 categorical blowup とその section ring の射影圏を結ぶ射が、
指定された移動後の点の一般位置条件の下で導来同値を誘導することを示す。
二点の場合の条件は、その二点が相異なることである。
RSS の section ring の比較 \cite[第5節、定義5.1の直後、命題5.4]{BlowupRSS2015} と合わせると、
この条件を確認した平面側には二つの例外対象を加えた SOD が得られる。
次の部分節では quadric 側の比較を証明し、
固定された一般の Sklyanin quadric の共通環にも
系\ref{blowup:common-product}を適用する。

% 2026-11-08 v7. Integration fragment for Japanese LuaLaTeX.
% The new comparison argument is proved below for the explicitly stated
% marked cubic Z-algebra presentation. See 01-research-notes.md for scope.
% Existing bibliography keys are reused. Add QuadricBR2020 and QuadricMU2021.

\subsection{楕円型 quadric の一点 blowup と切断環の比較}
\label{blowup:quadric-comparison}

以下では $k$ を標数零の代数閉体とする。
$A$ を三次元 cubic AS-regular $\mathbb Z$-代数とし、
$X=\operatorname{Proj}A$、$P_i=\pi(e_iA)$ と書く。
$A$ の楕円データを $(E,M_0,M_1,M_2)$ とし、$E$ は滑らかと仮定する。
この記法では $M_i$ は次数2の直線束であり、
$E\hookrightarrow\mathbb P^1\times\mathbb P^1$ を定める隣接する二つの pencil を持つ。
楕円曲線への包含を $i\colon E\hookrightarrow X$ とする。
標準的な同値 $\operatorname{qgr}(A/(g_i))\simeq\operatorname{coh}E$ を
$Li^*P_0=\mathcal O_E$ となるよう規格化すると
\begin{equation}\label{blowup:quadric-restrictions}
 F_0=\mathcal O_E,\qquad
 F_j=Li^*P_j=(M_0\otimes\cdots\otimes M_{j-1})^{-1}
 \quad(1\le j\le3).
\end{equation}
この同定は twisted homogeneous coordinate $\mathbb Z$-代数の定義から従い、
$\operatorname{Hom}(P_{j+1},P_j)=H^0(E,M_j)$ の制限は切断の評価写像である。
用いている cubic $\mathbb Z$-代数の分類と楕円データについては
\cite[命題5.1.2、\S5.5]{VdBQuadrics} を参照する。

\begin{lemma}[一点上の半単純な制限による消滅]
\label{blowup:quadric-kernel}
$T\in D^-_{\mathrm{coh}}(X)$ と $t\in E$ が
\[
 R\Gamma(X,T)=0,\qquad
 \mathcal H^j(Li^*T)\cong k(t)^{\oplus m_j}
 \quad(j\in\mathbb Z)
\]
を満たすとする。各 $m_j$ は有限とする。このとき $T=0$ である。
\end{lemma}

\begin{proof}
$(P_3,P_2,P_1,P_0)$ は full strong exceptional collection である。
最初の二対象を right mutation して
\[
 0\longrightarrow P_3\longrightarrow
 H^0(E,M_2)^\vee\otimes P_2\longrightarrow U\longrightarrow0
\]
を得る。Belmans--Raedschelders の計算
\cite[命題3.12]{QuadricBR2020} により、
$(P_2,U,P_1,P_0)$ は full strong exceptional collection であり、
$U$ と $P_1$ は完全に直交する。
従って $R\Gamma(X,-)$ の核は
\[
 \mathcal K=\langle P_2,U,P_1\rangle
\]
に対応する。
その tilting 代数 $H$ は、一つの頂点から他の二頂点へそれぞれ二本の矢が出る
関係式のない quiver の道代数であり、hereditary である。
ここで必要となるのは三頂点の部分代数だけである。

上の mutation 列を $E$ に制限する。
$M_2$ の二次元の全切断空間は basepoint-free なので、
制限した左の射は部分直線束の包含になる。
二次元ベクトル空間の行列式の基底を一つ固定すれば
\begin{equation}\label{blowup:quadric-mutated-restriction}
 Li^*P_2=F_2,\qquad Li^*U=F_2\otimes M_2,\qquad
 Li^*P_1=F_2\otimes M_1.
\end{equation}
二組の矢は、それぞれ $M_2$ と $M_1$ の全切断による写像である。

hereditary 性より、$D^-_{\mathrm{fd}}(H)$ の対象はそのコホモロジー加群の
シフトの直和に分解する。有限次元 $H$-加群の最小射影分解を
$\mathcal K$ に戻すと、その対応物は
\begin{equation}\label{blowup:quadric-two-term}
 W\otimes P_2\longrightarrow
 V\otimes U\oplus V'\otimes P_1
\end{equation}
という二項複体と $P_2$ の直和に分かれる。
最小性により、始域の $P_2$ から終域の $P_2$ への成分は零になる。
$Li^*P_2$ は直線束なので、$P_2$ の単独の直和因子は仮定から生じない。
また、各二項複体の制限のコホモロジーも $k(t)$ の直和である。
有限振幅の制限関手はこの各次数で有限な直和と可換する。

式\eqref{blowup:quadric-two-term}を制限し $F_2^{-1}$ を掛けると
\[
 \varphi\colon\mathcal O_E^{\oplus w}\longrightarrow
 M_2^{\oplus v}\oplus M_1^{\oplus v'}
\]
を得る。核は局所自由層の部分層かつ torsion なので零である。
余核が torsion であることから $w=v+v'$ である。
直線束の次数を比較すると
\[
 \operatorname{length}(\operatorname{coker}\varphi)=2(v+v')=2w.
\]
一方、余核が $k(t)^{\oplus m}$ なら、$t$ におけるファイバーへの全射から
$m\le v+v'=w$ となる。従って $2w\le w$ であり $w=0$ である。
全ての直和因子が零となり、$T=0$ を得る。
\end{proof}

\begin{theorem}[cubic 表示を持つ quadric の比較同値]
\label{blowup:quadric-section-equivalence}
上の $X$ が、構造層と $E$ の埋め込みを保つ同値
\[
 X\simeq\operatorname{Proj}Q
\]
によって、$A$ の二周期的な even Veronese の通常の斉次表示 $Q$ を持つとする。
さらに $Q$ は domain で、次数2の中心的正則元 $g$ を持ち
\[
 Q/gQ=B(E,\mathcal L,\alpha),\qquad \deg\mathcal L=4
\]
とする。$T_0=Q^{(2)}$ とし、
\[
 T_0/gT_0=B(E,\mathcal M,\tau),\qquad
 \mathcal M=\mathcal L\otimes\alpha^*\mathcal L,\qquad
 \tau=\alpha^2
\]
と書く。$|\tau|=\infty$ と仮定する。
任意の $t\in E$ に対して $R=T_0(t)$ とおけば、
Van den Bergh の一点 blowup $\alpha_t\colon\widetilde X_t\to X$ と切断環への射
$\delta_t\colon\widetilde X_t\to\operatorname{Proj}R$ は
\[
 R\delta_{t*}\colon D^b(\operatorname{coh}\widetilde X_t)
       \xrightarrow{\ \sim\ }D^b(\operatorname{qgr}R)
\]
という dg-enhanced 同値を与える。その逆は $L\delta_t^*$ である。
従って $\mathcal E_t=R\delta_{t*}\mathcal O_{L_t}(-1)$ とおけば
\begin{equation}\label{blowup:quadric-unconditional-sod}
 \operatorname{Perf}(\operatorname{qgr}R)
 =\left\langle
     \mathcal E_t\otimes_k\operatorname{Perf}(k),\quad
     R\delta_{t*}L\alpha_t^*\operatorname{Perf}(\operatorname{qgr}Q)
   \right\rangle
\end{equation}
\end{theorem}

\begin{proof}
Van den Bergh の構成に用いる仮定を確認する。
$E\subset X$ は滑らかな両側閉 Cartier divisor であり、
$\mathcal O_X(E)=\mathcal O_X(2)$ は ample である。
標準的な $P_i$ は $E$-flat で、制限は直線束であり、$X$ を生成する。
cubic AS-regular 性から $\operatorname{Mod}(X)$ は有限 homological dimension を持つ。
また $P_0=\mathcal O_X$ が exceptional なので
$R\Gamma(X,\mathcal O_X)=k$ である。
normal bimodule に点イデアルを掛けた $E$ 上の invertible bimodule は
$(\mathcal M(-t))_\tau$ であり、次数は7である。
その正の冪は $H^1=0$ を満たし ample である。
従って \cite[\S9.2、命題9.2.2、命題9.2.3]{BlowupVdB2001} が適用される。

切断環の同定もここで確認できる。
Van den Bergh の Rees bimodule algebra を $\mathcal D_t$ とし
$D_t=\Gamma(X,\mathcal D_t\mathcal O_X)$ とする。
この標準表示では
$Q_n=H^0(X,\mathcal O_X(n))$ for $n\ge0$ である。
従って Rees algebra の包含から $D_t\subset T_0$ を得る。
同文献の式(6.16)と命題9.2.3から
\[
 0\longrightarrow D_t(-1)\xrightarrow{g}D_t
 \longrightarrow B(E,\mathcal M(-t),\tau)\longrightarrow0
\]
を得る。次数1部分は
\[
 (D_t)_1/k g=H^0(E,\mathcal M(-t))\subset H^0(E,\mathcal M).
\]
この式は点の表示を固定する。
右加群の grading shift による点の移動はこの標識に含める。
$R=T_0(t)$ はこの次数1部分で生成される
\cite[定理5.3]{BlowupRSS2015}。
従って $R\subset D_t$ である。
両者の $g$ による商は同じ twisted homogeneous coordinate ring であり、
各次数の次元は
\[
 \dim_k R_n=\dim_k(D_t)_n=1+\frac{7n(n+1)}2.
\]
よって $R=D_t$ である。

命題9.2.3は $R\delta_{t*}L\delta_t^*=\mathrm{id}$ を $D^-$ 上で与える。
さらに $\delta_{t*}$ の cohomological dimension は高々1である。
$N\in D^-_{\mathrm{coh}}(\widetilde X_t)$ が $R\delta_{t*}N=0$ を満たすとする。
$T=R\alpha_{t*}N$ とおく。
blowup の分解 \cite[定理8.4.1]{BlowupVdB2001} により三角
\[
 L\alpha_t^*T\longrightarrow N\longrightarrow C\longrightarrow
\]
が存在し、$C$ は $\mathcal O_{L_t}(-1)$ のシフトの直和である。
pushforward の合成と $R\delta_{t*}N=0$ から
$R\Gamma(X,T)=R\Gamma(\widetilde X_t,N)=0$ となる。
同文献の式(9.11)と、$\delta_t$ が $E$ 上で恒等写像になることから
$Li^*N=0$ である。
$\mathcal O_{L_t}(-1)$ の $E$ への導来制限は一点の skyscraper である。
従って $Li^*T$ の各コホモロジーは同じ一点の skyscraper の有限直和であり、
補題\ref{blowup:quadric-kernel}により $T=0$ となる。
三角から $N\simeq C$ であり、その制限が零であることから $C=0$、従って $N=0$ である。
これで $R\delta_{t*}$ の核が $D^-_{\mathrm{coh}}$ 上で消滅した。

有界性は次の spectral sequence から従う。
$F\in\operatorname{coh}(\operatorname{Proj}R)$ とし $N=L\delta_t^*F$ とおく。
$R^j\delta_{t*}=0$ for $j>1$ より
\[
 0\longrightarrow R^1\delta_{t*}\mathcal H^{r-1}(N)
 \longrightarrow\mathcal H^r(R\delta_{t*}N)
 \longrightarrow\delta_{t*}\mathcal H^r(N)\longrightarrow0
\]
を得る。$R\delta_{t*}N=F$ なので、$r<-1$ に対して
$R\delta_{t*}\mathcal H^r(N)=0$ である。
核の消滅から $\mathcal H^r(N)=0$ となり、$L\delta_t^*$ の振幅は $[-1,0]$ に入る。
従って上の随伴同値は $D^b_{\mathrm{coh}}$ に制限する。
導来関手は元の bimodule functor から構成されるので dg enhancement と両立する。
最後に定理8.4.1の分解をこの同値で移すと
式\eqref{blowup:quadric-unconditional-sod}を得る。
\end{proof}

\begin{remark}[適用範囲と generic localization]
この定理の仮定は、$E$ と構造層を標識した cubic $\mathbb Z$-代数表示を含む。
この表示を既に持つ曲面族では、categorical blowup と切断環の比較を
さらに仮定する必要がなくなる。
RSS の任意の hypersurface 表示の quadric に適用する際には、
その標識付き表示との同定を確認する。

点イデアル bimodule は generic localization で単位 bimodule になる。
従って Rees bimodule algebra の局所化は一変数の斉次環となり、
$\alpha_t$ は generic quotient 上で恒等同値を誘導する。
また $D_t=R\subset T_0$ という上の同定は、$\delta_t$ が誘導する
generic quotient の同一視を、この包含が誘導する分数体の同一視に固定する。
従って、$D=Q_0(Q)\cong Q_0(R)$ とし両側の局所化関手を $\ell_Q,\ell_R$ と書けば
\[
 \ell_R\circ R\delta_{t*}L\alpha_t^*\simeq\ell_Q
\]
となる。例外対象 $\mathcal E_t$ の generic rank は零である。
ここで使う $Q_0(R)=Q_0(T_0)$ は
\cite[定理1.1(2)]{BlowupRSS2015}の分数環の保存とも一致する。
\end{remark}

\begin{proposition}[Sklyanin pencil の一般の quadric への適用]
\label{blowup:generic-sklyanin-quadric}
$S_4$ を、滑らかな楕円曲線 $E$、次数4の直線束 $\mathcal L$、
無限位数の平行移動 $\alpha$ に付随する四次元 Sklyanin 代数とする。
その中心的二次元部分空間から得られる pencil
$Q=S_4/(\Omega)$ の非空 Zariski 開集合上で、任意の $t\in E$ に対し
定理\ref{blowup:quadric-section-equivalence}の比較同値と
式\eqref{blowup:quadric-unconditional-sod}が成立する。
\end{proposition}

\begin{proof}
滑らかな $\operatorname{qgr}Q$ における二つの正規化された rank-one
maximal Cohen--Macaulay module を $X,Y$ とする。
各々は次数0の二元で生成される。
syzygy による置換 $M\mapsto\Omega M(1)$ が $X,Y$ を交換するとき、
$Q$ は \emph{standard} と呼ばれる。
この二つの可能性と standard の定義は
\cite[命題5.3、定義5.4]{QuadricMU2021} による。
$X,Y$ は次数の移動を除いて $Q^{\oplus2}$ の部分加群となるので、
$g$ は両者の上で正則である
\cite[命題5.4、補題5.5]{SmithVdB}。

standard 条件が pencil の一般点で満たされることを、$E$ への制限で確認する。
一つの ruling の line module を $L_D=Q/V(D)Q$ とし、
$D$ は $E$ 上の次数2の有効因子、
$V(D)=H^0(E,\mathcal L(-D))$ とする。
ここでは右加群の規約を用いる。
この secant の記述は \cite[Basic Facts 7.2]{RSS2023} の右加群版による。
line module と spinor module の対応により
\[
 0\longrightarrow M(-1)\longrightarrow Q\longrightarrow L_D\longrightarrow0
\]
があり、$M$ は $X,Y$ のいずれかである。
\cite[命題7.2]{SmithVdB}
残る中心元 $g$ は $L_D$ 上で正則である。
実際 $L_D$ は2-critical であり、$gL_D=0$ なら
$B(E,\mathcal L,\alpha)$-module としての Hilbert 多項式の一次係数は
$4$ の非負整数倍になるが、line module の一次係数は1である。
従って $gL_D\ne0$ であり、criticality より $g$ の核は零である。
よって上の列は $E$ への制限後も完全である。
$\mathcal L(-D)$ は次数2で globally generated なので
\[
 Li^*M(-1)=\mathcal O_E(-D).
\]
$N=\mathcal O_E(D)$ とおく。
右加群の grading shift は
$F\mapsto\alpha_*(F\otimes\mathcal L)$ に対応するので
\[
 Li^*M=\alpha_*(\mathcal L\otimes N^{-1}).
\]
他方、\cite[補題5.5(2)]{SmithVdB} の列
\[
 0\longrightarrow M(-1)\longrightarrow Q^{\oplus2}
 \longrightarrow\Omega M(1)\longrightarrow0
\]
を制限する。$\Omega M(1)$ も $X,Y$ の一方であり、$g$ はその上で正則である。
従って制限後も完全で、行列式から $Li^*\Omega M(1)=N$ を得る。
従って $\Omega M(1)\cong M$ なら
\begin{equation}\label{blowup:nonstandard-necessary}
 N\otimes\alpha^*N\cong\mathcal L.
\end{equation}
写像 $\operatorname{Pic}^2E\to\operatorname{Pic}^4E$、
$N\mapsto N\otimes\alpha^*N$ は、平行移動を除けば multiplication by 2 である。
式\eqref{blowup:nonstandard-necessary}の解は高々4個である。
secant の因子類 $N\in\operatorname{Pic}^2E$ を固定すると、
対応する line module を annihilate する central quadric は一意である。
\cite[\S10.3]{SmithVdB} の左加群での記述に対し、
右加群の記述は \cite[Basic Facts 7.2]{RSS2023} に明記されている。
従って non-standard な $Q$ は、上の高々4個の $N$ が定める有限個の
pencil の元に含まれる。これらと非 smooth な元を除けば $Q$ は standard である。

この開集合上で、Mori--Ueyama の full geometric helix
\cite[定理5.16]{QuadricMU2021}と
$0\to X(-2)\to Q(-1)^{\oplus2}\to Y(-1)\to0$ を用いると
\[
 \bigl(\mathcal O_Q(-1),\widetilde Y(-1),\widetilde X(-1),\mathcal O_Q\bigr)
\]
という full strong exceptional collection が得られる。
強例外性と二つの中間対象の完全直交性は
同文献の補題5.6、5.9、5.13からも直接従う。
$E$ への制限は全て直線束で、その次数は順に $-4,-2,-2,0$ である。
従って最初の三頂点は再び hereditary fork になり、
補題\ref{blowup:quadric-kernel}の rank と degree の証明がそのまま適用できる。
残りの仮定は RSS の quadric elliptic surface の設定で満たされ、
$g$ は中心的な次数2の正則元である。
$Q$ は AS-Gorenstein で depth が3なので、
$Q_n=H^0(\operatorname{qgr}Q,\mathcal O_Q(n))$ for $n\ge0$ も成り立つ。
定理\ref{blowup:quadric-section-equivalence}の証明を適用して結論を得る。
\end{proof}

\begin{corollary}[共通切断環の二つの分解と三次元の積]
\label{blowup:rss-common-sod}
命題\ref{blowup:generic-sklyanin-quadric}の開集合から $Q$ を取り、
$\tau=\alpha^2$、$|\tau|=\infty$ とする。
任意の $t\in E$ に対して、RSS の同型
\cite[定理7.3]{RSS2023}を
\[
 R=Q^{(2)}(t)
   \cong S^{(3)}(\tau^{-1}p+\tau^{-1}q),
 \qquad \sigma^3=\tau,\qquad p\ne q
\]
と書く。ここで $S=S(E,\sigma)$ は三次元 Sklyanin 代数である。
$\mathcal A_U=\operatorname{Perf}(\operatorname{qgr}U)$ とおくと、
dg-enhanced な fully faithful functor
$\Phi_S\colon\mathcal A_S\to\mathcal A_R$、
$\Phi_Q\colon\mathcal A_Q\to\mathcal A_R$ と例外対象
$E_1,E_2,F\in\mathcal A_R$ が存在して
\begin{align}
 \mathcal A_R
  &=\langle E_2\otimes_k\operatorname{Perf}(k),\,
          E_1\otimes_k\operatorname{Perf}(k),\,
          \Phi_S(\mathcal A_S)\rangle,
       \label{blowup:rss-plane-sod}\\
 \mathcal A_R
  &=\langle F\otimes_k\operatorname{Perf}(k),\,
          \Phi_Q(\mathcal A_Q)\rangle.
       \label{blowup:rss-quadric-sod}
\end{align}
両 embedding は、RSS の環包含と上の同型が定める
\[
 Q_0(S)\cong Q_0(R)\cong Q_0(Q)
\]
という generic quotient の同一視に整合する。

さらに $\mathcal P=\operatorname{Perf}(\mathbb P^1)$ とし、
テンソル積は dg Morita テンソル積の冪等完備化とする。
$\mathcal B=\mathcal P\otimes_k^{\mathrm L}\mathcal A_R$ は
\begin{align}
 \mathcal B
   &=\langle\mathcal P,\,\mathcal P,\,
            \mathcal P\otimes_k^{\mathrm L}\mathcal A_S\rangle,
        \label{blowup:rss-product-plane-sod}\\
 \mathcal B
   &=\langle\mathcal P,\,
            \mathcal P\otimes_k^{\mathrm L}\mathcal A_Q\rangle.
        \label{blowup:rss-product-quadric-sod}
\end{align}
を持つ。各成分は上の fully faithful functor の積による像を表す。
三次元の generic quotient でも、$u$ を中心的な不定元として
$Q_0(S)(u)\cong Q_0(Q)(u)$ が保たれる。
\end{corollary}

\begin{proof}
quadric 側は命題\ref{blowup:generic-sklyanin-quadric}から従う。
平面側では $p_1=\tau^{-1}p$、$p_2=\tau^{-1}q$ とおく。
RSS の反復公式 \cite[命題5.4]{BlowupRSS2015} により
$S^{(3)}(p_1+p_2)=S^{(3)}(p_1)(p_2)$ である。
各段階の切断環は、上の次数1部分と $g$ による商の比較によって
Van den Bergh の切断環と一致する。
$p_1\ne p_2$ なので $\sigma p_1\ne\sigma p_2$ である。
二点の場合、これは \cite[定義11.1.2]{BlowupVdB2001} の一般位置条件を満たす。
従って同文献の定理11.1.3による二段階の比較同値と、
各段階の blowup 分解である定理8.4.1を合成すれば
式\eqref{blowup:rss-plane-sod}を得る。
$E_1$ は第一段階の例外対象を第二段階の fully faithful functor で移した対象であり、
$E_2$ は第二段階の例外対象である。
generic quotient での整合性は、各 Rees bimodule と切断環包含の
generic localization についての先の議論を二回適用すれば従う。

dg-enhanced な admissible decomposition の各 inclusion と adjoint を
$\mathcal P$ とテンソルすると、その単位・余単位および射影三角もテンソルされる。
従って完全忠実性、半直交性、生成性が保たれ、積の二つの分解を得る。
最後に $\mathbb P^1$ の generic point に制限し、
中心的な変数 $u$ を加えた Ore 分数環を取れば最後の主張を得る。
この系は固定された無限位数パラメータでの主張である。
\end{proof}


\subsection{三次元多様体上の order の blowup}
\label{blowup:orders}

有限中心上の非可換幾何では、三次元に直接適用できる blowup 公式がある。
次の既知の定理には次元2という制限がない。

\begin{theorem}[Belmans--Presotto]\label{blowup:order-theorem}
$X$ を smooth quasiprojective variety、$\mathcal A$ を
$X$ 上の次数 $n$ の locally free order とし、
$\operatorname{gldim}\mathcal A<\infty$ と仮定する。
$Z\subset X$ を余次元 $c\geq2$ の smooth closed subvariety とし、
$Z$ は $\mathcal A$ の ramification locus と交わらず、
$\mathcal A|_Z\cong\operatorname{Mat}_n(\mathcal O_Z)$ と仮定する。
$p\colon\widetilde X=\operatorname{Bl}_Z X\to X$ とおくと
$p^*\mathcal A$ も有限 global dimension を持ち、
\begin{equation}\label{blowup:order-sod}
 D^b(\widetilde X,p^*\mathcal A)
 =\left\langle
   D^b(Z),\ldots,D^b(Z),
   Lp^*_{\mathcal A}D^b(X,\mathcal A)
 \right\rangle
\end{equation}
を得る。$D^b(Z)$ は $c-1$ 個である。
\end{theorem}

\begin{proof}[既知結果への帰着]
\cite[定理30、注意31]{BlowupBP2018} の Orlov 型公式を、
標準的な負の twist による順序で記述する。
$j\colon E\hookrightarrow\widetilde X$ と $q\colon E\to Z$ を
exceptional divisor とその射影とすると、前の成分の埋込みは
Morita 同値を介して
\[
 j_{\mathcal A,*}
 \bigl(q_{\mathcal A}^*(-)\otimes\mathcal O_E(a)\bigr),
 \qquad a=1-c,\ldots,-1
\]
で与えられる。ここで $\mathcal O_E(1)$ は $E=\PP(N_{Z/X})$ の相対的な
正の tautological twist で、$\mathcal O_E(E)=\mathcal O_E(-1)$ である。
この順序と twist は第\ref{sec:classical-factorizations}節と同じ Orlov の規約にそろえた。
定理の証明では derived projection formula で $Lp^*_{\mathcal A}$ の
fully faithfulness を示し、中心の近傍で order を split する
\'etale 基底変換を用いて Orlov の blowup 公式に帰着する。
この順序で $Lp^*$ から例外成分への Hom が消えることは
$Rq_*\mathcal O_E(a)=0$ $(1-c\le a\le-1)$ から直接確認できる。
例外成分同士の半直交性と生成性も、その分裂近傍での Orlov の計算から従う。
\end{proof}

特に $\dim X=3$ なら、滑らかな曲線を中心とする場合には $D^b(Z)$ が一つ、
点を中心とする場合には $\operatorname{Perf}(k)$ が二つ加わる。
中心の外側で $(\widetilde X,p^*\mathcal A)$ は $(X,\mathcal A)$ と同一になるので、
generic central simple algebra $\mathcal A\otimes k(X)$ も同一になる。
その division algebra 部分を選べば、関数斜体の保存も同時に得られる。

中心が ramification locus に入ると、pullback order の有限 global dimension は
改めて検証を要する。\cite[注意34]{BlowupBP2018} は周期的射影分解を持つ例を与える。
本研究への適用には、対象の量子化が有限中心上の order として表示され、
選んだ中心が定理\ref{blowup:order-theorem}の局所条件を満たすことが必要である。
三次元 Sklyanin 代数にこの方法を用いる場合、対応する楕円曲線の自己同型が
有限位数となる finite-over-center な部分族が適用範囲となる。
\cite[定理5]{BlowupBP2018}

\subsection{三次元 flop と noncommutative crepant resolution}
\label{blowup:nccr}

三次元の birational geometry における非可換環のもう一つの主要な用法は
noncommutative crepant resolution (NCCR) である。
標数零の代数閉体上の normal Gorenstein affine domain $R$ に対し、
非零 reflexive $R$-module $M$ の endomorphism algebra
$A=\operatorname{End}_R(M)$ が maximal Cohen--Macaulay かつ
有限 global dimension を持つとき、$A$ は NCCR を与える。
Van den Bergh は、$\dim R=3$ で $R$ が terminal singularities を持つ場合、
NCCR の存在と可換 crepant resolution の存在が同値であり、
これらの bounded derived category が同値になることを証明した。
\cite[定義4.1、補題4.2、定理6.6.3]{BlowupVdBNCCR2004}

滑らかな三次元多様体の間の flop についても、perverse coherent sheaf の projective generator の
endomorphism algebra を介して derived equivalence を構成できる。
\cite[定理C、定理4.4.2]{BlowupVdBFlops2004}
この構成は共通の可換中心 $R$ に有限な代数を用いる。
$K=\operatorname{Frac}(R)$、$r=\operatorname{rank}_R M$ とすると
$A\otimes_R K\cong\operatorname{Mat}_r(K)$ であり、
generic category は $K$ 上の線形代数として記述される。
本研究の量子化へ適用するには、対象をこのような中心付き解消として実現することが必要となる。

\subsection{非可換 \texorpdfstring{$\PP^1$}{P1}-束の SOD}
\label{blowup:p1-bundles}

三次元の SOD を与えるもう一つの一般的な構成は、曲面を底とする非可換
$\PP^1$-束である。
$X$ を smooth projective scheme とし、$X\times X$ 上の coherent sheaf
$\mathcal E$ の support が両射影に関して有限で、
両方の pushforward が rank2 の locally free sheaf であるとする。
Van den Bergh の非可換 symmetric sheaf $\mathbb Z$-代数を
$\mathbb S(\mathcal E)$ と書く。
Mori--Okawa--Ueda は
\begin{equation}\label{blowup:mou-sod}
 D^b\qgr\mathbb S(\mathcal E)
 =\left\langle
      f_1^*D^b\operatorname{coh}X,
      f_0^*D^b\operatorname{coh}X
   \right\rangle
\end{equation}
を証明している。dual gluing functor は sheaf bimodule
$\mathcal E$ による tensor functor で記述される。
\cite[定理5.1]{BlowupMOU2019}

$\dim X=2$ とすれば三次元の非可換束が得られる。
中心上の有限性を用いずに構成できる点で、
本研究の Ore 拡大による $\PP^1$-束モデルとも比較すべき結果である。
具体的な比較には、対応する sheaf bimodule と次数付き表示の同定を要する。
\eqref{blowup:mou-sod} は束の SOD であり、birational map の
blowup 因子分解を得るにはさらに中心と収縮 functor の構成が必要となる。

\subsection{任意次元の categorical resolution と refined blowup}
\label{blowup:categorical}

Kuznetsov--Lunts は標数零の体上の任意の separated finite-type scheme $X$ に
categorical resolution を構成した。
resolving category は smooth scheme の derived category を成分とする
有限 SOD を持ち、$X$ が proper のとき resolving category も proper である。
構成には可換 scheme の smooth-center blowup、非還元中心の厚みを記録する
Auslander algebra、dg gluing を用いる。
\cite[定理1.4]{BlowupKL2015}

Kaledin--Kuznetsov は、この局所的な一段階を filtered derived category により
refined blowup として記述した。
$X$ が noetherian で、任意の coherent sheaf が有限階数の locally free sheaf の商となると仮定し、
$Z=V(\mathcal I)$ と $Y=\operatorname{Bl}_{\mathcal I}X$ をとる。
十分大きい整数 $N$ に対して圏 $\operatorname{Bl}_N(X,Z)$ と fully faithful functor
\[
 \operatorname{Perf}(X)\longrightarrow\operatorname{Bl}_N(X,Z)
\]
が存在し、
\begin{equation}\label{blowup:refined-sod}
 \operatorname{Bl}_N(X,Z)
 =\left\langle D^b\operatorname{coh}Y,
   \underbrace{D^b\operatorname{coh}Z,\ldots,D^b\operatorname{coh}Z}_{N\text{ 個}}
   \right\rangle
\end{equation}
となる。$N$ は選んだ refinement の大きさである。
\cite[定義1.7、命題1.8、式(1.9)]{BlowupKK2015}

これらの結果は三次元の可換 scheme から構成する非可換解消にも適用できる。
本稿の量子化環からの適用には、その環または射影圏を
上記の中心付き構成に実現するデータが必要となる。
関数斜体の同型のみから、refined blowup の中心や gluing bimodule を復元する問題は残る。

\subsection{一般の Rees 構成と naive blowup}
\label{blowup:naive}

一般の非可換環では、両側イデアル $I$ の通常の Rees 環
$\bigoplus_{m\geq0}I^m$ の exceptional fiber の次元が変化し得る。
Van den Bergh の二変数の例を中央変数 $z$ で拡張すると
\[
 A=k\langle x,y,z\rangle/
 (xy-yx-y,\ xz-zx,\ yz-zy),\qquad I=(x,y,z)
\]
を得る。この環は PBW 基底を持ち $\operatorname{GKdim}A=3$ である。
$y=xy-yx$ により帰納的に $y\in I^m$ が全ての $m\geq1$ で成り立つ。
$A/(y)=k[x,z]$ なので
\[
 I^m/I^{m+1}\cong (x,z)^m/(x,z)^{m+1},\qquad
 \operatorname{gr}_I A\cong k[\bar x,\bar z].
\]
従って通常の Rees 構成の exceptional fiber は $\PP^1$ となる。
この計算は、三次元の点 blowup に対し $\PP^2$ 型の exceptional fiber を得るには、
中心と Rees 構成の追加条件が必要となることを示す。
元の二変数の例と twisted Rees 構成は
\cite[\S1.1]{BlowupVdB2001} を参照されたい。

任意次元の射影多様体から得られる環論的な別の構成として、
Keeler--Rogalski--Stafford の naive blowup がある。
$X$ を integral projective variety、$\dim X\geq2$ とし、
$\sigma\in\operatorname{Aut}(X)$、very ample かつ $\sigma$-ample な $\mathcal L$、
critically dense な無限 orbit を持つ点 $p$ を選ぶ。
$\mathcal I=\mathcal I_p$、
$\mathcal L_m=\bigotimes_{i=0}^{m-1}(\sigma^i)^*\mathcal L$ とおくと
\[
 R=\bigoplus_{m\geq0}
 H^0\!\left(X,
   \mathcal I\,\sigma^*\mathcal I\cdots(\sigma^{m-1})^*\mathcal I
   \otimes\mathcal L_m\right)
\]
に twisted multiplication を入れられる。
$R$ は noetherian となる一方で、
$\dim_k\operatorname{Ext}^1_{\qgr R}(\pi R,\pi R)=\infty$ が成り立つ。
\cite[定理1.1]{BlowupKRS2005}

この構成は $\dim X=3$ にも適用できるが、上の Ext により
$D^b(\qgr R)$ の properness は破れる。
また generic trivialization を選ぶと $R\subset k(X)[t;\sigma]$ となるため、
次数零分数斜体は可換体 $k(X)$ の部分体に入る。
従って本研究で現れる一般の非可換 $Q_0(A)$ と比較する際には、
この birationally commutative な性質を確認する必要がある。

さらに一般の quasi-scheme 上で中心を扱うため、Rogalski は
closed subcategory の well-closedness を研究した。
これは中心の ideal との tensor product に相当する右完全 functor の存在条件であり、
tiny simple object $P$ に対する
$\dim_k\operatorname{Ext}^1(M,P)<\infty$ が全ての noetherian $M$ で成り立つとき、
対応する closed point は very well-closed になる。
\cite[定理1.3、定理6.6]{BlowupRogalski2019}
この条件は Rees 型操作の定義に必要な圏論的基盤を与える。
本稿の三次元モデルへの応用では、その後に noetherianity、smoothness、
generic localization、exceptional component を検証する段階が残る。

\subsection{本研究における因子分解の範囲}
\label{blowup:scope}

中央 $\PP^1$ との積で表されるモデルでは、
曲面の blowup と blowdown を積に移して三次元の因子分解を構成できる。
SOD は命題\ref{blowup:product-proposition}の仮定の下で得られ、
有限中心上の order については定理\ref{blowup:order-theorem}が直接適用できる。
トーラス型では第\ref{sec:toric-nc-blowups}節の全 Rees blow-up 列がこの条件を満たす。
楕円型では第\ref{sec:relative-common-blowup}節の相対共通圏と、
第\ref{blowup:quadric-comparison}節の固定ファイバーの共通環を得た。
残る成分間の量子化で分数斜体の同型が判明している各対に対して、
中心付きの非可換 blowup を実際に構成し、射影圏と斜体の両方に整合する
functor を与えることが次の課題となる。
この課題は、Poisson blowup が存在すること、代数的量子化族が存在すること、
関数斜体が同型になることにそれぞれ追加される条件である。

\section{SOD と関数斜体を結ぶ局所化}
\label{sec:sod-functions}

\subsection{零ランク対象による商}

関数斜体を回復する際には、圏とともに零ランク対象の部分圏を指定する。
まず次数付き環について、この操作を明示する。

\begin{proposition}\label{prop:generic-serre}
$A$ を局所有限な連結次数付き noetherian domain とし、
非零斉次元についての左右 Ore 局所化 $Q_{\rm gr}(A)$ が存在するとする。
さらに $A_1\ne0$ とし、$D=Q_{\rm gr}(A)_0$ と置く。
$\mathcal T_\eta\subset\qgr A$ を、斉次局所化が零となる有限生成次数付き加群の像が
なす Serre 部分圏とする。このとき
\[
 (\qgr A)/\mathcal T_\eta\simeq\operatorname{mod}_{\rm fd}D.
\]
右辺は右 $D$-ベクトル空間として有限次元な加群の圏である。
\end{proposition}

\begin{proof}
Ore 局所化は完全であり、その核は各元がある非零斉次元で消える加群である。
局所化環上の有限生成加群とその射には有限個の生成元・関係・分母だけが現れるので、
Ore 条件による分母払いにより元の次数付き加群から得られる。
従って局所化はその核による Serre 商を与える。
有限次元次数付き $A$-加群はこの核に含まれるから、$\qgr A$ を経由できる。
非零 $t\in A_1$ は $Q_{\rm gr}(A)$ で可逆であり、
$Q_{\rm gr}(A)=D[t,t^{-1};\sigma]$ は強次数付き環となる。
次数零成分を取る操作と、$-\otimes_DQ_{\rm gr}(A)$ は互いに逆な同値である。
これを先の Serre 商と合成すればよい。
\end{proof}

この命題では二種類の torsion が現れる。
$\qgr A$ を定義する際には有限次元加群を除き、generic な商では斉次ランク零の
加群を全て除く。例えば非可換曲面の曲線に支えられる加群は後者に含まれる。

導来圏については、次の Ore 局所化の形が直接使える。

\begin{proposition}\label{prop:generic-perf}
$B$ を左右 Ore domain、$D=\Frac B$ とする。
$\mathcal N$ を $P\otimes_B^LD=0$ を満たす
$P\in\operatorname{Perf}(B)$ の thick 部分圏とすると、
\[
 \left(\operatorname{Perf}(B)/\mathcal N\right)^\natural
 \simeq\operatorname{Perf}(D).
\]
ここで $(-)^\natural$ は冪等完備化を表す。
\end{proposition}

\begin{proof}
$D$ は $B$ の非零元の右加群自己準同型 $B\to B$ を可逆化した環である。
左右 Ore 局所化は平坦なので $\operatorname{Tor}^B_i(D,D)=0$ $(i>0)$ である。
Neeman--Ranicki \cite[定理0.7]{NeemanRanicki} をこの可逆化に適用する。
同定される核はこれらの射の cone が生成する thick 部分圏であり、
同定後はちょうど $-\otimes_B^LD$ の核 $\mathcal N$ である。
\end{proof}

滑らかな整射影多様体 $X$ についても、generic point で零となる perfect complex の
部分圏を $\operatorname{Perf}_{\rm tors}(X)$ とすれば
\begin{equation}\label{eq:comm-generic-quotient}
 \left(\operatorname{Perf}(X)/\operatorname{Perf}_{\rm tors}(X)\right)^\natural
 \simeq\operatorname{Perf}(k(X))
\end{equation}
である。実際、generic point 上の有限次元ベクトル空間と射は稠密開集合に広がり、
有限個の射の分母はさらに小さい稠密開集合を取れば消える。
perfect complex もそのような開集合で cohomology の直和へ分裂する。
その補集合に支えられる cone を商で零とすれば、局所化した圏の対象は
$\mathcal O_X$ の有限直和とシフトから生成される。
その自己準同型は $k(X)$ となり、高次自己準同型は稠密なアフィン開集合への制限で消える。
これにより式\eqref{eq:comm-generic-quotient}を得る。

本稿の斉次環が normal な元の局所化からアフィン Ore チャートを持つ場合には、
このチャートと命題\ref{prop:generic-perf}を用いて generic な導来圏を計算できる。
一般の smooth proper dg 圏に対しては、同様の generic localization のデータを
追加することが関数斜体を論じるための条件となる。

\subsection{blow-up の例外部分が局所化で消える理由}

滑らかな曲面の一点 blow-up $f:\widetilde S\to S$ では
\[
 \operatorname{Perf}(\widetilde S)=
 \left\langle\mathcal O_E(-1),Lf^*\operatorname{Perf}(S)\right\rangle.
\]
例外対象 $\mathcal O_E(-1)$ は generic rank が零である。
$f$ は稠密開集合で同型であり、$Lf^*$ はその開集合上で恒等関手になる。
従って generic quotient に移ると $k(\widetilde S)=k(S)$ が回復される。
この記述は次の形で非可換の場合にも適用できる。

\begin{proposition}\label{prop:sod-compatible-generic}
冪等完備な pretriangulated dg 圏 $\mathcal C,\widetilde{\mathcal C}$ を取り、
$\widetilde{\mathcal C}$ に admissible な SOD
\[
 \widetilde{\mathcal C}=\langle\mathcal E,F(\mathcal C)\rangle
\]
があり、両側に generic localization
$q:\mathcal C\to\mathcal C_\eta$、
$\widetilde q:\widetilde{\mathcal C}\to\widetilde{\mathcal C}_\eta$ が指定されているとする。
ここで局所化先は対応する Verdier 商の冪等完備化とする。
$\mathcal E\subset\ker\widetilde q$ であり、$F$ が零ランク対象を保ち、
誘導される局所化関手が fully faithful ならば、
\[
 \mathcal C_\eta\simeq\widetilde{\mathcal C}_\eta.
\]
特にこの圏がそれぞれ関数斜体の perfect 圏ならば、両斜体は同型である。
\end{proposition}

\begin{proof}
SOD の生成性より、$\widetilde{\mathcal C}_\eta$ は
$\widetilde qF(\mathcal C)$ の像で生成される。
誘導関手の fully faithfulness と冪等完備性から同値が従う。
斜体 $D$ 上の perfect complex は、有限次元右 $D$-ベクトル空間のシフトの直和である。
したがって indecomposable な対象は $D[n]$ であり、その自己準同型環が $D$ になる。
同値はこの自己準同型環を保存する。
\end{proof}

Van den Bergh 型の曲面 blow-up では、適切な点の補集合上の比較がこの局所化との
両立性を与える。ring-theoretic な common blow-up については、
その圏とこの blow-up 圏との比較および smoothness を確認して適用する。
第\ref{blowup:quadric-comparison}節は一点 quadric の比較を与える。
相対的な共通圏では命題\ref{prop:relative-common-field}により、
二つの局所化が同じ kernel を持つことから斜体同型が従う。
トーリック共通模型では系\ref{cor:toric-nc-common}の単項式局所化が
同じ量子トーラスを与える。

\subsection{同じ SOD 成分と異なる関数斜体}

\begin{proposition}\label{prop:sod-does-not-determine-field}
$q\in\CC^*$ に対して
\[
 S_q=\CC\langle x,y,z\rangle/(xy-qyx,\ xz-zx,\ yz-zy)
\]
と置く。全ての $q$ に対して
\[
 D^b(\qgr S_q)=\langle\mathcal O,\mathcal O(1),\mathcal O(2)\rangle
\]
は full strong exceptional collection である。
その関数斜体は
\[
 D_q=\Frac\bigl(\CC\langle u,v\rangle/(uv-qvu)\bigr).
\]
$q=1$ の関数体は $\CC(u,v)$ であり、$q\ne1$ の関数斜体は非可換である。
従って SOD の各成分の同値型と個数から関数斜体を決定することはできない。
\end{proposition}

\begin{proof}
$S_q$ は三変数の skew-polynomial 環であり、順序付き単項式を基底に持つ
Koszul AS 正則環で、大域次元は三である。
三つの normal な変数による次数付き \v{C}ech 複体から
$H^1(\mathcal O(r))=0$、$H^2(\mathcal O(r))=0$ $(r\ge-2)$、
$H^0(\mathcal O(r))=(S_q)_r$ を得る。
これにより三つの対象の強例外性と半直交性が従う。
Koszul 分解は $\qgr S_q$ で全ての $m$ に対して
\[
 0\longrightarrow\mathcal O(m-3)\longrightarrow\mathcal O(m-2)^3
 \longrightarrow\mathcal O(m-1)^3\longrightarrow\mathcal O(m)\longrightarrow0
\]
を与えるので、この三対象は全ての twist を生成する。
有限生成次数付き加群の有限自由分解により full である。
$z$ は中心元であり、次数零の局所化は
$u=xz^{-1}$、$v=yz^{-1}$ を生成元とする量子平面になる。
$q\ne1$ なら $uv-vu=(q-1)vu\ne0$ であり、domain の分数環への埋込み後も非零である。
\end{proof}

この例で変化するのは SOD の貼合せである。
頂点 $0,1,2$ と、それぞれ $i\to i+1$ の三本の矢
$x_i,y_i,z_i$ $(i=0,1)$ を用いる Beilinson 代数の関係式は
\[
 x_1y_0=qy_1x_0,\qquad x_1z_0=z_1x_0,\qquad y_1z_0=z_1y_0.
\]
$z_0,z_1$ を可逆化して三頂点を同一視すると、その corner は量子平面となる。
さらに非零元を可逆化して $D_q$ を得る。
従って貼合せ代数と可逆化する射を保持すれば、関数斜体を回復できる。

generic quotient の後には、半直交性自体も通常は残らない。
例えば $\mathcal O$ と $\mathcal O(1)$ はどちらも generic point 上で一つの
自由 $D_q$-加群に移り、両方向の Hom が非零になる。

\begin{remark}\label{rem:generic-not-proper}
$D$ が $\CC$ 上無限次元の関数斜体ならば、$\operatorname{Perf}(D)$ を
proper な $\CC$-線形 dg 圏へ fully faithfully 埋め込むことはできない。
実際 $\operatorname{End}_D(D)=D$ は無限次元であり、properness の Hom の有限次元性と
両立しない。generic な圏は式\eqref{eq:comm-generic-quotient}や
命題\ref{prop:generic-perf}の商として現れる。
\end{remark}

曲面ではこの局所化、Van den Bergh の blow-up SOD、Sklyanin 曲面の明示的な
双有理変換が利用できる。Rains の非可換曲面論が与える factorization 型の結論には、
共通の反復 blow-up や quasi-ruled 模型に関する仮定が付く。
本稿ではトーラス型と楕円型に共通模型を構成した。
全成分の関数斜体同型をそのような共通模型で実現する問題は残る。

\section{反標準環からの導来圏の復元：正確な判定と障害}
\label{sec:restored-anticanonical}
第\ref{sec:anticanonical-common-ring}節は、反標準 orbit 環を
二点 blow-up 環と同定する。しかし環の同型とカテゴリーの同値は別の命題である。
その間に必要な条件を、既知の復元定理を用いて明示する。

\begin{theorem}[反標準 section functor の十分条件]\label{thm:restored-az}
$\mathcal H$ を $k$-線形、局所 noetherian な abelian category とし、
$O\in\mathcal H$ を noetherian object、$F:\mathcal H\to\mathcal H$ を
完全な自己同値とする。
任意の noetherian objects の間の Hom が有限次元であり、
$(O,F)$ が Artin--Zhang の意味で ample、すなわち
\begin{enumerate}
\item 任意の noetherian な $M\in\mathcal H$ に対し、有限個の $n_i\ge0$ と全射
$\bigoplus_iF^{-n_i}O\twoheadrightarrow M$ が存在する。
\item 任意の noetherian objects の全射 $M\twoheadrightarrow N$ について、
$n\gg0$ のとき
$\operatorname{Hom}(F^{-n}O,M)\to\operatorname{Hom}(F^{-n}O,N)$ が全射となる。
\end{enumerate}
と仮定する。
$R=\bigoplus_{n\ge0}\operatorname{Hom}_{\mathcal H}(O,F^nO)$ に
$u\cdot v=F^{\deg v}(u)\circ v$ で積を入れると、
$R$ は右 noetherian であり、自然な section functor は
$\mathcal H\simeq\operatorname{qgr}R$ を誘導する。
さらに、ある smooth proper dg 圏 $\mathcal C$ が
$D^b(\mathcal H)$ の指定 dg enhancement に同値ならば、
$\mathcal C\simeq D^b(\operatorname{qgr}R)$ である。
\end{theorem}
\begin{proof}
最初の結論は Artin--Zhang の ample pair による復元定理である
\cite{KeelerNCRings,ArtinZhang1994}。
後半は同じ abelian equivalence を有界導来圏へ適用する。
この定理は $F$ が \emph{指定した heart 上で t-exact} であることを前提とする。
三角圏上の自己同値であるというだけでは仮定を満たさない。
\end{proof}

\begin{corollary}[本研究に必要な検証項目]\label{cor:restored-az-criteria}
第\ref{sec:anticanonical-common-ring}節の階数23の共通圏 $\mathcal C_B$
について、反標準関手 $F=S_{\mathcal C_B}^{-1}[2]$ と
構造対象 $O$ が、ある noetherian な bounded heart $\mathcal H_B$ に対して
$F$ の t-exactness と定理\ref{thm:restored-az}の二つの ample 条件を満たし、
$D^b(\mathcal H_B)\simeq\mathcal C_B$ が成り立てば、
\[
\mathcal C_B\simeq D^b(\operatorname{qgr}\mathscr A),
\qquad
\mathscr A=\bigoplus_{n\ge0}\operatorname{Hom}(O,F^nO)
\]
となる。特に固定された無限位数ファイバーにおける
$\mathscr A\simeq\mathcal S^{(3)}(\tau^{-1}p+\tau^{-1}q)$ の
\emph{環としての}同定からこの結論を得るには、上記の heart と ample 性の検証が必要である。
\end{corollary}

\begin{proposition}[五つの section module の比較の十分条件]
$R$ を右 noetherian な次数付き代数とし、
$\mathcal C_B=\operatorname{Perf}(B)$ に五個の full exceptional objects
$E_1,\ldots,E_5$ を選ぶ。
$\operatorname{qgr}R$ 上の五個の objects $M_i$ が、
\[
\operatorname{RHom}_{\operatorname{qgr}R}(M_i,M_j)
 \simeq \operatorname{RHom}_{\mathcal C_B}(E_i,E_j)
\]
を合成と単位まで保つ dg 代数として実現し、
$\bigoplus_iM_i$ が $D^b(\operatorname{qgr}R)$ の classical generator
であるとする。このとき
$D^b(\operatorname{qgr}R)\simeq\mathcal C_B$ が成立する。
\end{proposition}
\begin{proof}
五個の対象の full dg endomorphism algebra を比較し、
導来 Morita 理論を適用する。
単に Euler 行列や $K_0$ 上の Serre 作用が一致するだけでは、
full dg endomorphism algebra の一致は従わない。
\end{proof}

\begin{remark}[再開時の有限入力]
もし定理\ref{thm:restored-az}の heart が得られれば、
各 $E_i$ に対して
\[
M_i=\bigoplus_{n\ge0}\operatorname{Hom}_{\mathcal H_B}(O,F^nE_i)
\]
を考えられる。$u\in\operatorname{Hom}(O,F^nE_i)$ と
$r\in R_m$ の右積は $u\cdot r=F^m(u)\circ r$ である。
十分大きな共通の次数シフトを選び、五つの加群の有限表示、
全高次 $\operatorname{Ext}$、生成性、各積 $M_i\otimes M_j\to M_k$
に相当する合成の比較を調べればよい。
第\ref{sec:anticanonical-common-ring}節に既にある
$\mathcal S_3$ の二点 evaluation の核と42本の二次関係式は、
この有限計算の入力となる。
ただし\emph{すべての次数についての生成性}と
$F$ の t-exactness は現時点で検証されていない。
\end{remark}

\begin{remark}[反標準環だけでは圏を復元できない例]\label{ex:restored-weak-dp}
$Y$ を $\mathbb P^2$ の一点 $p$ を blow-up し、
さらに最初の例外曲線上の点 $q$ を blow-up した滑らかな曲面とする。
$H,E_1,E_2$ を pullback basis とすると
\[
 -K_Y=3H-E_1-E_2,\qquad C=E_1-E_2,
 \quad C^2=-2,\quad (-K_Y)\cdot C=0.
\]
ここで $C$ は最初の例外曲線の狭義変換である。
$-K_Y$ は big かつ nef だが ample ではなく、反標準射は $C$ を収縮する。
反標準環 $R(Y,-K_Y)$ の $\operatorname{Proj}$ は $A_1$ 型特異点を持つ
反標準模型になる。従って
$\operatorname{coh}Y$ を $\operatorname{qgr}R(Y,-K_Y)$ と
反標準 section functor で同一視することはできない。
実際、右辺は特異な反標準模型の coherent sheaves であり、
$Y$ の滑らかな幾何をそのまま回復しない。
これは固定ファイバーの環の同型だけから圏の同値を
推論できないことを示す可換な障害例である。
\end{remark}

\section{2026年の一次文献の適用範囲}
\label{sec:restored-primary}
この節は2026年11月8日に公開原著の本文を確認した二つの結果を記録する。
先行結果は本研究の新しい定理として扱わない。

\begin{theorem}[Mori--Ueyama--Wu、既知]\label{thm:restored-muw}
$S$ を非退化な四次元 Sklyanin 代数とし、
$0\ne f\in k\Omega_1+k\Omega_2\subset Z(S)_2$ とする。
$S/(f)$ が smooth な central Sklyanin quadric ならば standard である。
また固定した $S$ について、四つの特異な central quadric の
方程式を明示し、それらが相互に次数付き代数として同型である。
\end{theorem}
\noindent
これは\cite[定理1.3、定理5.17]{MoriUeyamaWu2026}による。
無限位数条件なしでの standardness は、smooth な
\emph{central Sklyanin quadric} に対して成立する。
任意の normal quadratic quotient の standardness、
任意の非可換 quadric の smoothness、
あるいは本研究の導来 blow-up 比較を主張するものではない。
本書第\ref{sec:nc-blowups}節で扱う quadric が
この定理の central・nondegenerate の仮定を満たす場合に限り使える。

\begin{theorem}[De Deyn、既知]\label{thm:restored-dedeyn}
$k$ を標数零の体、$X$ を reduced separated finite type $k$-scheme、
$\mathcal A$ を $X$ 上の order とする。
このとき smooth separated quasi-compact Deligne--Mumford stack
$\mathcal X$、その上の Azumaya algebra $\mathcal B$、
および stacky blow-up の合成である固有双有理射 $\mathcal X\to X$
が存在し、$\pi^*\mathcal A\to\mathcal B$ は生成点で同型である。
\end{theorem}
\noindent
これは\cite[定理1.1]{DeDeyn2026}である。
本結果は任意次元での\emph{有限な order} を対象とし、
root stack と blow-up を組み合わせる。
一般の Sklyanin 代数の $\operatorname{qgr}$ や、中心が $k$ しかない
量子分数斜体に対して直ちに適用できるものではない。
さらに、Azumaya への解消は、そのまま指定した SOD や
非可換関数体の双有理同型を与えるわけではない。

本研究の Rees blow-up への直接適用を主張する前に、
まず各量子化の generic algebra が可換基礎スキーム上で
有限階数の order として表示できるかを判定する必要がある。
そうした有限中心条件がない場合には、
第\ref{thm:restored-categorical-gluing}定理の dg 圏論的構成と
第\ref{prop:restored-associated-graded}命題の斉次環模型を区別する。


\part{非可換有理性と双有理分類の展望}
\section{非可換有理性の定義と射影三空間への双有理性}\label{sec:nc-rationality}
\subsection{定義の区別}
可換な射影三次元多様体 $X$ の有理性は $\CC(X)\simeq\CC(t_1,t_2,t_3)$ で
定まるが、この右辺の非可換な標準模型は一意には存在しない。
本節では、連結 $\mathbb N$ 次数付き、有限生成な noetherian Ore domain $A$
に対する次数零分数斜体 $D(A)=Q_{\mathrm{gr}}(A)_0$ を用い、
同じ $\CC$-斜体を持つとき\emph{斜体的双有理同値}と呼ぶ。
一般の smooth proper dg 圏については、tilting 代数の有限次元性だけから
このような斜体が存在するとは限らない。従って圏量子化のみの場合に
「非可換関数体」を無条件に割り当てることはしない。

\begin{definition}[AS-有理性、研究上の定義案]\label{def:as-rational}
$A$ を上記の domain で GK 次元4のものとする。
$\qgr A$ が\emph{AS-有理的}であるとは、
ある連結次数付き noetherian AS-regular domain $B$ が存在し、
$\operatorname{gldim}B=4$、$H_B(z)=(1-z)^{-4}$、かつ
$D(A)\simeq D(B)$ が成り立つことである。
これは $A$ の明示した斉次表示における定義であり、
任意の dg 圏に同じ条件が内在的に定まると主張するものではない。
\end{definition}

$B$ を特定の量子多項式環に限定した定義は、上の存在条件より強い。
また古典的な rational 三次元多様体を特殊ファイバーに持つことを
\emph{変形有理性}と呼ぶこともできるが、これは斜体的双有理性とは別の性質である。
次数零斜体の構成が存在する場合、Veronese を取っても斜体は不変である。

\begin{lemma}[Veronese 不変性]\label{lem:veronese-function}
$A$ を連結次数付き domain とし、各非負次数に非零な斉次元が存在し、
斉次元からなる Ore 局所化が存在するとする。任意の $d\ge1$ について
$D(A^{(d)})\simeq D(A)$ となる。
\end{lemma}
\begin{proof}
$D(A)$ の元を、同次数の非零斉次元 $a,b$ による右分数 $ab^{-1}$ で表す。
$\deg a=\deg b=n$ として、$n+m\equiv0\pmod d$ となる $m\ge0$ の
非零な斉次元 $c$ を選べば、
$ab^{-1}=(ac)(bc)^{-1}$ であり、分子と分母はともに $A^{(d)}$ に属する。
逆の包含は明白である。
\end{proof}

\subsection{非可換射影三空間自身の異なる双有理型}
\begin{proposition}[量子型と Weyl 型の区別]\label{prop:distinct-nc-p3}
$k=\CC$ とし、$q\in k^\times$ は1の冪根でないとする。
次数1の生成元 $x_0,x_1,x_2,x_3$ からなる二つの次数付き代数を
\[
 B_q=\frac{k\langle x_0,x_1,x_2,x_3\rangle}
 {\bigl([x_0,x_i],[x_3,x_i],\;x_1x_2-qx_2x_1\bigr)},
 \qquad
 B_{\mathrm W}=\frac{k\langle x_0,x_1,x_2,x_3\rangle}
 {\bigl([x_0,x_i],[x_3,x_i],\;x_1x_2-x_2x_1-x_0^2\bigr)}
\]
で定める。ただし $[x_0,x_i]$ と $[x_3,x_i]$ は
その変数を中央にする全関係式を表す。
いずれも Hilbert 級数 $(1-z)^{-4}$ の4次元 AS-regular domain であるが、
$D(B_q)\not\simeq D(B_{\mathrm W})$ である。
\end{proposition}
\begin{proof}
$B_q$ は量子多項式環である。$B_{\mathrm W}$ は
$k[x_0,x_1,x_3]$ の導分 $\delta(x_1)=-x_0^2$、
$\delta(x_0)=\delta(x_3)=0$ による次数付き Ore 拡大である。
従っていずれも PBW 基底を持ち、Hilbert 級数は指定のもので、
標準的な量子多項式環／斉次 Weyl 環の AS-regular 性が適用される。
$x_0\ne0$ の次数零座標を $u=x_1x_0^{-1}$、$v=x_2x_0^{-1}$、
$w=x_3x_0^{-1}$ とすれば、
\[
 D(B_q)\simeq\Frac\bigl(k(w)\langle u,v\rangle/(uv-qvu)\bigr),
 \quad
 D(B_{\mathrm W})\simeq\Frac\bigl(k(w)\langle u,v\rangle/(uv-vu-1)\bigr).
\]
スカラー交換子不変量 $\AD_k(D)=[D^\times,D^\times]\cap k^\times$ は
第一の場合に $\langle q\rangle$、第二の場合に $\{1\}$ である。
後者は Weyl 斜体の標準的な Alev--Dumas 不変量であり、
前者は量子トーラスの同不変量である
（\cite{GoodearlYakimov,StaffordVanDenBergh2001} 参照）。
したがって二つの斜体は同型でない。
\end{proof}

\begin{remark}
上の命題は「非可換 $\PP^3$ の関数斜体」を一つに固定できないことを示す。
自由代数の universal free skew field を「非可換有理関数体」と呼ぶ
別の用語法もあるが、これは本書の AS-regular な射影三次元幾何とは
異なる概念である。
\end{remark}

\subsection{変形と Poisson 双有理性の必要条件}
\begin{proposition}[半古典極限の必要条件]\label{prop:semiclassical-necessity}
$\hbar$-torsion-free で完備な二つの可換極限を持つ量子化の
generic fraction algebra $\mathscr D_1,\mathscr D_2$ に、
$\hbar$ を固定し、指定した $\hbar$-整格子を相互に保つ同型 $\Phi$ が存在するとする。
このとき $\Phi$ の $\hbar$ についての還元は可換関数体の Poisson 同型を誘導する。
\end{proposition}
\begin{proof}
$[a,b]\in\hbar\mathscr D_i$ と
$\{\bar a,\bar b\}=\overline{\hbar^{-1}[a,b]}$ の定義を用いる。
$\Phi$ が格子および $\hbar$ を保つため、還元した同型は括弧を保つ。
\end{proof}
\begin{remark}
格子を保たない $k((\hbar))$-斜体の抽象的同型だけから
Poisson 双有理性を結論することはできない。また $X$ が rational でも、
その Poisson 括弧を $\PP^3$ の稠密開集合へ移すと通常は有理的な極を持ち、
$\PP^3$ 全体で正則な双ベクトルにはならない。
この極を解消してなお $H_B(z)=(1-z)^{-4}$ の AS-regular 模型を得ることは
追加の大域的問題である。
\end{remark}

\subsection{三次元における AS-有理性予想の正確な範囲}
本書のトーラス型では量子多項式環の関数斜体と一致し、
line--conic 型と Riccati 型にも、指定した非可換 $\PP^3$ の特殊族との明示的同型がある。
一方、垂直楕円型と $R(2,2)$ Sklyanin pencil との比較には
中心体上の有限拡大からの下降が残る。
三重積の最大 pencil 成分については、Segre 量子化と階数64の共通圏の
関数斜体を同定する問題が残る。
これらを解くことは、可換極限が rational な三次元量子化が
AS-有理的になるための具体的な十分条件を与えることに相当する。

\begin{quote}
\textbf{研究問題（量子有理性問題）.}
滑らかな rational 三次元多様体の integral な斉次量子化は、
高次の量子化データを適合させたとき、
何らかの4次元 AS-regular 非可換射影三空間に斜体的に双有理同値となるか。
さらに、古典的な双有理写像を保つ $\hbar$-整な同型として実現できるか。
\end{quote}
これは現段階では予想としても根拠が十分でないため、
肯定・否定いずれも未確定な研究問題として扱う。
Sklyanin に由来する楕円的付値、量子トーラスのスカラー交換子群、
中心の超越次数、Riccati 導分の微分体としての性質、
generic fiber 上の Galois 下降などを判別手段とする。
非可換曲面の研究では、Presotto--Van den Bergh の
Sklyanin 平面と quadric の同じ斜体を持つ構成
\cite{PVdB2016}、Rogalski--Sierra--Stafford の極小模型
\cite{RogalskiSierraStafford2020}、Artin の関数斜体の分類問題
\cite{StaffordVanDenBergh2001} が先行する。
これらは三次元に対する一般の肯定解を与えるものではない。

\section{トーリック量子化の AS-有理性と Casimir 体}\label{sec:toric-as-rational}
この節では、前節の AS-有理性の定義案を適用できる模型の範囲を
三次元トーリック多様体へ拡げる。量子トーリック多様体の張り合わせの
一般構成は Cirio--Landi--Szabo~\cite{CLS} による。
以下で新たに示す同型はその構成の関数斜体レベルの直接的な帰結であり、
既存文献に対する新規性は主張しない。

\subsection{双指標と任意のトーリックコンパクト化}
$k$ を標数零の体、$M\simeq\mathbb Z^3$ を指標格子とし、
$c:M\times M\to k^\times$ を双指標とする。
$k_c[M]$ は基底 $X^m$、積
\[
 X^m X^n=c(m,n)X^{m+n}
\]
をもつ twisted group algebra である。
$c$ の交代化 $q(m,n)=c(m,n)c(n,m)^{-1}$ は
$q(e_i,e_j)=q_{ij}$ をみたす。別の双指標で同じ交代化を持つものは、
$M$ が自由アーベル群であるため、基底の再規格化により同型な twisted group algebra を与える。
$T_c$ は $k_c[M]$ の Ore 分数斜体とする。

\begin{theorem}[トーリック AS-有理性の十分条件]\label{thm:toric-AS-rational}
$X$ を稠密トーラス $T=\operatorname{Spec}k[M]$ を持つ正規トーリック射影三次元多様体とする。
同じ双指標 $c$ による $X$ の量子トーリック変形 $X_c$ を、
各トーリック affine chart の twisted semigroup algebra により定める。
関数斜体を共通トーラスの Ore 局所化で定めると
\[
 D(X_c)\simeq \operatorname{Frac}k_c[M]\simeq D(B_q),
\]
ここで
\[
 B_q=k\langle x_0,x_1,x_2,x_3\rangle/
 (x_0x_i-x_ix_0,\ x_ix_j-q_{ij}x_jx_i\ (1\le i<j\le3))
\]
は四生成元の量子多項式環である。
したがって、$X_c$ の関数斜体を上記の通り定める限り、
$X_c$ は AS-有理的である。この主張はトーリック fan の滑らかさを要しない。
\end{theorem}
\begin{proof}
$X$ の任意の最大錐 $\sigma$ の semigroup algebra $k[\sigma^\vee\cap M]$ を
$c$ で twist すると、全て同じ $k_c[M]$ の部分代数になる。
稠密トーラス上で全ての指標を可逆化した Ore 局所化は $k_c[M]$ であり、
各 chart が同じ分数斜体 $T_c$ を与える。
$B_q$ は反復 skew polynomial extension だから、noetherian domain で
Hilbert 級数 $(1-z)^{-4}$ の四次元 AS-regular algebra である。
$x_0$ は中央な非零斉次元である。$u_i=x_ix_0^{-1}$ と置けば
\[
 (B_q[x_0^{-1}])_0
 =k\langle u_1,u_2,u_3\rangle/(u_i u_j-q_{ij}u_j u_i).
\]
これら $u_i$ を可逆化すると $k_q[M]$ になり、上述の交代化の一致から
$k_q[M]\simeq k_c[M]$ である。
全斉次非零元の Ore 局所化における次数零の斜体は
$\operatorname{Frac}((B_q[x_0^{-1}])_0)$ に一致する。
実際、同次数 $d$ の $a,b$ の右分数は
$(ax_0^{-d})(bx_0^{-d})^{-1}$ で表せる。よって $D(B_q)\simeq T_c$ となる。
\end{proof}

\begin{remark}[適用範囲]
この定理は指定した\emph{双指標 twist} に関するものであり、
トーリック三次元多様体の任意の非可換変形が AS-有理的という主張ではない。
また、この同型は一般に極小模型や SOD の同一視を意味しない。
具体的な $\PP^3$, $\PP^1\times\PP^2$, $(\PP^1)^3$ の共通 Rees blow-up は
第\ref{sec:toric-nc-blowups}節で扱った。
\end{remark}

\subsection{対数正準 Poisson 関数体の中心}
$\Lambda=(\lambda_{ij})\in\operatorname{Mat}_3(k)$ を歪対称行列とし、
$\{X^m,X^n\}=\lambda(m,n)X^{m+n}$ により $k[M]$ を Poisson algebra とする。
$L_\Lambda=\{m\in M: \lambda(m,n)=0\text{ for all }n\in M\}$ と置く。

\begin{proposition}[有理 Casimir 体]\label{prop:toric-Poisson-center}
$\operatorname{char}k=0$ とすると、
\[
 Z_{\mathrm{Pois}}(k(M))=\operatorname{Frac}(k[L_\Lambda]).
\]
とくに $\lambda_{12},\lambda_{13},\lambda_{23}$ が $\mathbb Q$ 上一次独立なら
有理 Casimir は定数のみである。
\end{proposition}
\begin{proof}
$D_i=\{X^{e_i},-\}/X^{e_i}$ は Laurent 多項式環上の重み導分であり、
$D_i(X^m)=\lambda(e_i,m)X^m$ である。
$f=a/b\in k(M)$ を既約分数で $D_i f=0$ とすると、
Laurent 多項式環 $k[M]$ は UFD なので
$a\mid D_i a$ および $b\mid D_i b$ となる。
$D_i a$ が零でない場合、その Newton polytope は $a$ のものに含まれる。
$D_i a=ac$ と Newton polytope の積の Minkowski 和の等式を合わせると、
$\operatorname{Newt}(c)=\{0\}$ であり、$D_i a=\alpha_i a$、
同様に $D_i b=\alpha_i b$ となる。零の場合も $\alpha_i=0$ とすればよい。
全ての $i$ についてこれを行うと、$a$ と $b$ の支持の任意の二つの指数の差は
$L_\Lambda$ に属し、さらに $a$ と $b$ は同じ重みを持つ。
したがって $f$ は $L_\Lambda$ に属する単項式の有理式である。
逆向きの包含は定義から従う。
\end{proof}

\begin{remark}[形式的量子化]
$k$ を標数零の体とし、$q_{ij}=\exp(\hbar\lambda_{ij})\in k((\hbar))$ とする。
量子トーラスの Laurent 代数の中心の基底は
$X^m\;(m\in L_\Lambda)$ である。なぜなら
$\exp(\hbar\lambda(e_i,m))=1$ と
$\lambda(e_i,m)=0$ が同値だからである。
分数斜体の中心をこの Laurent 中心の分数体と同定するには
量子トーラスの中心の分数体定理を用いる。
従って量子 Casimir 中心と Poisson 有理中心は同じ整数格子で制御される。
\end{remark}

\subsection{中央線形元を持つ非可換平面の安定化}
\begin{proposition}[中央な $\PP^1$ 方向の AS-有理性]\label{prop:central-stabilization}
$S$ を連結次数付き noetherian AS-regular domain で
$\operatorname{gldim}S=3$, $H_S(z)=(1-z)^{-3}$ とする。
$S_1$ に非零の中央元 $h$ が存在すると仮定する。
新しい次数1の中央変数 $v$ を付けた $B=S[v]$ は四次元 AS-regular domain で
$H_B(z)=(1-z)^{-4}$ を持ち、
\[
 D(B)\simeq D(S)(w),\qquad w=vh^{-1}\text{ は中央}.
\]
したがって $D(S)(w)$ を関数斜体に持つ三次元量子化は AS-有理的である。
\end{proposition}
\begin{proof}
中央多項式拡大の AS-regular 性と Hilbert 級数の積公式による。
$h$ は中央非零の次数1元だから
$R=(S[h^{-1}])_0$ は domain で、その分数斜体は $D(S)$ に一致する。
実際、同次数 $d$ の斉次元による分数 $ab^{-1}$ は
$(ah^{-d})(bh^{-d})^{-1}$ と書ける。
$w=vh^{-1}$ は $S[h^{-1}]$ と可換であり、
$(B[h^{-1}])_0=R[w]$ となる。
従って $D(B)=\operatorname{Frac}(R[w])=D(S)(w)$ である。
\end{proof}
\begin{remark}
一般の楕円 Sklyanin 平面には中央な次数1の元がないので、この命題だけでは
垂直楕円型や三重積最大 pencil 型の未解決比較を解決しない。
強い中央性の仮定を外した一般化は別の研究問題である。
\end{remark}

\part{未解決問題と検証}
\section{残る問題と検証可能性}

本稿で得た非可換分数体の同型は、成分間の対応を作る具体的な選択を伴う。
同じ Poisson 双有理同値を持つ任意の二つの量子化について、
高次係数を含めた分数体の分類を決定する問題が残る。
量子トーラスの係数群と Riccati の微分 Ore 型の障害は、
この分類のための直接計算できる不変量となる。

共通圏の二つの収縮と射影環の再構成は
第\ref{sec:relative-common-blowup}節で古典点を含む族として得られた。
第\ref{sec:anticanonical-common-ring}節では反標準切断環を計算し、
平面の二点 blow-up の相対部分環として同定した。
可換二点 blow-up では、二点が相異なると反標準類は ample、
infinitely-near では $(-2)$ 曲線の収縮で ample でなくなる。
両者の反標準 Hilbert 関数は一致するため、数値情報だけでは復元できない。
さらに、この section functor から共通圏全体を復元するには、
coherent heart、構造対象、ample 性または既知の derived comparison を照合する。
特定の quadric 側の反標準表示との比較にも、点と周期性の標識を保持する必要がある。
さらに次の三つを調べる。
第一に、Sklyanin pencil の有限中心拡大後の同型を
$\mathbb C(u)$ 上へ下降する条件を求める。
二つの ruling と点・直線束の descent datum を保持して計算する。
第二に、第\ref{sec:native-pencil}節の Segre 斉次環と
第\ref{pcat:sec}節の三重積の圏量子化を比較する。
斉次環側の八対象の Hom と合成を、命題\ref{pcat:kernel-bimodule}の
階数32の cross bimodule に照合し、導来同値と generic localization を同時に固定する。
この比較は、元の偏極と中心を指定した非可換 contraction を与えるための次の段階である。
第三に、$R(1,3)$ と一般の $S(2,3)$ の関数体を調べ、
Riccati 型の $S(2,3)$ 内への対応の像とファイバーを決定する。
微分 Ore 型については導分の有理第一積分、定数体、および微分体同型が主要なデータになる。

本稿の主要な式は添付コードで再現できる。
\texttt{check\_nc\_p1p2.py} は九成分の正規形の有理変換、
\texttt{check\_p3\_nc.py} は CGG 積と Weyl 変換、
\texttt{check\_nc\_links.py} は分数線形変換、GWA の半古典式、
量子トーラスの単項式変換を検算する。
\texttt{check\_riccati\_p3.py} は根の正規化、斉次化および normal なチャートを検算する。
\texttt{check\_relative\_blowup.py} は共通圏の二つの例外集の Euler 行列、
反標準関手の $K_0$ 作用と Hilbert 多項式を照合し、
\texttt{check\_toric\_blowup.py} は双指標の合成則と Rees チャートを検算する。
\texttt{check\_pencil\_quantization.py} は Segre の二次関係式、
楕円完全交叉の有限次数の Koszul 階数、ruling の根格子と
共通圏の有限階数を照合する。
これらの有限検算と、本文の生成性・局所化・全次数の平坦性の証明を併用する。
全成分の網羅性、AS 正則性、曲面双有理定理の依存関係と証明箇所は
本文および同梱の研究方針に示した。


\begin{thebibliography}{99}
\bibitem{ArtinZhang1994} M. Artin and J. J. Zhang,
Noncommutative projective schemes, Adv. Math. 109 (1994), 228--287.
\url{https://doi.org/10.1006/aima.1994.1087}.
\bibitem{KeelerNCRings} D. S. Keeler,
The rings of noncommutative projective geometry,
arXiv:math/0205005 (2002), Theorem 2.3.
\url{https://arxiv.org/abs/math/0205005}.
\bibitem{MoriUeyamaWu2026} I. Mori, K. Ueyama and W. Wu,
Irreducible noncommutative quadrics, arXiv:2606.24606 (2026),
Theorems 1.3 and 5.17.
\url{https://arxiv.org/abs/2606.24606}.
\bibitem{DeDeyn2026} T. De Deyn,
A stacky resolution of singularities for orders by Azumaya algebras,
arXiv:2606.05137v2 (2026), Theorem 1.1.
\url{https://arxiv.org/abs/2606.05137}.


\bibitem{Pym2015} B. Pym,
Quantum deformations of projective three-space,
Adv. Math. 281 (2015), 1216--1241.
\href{https://doi.org/10.1016/j.aim.2015.06.005}{doi:10.1016/j.aim.2015.06.005},
\href{https://arxiv.org/abs/1403.6444}{arXiv:1403.6444}.
\bibitem{PVdB2016} D. Presotto and M. Van den Bergh,
Noncommutative versions of some classical birational transformations,
J. Noncommut. Geom. 10 (2016), 221--244.
\href{https://doi.org/10.4171/JNCG/232}{doi:10.4171/JNCG/232},
\href{https://arxiv.org/abs/1410.5207}{arXiv:1410.5207}.
\bibitem{Presotto2018} D. Presotto,
Symmetric noncommutative birational transformations,
J. Noncommut. Geom. 12 (2018), 733--778.
\href{https://doi.org/10.4171/JNCG/290}{doi:10.4171/JNCG/290},
\href{https://arxiv.org/abs/1607.08383}{arXiv:1607.08383}.
本文の定理番号は出版版による。
\bibitem{RSS2023} D. Rogalski, S. J. Sierra and J. T. Stafford,
Ring-theoretic blowing down II: Birational transformations,
J. Noncommut. Geom. 17 (2023), 1045--1097.
\href{https://doi.org/10.4171/JNCG/510}{doi:10.4171/JNCG/510},
\href{https://arxiv.org/abs/2107.01991}{arXiv:2107.01991}.
本文の定理番号は出版版による。
\bibitem{VdBQuadrics} M. Van den Bergh,
Non-commutative quadrics, Int. Math. Res. Not. 2011 (17), 3983--4026.
\href{https://doi.org/10.1093/imrn/rnq234}{doi:10.1093/imrn/rnq234},
\href{https://arxiv.org/abs/0807.3753}{arXiv:0807.3753}.
\bibitem{SmithVdB} S. P. Smith and M. Van den Bergh,
Noncommutative quadric surfaces, J. Noncommut. Geom. 7 (2013), 817--856.
\href{https://doi.org/10.4171/JNCG/136}{doi:10.4171/JNCG/136}.
\bibitem{LiuWangWu} L.-Y. Liu, S.-Q. Wang and Q.-S. Wu,
Twisted Calabi--Yau property of Ore extensions,
J. Noncommut. Geom. 8 (2014), 587--609.
\href{https://doi.org/10.4171/JNCG/165}{doi:10.4171/JNCG/165},
\href{https://arxiv.org/abs/1205.0893}{arXiv:1205.0893}.
\bibitem{GoodearlLaunois} K. R. Goodearl and S. Launois,
The Dixmier--Moeglin equivalence and a Gel'fand--Kirillov problem for Poisson polynomial algebras,
Bull. Soc. Math. France 139 (2011), 1--39.
\href{https://doi.org/10.24033/bsmf.2598}{doi:10.24033/bsmf.2598}.
\bibitem{GoodearlYakimov} K. R. Goodearl and M. T. Yakimov,
Twist invariants of graded algebras, Contemp. Math. 715 (2018), 141--170.
\href{https://arxiv.org/abs/1706.06957}{arXiv:1706.06957}.
\bibitem{OrlovSing} D. Orlov,
Derived categories of coherent sheaves and triangulated categories of singularities,
\href{https://arxiv.org/abs/math/0503632}{arXiv:math/0503632v3}.
定理2.5および命題2.17を用いる。
\bibitem{BallardFarman} M. R. Ballard and B. A. Farman,
Kernels for noncommutative projective schemes,
\href{https://arxiv.org/abs/1709.06470v2}{arXiv:1709.06470v2}.
本文の定理番号はこの版による。
\bibitem{AKMW} D. Abramovich, K. Karu, K. Matsuki and J. W{\l}odarczyk,
Torification and factorization of birational maps,
J. Amer. Math. Soc. 15 (2002), 531--572.
\href{https://doi.org/10.1090/S0894-0347-02-00396-X}{doi:10.1090/S0894-0347-02-00396-X},
\href{https://arxiv.org/abs/math/9904135}{arXiv:math/9904135}.
\bibitem{KuznetsovSODBaseChange} A. Kuznetsov,
Base change for semiorthogonal decompositions,
Compos. Math. 147 (2011), 852--876.
\url{https://doi.org/10.1112/S0010437X10005166}.
\bibitem{OrlovBlowup} D. O. Orlov,
Projective bundles, monoidal transformations, and derived categories of coherent sheaves,
Russian Acad. Sci. Izv. Math. 41 (1993), 133--141.
\href{https://doi.org/10.1070/IM1993v041n01ABEH002182}{doi:10.1070/IM1993v041n01ABEH002182}.
\bibitem{NeemanRanicki} A. Neeman and A. Ranicki,
Noncommutative localisation in algebraic K-theory I,
Geom. Topol. 8 (2004), 1385--1425.
\href{https://doi.org/10.2140/gt.2004.8.1385}{doi:10.2140/gt.2004.8.1385},
\href{https://arxiv.org/abs/math/0410620}{arXiv:math/0410620}.
\bibitem{Polishchuk1997} A. Polishchuk, \emph{Algebraic geometry of Poisson brackets}, J. Math. Sci. 84 (1997), 1413--1444. \href{https://doi.org/10.1007/BF02399197}{doi:10.1007/BF02399197}.
\bibitem{p12-Cogolludo} J. I. Cogolludo-Agustín, \emph{Braid monodromy of algebraic curves}, Ann. Math. Blaise Pascal 18 (2011), 141--209. \href{https://doi.org/10.5802/ambp.295}{doi:10.5802/ambp.295}.
\bibitem{p12-Milnor} J. Milnor, \emph{Singular Points of Complex Hypersurfaces}, Annals of Mathematics Studies 61, Princeton University Press, 1968.
\bibitem{p12-LPT} F. Loray, J. V. Pereira, F. Touzet, \emph{Foliations with trivial canonical bundle on Fano 3-folds}, Math. Nachr. 286 (2013), 921--940. \href{https://doi.org/10.1002/mana.201100354}{doi:10.1002/mana.201100354}. \href{https://arxiv.org/abs/1112.3871v3}{arXiv:1112.3871v3}.
\bibitem{p12-GualtieriPym} M. Gualtieri, B. Pym, \emph{Poisson modules and degeneracy loci}, Proc. Lond. Math. Soc. 107 (2013), 627--654. \href{https://doi.org/10.1112/plms/pds090}{doi:10.1112/plms/pds090}.
\bibitem{p12-Pym} B. Pym, \emph{Constructions and classifications of projective Poisson varieties}, Lett. Math. Phys. 108 (2018), 573--632. \href{https://doi.org/10.1007/s11005-017-0984-5}{doi:10.1007/s11005-017-0984-5}.
\bibitem{p12-BFT} P. Belmans, E. Fatighenti, F. Tanturri, \emph{Polyvector fields for Fano 3-folds}, Math. Z. 304 (2023), article 12. \href{https://doi.org/10.1007/s00209-023-03261-2}{doi:10.1007/s00209-023-03261-2}.
\bibitem{p12-Okumura} K. Okumura, \emph{SNC Log Symplectic Structures on Fano Products}, Canad. Math. Bull. 63 (2020), 891--900. \href{https://doi.org/10.4153/S0008439520000120}{doi:10.4153/S0008439520000120}.

\bibitem{cube-Fano} P. Belmans, Fanography, family 3-27,
\url{https://fanography.info/3-27/} (2026年11月6日確認).
\bibitem{cube-KN} G. Kempf, L. Ness, The length of vectors in representation spaces,
in Algebraic Geometry, Lecture Notes in Math. 732 (1979), 233--243.
\href{https://doi.org/10.1007/BFb0066647}{doi:10.1007/BFb0066647}.
\bibitem{cube-LT} J.-G. Luque, J.-Y. Thibon, The polynomial invariants of four qubits,
Phys. Rev. A 67 (2003), 042303.
\href{https://arxiv.org/abs/quant-ph/0212069}{arXiv:quant-ph/0212069}.
\bibitem{cube-SM} C. Sanabria Malagón, An elementary proof of the symplectic spectral theorem,
Asian-Eur. J. Math., \href{https://arxiv.org/abs/1710.00697}{arXiv:1710.00697},
\href{https://doi.org/10.1142/S1793557119500335}{doi:10.1142/S1793557119500335}.

% 2026-11-08 v7. Insert inside thebibliography.
% RSS2023 is already present in the main manuscript and is not duplicated here.
\bibitem{BlowupVdB2001} M. Van den Bergh,
Blowing up of non-commutative smooth surfaces,
Mem. Amer. Math. Soc. 154 (2001), no. 734.
\url{https://arxiv.org/abs/math/9809116}.

\bibitem{BlowupRSS2015} D. Rogalski, S. J. Sierra and J. T. Stafford,
Noncommutative blowups of elliptic algebras,
Algebr. Represent. Theory 18 (2015), no. 2, 491--529.
\url{https://doi.org/10.1007/s10468-014-9506-7}.
\url{https://arxiv.org/abs/1308.2216}.

\bibitem{BlowupRSS2017} D. Rogalski, S. J. Sierra and J. T. Stafford,
Ring-theoretic blowing down. I,
J. Noncommut. Geom. 11 (2017), no. 4, 1465--1520.
\url{https://doi.org/10.4171/JNCG/11-4-9}.

\bibitem{BlowupPresotto2017} D. Presotto,
$\mathbb Z^2$-algebras as noncommutative blow-ups,
arXiv:1701.00413v2 (2017).
\url{https://arxiv.org/abs/1701.00413}.

\bibitem{BlowupBP2018} P. Belmans and D. Presotto,
Construction of non-commutative surfaces with exceptional collections of length 4,
J. London Math. Soc. (2) 98 (2018), no. 1, 85--103.
\url{https://doi.org/10.1112/jlms.12126}.
\url{https://arxiv.org/abs/1705.06943}.

\bibitem{BlowupMOU2019} I. Mori, S. Okawa and K. Ueda,
Moduli of noncommutative Hirzebruch surfaces,
arXiv:1903.06457 (2019).
\url{https://arxiv.org/abs/1903.06457}.

\bibitem{BlowupKL2015} A. Kuznetsov and V. A. Lunts,
Categorical resolutions of irrational singularities,
Int. Math. Res. Not. IMRN 2015 (2015), no. 13, 4536--4625.
\url{https://doi.org/10.1093/imrn/rnu072}.
\url{https://arxiv.org/abs/1212.6170}.

\bibitem{BlowupKK2015} D. Kaledin and A. Kuznetsov,
Refined blowups,
Math. Res. Lett. 22 (2015), no. 6, 1717--1732.
\url{https://doi.org/10.4310/MRL.2015.v22.n6.a9}.
\url{https://arxiv.org/abs/1410.7121}.

\bibitem{BlowupKRS2005} D. S. Keeler, D. Rogalski and J. T. Stafford,
Na\"ive noncommutative blowing up,
Duke Math. J. 126 (2005), no. 3, 491--546.
\url{https://arxiv.org/abs/math/0306244}.

\bibitem{BlowupRogalski2019} D. Rogalski,
Well-closed subschemes of noncommutative schemes,
Theory Appl. Categ. 34 (2019), no. 14, 375--404.
\url{https://www.tac.mta.ca/tac/volumes/34/14/34-14abs.html}.
\url{https://arxiv.org/abs/1707.07052}.

\bibitem{BlowupRains2025} E. M. Rains,
The (noncommutative) geometry of difference equations,
arXiv:2504.16187v1 (2025).
\url{https://arxiv.org/abs/2504.16187}.
\bibitem{BlowupVdBFlops2004} M. Van den Bergh,
Three-dimensional flops and non-commutative rings,
Duke Math. J. 122 (2004), 423--455.
\url{https://arxiv.org/abs/math/0207170}.

\bibitem{BlowupVdBNCCR2004} M. Van den Bergh,
Non-commutative crepant resolutions,
in \emph{The Legacy of Niels Henrik Abel}, Springer, Berlin, 2004, 749--770.
訂正版 arXiv:math/0211064v2 (2009).
\url{https://arxiv.org/abs/math/0211064}.


\bibitem{BlowupAbstract} M. Van den Bergh,
Abstract blowing down,
Proc. Amer. Math. Soc. 128 (2000).
\href{https://arxiv.org/abs/math/9806046}{arXiv:math/9806046}.

\bibitem{ToricCLS} L. Cirio, G. Landi and R. J. Szabo,
Algebraic deformations of toric varieties I. General constructions,
Adv. Math. \textbf{246} (2013), 33--88.
\url{https://doi.org/10.1016/j.aim.2013.06.025}.
\url{https://arxiv.org/abs/1001.1242}.

\bibitem{ToricElagin} A. Elagin,
Descent theory for semiorthogonal decompositions,
Sb. Math. \textbf{203} (2012), no.~5, 645--676.
\url{https://doi.org/10.1070/SM2012v203n05ABEH004238}.
\url{https://arxiv.org/abs/1206.2881}.
\bibitem{ToricCox}
D. A. Cox,
\emph{The homogeneous coordinate ring of a toric variety},
J. Algebraic Geom. \textbf{4} (1995), 17--50.
\href{https://arxiv.org/abs/alg-geom/9210008}{arXiv:alg-geom/9210008}.


\bibitem{RelativeArtin} The Stacks Project Authors,
The approximation property for G-rings,
Theorem 16.13.2, Tag 07QZ.
\url{https://stacks.math.columbia.edu/tag/07QZ}.

\bibitem{RelativeLowen} W. Lowen,
Obstruction theory for objects in abelian and derived categories,
Comm. Algebra 33 (2005), no.~9, 3195--3223.
\url{https://doi.org/10.1081/AGB-200066155}.
\url{https://arxiv.org/abs/math/0407019}.

\bibitem{RelativeBOR} P. Belmans, S. Okawa and A. T. Ricolfi,
Moduli spaces of semiorthogonal decompositions in families,
arXiv:2002.03303v3 (2025).
\url{https://arxiv.org/abs/2002.03303}.

\bibitem{QuadricBR2020} P. Belmans and T. Raedschelders,
Derived categories of noncommutative quadrics and Hilbert squares,
Int. Math. Res. Not. 2020 (19), 6042--6069.
\url{https://doi.org/10.1093/imrn/rny192}.
\url{https://arxiv.org/abs/1605.02795}.
\bibitem{QuadricMU2021} I. Mori and K. Ueyama,
A categorical characterization of quantum projective spaces,
J. Noncommut. Geom. 15 (2021), no.~2, 489--529.
\url{https://doi.org/10.4171/JNCG/403}.
\url{https://arxiv.org/abs/1708.00167}.

\bibitem{PencilRainsFiltered} E. M. Rains,
Filtered deformations of elliptic algebras,
arXiv:2107.13540v2 (2022).
\url{https://arxiv.org/abs/2107.13540}.
\bibitem{PencilRainsBirational} E. M. Rains,
The birational geometry of noncommutative surfaces,
arXiv:1907.11301 (2019).
\url{https://arxiv.org/abs/1907.11301}.
\bibitem{PencilOrlovGluing} D. Orlov,
Smooth and proper noncommutative schemes and gluing of DG categories,
Adv. Math. 302 (2016), 59--105.
\url{https://doi.org/10.1016/j.aim.2016.07.014}.
\url{https://arxiv.org/abs/1402.7364v5}.
\bibitem{ACAnnoLogvinenko} R. Anno and T. Logvinenko,
Spherical DG-functors,
J. Eur. Math. Soc. 19 (2017), 2577--2656.
\url{https://doi.org/10.4171/JEMS/724}.
\url{https://arxiv.org/abs/1309.5035v2}.
\bibitem{ACKeelerAmple} D. S. Keeler,
Ample filters of invertible sheaves,
J. Algebra 259 (2003), no.~1, 243--283.
Corrigendum, J. Algebra 507 (2018), 592--598.
訂正を含む arXiv:math/0108068v5 を用いる。
\url{https://doi.org/10.1016/S0021-8693(02)00557-4}.
\url{https://arxiv.org/abs/math/0108068v5}.
\bibitem{ACBocklandt} R. Bocklandt,
Graded Calabi Yau algebras of dimension 3,
J. Pure Appl. Algebra 212 (2008), 14--32.
\url{https://doi.org/10.1016/j.jpaa.2007.03.009}.
\url{https://arxiv.org/abs/math/0603558}.

\bibitem{VdBGlobal} M. Van den Bergh,
On global deformation quantization in the algebraic case,
J. Algebra 315 (2007), 326--395.
\href{https://arxiv.org/abs/math/0603200}{arXiv:math/0603200}.
\bibitem{BW2024} S. Barmeier and Z. Wang,
Deformations of categories of coherent sheaves via quivers with relations,
Algebraic Geometry 11 (2024), 1--36.
\href{https://doi.org/10.14231/AG-2024-001}{doi:10.14231/AG-2024-001},
\href{https://arxiv.org/abs/2107.07490}{arXiv:2107.07490}.
\bibitem{Schlessinger} M. Schlessinger,
Functors of Artin rings, Trans. Amer. Math. Soc. 130 (1968), 208--222.
\href{https://doi.org/10.2307/1994967}{doi:10.2307/1994967}.
\bibitem{StacksArtin} The Stacks Project Authors,
\'Etale neighbourhoods and Artin approximation, Lemmas 37.39.1--37.39.2.
\href{https://stacks.math.columbia.edu/tag/0CAT}{Tag 0CAT}.
\bibitem{CLS} L. Cirio, G. Landi and R. J. Szabo,
Algebraic deformations of toric varieties I. General constructions,
Adv. Math. 246 (2013), 33--88.
\href{https://arxiv.org/abs/1001.1242}{arXiv:1001.1242}.
\bibitem{MoriSmith} I. Mori and S. P. Smith,
The classification of 3-Calabi--Yau algebras with 3 generators and 3 quadratic relations,
Math. Z. 287 (2017), 215--241.
\href{https://doi.org/10.1007/s00209-016-1824-5}{doi:10.1007/s00209-016-1824-5},
\href{https://arxiv.org/abs/1502.07403}{arXiv:1502.07403}.
\bibitem{OkawaUedaQuadrics} S. Okawa and K. Ueda,
Noncommutative quadric surfaces and noncommutative conifolds.
\href{https://arxiv.org/abs/1403.0713}{arXiv:1403.0713}.
\bibitem{KontsevichAlgebraic} M. Kontsevich,
Deformation quantization of algebraic varieties,
Lett. Math. Phys. 56 (2001), 271--294.
\href{https://arxiv.org/abs/math/0106006}{arXiv:math/0106006}.
\bibitem{PolishchukQuant} A. Polishchuk, Algebraic geometry of Poisson brackets,
J. Math. Sci. 84 (1997), 1413--1444.
\href{https://pages.uoregon.edu/apolish/alg-geom-poisson.pdf}{著者公開PDF}.

\bibitem{StaffordVanDenBergh2001} J. T. Stafford and M. Van den Bergh,
Noncommutative curves and noncommutative surfaces,
Bull. Amer. Math. Soc. 38 (2001), 171--216.
\url{https://arxiv.org/abs/math/9910082}.
\bibitem{RogalskiSierraStafford2020} D. Rogalski, S. J. Sierra and J. T. Stafford,
Some noncommutative minimal surfaces,
Adv. Math. 369 (2020), 107151.
\url{https://doi.org/10.1016/j.aim.2020.107151}.
\end{thebibliography}
\end{document}
